\documentclass[11pt,a4paper]{article}
\usepackage[T1]{fontenc}
\usepackage[utf8]{inputenc}
\usepackage{lmodern}
\usepackage{microtype}
\usepackage{graphicx}
\usepackage{slashed}
\usepackage{tikz}
\usepackage{colortbl}
\usepackage{calc}



% Current II-D reference root for PapersMinimal v3.063; mathematical macros otherwise unchanged.
% Macros-v2_188: re-based current paper identities and A--G family; no change to geometry.
% Macros-v2_184: active Paper III-E v1.200/III-F v1.267 binding for framework v3.031; mathematical macros unchanged.
% Macros-v2_183: updated shared cross-reference bindings for Paper1B/III-E/III-F v3.030; mathematical macros unchanged.
% Macros-v2_182: active Paper II-F/V-A cross-references for framework v3.029; definitions unchanged.
% Paper-series v2.381: new I-B torsion/modular results, bibliography and current source roots.
% Paper-series v2.380: new exceptional bibliography and synchronized paper references.
% Paper-series v2.379 current 31-paper macro/aux registry.
% Paper-series v2.378: current 31-paper AUX/bibliography registry updated for Paper I-B v1.2 literature additions.
% Framework v2.371: current-source reference registry refreshed; old theorem labels retained.
% Framework v2.370: current 30-paper source/AUX root rebinding; native mathematical macros unchanged.
% Framework v2.368: synchronized merged 30-root registry; mathematical notation unchanged.
% Framework v2.358: active 30-paper source registry synchronized; existing mathematical macros unchanged.
% Framework v2.357: selective new source-owned certificates; active reference roots updated for five revised papers.
% Framework v2.356: active 30-root registry synchronized to the 30 papers in this release.
% Framework v2.275: exact positive dense-twin word cubature, finite repeatability and eighth-moment classification, nested sampling frames, joint action and POVM-jet exports; preserves v2.274.
% Framework v2.274: filename and root-source synchronization.
% Framework v2.273: restored missing Problem 3 theorems from v2.268; exact exceptional frames, twin geometry, convolution repeatability and projector/occurrence firewalls; merged on v2.272 without overwriting independent branches.
% Framework v2.272: merged v2.270 covariant native-readout mathematics with v2.271 Problem 4 finite-action Hopf hypercharge/holonomy corrections; no universal bulk-gauge uniqueness claim.
% Independent branches: v2.270 readout/guard exports and v2.270 K16 action-consumed charge closure; all prior scoped assumptions retained.
% Framework v2.269: synchronized 30 active AUX roots; mathematics and bibliography preserved from Macros-v2_161.
% Framework v2.266: Paper II-C converse-derived outcome resources; II-E, III-C, IV-A/B, V-A and VI-B downstream mathematical closures; 30-root registry synchronized.
% Framework v2.265: II-C coordinate-stage constitution and exact first-symbol host inheritance; prior mathematical domains retained.
% Framework v2.264: paper II-B reverse/audit and IV-A observable-minimal exchange; synchronized external references to 30 active roots.
% Framework v2.263: Paper I-C adds bridge-controlled root-carry parity and source-relative old-cap cancellation; other root mathematical bodies retained; active 30-root references synchronized.
% Framework v2.262: integrates Paper II-F radius, Casimir and readout results; active 30-root references synchronized.
% Framework v2.261: merged Paper II-B GP2-level/geometric repairs; preserves v2.260 II-D repairs; active references synchronized.
% Framework v2.260: Paper II-D signed RS, finite-symbol typing, and conditional curvature/Jordan repairs; references synchronized.
% Framework v2.259: conditional native SU(2) atlas and explicit later discharge; v2.258 mathematical bodies retained elsewhere.
% Framework v2.258: preserves parallel II-A/dynamics update and adds complete readout/frequency closures.
% Framework v2.257: exact minimal complete position readout and arbitrary-preparation frequency law.
% Framework v2.256: exact native inertia and sharp bounded ordered-source closure; preserves concurrent v2.255 work.
% Framework v2.255: preserves parallel detector integration and adds exact finite readout closures.
% Framework v2.254: Pasterski bibliography and synchronized 30-paper external-reference registry.
% Framework v2.253: typed raw coordinate defect in I-A and synchronized VI-A import; active references synchronized.
% Framework v2.252: synchronized source registry; optional III-G shell-frame calculation.
% Framework v2.251: forward common-action execution, joint source jets and scoped coupled control.
% Framework v2.251: exact full-heat variations, fermion source jets and explicit physical completion frontier.
% Framework v2.251: certified interacting source jets, pole stability and scoped completion proofs.
% Framework v2.251: synchronized active source registry; mathematical notation unchanged.
% Framework v2.251: synchronized active source registry; mathematical notation unchanged.
% Framework v2.251: synchronized active source registry; mathematical notation unchanged.
% Framework v2.251: synchronized source registry; no mathematical notation changes.
% Framework v2.244: synchronized 30-paper AUX registry and rooted phase-source action-order tests.
% Framework v2.242: signed rooted phase lift, complete transverse return and exact family cubic.
% Framework v2.241: proved rooted conormal geometry, conditional cap realization and scoped action tests.
% Framework v2.239: all 30 active source roots synchronized; no mathematical macro changes.
% Framework v2.238: Paper IV A/B/C split and 30-root reference registry.
% Framework v2.237: Paper VII-D technical split and 28-root reference registry.
% Framework v2.236: Paper VI-C technical-volume reference added; all current root bindings advanced.
% Framework v2.235: 26-root AUX and bibliography bindings synchronized; mathematical macros unchanged.
% BEGIN INLINE SOURCE: Paper00-TeXBuildCompatibility-v1_0.tex
% Rendering-only repairs for inherited long appendices and bitmap text fonts.
% Integer counters through 26 retain the original labels; later values use AA,
% AB, ... . This file changes no action coefficient, state or theorem premise.
\RequirePackage{alphalph}
\makeatletter
\let\BHPrimitiveAlph\@Alph
\newalphalph{\BHExtendedAlph}{\BHPrimitiveAlph}{26}
\let\@Alph\BHExtendedAlph
\AtBeginDocument{%
  \@ifpackageloaded{microtype}{\microtypesetup{expansion=false}}{}%
}
\makeatother

% END INLINE SOURCE: Paper00-TeXBuildCompatibility-v1_0.tex

% Framework v2.164: current reference/bibliography profile; mathematical definitions retained.
% Framework v2.162: finite metric, functional Hilbert and scoped Ward integration.
% Framework v2.161: inherits Macros-v2_57.tex exactly except the current external-reference binding.
% Macros v2.57 / framework v2.156: current reference registry only; mathematical definitions unchanged.
% Framework v2.143 / Macros v2.55: current external-reference bindings; mathematical notation unchanged.
% Framework v2.110 / Macros v2.54: current external-reference registry and bibliography; mathematical notation inherited.
% Framework v2.67 / Macros-v2_53: adds Born-action-cell and joint-EH readout symbols; absolute unit realization is metrological; MATCH_{J/scal}^{EH} remains the sole G24 normalization bridge.
% Framework v2.35 / Macros-v2_52: grouped paper numbering, active AUX registry, and Born-v6_20 bibliography binding synchronized; notation unchanged.
% Framework v2.26: active-source AUX and bibliography bindings synchronized; notation unchanged.
% Framework v2.09: reviewed early-reciprocity integration; scoped physical imports.
% Framework v2.08: Born--Euler mathematics and conditional complete-action potential criterion.
% Framework v2.07: active AUX registry includes Paper I-B v1.2; shared bibliography Born-v6_18.
% Macros-v2_46 / framework v2.04: mathematical atlas and downstream realization correction.
% Macros-v2_46 / framework v2.03: scoped structure integration and current dependency bindings.
% Macros-v2_46 / framework v2.01: shared source bindings; prior revision comments are historical.
% v2.42 / framework v1.99: Synchronized RG/AS owner and bibliography bindings.
% v2.41 / framework v1.99: active manuscript AUX bindings and geometric-closure bibliography synchronized.
% Release-v1.80 centralized source synchronization for scalar-circle/comparison-band integration.
% Release v1.79: Hurwitz hierarchy and kinetic-normalized action response integration; physical mass matching remains conditional.
% Release v1.78: approved M8 closure-batch integration; physical and analytic domains retained.
% Current v1.78 registry: all active TeX roots are synchronized below.
% Earlier release comments are historical scope records, not current status.
% Release-v1.77 synchronization: native same-owner Noether closure, generator--flow/history typing, analytic-certification firewalls, scalaron/K16, LF/QCD, quantum-Ward, and RAR boundary-owner transport are registered; shared mathematical notation is unchanged.
% Release-v1.76 synchronization: registry points to the v1.76 active paper roots.
% Release-v1.76 synchronization: registers the local A+Phi+g Noether/constraint closure and dynamics-frontier propagation; shared mathematical notation is unchanged.
% Release-v1.75 synchronization: registry points to the v1.75 active paper roots.
% Release-v1.75 synchronization: integrates the joint-Dirac Einstein bulk real form and the local odd-scalar heat-coefficient action on explicit analytic first-jet/interior domains; v1.74 Yang--Mills closure is retained. Fermionic reality, pure-metric curvature-squared/exact-nonlocal terms, noncanonical boundaries, physical vacuum selection, and actual solution-history tokens remain open.
% Release-v1.74 synchronization: registry points to the v1.74 active paper roots and retains the v1.73 theorem labels.
% Release-v1.74 synchronization: canonical Macros-v2_34.tex; generated local nonlinear Yang--Mills action, constrained connection-field phase space, source-free Gauss momentum map, and Legendre closure are integrated on their proved domains. The remaining non-gauge ACT_cl sectors and actual solution-history selection remain open.
% Release-v1.71 synchronization: closes the former JOINT_x/JOINT_5
% connection-occurrence gap by canonical structure-group extension of the
% nonflat Born connection, faithful associated representation, and the local
% cotangent--Clifford operator constructor.  Actual fluctuation members,
% actions, vacua, solutions, calibrations, and analytic QFT domains remain
% separately typed.  The active-AUX registry is advanced once.
% Release-v1.71 canonical macro and active-AUX registry.
% ======================================================================
% Macros-v2_39.tex -- canonical publication macro file
% Release v1.71 retains the v1.70 null-ray/QCD boundaries and additionally
% removes the duplicate structural connection actuality premise.  It does not
% claim global bundle topology, select a gauge/Higgs member, or promote the
% finite construction to an interacting continuum theory.
% Release v1.68 constructs the finite light-front total-fraction operators and
% parton number/momentum operator measures, proves their simplex and Hausdorff-limit
% theorems, and separates positive cut measures from scheme-renormalized
% light-ray distributions.  It also separates finite real Schur poles, stable
% norm-resolvent/Riesz pole-residue limits, and continued resonance
% pole/residue limits.  The Paper-VII-C scattering package is claim-indexed:
% local QFT and spectral claims import no S-matrix or infrared premise, while
% asymptotic amplitudes, LSZ, S-matrix poles, completeness, and cross sections
% import only their named massive or massless/long-range clauses.
% Release v1.67 factorizes BND48_H into a same-action Hessian/Schur
% descendant and still-open typed boundary-trace, critical-lift, and
% refinement-limit conditions.  It proves the finite-collar Dirichlet-to-
% Neumann law sqrt(y)tanh(L sqrt(y)), its nonuniform critical limit, and the
% finite-return ultra-infrared turnover.  The Paper-I-A Hardy constant mode is
% an exact geometric near-hit, not yet the physical metric boundary mode.
% The current local f(R) Hessian remains massive and therefore still refutes
% a RAR prediction on the minimal action branch.
% Release v1.66 derives the transverse depth-six normalizer bridge from the
% already generated C6 Reynolds image and Paper-I-A binary-dihedral normalizer.
% The bridge is the unique least 2D3-stable scalar-containing realization of
% the faithful quotient 2D3/C6=C2, has trace defect six and hence K=48, but is
% transverse to the nested active carrier.  Exact K=44 and K=52 contraction,
% semigroup, and energy-midpoint identities are registered as structural
% readouts rather than projector events.  The active E48 carrier gap remains
% unchanged.  A same-action boundary-symbol/Calderon intertwiner is isolated
% as the necessary and still unproved owner bridge to the static DtN sort;
% representation equivalence alone neither derives RAR nor creates a K=48
% action junction, pole, occupation, or cosmological transition.
% Release v1.65 derives the native E_48-centered carrier mass gap from the
% GP1--GP3-generated active spectrum, distinguishes its carrier-free interval
% from a particle located at E_48, and withdraws the unsupported promotion to
% a positive continuum Yang--Mills mass gap.  The latter remains a valid
% conditional transfer theorem on an enlarged continuum branch, but it is not
% an active framework prediction or an unfinished native-gap proof.
% Release v1.64 synchronizes Derivation Batches 14--22.  It registers the
% GP1--GP3 heat-profile orbit, same-Hessian electroweak response and pole
% ownership, K16/K48 action-response corrections, perturbation state and
% non-Gaussianity typing, finite-vacuum/gap/QCD limit repairs, factorized
% position-area ownership, the actual static metric response, and the ordered
% Wilson/string-tension limits.  Strict continuum convergence, positive
% physical gap, surviving higher cumulants, and positive tension remain open.
% Release v1.61 closes TIMEWIDTH16 as an independent direct-production datum.
% JUNC_{2|7} continues to own the finite intrinsic microscopic profile/width,
% while the actual same-action background/production history determines its
% physical-time pullback, support, and duration.  The current Paper-I-A--VII-C AUX
% registry is advanced once; shared mathematical notation is unchanged.
% Release v1.60 derives the final-octet primitive-remnant depletion, canonical
% Hopf-current representative, exact Hopf descent to dOmega/(4pi), and
% closed-FLRW a^{-3} neutral-current shape.  It retypes K=48 as the onset
% of the final one-unit geometric transfer and records cDM/RAR as distinct
% physical projections of that common neutral parent while preserving
% occurrence, Hilbert-owner, stress-matching, and static-kernel premises.
% Release v1.59 transports the v1.58 exceptional-shell/K16 closure through
% Papers III-F and VI-B--VII-C.  It replaces Paper VI-B's obsolete mixed-Hessian
% output typing by the third-variation pair vertex and Bogoliubov occupation
% support, derives the post-seed Higgs-pair feedback vertex and its exact
% alpha--delta kernel factors, propagates the unique K16/K48 action-junction
% split, identifies the QCD delta count with the upstream 2O shell atom only
% as a common-scalar numerical compatibility, and narrows THERM16 to dynamical
% approach/beta selection because the finite Gibbs/KMS target is already
% unique on a supplied simple observable algebra once its Hamiltonian and
% inverse temperature are supplied; direct-sum central weights remain typed.
% Release v1.58 registers the exceptional nearest-shell closure, the derived
% transverse K=16 action event, the unique interior action junction, the
% type-correct ordered pair vertex, the thin-layer squeeze law, and the
% radiation-like/matter-like kinematic split.  REG_0, transfer magnitude,
% physical layer width, and thermalization remain explicitly conditional.
% Release v1.57 registers the active-depth-six obstruction, the native
% K=48 carrier desert, the typed 7|9 owner reduction, and the dynamical
% reheating handoff across Papers III-B--VI-B.
% Release v1.56 registers the least-support n_2 inflation dressing of the
% persistent n_7 Laurent response, the analytic global-convexity and positive
% endpoint-measure proofs, the derived inflation core and K=16 coefficient
% freeze, and the no-particle dark-sector split into static RAR plus the
% conditional homogeneous effective-curvature response at reflected K=48.
% There is no n_6 action, scalaron-DM mass, or second Einstein--Hilbert term.
% Release v1.55 registers the protected-octet resolved/residual reflection
% q_sp+Z_W=8, its j<->8-j and K<->64-K scale covariance, and the derived
% rank-one-exit quadratic grade C_1^2/U=C_2 in Papers III-C, III-E, and VIII.
% The reflection is a typed trace/scale theorem, not a time reversal or a
% physical event-occurrence theorem; reheating transfer remains conditional.
% Release v1.54 registers the exact positive-part Laurent indices, the
% active-domain/terminal-endpoint action split, the no-second-EH firewall, and
% the q_DM=6/K_DM=48 scalaron readout typing across Papers III-D/E/F, VI,
% VIII, and IX.
% Release v1.53 registers the common nonnegative Laurent obstacle projection
% for n_1^+ and n_7^+, discharges the three former exit/freeze assumptions,
% and advances the Paper-III-D/E/F, VIII, and IX consumers.  Shared notation
% and bibliography ownership are unchanged.
% Release v1.51 registers the PQTI derivation batch and its updated
% Paper-III-A/C/D/E/F, VIII, and IX consumers.  Shared notation and
% bibliography ownership are unchanged.
% Release v1.50 advances the Paper-III-B AUX target to v1.40 after the
% rigorous audit repair; shared notation and bibliography ownership are unchanged.
% Release v1.49 advances the Paper-II-E/F AUX target to v3.14 after the
% minimal audit repair; shared notation and bibliography ownership are unchanged.
% Release v1.46 synchronizes Paper II-A's minimal frozen constructor-origin
% theorem and its five dependent AUX targets.  Shared notation is unchanged.
% Release v1.45 updates the Paper-II-B AUX target to v2.23 after the
% GP3-rigidity and appendix-placement repair; shared notation is unchanged.
% Release v1.44 advances the centralized Paper-I-A--VII-C AUX registry after the
% Paper-I-A v5.16 protected-completion scope repair.  Every current
% paper imports this one file; no paper-level opt-in or \externaldocument list
% is permitted.  The bibliography remains centralized here through Born-v6_13.
% Shared bibliography, theorem infrastructure, colors, notation alignment,
% and current Paper-I-A--VII-C external-reference registration.
% This file centralizes common imports, bibliography sorting, colors, theorem
% infrastructure, and the notation shared by the complete series.
% v2.9 (release v1.43) made the Paper-I-A--VII-C AUX registry automatic and
% centralized the bibliography and all external-reference ownership.
% v2.8 release-registry refresh (v1.33) registers the rooted generated-parent
% cascade and the POL/REG/free-law/11+4/Born-group closure revisions.  No
% shared mathematical notation or Paper-II-A fundamental principle is changed.
% v2.8 release-registry refresh (v1.32) registers the universal generated-
% continuation revisions in Papers II-A, III-A, III-C, IV, and VIII.  No
% shared mathematical notation or Paper-II-A fundamental principle is changed.
% v2.8 release-registry refresh (v1.30) registers the CRP4/projective-state
% source revisions and their direct registry-only consumers.  No shared
% mathematical notation or Paper-II-A axiom is changed.
% v2.7 release-registry refresh (v1.27) updates only the current external
% Paper-I-A--VII-C basenames after the placement/assumption synchronization; no
% shared mathematical notation or Paper-II-A axiom is changed.
% v2.7 synchronizes the geometric defect-repair revision: Paper-II-A v3.42,
% Paper-III-C v1.26, Paper-III-D v1.27, Paper-III-E v1.32, and Paper-III-F
% v1.34.  It registers the common trace-defect/spectral-filtration theorem,
% its owner-typed response pullback, and the reduced GRED and TERM interfaces.
% No shared mathematical notation is changed.
% v2.6 synchronizes Paper-III-E v1.31's cone-type repair, minimal-response
% classification, and explicit classical action-role boundary, together with
% Paper-III-F v1.32's downstream response-atlas synchronization.  No shared
% mathematical notation is changed.
% v2.5 synchronizes the Paper-III-D v1.26 faithful-trace-duality proof,
% exact ordered-owner joint locus, and typed alternative-presentation
% bookkeeping repair, together with its Paper-III-E and closure-frontier
% references.  No shared mathematical notation is changed.
% v2.4 synchronizes the Paper-III-C v1.25 orbit/process/owner repair, its
% downstream Papers III-B--F, V, VI, VIII, and IX revisions, and Born-v6_13.
% v2.3 synchronizes the Paper-II-A v3.41 witness/support/principal inheritance
% layer, the Paper-III-A v1.24 joint-realizability/readout firewall, all dependent
% host-rate revisions, and the current Paper-I-A--VII-B source registry.
% v2.2 synchronizes the generated active-realization, determinant-depth,
% dimension-weighted hierarchy, and typed cosmological-event interfaces.
% v2.1 synchronizes the exceptional-support interface revisions: Paper I-A
% v5.9, Paper II-A v3.30, Paper II-E/F v3.7, and Paper IV v4.16.
% v2.0 introduces the stage-safe G0 occurrence notation, synchronizes the
% active Paper-I-A--IV source registry (current rigorous Paper-III closure:
% III-A v1.17, III-B v1.22, III-C--E v1.17, III-F v1.14, IV v4.14), and makes
% external-reference self-skip depend on logical paper role rather than a
% versioned job name.
% ======================================================================

% ======================================================================
% macros.tex — Shared macros for the Born Geometry series
% ======================================================================

% ======================================================================
% 1. PACKAGES, BIBLIOGRAPHY, COLORS, AND UNIT DEFINITIONS
% ======================================================================
% Core packages shared across the complete Paper-I-A--VII-C series.  Loading them here makes
% the macro file a genuine series preamble while remaining compatible with existing
% paper preambles that load the same packages before \input.
\usepackage{amsmath}
\usepackage{amssymb}
\usepackage{amsthm}
\usepackage{mathrsfs}
\usepackage{array}
\usepackage{booktabs}
\usepackage{xcolor}
\makeatletter
\@ifpackageloaded{geometry}{}{\usepackage{geometry}\geometry{margin=1in}}
\makeatother
\usepackage{mdframed}

% Shared support packages required by the macros below and by most papers.
\usepackage{ragged2e}
\usepackage{longtable}
\usepackage{tabularx}
\usepackage{enumitem}
\usepackage[most]{tcolorbox}

% Series-wide numeric bibliography convention using the portable BibTeX path.
% All Papers I-A--VII-C use the ordinary \cite interface.
\usepackage[numbers,sort&compress]{natbib}
\providecommand{\printbibliography}{%
  \bibliographystyle{plainnat}%
  \bibliography{Born-v6_59}%
}

% Shared colors.  \providecolor avoids clobbering a paper-specific override.
\providecolor{postulatecolor}{RGB}{230,242,255}
\providecolor{compare}{RGB}{230,245,230}
\providecolor{definitioncolor}{RGB}{230,242,255}

\providecolor{openriskcolor}{RGB}{255,235,235}
\providecolor{bridgecolor}{RGB}{240,240,255}
\providecolor{gold}{RGB}{255,215,0}

% ======================================================================
% 2. CUSTOM ENVIRONMENTS
% ======================================================================
% Common table column types used in all papers.  Guarded so that a paper may
% define its own versions before loading this file.
\makeatletter
\@ifundefined{NC@find@L}{%
  \newcolumntype{L}[1]{>{\RaggedRight\arraybackslash}p{#1}}%
}{}
\@ifundefined{NC@find@C}{%
  \newcolumntype{C}[1]{>{\Centering\arraybackslash}p{#1}}%
}{}
\@ifundefined{NC@find@Y}{%
  \newcolumntype{Y}{>{\RaggedRight\arraybackslash}X}%
}{}
\makeatother

% ======================================================================
% 3. THEOREM STYLES AND ENVIRONMENTS
% ======================================================================
% Custom underlined theorem style (useful for remarks or notes)
\newtheoremstyle{underlined}
  {3pt}   % Space above
  {3pt}   % Space below
  {}      % Body font (empty means normal/roman text, typical for remarks)
  {}      % Indent amount
  {}      % Theorem head font (can be \itshape, \bfseries, etc.)
  {.}     % Punctuation after theorem head
  { }     % Space after theorem head
  {\underline{\thmname{#1}\thmnumber{ #2}}\thmnote{ (#3)}} % Theorem head formatting

% Standard plain theorem style environments.  Each environment is guarded so
% the shared macro file can be loaded after a paper-specific theorem setup.
\theoremstyle{plain}
\makeatletter
\@ifundefined{theorem}{\newtheorem{theorem}{Theorem}[section]}{}
\@ifundefined{lemma}{\newtheorem{lemma}{Lemma}[section]}{}
\@ifundefined{proposition}{\newtheorem{proposition}{Proposition}[section]}{}
\@ifundefined{corollary}{\newtheorem{corollary}{Corollary}[section]}{}
\makeatother

% Definition style environments (typically non-italicized body)
\theoremstyle{definition}
\makeatletter
\@ifundefined{definition}{\newtheorem{definition}{Definition}[section]}{}
\@ifundefined{postulate}{\newtheorem{postulate}{Postulate}[section]}{}
\@ifundefined{closure}{\newtheorem{closure}{Closure}[section]}{}
\makeatother

% Applying the custom underlined style
\theoremstyle{underlined}
\makeatletter
\@ifundefined{remark}{\newtheorem{remark}{Remark}[section]}{}
\@ifundefined{convention}{\newtheorem{convention}{Convention}[section]}{}
\@ifundefined{principle}{\newtheorem{principle}{Principle}[section]}{}
\makeatother

% Framed proof environment.
% Every \begin{proof} ... \end{proof} is automatically framed/shaded.
% This avoids mdframed/trivlist nesting, which can drop the top border after
% theorem-like environments.
\makeatletter
\renewenvironment{proof}[1][\proofname]{%
  \par\addvspace{6pt}%
  \noindent
  \begin{tcolorbox}[%
    enhanced,
    breakable,
    colback=black!4,
    colframe=black!35,
    boxrule=0.35pt,
    arc=1pt,
    left=6pt,
    right=6pt,
    top=4pt,
    bottom=4pt,
    boxsep=0pt,
    before skip=0pt,
    after skip=6pt,
    parbox=false
  ]%
  \pushQED{\qed}%
  \normalfont\itshape #1\@addpunct{.}\enspace
  \normalfont\ignorespaces
}{%
  \unskip\nobreak\hfill\popQED
  \end{tcolorbox}%
  \par\addvspace{2pt}%
}
\makeatother

\mdfdefinestyle{postulatebox}{
  backgroundcolor=yellow!10,
  linecolor=black,
  linewidth=0.4pt,
  skipabove=\baselineskip,
  skipbelow=\baselineskip,
  nobreak=true
}

% ======================================================================
% PROOF-READABILITY HELPERS
% ======================================================================

\newcommand{\ProofPlan}[1]{%
  \par\smallskip
  \noindent\textbf{Proof plan. }#1%
  \par\smallskip
}

\newlist{proofsteps}{enumerate}{2}
\setlist[proofsteps,1]{
  label=\textbf{(\alph*)},
  ref=\alph*,
  leftmargin=2.2em,
  itemsep=0.45\baselineskip,
  topsep=0.35\baselineskip
}
\setlist[proofsteps,2]{
  label=\textbf{(\roman*)},
  ref=\roman*,
  leftmargin=2.0em,
  itemsep=0.25\baselineskip,
  topsep=0.2\baselineskip
}


% ======================================================================
% 4. NUMBER SYSTEMS AND ALPHABETS
% ======================================================================
\providecommand{\RR}{\mathbb{R}} % Real numbers
\providecommand{\PP}{\mathbb{P}} % Projective space / Probability
\providecommand{\CC}{\mathbb{C}} % Complex numbers
\providecommand{\ZZ}{\mathbb{Z}} % Integers
\providecommand{\NN}{\mathbb{N}} % Natural numbers
\providecommand{\HH}{\mathbb{H}} % Quaternions / Hyperbolic space
\providecommand{\LFRayGeom}{\mathsf{LF\!\mbox{-}RAY\!\mbox{-}GEOM}}
\providecommand{\LFRep}{\mathsf{LF\!\mbox{-}REP}}
\providecommand{\LFState}{\mathsf{LF\!\mbox{-}STATE}}
\providecommand{\LFLim}{\mathsf{LF\!\mbox{-}LIM}}
\providecommand{\LFScat}{\mathsf{LF\!\mbox{-}SCAT}}
% Number systems / projective spaces
\providecommand{\CP}{\mathbb{CP}}
\providecommand{\RP}{\mathbb{RP}}

% ======================================================================
% 5. GROUPS AND LIE THEORY
% ======================================================================
\providecommand{\SO}{\ensuremath{\mathrm{SO}}}     % Special Orthogonal group
\providecommand{\SU}{\ensuremath{\mathrm{SU}}}     % Special Unitary group
\providecommand{\Spin}{\ensuremath{\mathrm{Spin}}} % Spin group
\providecommand{\CS}{\ensuremath{\mathrm{CS}}}     % Chern-Simons (capitalized)
\providecommand{\cs}{\ensuremath{\mathrm{cs}}}     % Chern-Simons (lowercase)

% ======================================================================
% 6. COMMON MANIFOLDS AND TOPOLOGY
% ======================================================================
\providecommand{\Sone}{S^1}                 % Standard circle

% ======================================================================
% 7. MATH OPERATORS (Declared via \DeclareMathOperator)
% ======================================================================
\DeclareMathOperator{\End}{End}   % Endomorphisms
\DeclareMathOperator{\Aut}{Aut}   % Automorphisms
\DeclareMathOperator{\spec}{spec} % Lowercase spectrum
\DeclareMathOperator{\tr}{tr}     % Lowercase trace (if you already have \Tr)
\DeclareMathOperator{\supp}{supp} % Support
\DeclareMathOperator{\TrA}{Tr_{A}} % Partial trace over subsystem A
\DeclareMathOperator{\TrB}{Tr_{B}} % Partial trace over subsystem B
% Lie/algebra operators
\DeclareMathOperator{\Ad}{Ad}
\DeclareMathOperator{\Sym}{Sym}
\DeclareMathOperator{\su}{su}
% ======================================================================
% 8. STANDARD OPERATOR-LIKE SYMBOLS (Safe definitions via \providecommand)
% ======================================================================
\providecommand{\vol}{\operatorname{vol}}   % Volume form
\providecommand{\Vol}{\operatorname{Vol}}   % Total volume
\providecommand{\Ric}{\operatorname{Ric}}   % Ricci tensor
\providecommand{\Rm}{\operatorname{Rm}}     % Riemann tensor
\providecommand{\Tr}{\operatorname{Tr}}     % Uppercase trace
\providecommand{\Spec}{\operatorname{Spec}} % Uppercase spectrum
\providecommand{\diag}{\operatorname{diag}} % Diagonal matrix
\providecommand{\rank}{\operatorname{rank}} % Rank
\providecommand{\scal}{\mathcal R} % Scalar curvature
\providecommand{\Cl}{\operatorname{Cl}}     % Clifford algebra
\providecommand{\diam}{\operatorname{diam}} % Diameter

% ======================================================================
% 9. FUNCTIONS, TRIGONOMETRY, AND CALCULUS
% ======================================================================

% Codifferential naming conventions
% (Avoid \ddag: commonly defined by physics.sty and symbol packages)
\providecommand{\codiff}{\mathrm{d}^{\dagger}} % Codifferential operator

% Quantum mechanics bra-ket notation (needed for basis tables)
\providecommand{\ket}[1]{\lvert #1 \rangle}
\providecommand{\bra}[1]{\langle #1 \rvert}
\providecommand{\braket}[2]{\langle #1 \mid #2 \rangle}

% ======================================================================
% 10. PHYSICAL SCALES AND CONSTANTS (Paper 1c / Born Geometry)
% ======================================================================
\providecommand{\Rp}{\ensuremath{R_{\mathrm P}}}     % matched common EH curvature scale (not defined from observed G)
\providecommand{\hBorn}{\ensuremath{h_{\mathrm B}}} % physical representative 2\pi\hbar of one primitive Born scalar action cell
\providecommand{\UEHjoint}{\ensuremath{U_{J}}} % joint-Dirac EH curvature-coordinate scale 8 kappa_J
\providecommand{\Gread}{\ensuremath{G_{\rm scal}}} % scalar-action coordinate c_op^3/(h_B U); physical matching explicit
\providecommand{\GreadJoint}{\ensuremath{G_{J,\mathrm{raw}}}} % raw ordinary-trace coordinate = G_scal/96
\providecommand{\Rds}{R_{\mathrm{dS}}}    % de Sitter Curvature
\providecommand{\Mpl}{M_{\mathrm{Pl}}}    % Planck mass
\providecommand{\lp}{\ell_{\mathrm{Pl}}}  % Planck length
\providecommand{\Pmax}{P_{\mathrm{max}}}  % Maximum momentum
\providecommand{\Mplbar}{\bar{M}_{\mathrm{Pl}}} % Reduced Planck mass
\providecommand{\eff}{\mathrm{eff}}       % Subscript for effective quantities
\providecommand{\Cosmo}{\ensuremath{\Lambda}}                % Cosmological constant (use \Cosmo in source)
\providecommand{\LCDM}{\ensuremath{\Cosmo\mathrm{CDM}}}  % \LambdaCDM (typesets with \Cosmo)
\providecommand{\Lcut}{\ensuremath{\Lambda_{\mathrm{cut}}}}  % Cutoff scale (brace when used as script)
\providecommand{\Rcs}{\ensuremath{R_{\mathrm{cs}}}}      % Configuration-space radius
\providecommand{\CosmoEff}{\ensuremath{\Cosmo_{\eff}}}   % Effective cosmological constant

% Disambiguation: \Rmom is the momentum-space (S^3) radius. 
% Reserve plain R for the 4D curvature scalar in f(R).
\providecommand{\Rmom}{\ensuremath{R_{\mathrm{mom}}}} 
\providecommand{\Punit}{P_{0}} % Momentum unit set by \Rmom (= \hbar/\Rmom)
\providecommand{\Rvec}{\ensuremath{R_{\mathrm{vec}}}} % invariant vector-space radius in the reduced round sector
% Specific numeric values and expressions
\providecommand{\Vlogcap}{L_{\mathrm{Hopf}}^{(1)}}
\providecommand{\Hopf}{\operatorname{Hopf}}
% Scalar phase (S^1) sector + universal cover notation (Paper I-A)
\providecommand{\tper}{T_\star}      % Time period/scale
\providecommand{\Eper}{E_\star}      % Energy period/scale
\providecommand{\ttilde}{\tilde t}   % Universal cover time
\providecommand{\Etilde}{\tilde E}   % Universal cover energy

% ======================================================================
% 10A. SERIES-WIDE LOGICAL / NOTATIONAL DISAMBIGUATIONS
% ======================================================================
% Use these symbols to avoid collisions among: (i) the Paper--III
% no-insertion/geometric-closure class, (ii) Paper--I-A trimer diagonals,
% (iii) Paper--V-A/V-C family-space diagonals, (iv) scalar-index running,
% and (v) protected shell-clock labels.
\providecommand{\Rprim}{\ensuremath{\mathsf R_{\mathrm{prim}}^{\Sigma}}} % primitive finite structural readouts
\providecommand{\Rloop}{\ensuremath{\mathsf R_{\circ}^{\Sigma}}} % closed-loop readout role
\providecommand{\Rbdry}{\ensuremath{\mathsf R_{\partial}^{\Sigma}}} % boundary-conditioned readout role
\providecommand{\Rgen}{\ensuremath{\mathsf R_{\mathrm{fin}}^{\Sigma}}} % finite generated closure of primitive roles
\providecommand{\PhysOcc}{\operatorname{PhysOcc}}
\providecommand{\Actual}{\operatorname{Actual}}
\providecommand{\GeomLic}{\operatorname{GeomLic}}
\providecommand{\InvDist}{\operatorname{InvDist}}
\providecommand{\Dcl}{\ensuremath{\mathsf{D}_1}}      % unmarked positive Dirac-even closed spectral-readout form class
\providecommand{\Dheat}{\ensuremath{\mathsf D_{\mathrm{heat}}}} % one-parameter heat-semigroup subclass
\providecommand{\DtriI}{\ensuremath{D^{\mathrm{tri}}_1}} % Paper--I-A cyclic-trimer diagonal
\providecommand{\DtriII}{\ensuremath{D^{\mathrm{tri}}_2}}
\providecommand{\DfamI}{\ensuremath{D^{\mathrm{fam}}_1}} % Family-space diagonal selector
\providecommand{\DfamII}{\ensuremath{D^{\mathrm{fam}}_2}}
\providecommand{\alpharun}{\ensuremath{\alpha_{\mathrm{run}}}} % scalar spectral-index running
\providecommand{\Ksh}{\ensuremath{K_{\mathrm{sh}}}} % protected shell-clock coordinate
\providecommand{\ddef}{\ensuremath{\delta_{\mathrm{def}}}} % universal defect scalar after Paper--III realization

% ======================================================================
% 10B. SERIES-WIDE LOGICAL STATUS TAGS
% ======================================================================
% Upright small-caps wrapper for audit/status labels.  It prevents LaTeX from
% trying to use unavailable small-caps italic font shapes inside theorem bodies
% such as TU/lmr/m/scit.
\providecommand{\AuditLabel}[1]{}
\renewcommand{\AuditLabel}[1]{{\normalfont\scshape #1}}

% These tags are bookkeeping labels for consistency audits. They do not
% create new mathematical hypotheses; they record how a datum is being used
% in the proof chain.
\providecommand{\StatusPrimitive}{}
\renewcommand{\StatusPrimitive}{\AuditLabel{Primitive}}
\providecommand{\StatusDerived}{}
\renewcommand{\StatusDerived}{\AuditLabel{Derived}}
\providecommand{\StatusImported}{}
\renewcommand{\StatusImported}{\AuditLabel{Imported}}
\providecommand{\StatusConditional}{}
\renewcommand{\StatusConditional}{\AuditLabel{Conditional}}
\providecommand{\StatusBridge}{}
\renewcommand{\StatusBridge}{\AuditLabel{Bridge-derived}}
\providecommand{\StatusComparison}{}
\renewcommand{\StatusComparison}{\AuditLabel{Comparison-only}}

% ======================================================================
% 10C. SERIES-WIDE CLAIM-WORDING AUDIT LABELS
% ======================================================================
% These are textual audit labels used to keep modal language consistent
% across the series.  They introduce no new mathematical assumptions.
\providecommand{\ClaimForced}{}
\renewcommand{\ClaimForced}{\AuditLabel{forced}}
\providecommand{\ClaimUnique}{}
\renewcommand{\ClaimUnique}{\AuditLabel{unique}}
\providecommand{\ClaimCanonical}{}
\renewcommand{\ClaimCanonical}{\AuditLabel{canonical}}
\providecommand{\ClaimMinimal}{}
\renewcommand{\ClaimMinimal}{\AuditLabel{minimal}}
\providecommand{\ClaimClosed}{}
\renewcommand{\ClaimClosed}{\AuditLabel{closed}}
\providecommand{\ClaimPrediction}{}
\renewcommand{\ClaimPrediction}{\AuditLabel{prediction}}

% ======================================================================
% 10D. SERIES-WIDE THEOREM-CHAIN AUDIT LABELS
% ======================================================================
% These labels mark the direction of load-bearing theorem chains. They are
% audit annotations only; they introduce no new mathematical assumptions.
\providecommand{\ChainInput}{}
\renewcommand{\ChainInput}{\AuditLabel{Chain input}}
\providecommand{\ChainOutput}{}
\renewcommand{\ChainOutput}{\AuditLabel{Chain output}}
\providecommand{\ChainGuard}{}
\renewcommand{\ChainGuard}{\AuditLabel{Chain guard}}
\providecommand{\ChainNoBackflow}{}
\renewcommand{\ChainNoBackflow}{\AuditLabel{No backflow}}

% ======================================================================
% 10E. SERIES-WIDE LOCAL THEOREM-AUDIT LABELS
% ======================================================================
% These labels mark a statement-level dependency audit around individual
% theorem/definition nodes. They are bookkeeping labels only.
\providecommand{\LocalAudit}{}
\renewcommand{\LocalAudit}{\AuditLabel{Local audit}}
\providecommand{\UsesOnly}{}
\renewcommand{\UsesOnly}{\AuditLabel{Uses only}}
\providecommand{\Exports}{}
\renewcommand{\Exports}{\AuditLabel{Exports}}
\providecommand{\DoesNotUse}{}
\renewcommand{\DoesNotUse}{\AuditLabel{Does not use}}

% ======================================================================
% 10F. SERIES-WIDE NUMERICAL-ROW AUDIT LABELS
% ======================================================================
% These labels classify individual rows in numerical, benchmark, status,
% and comparison tables.  They are audit annotations only and do not change
% any theorem, value, formula, or comparison number.
\providecommand{\RowAudit}{}
\renewcommand{\RowAudit}{\AuditLabel{Numerical-row audit}}
\providecommand{\RowTheorem}{}
\renewcommand{\RowTheorem}{\AuditLabel{Theorem-fixed row}}
\providecommand{\RowBridge}{}
\renewcommand{\RowBridge}{\AuditLabel{Bridge-fixed row}}
\providecommand{\RowRepresentative}{}
\renewcommand{\RowRepresentative}{\AuditLabel{Representative row}}
\providecommand{\RowComparison}{}
\renewcommand{\RowComparison}{\AuditLabel{Comparison row}}
\providecommand{\RowCalibrated}{}
\renewcommand{\RowCalibrated}{\AuditLabel{Calibrated row}}
\providecommand{\RowEstimate}{}
\renewcommand{\RowEstimate}{\AuditLabel{Estimate row}}
\providecommand{\RowOpen}{}
\renewcommand{\RowOpen}{\AuditLabel{Open row}}
\providecommand{\RowNoSelector}{}
\renewcommand{\RowNoSelector}{\AuditLabel{No selector}}

% ======================================================================
% 10G. TABLE HELPERS
% ======================================================================
% Table helpers used by Paper V-C.  The unused readout-label preset has been
% removed; only active table helpers remain.
\providecommand{\gpsection}[1]{%
  \addlinespace[0.55em]
  \multicolumn{3}{@{}l}{\textsc{#1}}\\[-0.25em]
  \midrule
}
\providecommand{\mathcell}[1]{%
  \(\begin{gathered}#1\end{gathered}\)%
}

\newcommand{\DerBy}[1]{%
  \begin{tabular}[t]{@{}l@{}}#1\end{tabular}%
}

\newcommand{\TableNotes}[1]{%
  \par\medskip
  \begin{minipage}{0.94\linewidth}
  \footnotesize\emph{Notes.} #1
  \end{minipage}\par
}
% ======================================================================
% 10H. PAPER VI-A / QCD NOTATION
% ======================================================================
% Shared symbols used by Paper VI-A and downstream QCD/Yang--Mills handoffs.
\providecommand{\Pqcd}{\mathfrak P_{\rm VIII}^{\rm QCD}}
\providecommand{\Iqcd}{\mathfrak I_{\rm QCD}^{(7)}}
\providecommand{\muQCD}{\mu_{\rm QCD}^{\rm int}}
\providecommand{\LambdaQCD}{\Lambda_{\rm QCD}}
\providecommand{\LambdaBH}{\Lambda_{\rm BH}^{\rm QCD}}
\providecommand{\HQCD}{\mathcal H_{\rm QCD}^{\rm BH}}
\providecommand{\HQCDfin}{H_{\rm QCD}^{\rm fin}}
\providecommand{\Hraw}{\mathcal H_{\rm raw}}
\providecommand{\HphysQCD}{\mathcal H_{\rm phys}^{\rm QCD}}
\providecommand{\Vcol}{V_C}
\providecommand{\Ps}{P_{\rm sing}}
\providecommand{\PiOct}{\Pi_8}
\providecommand{\MSbar}{\overline{\rm MS}}
\providecommand{\YM}{\mathrm{YM}}
\providecommand{\Ktri}{K_{\triangle}}
\providecommand{\KYM}{K_{\YM}^{\rm BH}}
\providecommand{\aBH}{a_{\rm BH}}
\providecommand{\epsBH}{\varepsilon_{\rm BH}}
\newcommand{\RowConstructive}{\textnormal{\textsc{Constructive}}}
\newcommand{\RowObligation}{\textnormal{\textsc{Bridge obligation}}}
\newcommand{\RowTarget}{\textnormal{\textsc{Target row}}}
\newcommand{\RowValidation}{\textnormal{\textsc{Validation row}}}
\providecommand{\slashed}[1]{\not{#1}}

% ======================================================================
% 10I. PAPER VI-B / CONTINUOUS-QFT NOTATION
% ======================================================================
% Shared symbols used by Paper VI-B and any downstream continuous-QFT or
% Yang--Mills handoff paper.  All definitions are guarded so individual
% papers may override them locally before loading this file if needed.
\providecommand{\A}{\mathfrak A}
\providecommand{\Hphys}{\mathcal H_{\rm phys}}
\providecommand{\QCD}{\mathrm{QCD}}
\providecommand{\BH}{\mathrm{BH}}
\providecommand{\loc}{\mathrm{loc}}

\providecommand{\RBH}{\mathcal R_{\rm BH}}
\providecommand{\DG}{\mathfrak D_G}
\providecommand{\PW}{\mathcal{P\!W}} % Peter--Weyl spectral-symbol subspace
\providecommand{\Toe}{\operatorname{Toe}} % Toeplitz compression/representative
\providecommand{\Hol}{\operatorname{Hol}} % Holonomy

% ======================================================================
% 11. LATE-LOADED PACKAGES (Must be at the end of the preamble)
% ======================================================================
% Load late packages only if a paper has not already loaded them, avoiding
% option clashes in paper-specific preambles.
\makeatletter
\@ifpackageloaded{hyperref}{}{\usepackage{hyperref}}
\makeatother
% --- Series-wide disambiguations (Pass 3) ---
\providecommand{\diracdelta}{\delta_{\rm D}} % Dirac delta distribution (delta reserved for hierarchy)
\providecommand{\Hg}{H_{\mathrm g}} % Hodge harmonic projector (keep H for Hubble)
\providecommand{\Gr}{G_{\mathrm r}} % Green operator (keep G for Newton)
\providecommand{\FD}{\mathrm{D}} % functional/variational derivative symbol (delta reserved)
\providecommand{\Id}{\mathbf{1}} % identity/Kronecker tensor (avoid \delta_{ij})


% BEGIN INLINE SOURCE: Paper00-ActiveExternalReferences-v2_130.tex
% Framework v2.191: current reference binding; existing definitions retained.
% Framework v2.190: current22-root registry with native clock-evolution proofs.
% Framework v2.189: exact current22-root external registry, including retained v2.188 hardening owners.
% Framework v2.187: current registry synchronization only; source release v2.186 preserved.
% Framework v2.186: authoritative current 22-paper registry.
% Framework v2.182: all current 22 logical paper AUX bindings.
% Framework v2.170: current I-B v1.43 external AUX binding.
% Framework v2.167: exact current 22-root registry after CP--BL2 integration.
% Framework v2.166: exact current 22-root compact-release AUX registry.
% Framework v2.165: authoritative current root bindings.
% Framework v2.164: exact current 22 manuscript AUX bindings.
% Framework v2.162: finite metric, functional Hilbert and scoped Ward integration.
% Framework v2.161: current 22 manuscript AUX bindings; source only.
% Framework v2.156 / registry v2.113: synchronized current 22 manuscript roots; source-only registry.
% Framework v2.143 / registry v2.111: active manuscript AUX basenames; no generated AUX or PDF included.
% ======================================================================
% Central external-reference imports, numbers only, no PDF jumping
% ======================================================================

% xr-hyper is loaded centrally so that papers may retain their existing
% \usepackage{xr-hyper} lines without redefining \externaldocument after this
% registry is installed.  If plain xr was loaded earlier, retain it.
\makeatletter
\@ifpackageloaded{xr-hyper}{}{%
  \@ifpackageloaded{xr}{}{\usepackage{xr-hyper}}%
}
\makeatother

% Make external-document registration idempotent.  The registry identifies a
% logical paper by its Paper<number><optional letter> prefix.  Prefix-qualified
% imports remain independent.  Current paper roots contain no local
% \externaldocument calls; this guard also makes accidental duplicates harmless.
\ExplSyntaxOn
\prop_new:N \g_bh_external_refs_prop
\cs_new_eq:NN \BHOriginalExternalDocument \externaldocument

\cs_new_protected:Npn \bh_external_document:nnnn #1#2#3#4
 {
  \tl_set:Nn \l_tmpa_tl {#3}
  \regex_replace_once:nnN
   { \A (?: xref/ )? (Paper[0-9]+[A-G]?) .* \Z }
   { \1 }
   \l_tmpa_tl
  \tl_set:Nx \l_tmpb_tl { #1 | \tl_use:N \l_tmpa_tl }
  \prop_if_in:NVTF \g_bh_external_refs_prop \l_tmpb_tl
   {
    \PackageInfo{BornHopf}
     {Skipping~duplicate~external~refs~for~#3}
   }
   {
    \prop_gput:NVn \g_bh_external_refs_prop \l_tmpb_tl {#3}
    \IfNoValueTF{#4}
     { \BHOriginalExternalDocument[#1][#2]{#3} }
     { \BHOriginalExternalDocument[#1][#2]{#3}[#4] }
   }
 }

\RenewDocumentCommand{\externaldocument}{O{} O{} m o}
 { \bh_external_document:nnnn {#1}{#2}{#3}{#4} }

\NewDocumentCommand{\BHExternalRefs}{m}
 {
  % Do not import the aux file of the logical paper currently being compiled,
  % even when the job and registry use different version suffixes.
  \tl_set:Nx \l_tmpa_tl { \c_sys_jobname_str }
  \regex_replace_once:nnN
   { \A (Paper[0-9]+[A-G]?) .* \Z }
   { \1 }
   \l_tmpa_tl
  \tl_set:Nn \l_tmpb_tl { #1 }
  \regex_replace_once:nnN
   { \A (Paper[0-9]+[A-G]?) .* \Z }
   { \1 }
   \l_tmpb_tl
  \str_if_eq:eeTF { \tl_use:N \l_tmpa_tl } { \tl_use:N \l_tmpb_tl }
   {
    \PackageInfo{BornHopf}
     {Skipping~external~refs~for~#1~because~this~is~the~current~document}
   }
   {
    \IfFileExists{xref/#1.aux}
     { \externaldocument[][nocite]{xref/#1}[#1.pdf] }
     {
      \IfFileExists{#1.aux}
       { \externaldocument[][nocite]{#1} }
       {
        \PackageInfo{BornHopf}
         {External~aux~file~#1.aux~not~found;~no~labels~imported}
       }
     }
   }
 }
\ExplSyntaxOff

% External references are automatic.  This file is the sole owner of both the
% bibliography and the current Paper-I-A--VII-D AUX registry (30 roots including I-C).  Missing auxiliaries
% are informational on the first pass and are populated by the source-only
% two-pass series build.  Basenames deliberately omit temporary download/copy
% suffixes such as ``(1)'' or ``-1''.
  % Paper I-A: round S3 compendium.
  \BHExternalRefs{Paper1A-Geometry-v5_180}

  % Paper I-B: complete exceptional geometry and finite spectral structures.
  \BHExternalRefs{Paper1B-ExceptionalGeometry-v1_22}
  % Paper I-D: technical companion of I-A through I-C.
  \BHExternalRefs{Paper1D-TechnicalMaterial-v1_69}
  % Paper I-C: reciprocal M8 compendium.
  \BHExternalRefs{Paper1C-M8Geometry-v1_144}

\ifdefined\BHMathOnly
  % Papers I-A and I-B import only the mathematical atlas labels.
\else
  % Paper II: five physical-mathematical main papers plus one technical proof volume.
  \BHExternalRefs{Paper2A-GeometricConstitution-v3_227}
  \BHExternalRefs{Paper2B-MomentumCarrier-v2_147}
  \BHExternalRefs{Paper2C-ReciprocalLocalization-v2_174}
  \BHExternalRefs{Paper2D-SymmetrySpectra-v2_145}
  % The former II-F export labels remain in II-E; current II-F is the technical volume.
  \BHExternalRefs{Paper2E-KinematicExports-v3_140}
  % Paper II-F: complete technical proof and calculation owner for II-A--II-E.
  \BHExternalRefs{Paper2F-TechnicalMaterial-v1_71}

  % Paper III: Dynamics and Intrinsic Geometric Constitution (six components).
  \BHExternalRefs{Paper3A-GeneratedDynamics-v1_180}
  \BHExternalRefs{Paper3B-SpectralAction-v1_191}
  \BHExternalRefs{Paper3C-ProtectedHierarchy-v1_193}
  \BHExternalRefs{Paper3D-RunningResponse-v1_167}
  \BHExternalRefs{Paper3E-GravityAction-v1_203}
  \BHExternalRefs{Paper3F-CompleteDynamics-v1_276}
  % Paper III-G: complete preserved technical proof owner.
  \BHExternalRefs{Paper3G-TechnicalMaterial-v1_57}

  % Paper IV-A: generated finite quantum mechanics.
  \BHExternalRefs{Paper4A-FiniteQuantumMechanics-v4_216}
  % Paper IV-B: complete finite quantum processes and physical response.
  \BHExternalRefs{Paper4B-QuantumResponse-v1_54}
  % Paper IV-C: extended proofs and finite technical calculations.
  \BHExternalRefs{Paper4C-TechnicalMaterial-v1_50}

  % Paper V-A: The Standard Model.
  \BHExternalRefs{Paper5A-FiniteMatterGeometry-v3_218}

  % Paper V-B: Finite Spectrum and Flavor Closure.
  \BHExternalRefs{Paper5B-TerminalSpectrum-v5_202}

  % Paper V-C: Finite Interaction Boundaries of the Standard Model.
  \BHExternalRefs{Paper5C-InteractionResponse-v2_182}
  % Paper V-D: preserved technical proof and calculation owner.
  \BHExternalRefs{Paper5D-TechnicalMaterial-v1_58}

  % Paper VI-A: Gravity and Closed Cosmology.
  \BHExternalRefs{Paper6A-GravityCosmology-v2_267}

  % Paper VI-B: Geometric Dark Sector and Cosmological Readouts.
  \BHExternalRefs{Paper6B-CosmologyReadouts-v1_214}

  % Paper VI-C: complete technical proof and comparison owner.
  \BHExternalRefs{Paper6C-TechnicalMaterial-v1_53}

  % Paper VII-A: Finite QCD Kernel and Nonperturbative Closure.
\BHExternalRefs{Paper7A-QCDKernel-v1_180}

  % Paper VII-B: Hadronic, Thermal, and Nuclear QCD Readouts.
  \BHExternalRefs{Paper7B-QCDReadouts-v1_165}

  % Paper VII-C: Continuous Quantum Field Theory.
  \BHExternalRefs{Paper7C-ContinuousQFT-v2_208}

  % Paper VII-D: complete technical material for VII-A through VII-C.
  \BHExternalRefs{Paper7D-TechnicalMaterial-v1_60}

\fi

% END INLINE SOURCE: Paper00-ActiveExternalReferences-v2_130.tex


\hypersetup{hidelinks}





\usetikzlibrary{positioning, calc}

\title{\textbf{Born--Hopf Geometry II-F: Technical Proofs and Calculations}\\[0.5em]
\large Compact Parent, Finite Momentum, Reciprocal Operators, Spectra, and Exports}
\author{Peter Kugelmann}
\date{Version 1.71 --- 20 November 2026}
\setlength{\emergencystretch}{3em}


\begin{document}
\maketitle
\begin{abstract}
Paper~II-F is the technical proof and calculation companion to the
constitution, momentum-carrier, reciprocal-localization, spectral-kinematic
and geometric-export results of Papers~II-A--II-E.  It gives detailed
compact-parent and elliptic arguments, exactness countermodels, conditional
finite-discreteness and local-regularity results, optional momentum towers,
finite operator matrices, coordinate-completion comparisons, covariant
positive-operator measures, sharp moment and rank certificates, spectral and
Hopf/Berry calculations, higher-spin checks, equivariance obstructions, and
independent exceptional-incidence and rank-two certificates.

The statements retain their individual hypotheses.  Optional constructions
and non-implication countermodels do not change the generated GP2 momentum
support or discharge the separate universal GP3 whole-candidate comparison.
The existence of a finite operator, frame, spectral representative or
measurement model does not itself establish physical occurrence, a selected
action, an experimental instrument, or a particle spectrum.  Exceptional
subgroup foundations are imported from Paper~I-B; general round-$S^3$ and
Hopf geometry are imported from Paper~I-A.  Detailed calculations herein
supply technical evidence on those specified mathematical domains.
\end{abstract}
\tableofcontents
\clearpage

\section{Constitutional exactness and non-implication tests}
\label{p2ftech:sec:constitution}

\subsection{Exact field realization and stationary-path typing}
\label{p2a:app:v2190-exactness-stationarity}

This section records a diagnostic, not an alternative physical action or
a new constitutional assumption.  A generated diagnostic residual and a
GP2 realization defect have different roles by
Definitions~\ref{p2:def:intrinsic-generated-defect}
and~\ref{p2:def:gp2-realization-defect}.  A faithful realization preserves
the former even when it is nonzero.  GP3's structural preorder is not
numerical minimization of an action.

In the later native odd-scalar field family, where a nonzero compactly
supported Hermitian section \(\Phi_1\) is allowed, linearity licenses
\(\Phi=\lambda\Phi_1\).  On the Euclidean local scalar action domain,
\[
 I_E(\lambda\Phi_1)=a\lambda^2+b\lambda^4,
 \qquad b=\frac1{16\pi^2}\int\operatorname{Tr}_F(\Phi_1^4)\,d\mu_E>0.
\]
The retained divergence integrates to zero there.  Therefore
\(2a\lambda+4b\lambda^3\) does not vanish identically.  Exact realization
of this generated field-role family cannot require every member to be
stationary.  This tests the inference from field-role exactness to
stationarity before a stationary-history target is imposed; it is not a
nonstationary member of an already constituted on-shell history role and
does not prove formal independence of all dynamics from the axioms.

Likewise, same-owner symmetry invariance gives
\(dS(R_\xi)=0\) along actual symmetry directions by
Corollary~\ref{p2a:cor:v368-native-noether-antecedent}.  It does not supply
\(dS(v)=0\) for every physical variation.  Mere preservation of a phase
form is weaker than tangency to the characteristic kernel of its exterior
derivative.  The positive theorem uses the licensed complete stationary
relation and the exact operation-identification certificate, rather than
any of these stronger inferences.

For the unmarked clock-comparison theorem, in a tangent basis consisting
of transverse vectors and the characteristic vector \(X=(Y,1)\), inversion
is \(\operatorname{diag}(-I,1)\).  In the original clocked coordinates its
matrix is
\[
 g=\begin{pmatrix}-I&2Y\\0&1\end{pmatrix},
 \qquad g^T\beta g=\beta,\quad dt\,g=dt,\quad gX=X.
\]
Those tensor equalities do not establish preservation of all parent marks.
For example a separately retained nonzero energy-gradient covector need
not be fixed by this inversion.  Its action cannot be omitted from the
isotropy test merely because the curvature/clock tensors are preserved.
The theorem's zero fixed-space certificate must therefore use the complete
actual pre-generator owner, or an independently licensed operation
factorization through the bare comparison descendant.

\begin{proposition}[A retained transverse covector obstructs zero fixed space]
\label{p2a:prop:v2190-retained-mark-fixed-vector}
At a regular comparison point let \(V=\ker dt\), with nondegenerate
\(\Omega_V=\beta|_{V\times V}\).  If a nonzero covector
\(\ell\in V^*\) is retained by every actual isotropy element, then the
unique vector \(Z\in V\) with \(\iota_Z\Omega_V=\ell\) is nonzero and
fixed by that isotropy.  Hence \(V^G\ne0\).  In particular this applies
to \(\ell=d\mathcal H|_V\) when the complete owner retains
\(\mathcal H\) as an invariant scalar and that differential is nonzero.
\end{proposition}
\begin{proof}
Every isotropy element preserves \(V\), \(\Omega_V\), and \(\ell\).
It therefore sends \(Z\) to another solution of the same nondegenerate
equation.  Uniqueness fixes \(Z\), and \(\ell\ne0\) makes it nonzero.
Preservation of a scalar at the owner level preserves its differential
at a fixed point; restriction to \(V\) gives the stated example.
\end{proof}

In the time-independent canonical comparison chart this fixed vector is
\(Z=(Y,0)\), with \(\iota_Y\Omega=d\mathcal H\).
The affine invariant directions include \(X+aZ\) for scalar \(a\),
all with \(dt=1\).  This exhibits a pointwise ambiguity of the naturality
test, not admitted alternate physical laws or a proof of constitutional
independence.  Preserving additional parent marks can restrict the test
further; the full owner must still be used.



\subsection{Selector non-implication theorems and countermodels}
\label{p2:app:selector-limitations}
This section collects negative or scope-delimiting results whose outcomes are
stated in the main text but whose detailed proofs are not part of the positive
constitutional chain.  They show, respectively, that one finite static algebra
does not reconstruct its parent spectral grammar, that GP1-TB does not imply
the optional comparison regularity conditions, that boundedness alone does not
select a realization label, and that the legacy finite-level capacity formulas
remain distinct from physical momentum-seed selection.

\subsubsection{Generic-net obstruction to exceptional-core ancestry}
\label{p2a:app:generic-net-exceptional-obstruction}

\begin{proof}[Proof of Proposition~\ref{p2a:prop:generic-resolution-not-exceptional}]
Fix \(\varepsilon>0\) and choose a finite \(\varepsilon/3\)-net
\(Z_0=\{z_1,\ldots,z_N\}\subset S^3\).  The covering-radius functional
\[
 \rho(Z):=\sup_{x\in S^3}\min_{z\in Z}d(x,z)
\]
is one-Lipschitz in Hausdorff distance.  Hence every pointwise perturbation
\(Z'=\{z'_1,\ldots,z'_N\}\) satisfying
\[
 \max_i d(z_i,z'_i)<\varepsilon/3
\]
is still an \(\varepsilon\)-net.

After shrinking to a connected product of geodesically convex coordinate
balls if necessary, this perturbation neighborhood is a connected real-
analytic manifold.  Each exact multiplication relation
\(z'_iz'_j=z'_k\) is the zero locus of a nontrivial real-analytic map.  For each
\(24\)-element index subset and each ordering of the Hurwitz vertices, equality
of its Gram matrix with the \(2T\) Gram matrix is likewise the zero locus of
a nontrivial real-analytic map.  None of these maps vanishes identically:
one may vary a single point while holding the remaining points fixed.  Each
excluded locus is therefore proper and nowhere dense.  Only finitely many
such loci occur, so their union has empty interior.  One may consequently
choose \(Z'\) outside all of them while retaining the \(\varepsilon\)-net
property.  The resulting support is not group closed and contains no exact
regular-\(24\)-cell subset.

There is also a direct cardinality obstruction at coarse resolution.  Whenever
the declared resolution is at least the diameter, a singleton is an allowed
finite net.  Thus no subset-valued constructor defined on the complete generic
atlas can always extract \(24\) Hurwitz vertices.  A constructor that ignores
the input support and inserts an ambient \(24\)-cell is precisely an additional
exceptional-core constructor, not an output of point resolution.  These two
countermodels prove every clause of the proposition.
\end{proof}

\subsubsection{Proof of the boundedness nonselection theorem}
\begin{proof}[Proof of Theorem~\ref{p2:thm:no-cutoff-from-finiteness}]
The stated FRC admissibility and boundedness conditions are satisfied by every
member of the displayed family and therefore contain no discriminator among
them.  An order on
\(\Lambda\), including the existence of an endpoint, is not a natural section
from the previously constituted geometry.  The global GP3 comparison has
already allocated the operation role; it supplies no order on an unrelated
spectral or truncation family and cannot manufacture one.
\end{proof}

\begin{proposition}[Scalar-orientation inversion countermodel]
\label{p2:prop:scalar-orientation-inversion-countermodel}
Remove the typed provenance occurrence \(\mathsf{HS}_t\), while retaining an
oriented Paper-I-A Hopf circle, an independently allocated scalar circle
\(H_{\rm succ}\cong U(1)\), its full dual, and the scalar Poincar\'e pairing.
Then those retained data do not determine a positive scalar-succession
primitive.
\end{proposition}

\begin{proof}
On the unmarked scalar circle, inversion
\[
 I:H_{\rm succ}\longrightarrow H_{\rm succ},\qquad z\longmapsto z^{-1},
\]
is a continuous group automorphism and sends any primitive lattice generator
\(e_t\) to \(-e_t\).  Simultaneously sending its full-dual primitive
\(e_E\mapsto-e_E\) preserves
\[
 e_E(e_t)=1,
 \qquad
 \Omega(e_E,e_t)=1,
 \qquad
 dE\wedge dt.
\]
Thus the abstract scalar topology, compact group law, full duality, primitive
phase class, and symplectic orientation survive simultaneous inversion.
Without a typed arrow carrying the already fixed Paper-I-A sign into the scalar
role, the two scalar-succession interpretations satisfy the same frozen data.
Selecting one would insert an ungenerated sign, contrary to GP2.
\end{proof}


\begin{theorem}[Finite spectral algebra under injective reparametrization]
\label{p2:thm:finite-spectral-reparametrization-no-go}
Let \(H=H^*\) act on a finite-dimensional carrier, and let \(f\) be injective
on the finite spectrum \(\sigma(H)\).  Then a polynomial \(g\) exists with
\[
        g(f(H))=H.
\]
Consequently \(H\) and \(f(H)\) have the same spectral projections,
commutant, invariant subspaces, symmetry stabilizer, and generated unital
finite functional-calculus algebra.  Their numerical spectra need not agree:
the common algebra contains many possible distinguished self-adjoint
generators and does not select one of them.
\end{theorem}

\begin{proof}
Interpolate the inverse function \(f^{-1}\) on the finite set
\(f(\sigma(H))\) by a polynomial \(g\).  Functional calculus gives
\(g(f(H))=H\).  Conversely \(f(H)\) belongs to the functional calculus of
\(H\), so the two generated finite algebras coincide.  The listed invariants
follow at once from that coincidence.  Injectivity relabels eigenvalues without
merging or splitting eigenspaces, so the spectral projections agree, while the
sets of numerical eigenvalues generally do not.
\end{proof}


\begin{corollary}[Static-algebra selector boundary]
\label{p2:cor:static-algebra-selector-boundary}
Let a selector receive as its entire input only the represented unital
\(*\)-algebra generated by one finite static self-adjoint realization, with no
distinguished generator, parent lift, filtration, compression map, or symbol
data.  Such a selector cannot determine which injective spectral
parametrization is intended and therefore cannot determine a parent
differential order, locality class, principal symbol, or dispersion law from
that algebra alone.  A selector supplied with
\(\operatorname{Fil}^{\bullet}\operatorname{Op}(\mathfrak P_R)\),
\(Q_{R,m}\), or a specified parent lift lies outside this conclusion.
\end{corollary}

\begin{proof}
Theorem~\ref{p2:thm:finite-spectral-reparametrization-no-go} gives the same
represented algebra for every injective relabeling of the finite spectrum, so
any function of that algebra alone has the same value on all such relabelings.
Parent differential order, locality, symbol, and dispersion belong instead to
the filtered parent datum recorded in
Definition~\ref{p2:def:paired-realization}.  Supplying that datum changes the
selector input and removes the hypothesis of this corollary.
\end{proof}


\begin{proposition}[GP1-TB alone does not imply either parent-comparison regularity condition]
\label{p2:prop:gp1-clause-sharpness}
Total boundedness, a nonisolated identity, and the full separating
metric-profile family imply neither (FD) nor (LR) for an abstract parent
comparison:
\begin{enumerate}[label=(\roman*), leftmargin=2em]
\item The \(p\)-solenoid
\(\Sigma_p:=\varprojlim(S^1\xleftarrow{\,z\mapsto z^p\,}S^1)\), with a
standard compatible invariant metric, is compact and finite-dimensional but
not locally connected.  Its proper dense arc component shows that (LR) fails
while (FD) holds.
\item The weighted torus
\(T^\infty=\prod_{k\ge1}S^1\), with
\(d(\theta,\theta')=\sum_{k\ge1}2^{-k}\rho(\theta_k,\theta'_k)\), satisfies
(LR) with \(C_R=1\) by coordinatewise minor-arc interpolation, but contains
closed copies of \(T^n\) for every \(n\); hence (FD) fails.
\end{enumerate}
Both groups act on themselves by translations, so \(D_\infty=d\), and both
have canonical separating metric profiles.  They are parent/presentation
countermodels only: their regular translation groups are not faithful exact
reversible operation parents of a finite FRC carrier.  Thus they do not
contradict automatic compact Lie closure, but they show that (FD) and (LR)
cannot be inferred from GP1-TB or metric separation alone.
\end{proposition}

\begin{proof}
The standard inverse-limit description \cite{HewittRoss1979} gives
\(\dim\Sigma_p=1\) and a proper dense arc component.  (LR) fails directly: the
arc components are dense, so every \(\kappa D_X\)-ball around the identity
contains points of other arc components, and such a point admits no path to
\(e\) at all, rectifiable or otherwise.  Equivalently, by the contrapositive of
Theorem~\ref{p2:thm:gp1-reachability-local-connectedness}, a space that is not
locally connected cannot satisfy (LR).  For \(g=(\theta_k)_k\in T^\infty\), the path
\(\gamma(t)=(t\theta_k)_k\) along coordinatewise minor arcs has length
\(L(\gamma)=d(g,e)\), proving (LR).  Monotonicity of covering dimension on
closed subspaces \cite{HurewiczWallman1941} gives \(\dim T^\infty\ge n\) for
every \(n\).  Compactness
implies GP1-TB in both cases.  If either regular translation group acted
faithfully by exact reversible operations on a finite carrier, its closure
would be compact Lie by
Theorem~\ref{p2:thm:frc-automatic-compact-lie-closure}, contrary to the stated
non-Lie structure.
\end{proof}


\begin{corollary}[Boundedness plus origin does not select a family member]
\label{p2:cor:gp1-gp2-no-family-member}
If a whole family of paired FRC realizations has the same GP2 origin and no
member is uniquely distinguished by the prior geometry, FRC+GP2 retain the
family and select no cutoff, extremal label, or numerical scale.
\end{corollary}

\begin{proof}
FRC is an admissibility condition: it removes nonfinite or unbounded
realizations but supplies no numerical minimization among the surviving
family.  In a registered generated-boundary domain the intrinsic construction
filtration may order candidate bounds, but GP1 still contains no rule choosing
the least bound.  GP2 excludes ungenerated prior marks and candidates lacking
exact host typing.  If all displayed members nevertheless realize the same
role geometry exactly and faithfully, choosing a label, extremum, or scale
would add an ungenerated datum unless one member is placed below every other
by an already justified role-faithful GP3 structural comparison.
\end{proof}


\begin{corollary}[Exact capacity of the admissible Paper-I-A module]
\label{p2:cor:admissible-module-capacity}
Invoke the local \(\mathsf P_{\rm sli}\) interface and put
\[
 x:=P_{\rm ceil}/\Punit,\qquad
 n:=m_{\rm adm}=\lfloor x\rfloor.
\]
Then the maximal admissible integer member of the declared Paper-I-A module
family has
\[
 \dim\mathcal H_{m_{\rm adm}}
 =\binom{2n+4}{4},
 \qquad
 \#\{(j_L,j_R):j_L+j_R\le n\}
 =\binom{2n+2}{2},
\]
and
\[
 m_{\rm adm}=n
 \quad\Longleftrightarrow\quad
 n\Punit\le P_{\rm ceil}<(n+1)\Punit.
\]
If one uses only the realized label \((\mathrm{LVL}_p^1)\), whether proved by
the revised host-closure theorem or otherwise, without separately asserting an
admissibility ceiling \(m_{\rm adm}=1\), then \(m_p\le s_\rho\) implies merely
\(P_{\rm ceil}\ge\Punit\), with no upper bound.  The stronger upper bound
\(P_{\rm ceil}<2\Punit\) follows only from the different statement
\(m_{\rm adm}=1\).  Neither statement selects a member of
\(\mathfrak M_{\rm coord}\) or constitutes a point-resolution occurrence.
When such a resolution occurrence is present, support finiteness follows
instead from Theorem~\ref{p2:thm:derived-finite-point-support-role}.
\end{corollary}

\begin{proof}
Paper~I-A's declared multiplicity-one module family gives
\[
 \dim\mathcal H_n
 =\frac{(2n+1)(2n+2)(2n+3)(2n+4)}{24}
 =\binom{2n+4}{4}.
\]
Writing \(a=2j_L\) and \(b=2j_R\), the retained sectors correspond exactly
to \(a,b\in\mathbb Z_{\ge0}\) with \(a+b\le2n\), whose number is
\(\binom{2n+2}{2}\).  The interval equivalence is the defining property of
the floor function.  Finally, \(1\le s_\rho=\tfrac12\lfloor2x\rfloor\)
is equivalent to \(x\ge1\) and imposes no upper bound.  The attachment data
remain separate by Paper~I-A's coordinate-completion classification and the
finite-resolution support theorem.
\end{proof}



\subsection{GP2 exactness and scope tests}
\label{p2:app:gp2-scope-diagnostics}
The following statements delimit what GP2 requires and exactifies on its own.  They are kept
for logical transparency but are not premises of the positive parent-allocation
chain in the main text.

This subsection is scope-delimiting.  Its statements record what GP2
exactifies or requires, and what lies outside GP2; no later result in this paper consumes any of them.

\begin{remark}[Exact content of the GP2 exact/faithful filter]
\label{p2:rem:gp2-origin-filter-alone}
The primitive comparison occurrence and local-reapplication constructor belong
to the strong role predicate \(\operatorname{Geo}_R\) asserted by
\(\mathsf E_R\); they are not products of GP2.  A proposed datum
\(e\) passes GP2 only if it has prior geometric origin or is generated by the
constituted typed constructor grammar, is represented in the correct sort, and
is neither inserted relative to the frozen inventory nor omitted from the
candidate in which it is generated.  Canonical finite images are objectified only under the typed
antecedents of Postulate~\ref{p2:post:GP2}; arbitrary theorem consequences are
not thereby constituted.  GP2 alone neither constitutes \(\mathsf E_R\),
asserts GP1-TB total boundedness, nor performs the GP3 irreducible/least
comparison.
\end{remark}

\begin{remark}[GP2 chronological origin]
\label{p2:rem:gp2-chronological-origin-alone}
Unfolding the chronology clause of GP2: a proposed datum cannot supply its own
geometric origin, and a later datum cannot retroactively originate an earlier
branch choice.  Naturality is evaluated under the equivalences present
immediately before the proposed extension.  Mere mathematical existence,
intrinsic residence, or membership in a family is not a GP2 origin.  The
assertion \(\mathsf E_R\) is different: its strong
\(\operatorname{Geo}_R\) condition has already been constituted by Geometric
Primacy before GP2 is applied.  This is the constitutional content of
Geometric Primacy rather than a GP2-derived result; the chronology is
load-bearing because the
alternatives are vacuous.  Allowing \(e\) among the inputs used to construct
\(e\) makes the origin test vacuous, and allowing a later refinement to reduce
the earlier equivalence group would permit every arbitrary selection to certify
itself after the fact.  Stage-indexed evaluation prevents both circularities.
\end{remark}

\begin{theorem}[GP2 local generated closure is the identity component]
\label{p2:thm:gp2-local-closure-connectedness}
If a supplied operation parent is a Lie group \(G\) with Lie algebra
\(\mathfrak g\), the subgroup generated by its local reapplications is
\[
 \langle\exp(\mathfrak g)\rangle=G^0.
\]
Therefore every \(\mathfrak M\in\operatorname{Adm}_R\) realizing the primitive
assertion \(\mathsf E_R\) has connected generated parent.  A component
outside \(G^0\) is admissible only when the prior operation grammar contains an
independently generated discrete role; it is not generated by local transfer.
This theorem does not assert that the full automorphism group of a geometry is
connected.
\end{theorem}

\begin{proof}
Every exponential lies in \(G^0\), while an exponential neighborhood of the
identity generates the connected identity component.  No finite product of
identity-neighborhood operations reaches another component.  Candidate-internal
GP2 exactness retains the generated local closure, while chronological
no-insertion forbids an independently ungenerated discrete seed.
\end{proof}

Under (LR), Theorem~\ref{p2:thm:gp1-reachability-local-connectedness} makes
the identity component of the operation closure open before Lie regularity is
invoked; germ profiles sufficiently close to the identity therefore lie in
\(G^0\), and the containment of the primitive parent in \(G^0\) does not
depend on the Lie hypothesis of this theorem.

\begin{theorem}[Candidate-internal and chronological infinitesimal retention]
\label{p2:thm:gp2-alone-frame-natural-alternatives}
Let \(G_{\mathfrak M}\) be the Lie operation group of an admissible candidate
\(\mathfrak M\).  Candidate-internal GP2 exactness retains its complete
sequential multiplication in \(\Sigma_R[\mathfrak M]\), and differentiation
therefore retains the induced Lie bracket on
\(V_{\mathfrak M}=T_eG_{\mathfrak M}\).  These data cannot be deleted inside
\(\mathfrak M\), but before GP3 they are not imposed on a different complete
candidate unless \(\operatorname{Geo}_R\) itself requires the corresponding
diagram.  If \(\mathfrak M\) is selected, its multiplication and bracket enter
\(P_R\) and become undeletable in every later closure.  The same
cross-candidate preservation already applies when a richer geometrically
primary role explicitly requires the full A-profile family.

Under (FD)+(LR), the displacement construction also determines an
Ad-invariant norm \(F_{\mathfrak M}\) on \(V_{\mathfrak M}\).  GP2 excludes an
ungenerated basis, orientation, lattice, root datum, representation, or higher
tensor, but does not delete the multiplication, bracket, or
\(F_{\mathfrak M}\) from the candidate in which they are generated and does
not enlarge their symmetry group to all of
\(O(V_{\mathfrak M},g_F)\).
\end{theorem}

\begin{proof}
Inside every candidate, composition and equality of actual global profiles
belong to \(\operatorname{Gen}_R(\mathfrak M)\), hence to
\(\Sigma_R[\mathfrak M]\) by \((\mathrm{GP2\text{-}exact})\).  The Lie bracket
is the derivative of their commutator and is therefore fixed by that
candidate's multiplication.  The norm \(F_{\mathfrak M}\) is constructed
canonically from its first-order displacement in
Corollary~\ref{p2:cor:gp1-first-order-norm-derived}.  Removing a proof-witness
basis leaves these intrinsic data unchanged; removing the data themselves
would violate candidate-internal no-omission.  GP3 may select a different
complete candidate, but it does not truncate this one.  Once
\(\mathfrak M_R\) is selected, its entire generated presentation joins
\(P_R\), so every later comparison must preserve it.  Before primitive-role
allocation, by contrast, only \(P_{<R}\) and the strong role predicate asserted
in Definition~\ref{p2:def:geometric-role-existence} are prior content; the
branch status is an output of GP1--GP2 exactification, and no candidate's full
self-presentation is frozen.
\end{proof}

\begin{corollary}[Exact scope of GP2 alone]
\label{p2:cor:gp2-no-existence-alone}
Given the geometrically constituted primitive assertion, GP2 defines the
constructor-relative exact-and-host-faithful admissible class: it enforces no
insertion relative to \(P_{<R}\), no omission inside each candidate, sort
faithfulness, and preservation of every frozen or candidate-generated
operation relation and typed canonical image in its proper chronological
scope.  Together with GP1 it
derives the nontrivial connected exact-operation parent and reads its
commutator branch as in
Theorem~\ref{p2:thm:geometric-existence-exactification}.  It does not by itself supply
total boundedness, compactness, (FD), (LR), Lie
regularity, a numerical scale, Cartan rank one, or a global-form choice.  On
the geometric branch, the canonical mathematical metric-profile family
already detects every nontrivial commutator.  If that family is invoked as
the physical B-test separator, metric-curry closure is additionally required.
No GP3 step is required for this candidate-level kernel detection; GP3 is
nevertheless required to compare distinct faithful candidates globally.
\end{corollary}

\begin{proof}
These statements are the coupled consequences of the no-insertion,
no-omission, typed-host-faithfulness, constructor-relative-closure, and
canonical-image clauses of Postulate~\ref{p2:post:GP2}.  None of the excluded
conclusions occurs in that postulate: compactness and Lie regularity come from
FRC together with finite faithful linear realization and automatic
unitarization, (FD)+(LR) are the optional hypotheses of
Definition~\ref{p2:def:lie-regularity-conditions}, and global comparison of
distinct candidates is the GP3 step.  Candidatewise kernel detection is
Theorem~\ref{p2:thm:canonical-metric-profile-separation}; promoting its
profiles to physical tests consumes metric-curry closure exactly as the
statement records.
\end{proof}

\begin{remark}[Why chronology must precede candidate self-data]
\label{p2:rem:mode-neutrality-load-bearing}
The groups \(T^2\) and \(SU(3)\), acting on themselves by left translation
with their full distance-profile tests, satisfy compact connected Lie
regularity and faithful geometric separation.  They show that GP1 and
effectiveness alone do not imply an \(S^1/S^3\)-only conclusion.  They do not
defeat the cumulative chain: before the rank, dimension, and multiplication
table of either candidate are frozen, a nontrivial character or circle
subrole of \(T^2\) witnesses comparison with a fresh complete exact \(U(1)\)
candidate, while a compact root subrole of \(SU(3)\) witnesses comparison with
a fresh relation-exact \(SU(2)\) candidate.  Neither comparison truncates the
source model or claims that its full self-presentation survives in the target.
If \(\operatorname{Geo}_R\) requires the second toral direction or an
\(SU(3)\)-specific direction, bracket, relation, or test, the corresponding
comparison is inadmissible and the endpoint theorem abstains.  This
role-local statement does not
exclude \(T^2\) formed later as the product of two separately constituted
scalar roles.  Similarly \(SO(3)\) is locally irreducible at the common Lie-algebra level but carries the
additional quotient relation excluded by GP2 at globalization.
\end{remark}



\section{Compact-Lie classification and elliptic consequences}
\label{p2f:sec:compact-lie-elliptic}
\subsection{Compact-Lie and elliptic technical consequences}
\label{p2:app:compact-lie-consequences}

\subsubsection{Cocharacter proof for the rank-one parent}
\label{p2:app:v193-u1-su2-cocharacters}
\begin{proof}[Proof of Corollary~\ref{p2:cor:v193-u1-su2-cocharacters}]
Apply Paper~1B, Proposition~\ref{p1b:result:circle-cocharacters}, to the
already derived $A_{\rm v}\cong SU(2)$. Its nonnegative weight is $|n|$;
retaining the oriented primitive gives the signed degree $n$ used here.
The classification supplies conjugacy classes of inclusions, while the
specific faithful Hopf inclusion is the occurrence already constructed in
Lemma~\ref{p2:lem:chosen-hopf-vector-parent-inclusion}.
\end{proof}

The results in this technical section are standard consequences of the compact connected
least parent-role skeleton already derived in the main text.  They fix hypotheses and provide
a shared technical reference for later papers, but none is an additional premise of
the constitutional allocation theorem.

The following compact-Lie, homogeneous-metric, representation, and
elliptic consequences are recorded with their hypotheses and proofs.
They do not add premises to the constitutional classification.  Throughout
its first tier, \(H\) denotes either a finite-dimensional compact connected
effective candidate or a GP3-derived least parent-role skeleton, with retained Lie algebra
\(\mathfrak h\), bracket, and generated test action
\(\rho_{\mathcal B}\).  Where (FD)+(LR) are additionally imposed on the
primitive comparison, the optional operational norm \(F\) and John metric
\(g_F\) are also available.  These general compact-Lie results are
consistency tools; the master allocation theorem does not depend on them as
extra physical hypotheses.

\subparagraph{Metric-domain convention.}
Every occurrence of \(F\) or \(g_F\) below carries the effective hypotheses
(FD)+(LR) from
Theorems~\ref{p2:thm:gp1-propagation-derived} and
\ref{p2:thm:operational-john-metric}.  A statement formulated for an
arbitrary retained bi-invariant Riemannian metric \(g\) may instead be
instantiated by the canonical \(g_K\) on the selected \(SU(2)\) branch or by
a displayed finite-representation proof metric \(g_\rho\).  No such
instantiation identifies that metric with the primitive comparison
\(D_\infty\).

\begin{theorem}[Finite-representation proof metric and curvature sign]
\label{p2:thm:finite-mode-trace-metric}
For any faithful finite-dimensional unitary realization \(\rho\) supplied by
Theorem~\ref{p2:thm:finite-mode-lie-regularity},
\[
 g_\rho(X,Y):=-\operatorname{Re}\operatorname{Tr}
 \bigl(d\rho(X)d\rho(Y)\bigr)
\]
is a positive-definite \(\operatorname{Ad}(G)\)-invariant inner product on
\(\mathfrak g\).  Its translated metric on \(G\) is bi-invariant and satisfies
\[
 K(X,Y)=\frac14
 \frac{\lVert[X,Y]\rVert^2}{\lVert X\wedge Y\rVert^2}\ge0.
\]
It is a proof metric, not a GP2 authorization to retain every higher trace
tensor of a chosen representation.
\end{theorem}

\begin{proof}
Faithfulness makes \(d\rho\) injective, and skew-Hermiticity makes the displayed
form positive definite.  Trace cyclicity proves Ad-invariance, which is
equivalent to bi-invariance of the translated metric.  The standard curvature
formula for a bi-invariant metric \cite{Milnor1976} gives the displayed
expression.  It gives
strictly positive sectional curvature on the nonabelian three-dimensional
branch, but compactness by itself gives only the nonnegative formula.
\end{proof}

\begin{theorem}[Consequences after comparison compactness and finite-operation Lie closure]
\label{p2:thm:gp1-compact-geometric-consequences}
For the compact completion \(\widehat X\), and for the compact Lie operation
group \(G\) when the finite-representation theorem is invoked:
\begin{enumerate}[label=(\roman*), leftmargin=2em]
\item every continuous scalar quantity constituted on \(\widehat X\) has
compact range and attains its extrema;
\item on any subsequently generated compact smooth Riemannian parent without
boundary, every smooth tensor and each fixed-order covariant derivative is
bounded, volume is finite, and injectivity radius is positive;
\item every smooth vector field on such a boundaryless compact parent is
complete;
\item a compact operation group has normalized Haar measure, is unimodular,
and every continuous homomorphism from it to \((\mathbb R^n,+)\) is trivial;
\item for every operation \(a\) in a compact operation group there are positive
integers \(n_k\to\infty\) with \(a^{n_k}\to e\) along a subsequence.
\end{enumerate}
Thus bounded global composition permits periodic or almost-periodic
reapplication but no nontrivial globally additive unwrapped real coordinate.
\end{theorem}

\begin{proof}
Item (i) is the extreme-value theorem.  Item (ii) follows by applying it to
continuous tensor norms on a compact smooth manifold without boundary and by
the positive injectivity-radius theorem on that domain
\cite{do1992riemannian}.  A smooth density integrates to a finite
value over a compact manifold, giving the finite volume in (ii).  Smooth vector
fields on compact manifolds without boundary are complete, proving (iii).
Compact groups possess finite Haar measure, which can be normalized; the
modular function is a continuous homomorphism \(G\to(\mathbb R_{>0},\times)\),
so its image is a compact subgroup of \(\mathbb R_{>0}\) and therefore
trivial, giving unimodularity; and a continuous image in \(\mathbb R^n\) is a
compact subgroup and hence is \(\{0\}\), proving (iv).  For (v), compactness
gives two arbitrarily late powers in the same sufficiently small neighborhood;
their quotient is a positive power arbitrarily close to the identity.  A
diagonal subsequence gives the claim.
\end{proof}

\begin{theorem}[FRC-b bounded spectral range gives finite spectral subspaces]
\label{p2:thm:gp1-bounded-spectrum-finite-carrier}
Let \(\Omega\) be a compact smooth parent manifold without boundary --- an ambient domain for
this import, unrelated to the comparison carrier \(X\) of
Definition~\ref{p2:def:intrinsic-operational-boundedness} --- let
\(E\to\Omega\) be a finite-rank
Hermitian vector bundle, fix a smooth positive density \(\mu\), and let \(L\)
be a nonnegative self-adjoint elliptic geometric differential operator of
positive order on
\(L^2(\Omega,E;\mu)\).  Then \(L\) has compact resolvent, discrete spectrum of
finite multiplicity, and eigenvalues tending to infinity.  Consequently
every FRC-b bounded spectral interval defines a finite-dimensional spectral
subspace
\[
 \mathcal E_{\leq\Lambda}(L)
 :=\mathbf1_{[0,\Lambda]}(L)L^2(\Omega,E;\mu),
 \qquad \dim\mathcal E_{\leq\Lambda}(L)<\infty.
\]
Finiteness of the spectral subspace therefore follows from compact geometry,
the positive-order elliptic (hence compact-resolvent) operator, and the bounded
spectral range.  A bound on a single nonelliptic component, without a joint
positive-order elliptic or Casimir bound, need not control multiplicity and is
not sufficient: on \(T^2\) the operator \(L_1=-\partial^2/\partial x_1^2\) has
eigenvalue \(m^2\) with infinite multiplicity, indexed by the free second
mode, so every bounded interval of its spectrum already carries an
infinite-dimensional space.  The notation
\(\mathcal E_{\leq\Lambda}(L)\) is deliberately distinct from Paper~I-A's
finite module \(\mathcal H_m\).  On a compact manifold with boundary the same
conclusion may be used only after a particular elliptic self-adjoint boundary
realization with compact resolvent and its operator domain have been supplied;
no boundary branch is invoked in this paper.
\end{theorem}

\begin{proof}
The compact-resolvent theorem for positive-order elliptic differential
operators on closed compact manifolds \cite{gilkey1995invariance} gives
an orthonormal eigenbasis, finite multiplicities, and no finite accumulation
point of eigenvalues.  Hence only finitely many eigenspaces occur below a fixed
\(\Lambda\), and their finite direct sum is finite-dimensional.
\end{proof}

\subparagraph{Tier I: compact-Lie and generated-test closure \([{\rm CL}]\).}

\begin{theorem}[Global structure of the compact connected operation parent]
\label{p2:thm:compact-connected-global-structure}
There exist a torus \(T^k\), compact simply connected simple Lie groups
\(H_1,\ldots,H_m\), and a finite central subgroup \(F_0\) such that
\[
 H\cong
 \frac{T^k\times H_1\times\cdots\times H_m}{F_0}.
\]
More intrinsically,
\[
 H=Z(H)^0[H,H],
 \qquad
 Z(H)^0\cap[H,H]\ \text{is finite},
\]
the commutator subgroup \([H,H]\) is closed, connected, and semisimple, and
\[
 \operatorname{rank}H
 =k+\sum_{i=1}^m\operatorname{rank}H_i.
\]
Thus (CL) determines a finite toral--simple decomposition type, although it
does not force \(k+m=1\) or select any simple type.
\end{theorem}

\begin{proof}
The Lie algebra of a compact connected Lie group has the orthogonal ideal
decomposition
\[
 \mathfrak h=\mathfrak z(\mathfrak h)
 \oplus\mathfrak h_1\oplus\cdots\oplus\mathfrak h_m
\]
with compact simple ideals \(\mathfrak h_i\).  The central algebra integrates
to \(Z(H)^0\cong T^k\), while the simply connected integrations of the simple
ideals form \(H_1\times\cdots\times H_m\).  Multiplication induces a
surjective homomorphism from their product to \(H\); its kernel is discrete
because the differential is an isomorphism, compact because the domain is
compact, and central because it lies in the overlap of commuting factors.
It is therefore finite.  The same construction proves the intrinsic
factorization and the stated properties of \([H,H]\).  Maximal toral
dimensions add across the Lie-algebra direct sum, giving the rank formula;
see, for example, \cite{chevalley1946theory,Hall2015lie}.
\end{proof}

\begin{theorem}[Maximal tori and surjectivity of the exponential map]
\label{p2:thm:maximal-torus-exponential}
Every element of \(H\) lies in a maximal torus, every torus is contained in a
maximal torus, and all maximal tori are conjugate.  Consequently
\[
 \exp:\mathfrak h\longrightarrow H
\]
is surjective.  The common dimension of the maximal tori is
\(\operatorname{rank}H\).  These statements supply toral representatives but
do not select one maximal torus or one logarithm of an operation.
\end{theorem}

\begin{proof}
The compact connected maximal-torus theorem states that every element lies in
a maximal torus, every torus extends to one, and all maximal tori are
conjugate.  If \(h\in H\) lies in such a torus \(T\), the exponential map of
\(T\cong\mathbb R^r/\Lambda\) is the quotient map and hence surjective, so
\(h=\exp X\) for some \(X\in\mathfrak t\subseteq\mathfrak h\).
The remaining assertions are the standard maximal-torus consequences.
\end{proof}

\begin{theorem}[Closures of one-parameter operation subgroups]
\label{p2:thm:one-parameter-operation-closure}
For \(X\in\mathfrak h\), the closure
\[
 T_X:=\overline{\{\exp(tX):t\in\mathbb R\}}
\]
is a compact connected abelian Lie subgroup, hence a torus; it is trivial
exactly when \(X=0\).  For \(X\ne0\), the following are equivalent:
\[
\begin{split}
 &\exp(TX)=e\ \text{for some }T>0,\\
 &\ker\bigl(t\mapsto\exp(tX)\bigr)\ne\{0\},\\
 &T_X\cong U(1).
\end{split}
\]
If \(\dim T_X>1\), the one-parameter subgroup is nonperiodic and dense in
that higher-dimensional torus.  Thus compactness implies recurrence, not
universal periodicity.
\end{theorem}

\begin{proof}
The image of \(\mathbb R\) is connected and abelian, so its closure has the
same properties.  Every compact connected abelian Lie group is a torus.
For \(X\ne0\), the homomorphism has either zero kernel or a lattice
\(T\mathbb Z\).  A nonzero kernel makes its image a compact circle and hence
closed.  Conversely, a nontrivial one-parameter subgroup whose closure is a
circle is a nontrivial continuous homomorphism
\(\mathbb R\to U(1)\), so it has a nonzero lattice kernel.  A zero kernel
therefore occurs exactly when the image is dense in a torus of dimension at
least two.
\end{proof}

\begin{theorem}[Regular homogeneity of the operation parent]
\label{p2:thm:regular-homogeneity-closure}
\label{p2:cor:homogeneity-from-transfer-group}
\label{p2:prop:homogeneity-derived}
The left regular action
\[
 L:H\times H\longrightarrow H,\qquad (a,x)\longmapsto ax
\]
is smooth, free, and transitive.  It preserves the translated operational
Finsler metric and the John metric \(g_F\), when those optional
primitive-comparison regularity structures are available, and it preserves
the canonical negative-Killing metric \(g_K\) on the selected \(SU(2)\)
branch.  More generally it preserves any retained left-invariant metric,
including \(g_{\rm rk}\) in the rank-one catalogue.  Hence the operation
parent is intrinsically homogeneous and every point has the same metric and
differential invariants.  The identity remains algebraically distinguished
as the neutral operation, but it is not a metrically preferred point.
This is the canonical operation-value self-realization \((H,L)\).  It
concerns the operation parent \(H\), not an external state space on which
\(H\) may act.
\end{theorem}

\begin{proof}
For \(x,y\in H\), the unique left translation taking \(x\) to \(y\) is
\(L_{yx^{-1}}\), proving freeness and transitivity.  Multiplication and
inversion are smooth in the derived Lie structure.  Every left-invariant
structure is preserved by definition; this applies to the optional
\(F,g_F\), to the canonical \(g_K\) on the selected branch, and to
\(g_{\rm rk}\) in Tier~II.  At the same
time, only the identity satisfies
\(ex=xe=x\) for every \(x\), which proves the stated algebraic--metric
distinction.
\end{proof}

\begin{theorem}[Intrinsic left--right isometry action and inversion]
\label{p2:thm:left-right-isometry-action}
Let \(g\) be any retained bi-invariant Riemannian metric on \(H\).
The formula
\[
 (a,b)\cdot x:=axb^{-1}
\]
defines an isometric action of \(H\times H\) on
\((H,g)\) with kernel
\[
 \Delta Z(H):=\{(z,z):z\in Z(H)\}.
\]
It is transitive, and its stabilizer of the identity is the diagonal subgroup
modulo this kernel:
\[
 (H\times H)_e/\Delta Z(H)
 \cong H/Z(H)\cong\operatorname{Inn}(H).
\]
The induced tangent action at the identity is \(\operatorname{Ad}(H)\).
Inversion
\[
 \iota(x)=x^{-1}
\]
is an additional involutive isometry; it interchanges the left and right
factors and is a group automorphism only when \(H\) is abelian.
\end{theorem}

\begin{proof}
Bi-invariance of \(g\) makes both left and right translations isometries.
If \(axb^{-1}=x\) for every \(x\), then evaluation at \(e\) gives \(a=b\),
and the remaining equality says \(a\in Z(H)\); this proves the kernel.
Transitivity follows already from the left factor.  A pair fixes \(e\)
exactly when \(a=b\), and its action is conjugation, whose kernel is the
center and whose differential is the adjoint representation.  Finally,
bi-invariance gives
\(d(x^{-1},y^{-1})=d(x,y)\), and
\((xy)^{-1}=y^{-1}x^{-1}\) proves the last assertion.
\end{proof}

\begin{theorem}[Canonical parallelism and universal topological consequences]
\label{p2:thm:parallelism-topology-closure}
Left translation gives the intrinsic vector-bundle trivialization
\[
 TH\longrightarrow H\times\mathfrak h,\qquad
 v_x\longmapsto
 \bigl(x,(dL_{x^{-1}})_xv_x\bigr).
\]
Consequently \(H\) is parallelizable, orientable, spin-admissible,
boundaryless, and of finite CW type.  If \(\dim H>0\), then
\[
 \chi(H)=0.
\]
These are existence statements.  The constitutional chain selects no basis of
\(\mathfrak h\), orientation convention, framing, or spin structure.
\end{theorem}

\begin{proof}
The displayed map and its inverse
\((x,X)\mapsto(dL_x)_eX\) are smooth, so the tangent bundle is trivial.
All Stiefel--Whitney classes of a trivial bundle vanish; in particular
\(w_1=w_2=0\), which gives orientability and existence of a spin structure.
A finite-dimensional Lie group is a manifold without boundary, and a compact
smooth manifold has the homotopy type of a finite CW complex.  When
\(\dim H>0\), any nonzero \(X\in\mathfrak h\) gives a nowhere-zero
left-invariant vector field.  Poincaré--Hopf then yields \(\chi(H)=0\).
Choosing a global frame or its orientation requires a basis and is not part
of the bundle trivialization itself.
\end{proof}

\begin{theorem}[Finite-resolution distance-profile tomography]
\label{p2:thm:finite-resolution-profile-tomography}
Let \(Z_\varepsilon\subset X\) be a finite \(\varepsilon\)-net for the
canonical completion \(\widehat X\), as supplied by GP1, and define
\[
 \kappa_{Z_\varepsilon}:\widehat X\longrightarrow
 \ell_\infty(Z_\varepsilon),\qquad
 \kappa_{Z_\varepsilon}(x)
 :=\bigl(\widehat d(x,z)\bigr)_{z\in Z_\varepsilon}.
\]
Then \(\ell_\infty(Z_\varepsilon)\) is finite-dimensional and, for every
\(x,y\in\widehat X\),
\[
 \widehat d(x,y)-2\varepsilon
 \leq
 \left\|\kappa_{Z_\varepsilon}(x)
       -\kappa_{Z_\varepsilon}(y)\right\|_\infty
 \leq \widehat d(x,y).
\]
Thus finitely many generated distance tests reconstruct every pairwise
distance with additive error at most \(2\varepsilon\).  GP1 supplies finite
approximate tomography at every positive resolution, while exact profile
separation is recovered as \(\varepsilon\downarrow0\).
\end{theorem}

\begin{proof}
The upper bound is the reverse triangle inequality.  Choose
\(z\in Z_\varepsilon\) with \(\widehat d(x,z)\leq\varepsilon\).  Then
\[
 |\widehat d(y,z)-\widehat d(x,z)|
 \geq \widehat d(y,z)-\widehat d(x,z)
 \geq \widehat d(x,y)-2\widehat d(x,z)
 \geq \widehat d(x,y)-2\varepsilon.
\]
Corollary~\ref{p2:cor:gp1-resolution-completeness-derived} supplies such a
finite net with centres in the dense original carrier.
\end{proof}

\begin{theorem}[Derived finite point-resolution support role]
\label{p2:thm:derived-finite-point-support-role}
Let \((X,d)\) be an actual distinguishable comparison role satisfying GP1-TB,
and assume its already constituted readout grammar is metric-curry closed.
For a realized positive resolution \(\varepsilon\), let
\(\mathsf{RES}^{\rm pt}_\varepsilon\) denote the G0 occurrence of the local
point-reference readout: finite partial evaluations are object-typed as one
readout support rather than retained only as proof witnesses.  Then the
unmarked support output
\[
 \mathsf{Supp}_\varepsilon(X)
 :=
 \left\{
 [Z,\kappa_Z]:
 Z\subset X\text{ is a finite }\varepsilon\text{-net}
 \right\}
\]
is nonempty, every representative has finite point support, and every
representative reconstructs the completed metric with additive error at most
\(2\varepsilon\).

Consequently, once distinguishable point-resolution readout occurs, the
support role and its finiteness require no independent finite-support datum:
metric currying supplies the reference-test constructor, GP1 supplies finite
resolution, and GP2 objectification and no-omission retain the complete
unmarked output class.  GP1--GP2 do not select a labelled net.  An actual
local readout may mark one representative, but that marking is a contingent
G0 occurrence and does not alter the support atlas.
\end{theorem}

\begin{proof}
Corollary~\ref{p2:cor:gp1-resolution-completeness-derived} makes the displayed
class nonempty and makes every \(Z\) finite.  Metric-curry closure generates
the tests \(b_z=d(z,\cdot)\), and the point-reference occurrence gives their
finite tuple the required object typing.  Theorem~\ref{p2:thm:finite-resolution-profile-tomography}
gives the \(2\varepsilon\) reconstruction bound.  GP2 typed host faithfulness
retains the object-valued image and its relations, while
Lemma~\ref{p2:lem:gp2-canonical-construction} forbids promoting a particular
representative when the prior geometry distinguishes no such choice.
\end{proof}

\begin{proposition}[Generic point resolution does not force exceptional support]
\label{p2a:prop:generic-resolution-not-exceptional}
Under the hypotheses of
Theorem~\ref{p2:thm:derived-finite-point-support-role}, the generated atlas
\(\mathsf{Supp}_\varepsilon(X)\) imposes finite \(\varepsilon\)-coverage and
the stated metric-profile error, but it does not imply any of
\[
 |Z|=24,\qquad Z=-Z,\qquad ZZ=Z,\qquad -1\in Z,
 \qquad \delta(Z)=\frac\pi3,
\]
nor the spherical-design or optimal-packing predicates of the Hurwitz shell.
In particular,
\[
 \boxed{
 \mathsf{RES}^{\rm pt}_\varepsilon
 \Longrightarrow\mathsf{Supp}_\varepsilon(S^3)
 \quad\not\Longrightarrow\quad
 \operatorname{CoReal}(\mathsf{Supp}_{2T}).}
\]
For every \(\varepsilon>0\), an allowed \(\varepsilon\)-net exists that is
not group closed and contains no subset orthogonally congruent to the regular
\(24\)-cell.  Thus exceptional-domain entry is a genuine additional ancestry
condition, not omitted Paper-I-A geometry.  The perturbation proof and the
coarse-cardinality countermodel are recorded in
Section~\ref{p2a:app:generic-net-exceptional-obstruction}.
\end{proposition}

\begin{proposition}[No natural embedded finite representative selector from unmarked homogeneous data]
\label{p2a:prop:no-natural-finite-exceptional-selector}
The following representative-selector obstructions hold.
\begin{enumerate}[label=(\roman*),leftmargin=2.5em]
\item Let \(Z\subset S^3\) be determined solely from the unmarked oriented
round metric by a rule natural under every orientation-preserving round
isometry.  Then \(Z=\varnothing\) or \(Z=S^3\).  In particular no such rule
selects a nonempty finite embedded support.
\item Let \(Z\subset SU(2)\) be determined solely from the unmarked Lie-group
structure by a rule natural under every Lie-group automorphism.  If \(Z\) is
finite, then
\[
 Z\subset Z(SU(2))=\{\pm I\}.
\]
In particular no such rule selects a 24-point embedded subset or an embedded
copy of \(2T\).
\end{enumerate}
Consequently neither unmarked homogeneous object supplies a natural embedded
24-point representative.  This conclusion does not forbid a natural orbit or
family of such representatives, nor an invariant functional or physical
process whose unique minimizer orbit is \([2T]\); any theorem using that
orbit-level route must separately constitute and license the role or
mechanism.
\end{proposition}

\begin{proof}
For (i), naturality on the same unmarked round object gives
\(hZ=Z\) for every \(h\in\operatorname{Isom}^{+}(S^3)\).  If
\(z\in Z\) and \(y\in S^3\), transitivity supplies an \(h\) with
\(h(z)=y\), hence \(y\in Z\).  Thus every nonempty invariant subset is all
of \(S^3\), which is not finite.

For (ii), automorphism naturality implies invariance under every inner
automorphism, so \(Z\) is a union of conjugacy classes.  Every noncentral
conjugacy class in \(SU(2)\) is a positive-dimensional two-sphere and is
therefore infinite.  A finite conjugacy-invariant subset can consequently
contain only central elements, proving \(Z\subset\{\pm I\}\).  Such a set
cannot have cardinality \(24\) and cannot be an embedded \(2T\).
\end{proof}

\begin{corollary}[Exceptional-support status after the representative no-go]
\label{p2a:cor:exceptional-support-status-after-no-go}
The representative no-go does not leave one universal ``\(2T\) occurrence
root.''  It leaves a conditional support-entry atlas.  A physical theorem
which actually consumes exceptional point/group geometry must enter through a
typed branch such as
\begin{enumerate}[label=(\alph*),leftmargin=2.5em]
\item a direct occurrence of the required exceptional support role;
\item an occurring optimal \(24\)-point support or a bridged point-support
constructor satisfying the corresponding Paper-I geometric hypotheses;
\item a finite operation/comparison role satisfying
Theorem~\ref{p2a:thm:v236-qop-exceptional-frontier};
\item an occurring stronger exceptional parent whose licensed descendant
constructor produces the required core while retaining the parent role; or
\item an abstract/torsor realization followed by an independently licensed
same-owner action or embedding bridge.
\end{enumerate}
Once such a branch occurs, all nonredundant unmarked descendants are
\(\mathrm{GEN}\) and need no duplicate occurrence certificate.  A labelled
embedded representative, performed support readout, orientation or trimer
remains a member token only when a physical claim consumes it.  Claims using
only \(\mathrm{Cmp}_m\) do not thereby require a physical exceptional support.
\end{corollary}

\begin{corollary}[Core-support inheritance from point resolution]
\label{p2:cor:core-support-inheritance-no-duplicate-witness}
Assume the hypotheses of
Theorem~\ref{p2:thm:derived-finite-point-support-role}.  Suppose in addition
that the frozen support grammar contains a candidate-independent core
constructor
\[
 C_{\rm core}:\mathsf{Supp}_\varepsilon(X)\longrightarrow
 \mathsf{CoreSupp}(X)
\]
whose nonempty output is canonical up to the unmarked support automorphisms.
Then the complete output orbit is a generated-role descendant of the occurring
point-reference composite and is co-realized without a second support
occurrence certificate.  If a further bridge theorem identifies this
point-resolution support domain with the domain of the Paper-I-B exceptional
core constructor, whose unique output orbit is \([2T]\), this gives
\[
 \mathsf{RES}^{\rm pt}_\varepsilon\text{ occurrence}
 \Longrightarrow
 \operatorname{CoReal}(\mathsf{Supp}_{2T}\mid
                        \mathsf{RES}^{\rm pt}_\varepsilon).
\]
A labelled Hurwitz vertex, body diagonal, orientation, or trimer remains a
representative choice or a further token; none is selected by this corollary.
The Paper-I-B exceptional atlas by itself does not prove the required domain
bridge.  Until that bridge is generated or its constructor is actually
applied under G0, the exceptional branch has not been entered.  Its remaining
status is an optional support occurrence or the output of a separately
constituted orbit-level support-selection mechanism, not an unresolved
natural embedded selector of the unmarked round/group data.  Once the
antecedent holds, the occurrence of its \([2T]\) output is redundant; a
labelled embedded representative is not thereby selected.
\end{corollary}

\begin{proof}
The point-reference occurrence generates the finite metric support atlas by
Theorem~\ref{p2:thm:derived-finite-point-support-role}.  The core constructor
is frozen, invariant, candidate-independent, and typed on that atlas only
under the additional domain bridge, so its complete output orbit
is canonical and nonredundant in its declared support sort.  It therefore
satisfies Definition~\ref{p2:def:finitely-presented-generated-role-descendant},
and Theorem~\ref{p2:thm:generated-role-corealization} gives co-realization.
Passing to a labelled member would add a mark not fixed by the orbit, proving
the representative clause.  If the domain bridge or constructor application
is absent, GP2 no-insertion forbids the conclusion, proving the final scope
clause.
\end{proof}

\begin{corollary}[Two exact conditional discharge routes to the unmarked \(2T\) core]
\label{p2a:cor:two-conditional-2t-discharge-routes}
Assume the occurring point-resolution support of
Theorem~\ref{p2:thm:derived-finite-point-support-role}.  Either of the
following independently generated antecedents is sufficient to discharge the
unmarked exceptional-core output without a second support occurrence.
\begin{enumerate}[label=(\roman*),leftmargin=2em]
\item The occurring representative \(Z\subset S^3\) satisfies
\[
 |Z|=24,
 \qquad
 \min_{z\ne z'}d_{S^3}(z,z')=\frac\pi3.
\]
Then Paper~I-B,
Theorem~\ref{p1:thm:s3one-classical-completions}, identifies \([Z]\), up to
orthogonal equivalence, with the unique optimal four-dimensional kissing
configuration \([2T]\).
\item A prior bridge proves that the occurring support is a finite subgroup
\(\Gamma\subset SU(2)\), contains \(-1\), lies in one common axis-free
exceptional role, and is compared within that role by subgroup leastness.
Paper~I-B then restricts the candidates to \(2T,2O,2I\), and
Corollary~\ref{p1:cor:conditional-exceptional-role-order} makes \(2T\) the
unique least class.
\end{enumerate}
Neither antecedent follows from finite \(\varepsilon\)-coverage, from the
label \(m=1\), from \(\dim\mathcal H_1=15\), or from the identity
\(24=15+9\).  Stipulating either antecedent solely in order to obtain \(2T\)
rephrases rather than removes the exceptional-domain root.
\end{corollary}

\begin{proof}
Part~(i) is the Paper-I-B optimal-kissing identification transported by
Corollary~\ref{p2:cor:core-support-inheritance-no-duplicate-witness}.
For part~(ii), Paper I-B supplies the three-class exceptional range and the
inclusions \(2T\subset2O\) and \(2T\subset2I\); GP3 can use that order only
after the three groups have been typed as candidates for the same physical
role.  Proposition~\ref{p2a:prop:generic-resolution-not-exceptional} proves
that neither antecedent is generated by the generic support atlas.
\end{proof}

The full-quadratic specialization is now proved in
Theorem~\ref{p2a:thm:v233-quadratic-support-discharge}. It derives
exceptionality and the central element from exactness once finite
group closure and role-faithful comparison are licensed. The
72-point counterexample in Paper~I-C shows why those remaining
requirements cannot be replaced by exact sampling alone.

\begin{corollary}[Conditional exhaustion of the trimer/axis binary-support branch]
\label{p2:cor:conditional-trimer-axis-support-exhaustion}
Assume an occurring point-resolution role as in
Theorem~\ref{p2:thm:derived-finite-point-support-role}, and let
\(\Gamma\subset SU(2)\) be a group-closed support descendant for which an
additional bridge establishes \(-1\in\Gamma\) and the Paper-I-B
trimer/axis-admissibility disjunction.  Then its complete unmarked binary
support range is exactly
\[
 \left\{
   \mathsf{Supp}_{C_6^{\rm grp}},\quad
   \mathsf{Supp}_{2T},\quad
   \mathsf{Supp}_{2O},\quad
   \mathsf{Supp}_{2I}
 \right\}.
\]
This is an exhaustion statement, not a GP3 selection: it does not force any
member to occur, and if distinct role constructors generate several members,
GP2 no-omission retains all of them.  The generic metric support atlas
\(\mathsf{Supp}_\varepsilon(X)\) remains unrestricted by this four-class
range.

On the core branch, once the bridge in
Corollary~\ref{p2:cor:core-support-inheritance-no-duplicate-witness}
co-realizes \(\mathsf{Supp}_{2T}\), the Paper-I-B equivariant constructor
generates the complete unmarked families of cyclic diagonal stabilizers
\(C_6^{\rm grp}\) and link graphs \(C_6^{\rm link}\) without a duplicate
role-level occurrence witness.  A particular labelled vertex, diagonal, or
orientation remains a further G0 token if it is used physically.  The
two-coset configuration \(2T\sqcup g2T\) is the \(2O\) normalizer support,
not a fifth class and not the direct product \(C_2\times2T\).
\end{corollary}

\begin{proof}
Paper-I-B Theorem~\ref{p1:thm:trimer-axis-binary-support-range} proves the exact
four-class range from the stated bridge hypotheses.  The first paragraph's
status clauses follow from GP2 exact/no-omission closure and GP3 no-selection.
For the core branch, apply the preceding corollary and then Paper-I-B
Corollary~\ref{p1:cor:generated-c6-stabilizer-link-bridge}; the latter
constructor is equivariant on the complete unmarked orbit, so
Theorem~\ref{p2:thm:generated-role-corealization} removes any duplicate
role-level witness but cannot select a marked representative.  Paper-I-B
Proposition~\ref{p1:prop:2O-two-cosets-not-product} proves the final clause.
\end{proof}

\begin{theorem}[Completion of the generated distance-test algebra]
\label{p2:thm:generated-test-algebra-completion}
Let
\[
 b_z(x):=\widehat d(x,z),\qquad z,x\in\widehat X,
\]
and let \(\mathcal A_d\) be the unital real algebra generated by these
profiles.  Then
\[
 \overline{\mathcal A_d}^{\|\cdot\|_\infty}
 =C(\widehat X,\mathbb R).
\]
Hence finite sums and products of the canonical generated profiles are
uniformly dense in all continuous real tests on the compact comparison parent.  This
is a functional completion theorem, not a claim that every continuous
function is a primitive B-operation.
\end{theorem}

\begin{proof}
The algebra contains the constants.  The profile family separates points by
Theorem~\ref{p2:thm:canonical-metric-profile-separation}, so the algebra it
generates also separates points.  The real Stone--Weierstrass theorem on the
compact Hausdorff space \(\widehat X\) gives the stated uniform density.
\end{proof}

\begin{theorem}[Universal separated effective-image factorization]
\label{p2:thm:universal-gp3-reduction}
Let \(H\) be a compact group acting continuously on the canonical compact GP1
completion \(\widehat X\), and let
\(\mathcal B\subset C(\widehat X,\mathbb R)\) be a generated test doctrine
invariant under pullback by \(H\).  Define
\[
 x\sim_{\mathcal B}y
 \quad\Longleftrightarrow\quad
 b(x)=b(y)\quad\text{for every }b\in\mathcal B
\]
and let
\[
 K_{\mathcal B}
 :=\{h\in H:b(hx)=b(x)\text{ for every }b\in\mathcal B,
                         x\in\widehat X\}.
\]
Then \(\sim_{\mathcal B}\) is a closed \(H\)-congruence and
\(K_{\mathcal B}\) is a closed normal subgroup.  The evaluation map
\[
 \operatorname{ev}_{\mathcal B}:\widehat X\longrightarrow\mathbb R^{\mathcal B},
 \qquad
 x\longmapsto(b(x))_{b\in\mathcal B}
\]
identifies \(\widehat X/{\sim_{\mathcal B}}\) canonically with its compact
image.
The induced operation group is canonically
\[
 H_{\mathrm{eff}}=H/K_{\mathcal B},
\]
and its action on the separated test image is faithful.  Every quotient
factorization in the universal statement is, by definition, a continuous
equivariant quotient pair
\[
 \pi:\widehat X\twoheadrightarrow Y,
 \qquad q:H\twoheadrightarrow L,
 \qquad \pi(hx)=q(h)\pi(x),
\]
with \(Y\) compact Hausdorff and \(L\) Hausdorff.  If
\[
 x\sim_{\mathcal B}y\Longrightarrow\pi(x)=\pi(y),
 \qquad K_{\mathcal B}\subseteq\ker q,
\]
then there are unique continuous equivariant factorizations
\[
 \bar\pi:\widehat X/{\sim_{\mathcal B}}\longrightarrow Y,
 \qquad
 \bar q:H/K_{\mathcal B}\longrightarrow L.
\]
If also \(\pi(x)=\pi(y)\Rightarrow x\sim_{\mathcal B}y\) and
\(\ker q\subseteq K_{\mathcal B}\), the factor maps are isomorphisms.  These
quotients are canonical; a smallest separating
subfamily of tests, when one exists, need not be unique.  This is a
candidate-internal effective-image theorem, not by itself a GP3 selection.  It
licenses a GP3 comparison only when the quotient congruence is generated and
both the source and a separately completed target are GP1-admissible,
candidate-internally GP2-exact, and faithful to the same role geometry.
\end{theorem}

\begin{proof}
Two points have the same evaluation image exactly when they are
\(\sim_{\mathcal B}\)-equivalent.  Since \(\widehat X\) is compact and the
product is
Hausdorff, the induced continuous bijection from the quotient onto the
evaluation image is a homeomorphism.  Pullback invariance gives
\(b\circ h\in\mathcal B\) for every \(b\in\mathcal B\) and \(h\in H\).
Hence, for \(x\sim_{\mathcal B}y\),
\[
 b(hx)=(b\circ h)(x)=(b\circ h)(y)=b(hy),
\]
so the state equivalence is an \(H\)-congruence.  The subgroup
\(K_{\mathcal B}\) is the kernel of the pullback action on the tests; it is
normal, and it is closed because it is the intersection of the closed sets
defined by \(b(hx)=b(x)\).  Hence \(H/K_{\mathcal B}\) acts on the state
quotient.  If a coset acts trivially there, then \(hx\sim_{\mathcal B}x\) for
every \(x\), so \(h\in K_{\mathcal B}\); the action is faithful.

The first implication for \(\pi\) says exactly that \(\pi\) is constant on
\(\sim_{\mathcal B}\)-classes, and the inclusion for \(q\) says exactly that
\(q\) kills \(K_{\mathcal B}\).  The two elementary quotient universal
properties give unique \(\bar\pi\) and \(\bar q\), and equivariance follows
from that of \((\pi,q)\).  Under the converse implications the induced maps
are bijective.  They are then homeomorphisms because their domains are compact
and their codomains Hausdorff.  Minimal generating or separating subfamilies
remain choices inside \(\mathcal B\) and are not supplied by either quotient.
\end{proof}

\begin{theorem}[Canonical faithful Banach realization of the completed geometric pairing]
\label{p2:thm:canonical-banach-realization}
Let \((\widehat X,\widehat d)\) be the canonical GP1 completion and let
\[
 \kappa:\widehat X\longrightarrow C(\widehat X,\mathbb R),
\qquad
\kappa(x)=b_x,
\qquad
 b_x(y)=\widehat d(x,y).
\]
Then
\[
 \lVert\kappa(x)-\kappa(y)\rVert_\infty=\widehat d(x,y),
\]
so \(\kappa\) is an isometric embedding with compact, hence closed, image.
Every completed generated isometry \(g\) acts linearly on
\(C(\widehat X,\mathbb R)\) by
\[
 (U_gf)(y):=f(g^{-1}y),
\]
and this action satisfies
\[
 U_g\kappa(x)=\kappa(gx),
 \qquad
 \lVert U_gf\rVert_\infty=\lVert f\rVert_\infty.
\]
The representation \(g\mapsto U_g\) is faithful.  Thus the complete
geometric A/B pairing has a canonical coordinate-free linear realization;
no basis, preferred point, or selected finite representation is required.
\end{theorem}

\begin{proof}
The equality is Theorem~\ref{p2:thm:canonical-metric-profile-separation}.
Continuity of \(\kappa\) and compactness of \(\widehat X\) make its image compact and
therefore closed in the Banach space \(C(\widehat X,\mathbb R)\).  Since \(g\) is a
surjective isometry,
\[
 \lVert U_gf\rVert_\infty
 =\sup_{y\in\widehat X}|f(g^{-1}y)|
 =\lVert f\rVert_\infty,
\]
and
\[
 (U_g\kappa(x))(y)
 =\widehat d(x,g^{-1}y)
 =\widehat d(gx,y)
 =\kappa(gx)(y).
\]
If \(U_g\) is the identity, then \(\kappa(gx)=\kappa(x)\) for every \(x\).
Injectivity of \(\kappa\) gives \(gx=x\) for every \(x\), hence \(g=e\) in
the generated isometry group.
\end{proof}

\begin{theorem}[Faithful quotient-kernel criterion]
\label{p2:thm:faithful-quotient-kernel}
Let a compact generated operation group \(G\) act on a generated test family
\(\mathcal B\subset C(X,\mathbb R)\), invariant under pullback, through
\[
 (\rho(g)b)(x):=b(g^{-1}x),
 \qquad
 K_{\mathcal B}:=\ker\rho.
\]
For a closed normal subgroup \(N\triangleleft G\), the action factors through
\(G/N\) if and only if \(N\subseteq K_{\mathcal B}\).  When it factors, the
induced action of \(G/N\) is faithful if and only if
\[
 N=K_{\mathcal B}.
\]
For the canonical distance-profile family of
Theorem~\ref{p2:thm:canonical-metric-profile-separation},
\[
 K_{\mathcal B}=\{e\}.
\]
Consequently no nontrivial quotient of the same generated geometric pairing
preserves its full transformation--test content faithfully.
\end{theorem}

\begin{proof}
An action descends to \(G/N\) exactly when every element of \(N\) acts
trivially, which is \(N\subseteq K_{\mathcal B}\).  The kernel of the induced
action is then \(K_{\mathcal B}/N\), so it is trivial exactly when
\(K_{\mathcal B}=N\).  Canonical metric profiles separate distinct generated
isometries by Theorem~\ref{p2:thm:canonical-metric-profile-separation}; hence
their action kernel is trivial.  The final statement concerns quotients of the
same constituted pairing and does not prohibit a different branch whose
primitive comparison parent was already a quotient.
\end{proof}

\begin{theorem}[Unique normalized invariant measure and averaging]
\label{p2:thm:unique-haar-averaging}
Every compact operation parent obtained from a compact comparison action or
from the finite reversible closure theorem carries a unique normalized
left-invariant regular Borel measure \(\mu_G\).  It is also right invariant.
For any derived bi-invariant Riemannian metric \(g\),
\[
 \mu_G
 =\frac{\operatorname{vol}_g}{\operatorname{vol}_g(G)}.
\]
The normalized measure is unchanged by a common positive rescaling of \(g\).
For every continuous finite-dimensional representation
\(\pi:G\to\operatorname{GL}(W)\),
\[
 \mathsf P_W:=\int_G\pi(g)\,d\mu_G(g)
\]
is the canonical projection onto the invariant subspace \(W^G\).  This
constructs invariant averaging but does not by itself assign a probability,
Born, or statistical interpretation to \(\mu_G\).
\end{theorem}

\begin{proof}
Existence and uniqueness of normalized left Haar measure hold for every
compact group.  Compact groups are unimodular, so the normalized measure is
also right invariant.  Riemannian volume of a bi-invariant metric is both left
and right invariant and therefore agrees with Haar measure after
normalization.  Rescaling the metric by a positive constant multiplies both
numerator and denominator by the same total-volume factor.  Haar invariance
gives \(\pi(h)\mathsf P_W=\mathsf P_W=\mathsf P_W\pi(h)\); idempotence follows
by a second integration, and the image is exactly \(W^G\).
\end{proof}

\begin{theorem}[Lie closure derives intrinsic real-analytic regularity]
\label{p2:thm:derived-real-analyticity}
Every finite-dimensional Lie operation parent supplied either by the direct
finite exact-operation closure or by the optional GP1-TB+(FD)+(LR) comparison
route carries its canonical
compatible real-analytic Lie-group structure.  Multiplication, inversion,
left and right translations, conjugation, and the exponential map are real
analytic.  Every continuous homomorphism between such
finite-dimensional Lie groups is real analytic.  No preferred analytic chart
or coordinate atlas is thereby selected.
\end{theorem}

\begin{proof}
Every finite-dimensional Lie group admits a unique compatible real-analytic
Lie-group structure, and continuous homomorphisms between finite-dimensional
Lie groups are analytic; see, for example, \cite{Hall2015lie}.  Applying this
standard structure theorem after either
Theorem~\ref{p2:thm:frc-automatic-compact-lie-closure} or
Theorem~\ref{p2:thm:finite-mode-lie-regularity} proves the claim.  The
analytic structure is intrinsic, whereas selecting a particular chart would
be additional coordinate data.
\end{proof}

\begin{theorem}[Peter--Weyl and complete representation-theoretic closure]
\label{p2:thm:peter-weyl-complete-reducibility}
Every finite-dimensional continuous complex representation of \(H\) admits an
invariant Hermitian inner product and is completely reducible.  If
\(\widehat H\) is the unitary dual, then the left--right regular
representation has the Hilbert decomposition
\[
 L^2(H,\mathbb C)
 \cong
 \widehat{\bigoplus}_{\pi\in\widehat H}
 V_\pi\otimes V_\pi^*.
\]
Finite linear combinations of matrix coefficients are uniformly dense in
\(C(H,\mathbb C)\), and the real linear span of their real and imaginary
parts is uniformly dense in \(C(H,\mathbb R)\).  Thus the operation parent has a complete
harmonic representation theory.  No representation, basis, finite subset of
\(\widehat H\), or physical mode interpretation is selected.
\end{theorem}

\begin{proof}
Averaging any Hermitian inner product over the normalized Haar measure of
Theorem~\ref{p2:thm:unique-haar-averaging} makes it \(H\)-invariant.
Orthogonal complements of invariant subspaces are then invariant, proving
complete reducibility by induction.  The Peter--Weyl theorem supplies the
displayed orthogonal sum and uniform density of matrix coefficients on every
compact group.  The real statement follows by taking real and imaginary
parts.  These constructions use all irreducibles and contain no canonical
finite truncation.
\end{proof}

\begin{theorem}[Canonical scalar Laplace--Casimir and heat-kernel package]
\label{p2:thm:laplace-casimir-heat-closure}
Let \(d=\dim H\) and let \(\Delta_{g_F}\ge0\) be the scalar
Laplace--Beltrami operator of the derived bi-invariant John metric.  Then
\(\Delta_{g_F}\) is the quadratic Casimir of the left and right regular
actions, commutes with both actions, and has a unique nonnegative
self-adjoint extension with compact resolvent.  Its spectrum is discrete,
each eigenspace is finite-dimensional and smooth, and
\[
 e^{-t\Delta_{g_F}}f(x)
 =\int_H K_t(x,y)f(y)\,d\operatorname{vol}_{g_F}(y)
 \qquad(t>0)
\]
for a smooth positive heat kernel \(K_t\).  Define
\[
 N_\Delta(\Lambda)
 :=\#\{\text{eigenvalues of }\Delta_{g_F}\leq\Lambda\},
\]
counted with multiplicity, and let
\[
 \omega_d
 :=
 \operatorname{vol}_{\mathbb R^d}
 \{\xi\in\mathbb R^d:\lvert\xi\rvert\leq1\}
 =
 \frac{\pi^{d/2}}{\Gamma(1+d/2)}
\]
be the Euclidean unit-ball volume.  Then
\[
 N_\Delta(\Lambda)
 \sim
 \frac{\omega_d\,\operatorname{vol}_{g_F}(H)}
 {(2\pi)^d}\,\Lambda^{d/2}
 \qquad(\Lambda\to\infty).
\]
The operator and its spectral asymptotics are canonical consequences of
\(g_F\); a Hamiltonian interpretation, spectral ceiling, or selected finite
mode space is not.
\end{theorem}

\begin{proof}
For an orthonormal basis \(X_1,\ldots,X_d\) of \(\mathfrak h\), interpreted as
left-invariant vector fields,
\[
 \Delta_{g_F}=-\sum_{a=1}^dX_a^2.
\]
Ad-invariance of \(g_F\) makes this element basis-independent and central in
the universal enveloping algebra, so the resulting operator commutes with
both regular actions.  Standard elliptic theory on the compact manifold \(H\)
gives essential self-adjointness, compact resolvent, smooth finite-multiplicity
eigenspaces, and the heat kernel.  Weyl's law gives the displayed asymptotic;
see, for example, \cite{berline1992heat,Vassilevich2003HeatKernel}.
\end{proof}

\begin{theorem}[Bi-invariant geodesic, curvature, and symmetric-space closure]
\label{p2:thm:bi-invariant-global-symmetry}
Let \(H\) carry any retained bi-invariant Riemannian metric \(g\), including
\(g_F\) in its optional domain, a finite-representation metric \(g_\rho\),
or \(g_K\) on the selected noncommutative branch.  Then every geodesic has
the form
\[
 \gamma(t)=h\exp(tX),
\]
and for left-invariant vector fields
\[
 \nabla_XY=\frac12[X,Y].
\]
With the curvature convention used in this paper,
\[
 R(X,Y)Z=-\frac14[[X,Y],Z],
\]
\[
 K(X,Y)=\frac14
 \frac{\|[X,Y]\|^2}{\|X\wedge Y\|^2}\ge0,
\qquad
 \operatorname{Ric}(X,Y)=-\frac14B_{\mathfrak h}(X,Y).
\]
For each \(h\in H\),
\[
 s_h(x):=hx^{-1}h
\]
is an involutive isometry fixing \(h\) with differential
\((ds_h)_h=-\operatorname{id}\).  Hence \(H\) is geodesically complete,
compact and globally symmetric, with \(\nabla R=0\), nonnegative sectional
curvature, zero curvature on every commuting two-plane, zero Ricci curvature
on the central directions, and positive Ricci curvature on every nonzero
semisimple direction.
\end{theorem}

\begin{proof}
For a bi-invariant metric, the Koszul formula gives
\(\nabla_XY=\frac12[X,Y]\), so one-parameter subgroups and their translates
are geodesics for all time.  A direct connection computation gives
\[
 R(X,Y)Z
 =\nabla_X\nabla_YZ-\nabla_Y\nabla_XZ-\nabla_{[X,Y]}Z
 =-\frac14[[X,Y],Z].
\]
Ad-invariance then gives the sectional-curvature formula, and tracing the
curvature operator gives
\(\operatorname{Ric}=-\frac14B_{\mathfrak h}\).
The Killing form vanishes on the center and is negative definite on the
compact semisimple ideal.

Inversion is an isometry fixing the identity and left and right translations
are isometries; their composition gives \(s_h\).  Direct calculation gives
\(s_h^2=\operatorname{id}\), and inversion differentiates to
\(-\operatorname{id}\) at the identity.  The point symmetries prove global
symmetry and \(\nabla R=0\); the remaining conclusions follow from the
displayed formulas.
\end{proof}

\subparagraph{Tier II: least-skeleton rank-one consequences
\([{\rm MR}]\).}

The remaining theorems in this tier are standard consistency refinements of
operational rank one.  On the GP3-least roles, (OR1) is already a corollary
of Theorem~\ref{p2:thm:master-compact-role-allocation}; it is not supplied as
another principle.  The \(SO(3)\) entries are retained only to display the
local quotient comparison excluded by GP2 relation exactness.

Before imposing rank one, fix on the arbitrary nontrivial compact connected
parent \(H\) a positive bi-invariant Riemannian metric \(g_H\), obtained from
any positive Ad-invariant inner product.  This is retained proof data for the
metric statements, not a selected primitive comparison.  Only after one of
the equivalent rank-one conditions below holds do we abbreviate the same
metric by \(g_{\rm rk}:=g_H\).  On the selected noncommutative branch it may be
the canonical negative-Killing metric \(g_K\); in the optional (FD)+(LR)
domain it may be \(g_F\).  The \(SO(3)\) row uses the quotient metric only as a
comparison.  The convention retains the free radius and does not assert that
the primitive comparison selected one of these scales.

\begin{theorem}[Equivalent intrinsic formulations of operational rank one]
\label{p2:thm:rank-one-equivalence-package}
For the nontrivial compact connected parent \(H\), conditions
\emph{(i)}, \emph{(ii)}, \emph{(iii)}, and \emph{(v)} below are equivalent.
If the optional (FD)+(LR) domain is in force, so that \(F\) exists, condition
\emph{(iv)} is equivalent to them as well.
\begin{enumerate}[label=(\roman*), leftmargin=2em]
\item \(H\) satisfies operational Cartan rank one:
\[
 [X,Y]=0\Longrightarrow X\wedge Y=0.
\]
\item \(\operatorname{rank}H=1\), equivalently every maximal torus is a
circle.
\item Every nonzero one-parameter subgroup of \(H\) is periodic.
\item For the retained Ad-invariant norm \(F\), the transformations
\(\{\pm\operatorname{Ad}_h:h\in H\}\) act transitively on the \(F\)-unit
sphere.
\item Either \(\dim H=1\), or every tangent two-plane has strictly positive
sectional curvature for \(g_{\rm rk}\).
\end{enumerate}
In the optional domain, these conditions force \(F\) to be Euclidean and
determined up to scale by \(g_{\rm rk}\) (equivalently by \(g_F\) there).
\end{theorem}

\begin{proof}
Every commuting family in a compact Lie algebra lies in a maximal toral
algebra, proving (i)\(\Leftrightarrow\)(ii).  Theorem~
\ref{p2:thm:one-parameter-operation-closure} proves
(ii)\(\Leftrightarrow\)(iii): if a maximal torus had dimension at least two,
an irrational direction would generate a dense nonperiodic flow.

Condition (iv) implies rank one because, inside a maximal toral algebra,
conjugate points meet only a finite Weyl-group orbit; an \(F\)-unit sphere of
dimension at least one cannot be one finite
\(\{\pm1\}W\)-orbit.  Conversely the rank-one classification leaves
\(U(1)\), \(SU(2)\), and \(SO(3)\).  The sign interchanges the two unit
directions of \(\mathfrak u(1)\), while in the nonabelian cases
\(\operatorname{Ad}(H)=SO(3)\) is transitive on directions.  Its invariance
also forces the \(F\)-unit ball to be a Euclidean ball.

Finally, the curvature formula in
Theorem~\ref{p2:thm:bi-invariant-global-symmetry} shows that a two-plane has
zero curvature exactly when it is spanned by independent commuting vectors.
This proves (i)\(\Leftrightarrow\)(v).
\end{proof}

\begin{theorem}[Actual metric isotropy after operational rank one]
\label{p2:thm:rank-one-actual-isotropy}
Every operational-rank-one parent with \(g_{\rm rk}\) is homogeneous and
isotropic in the global Riemannian sense: the isometry stabilizer of each
point acts transitively on its unit tangent sphere.

For either nonabelian global form,
\[
 \operatorname{Ad}(H)=SO(3)
\]
is the connected isotropy action at the identity, realized by actual global
conjugation isometries.  Together with inversion it generates the full tangent
isotropy \(O(3)\).  Translation carries this action to every point.  For
\(U(1)\), inversion interchanges the two unit tangent directions.
Thus isotropy is a theorem on the GP3-least roles (equivalently after the
rank-one condition), not a consequence of a quadratic form alone and not an
additional hypothesis on the classified branches.
\end{theorem}

\begin{proof}
Regular homogeneity is
Theorem~\ref{p2:thm:regular-homogeneity-closure}.  In the nonabelian cases the
Lie algebra is \(\mathfrak{su}(2)\), and the adjoint image is the standard
rotation group \(SO(3)\).  Theorem~\ref{p2:thm:left-right-isometry-action}
realizes it by conjugation isometries fixing the identity.  Since the
\(g_{\rm rk}\)-norm is Euclidean, \(SO(3)\) is transitive on the unit
two-sphere.
The derivative of inversion is \(-I\), so adjoining it gives \(O(3)\).
Conjugating the identity stabilizer by a left translation gives the
stabilizer at any point.  In dimension one the unit sphere has two points,
which inversion exchanges.
\end{proof}

\begin{theorem}[Rank-one space forms and Einstein curvature]
\label{p2:thm:rank-one-space-form-einstein}
With \(g_{\rm rk}\), the abelian branch is a flat round circle
\[
 U(1)\cong S^1(r)
\]
for one unfixed \(r>0\).  For each nonabelian parent there is one unfixed
curvature radius \(R>0\); its simply connected and adjoint global forms are,
respectively,
\[
 SU(2)\cong S^3(R),
 \qquad
 SO(3)\cong\mathbb{RP}^3(R)
\]
isometrically, the latter with the quotient of the round metric.  Locally,
\[
 K=R^{-2},\qquad
 \operatorname{Ric}=2R^{-2}g_{\rm rk},\qquad
 \operatorname{Scal}=6R^{-2}.
\]
Hence both nonabelian global forms are Einstein and locally round.  These
local tensors do not distinguish the cover; GP2 relation exactness does.
No value of \(r\) or \(R\) is selected by the constitutional chain.
\end{theorem}

\begin{proof}
Every invariant metric on a connected one-dimensional compact group is a
constant multiple of the standard circle metric.  On
\(\mathfrak{su}(2)\), every positive Ad-invariant inner product is a positive
multiple of the negative Killing form.  For an oriented orthonormal basis,
\[
 [X,Y]=\lambda\,X\times Y
\]
with \(\lambda\ne0\); the curvature formula gives
\(K=\lambda^2/4\) on every two-plane.  Put \(R=2/|\lambda|\).  A complete
simply connected three-manifold of constant curvature \(R^{-2}\) is
\(S^3(R)\), which identifies the universal cover \(SU(2)\).  The only
nontrivial connected quotient is by its central antipodal element, giving
\(\mathbb{RP}^3(R)\cong SO(3)\).  The Ricci and scalar formulas are the
standard contractions of a three-dimensional constant-curvature tensor.
Positive rescaling changes \(r\) or \(R\) without changing any structural
hypothesis.
\end{proof}

\begin{theorem}[Exact automorphism and metric-isometry groups]
\label{p2:thm:rank-one-exact-symmetry-groups}
For the rank-one metrics \(g_{\rm rk}\), the operation automorphism group,
connected
metric-isometry group, and full metric-isometry group are respectively
\[
\begin{array}{c|c|c|c}
H&
\operatorname{Aut}(H)&
\operatorname{Isom}_0(H,g_{\rm rk})&
\operatorname{Isom}(H,g_{\rm rk})\\ \hline
U(1)&O(1)&SO(2)&O(2)\\
SU(2)&SO(3)&SO(4)&O(4)\\
SO(3)&SO(3)&SO(3)\times SO(3)&O(4)/\{\pm I\}.
\end{array}
\]
In particular, for the nonabelian branches the operation-law stabilizer on
\(\mathfrak h\) is \(SO(3)\), while the full metric stabilizer at a point is
\(O(3)\).  Reflections preserve the round metric but do not preserve the Lie
bracket; they arise from composing an operation automorphism with inversion.
\end{theorem}

\begin{proof}
The automorphisms of \(U(1)\) are \(z\mapsto z^{\pm1}\).  Every automorphism
of the compact \(A_1\) Lie algebra, and hence of either connected nonabelian
global form, is inner, giving \(SO(3)\).

Theorem~\ref{p2:thm:rank-one-space-form-einstein} identifies the metric
spaces.  The round circle has isometry groups \(SO(2)\) and \(O(2)\).
The round three-sphere has connected and full isometry groups \(SO(4)\) and
\(O(4)\).  Isometries of the antipodal quotient are the projective orthogonal
group \(O(4)/\{\pm I\}\); its identity component is
\[
 SO(4)/\{\pm I\}\cong SO(3)\times SO(3).
\]
At one point their connected and full stabilizers are \(SO(3)\) and \(O(3)\).
An orientation-reversing linear isometry sends the cross-product bracket to
its negative, so it is not a Lie-algebra automorphism.  Inversion supplies the
required anti-automorphism component.
\end{proof}

\begin{theorem}[Complete topology and spin-admissibility of the rank-one list]
\label{p2:thm:rank-one-topology-spin}
The complete basic topological data of the three effective rank-one parents
are
\[
\begin{array}{c|c|c|c|c}
H&\text{manifold}&\pi_1(H)&Z(H)&H^\bullet(H;\mathbb Z)\\ \hline
U(1)&S^1&\mathbb Z&U(1)&H^0=H^1=\mathbb Z\\
SU(2)&S^3&0&\mathbb Z_2&H^0=H^3=\mathbb Z\\
SO(3)&\mathbb{RP}^3&\mathbb Z_2&1&
H^0=H^3=\mathbb Z,\ H^2=\mathbb Z_2 .
\end{array}
\]
All omitted integral cohomology groups vanish.  All three manifolds are
parallelizable and spin-admissible.  The numbers of spin structures are,
respectively,
\[
 2,\qquad 1,\qquad 2.
\]
The theorem counts admissible spin structures but selects none.
\end{theorem}

\begin{proof}
The manifold and fundamental-group identifications follow from the
classification and universal-cover kernel.  The centers are standard.
The cellular chain complexes of \(S^1\), \(S^3\), and
\(\mathbb{RP}^3\), together with the universal coefficient theorem, give the
displayed integral cohomology; in particular
\(H^2(\mathbb{RP}^3;\mathbb Z)=\mathbb Z_2\) while
\(H^1(\mathbb{RP}^3;\mathbb Z)=0\).  Parallelizability and spin-admissibility
are Theorem~\ref{p2:thm:parallelism-topology-closure}.  On a connected spin
manifold the spin structures form a torsor for \(H^1(H;\mathbb Z_2)\), whose
cardinalities for the three spaces are \(2,1,2\); see
\cite{hatcher2002algebraic,lawson1989spin}.
\end{proof}

\begin{theorem}[Global geodesic periods, radii, and cut loci]
\label{p2:thm:rank-one-geodesic-cut-locus}
Use the radii \(r\) and \(R\) of
Theorem~\ref{p2:thm:rank-one-space-form-einstein}.  Every nonconstant
unit-speed geodesic is closed, and the global data are
\[
\begin{array}{c|c|c|c|c}
H&\text{primitive period}&\operatorname{diam}&\operatorname{inj}&
\operatorname{Cut}(e)\\ \hline
U(1)&2\pi r&\pi r&\pi r&\text{antipodal point}\\
SU(2)&2\pi R&\pi R&\pi R&\text{antipodal point}\\
SO(3)&\pi R&\pi R/2&\pi R/2&\mathbb{RP}^2 .
\end{array}
\]
The \(SU(2)\) and \(SO(3)\) rows give global metric diagnostics of the
central quotient, but they do not replace the generated-test kernel that
decides which global form was originally constituted.
\end{theorem}

\begin{proof}
On \(S^1(r)\) and \(S^3(R)\), unit-speed geodesics are great circles of
period \(2\pi r\) and \(2\pi R\).  On \(S^1(r)\), the antipode is the cut
point from the base point but is not a conjugate point: the differential of
the one-dimensional exponential map never vanishes.  On \(S^3(R)\), the
antipode is both the first conjugate point and the cut point.  These facts
give the displayed diameter and injectivity radius.  The antipodal quotient
halves the primitive geodesic
period and identifies distances \(\theta\) with \(\pi-\theta\); hence
\(\operatorname{diam}=\operatorname{inj}=\pi R/2\).  Directions at distance
\(\pi R/2\), modulo antipodal identification in the tangent sphere, form
\(\mathbb{RP}^2\), the cut locus in \(\mathbb{RP}^3\).
\end{proof}

\begin{theorem}[Generated Cartan orientation on the nonabelian branch]
\label{p2:thm:rank-one-cartan-orientation}
On either nonabelian rank-one global form, define at the identity
\[
 \eta_e(X,Y,Z):=g_{\rm rk}(X,[Y,Z])
\]
and extend by left translation.  Then \(\eta\) is a nonzero closed
bi-invariant three-form.  It is a constant nonzero multiple of the Riemannian
volume form and therefore determines an orientation from the retained
multiplication and metric.  Every operation automorphism preserves \(\eta\),
whereas inversion satisfies
\[
 \iota^*\eta=-\eta.
\]
This generated Cartan orientation is neither a time orientation nor a
selected external spatial-orientation convention.
\end{theorem}

\begin{proof}
Ad-invariance of \(g_{\rm rk}\) makes
\(\eta_e\) alternating and Ad-invariant, so its left translate is also right
invariant.  The standard Chevalley--Eilenberg differential of the Cartan
three-form vanishes by invariance and the Jacobi identity.  On
\(\mathfrak{su}(2)\), the bracket is nonzero and the dimension is three, so
\(\eta_e\ne0\) and spans the one-dimensional space of alternating
three-forms.  Automorphisms preserve both the bracket and the unique
Ad-invariant metric scale, hence preserve \(\eta\).  Since
\((d\iota)_e=-I\), inversion changes the sign of a three-form.
\end{proof}

\begin{theorem}[Rank-one irreducibles, scalar spectra, and central descent]
\label{p2:thm:rank-one-representation-descent}
The irreducible complex representations and scalar Laplace spectra of the
rank-one parents are as follows.
\begin{enumerate}[label=(\roman*), leftmargin=2em]
\item For \(U(1)\), the irreducibles are the characters
\[
 \chi_n(e^{i\theta})=e^{in\theta},\qquad n\in\mathbb Z,
\]
and on \(S^1(r)\) their Laplace eigenvalues are \(n^2/r^2\).
\item For \(SU(2)\), the irreducibles are \(V_j\) with
\[
 j\in\tfrac12\mathbb Z_{\ge0},\qquad \dim V_j=2j+1.
\]
Writing \(\ell=2j\), the scalar spectrum on \(S^3(R)\) is
\[
 \lambda_\ell=\frac{\ell(\ell+2)}{R^2},
 \qquad
 \operatorname{mult}(\lambda_\ell)=(\ell+1)^2,
 \qquad \ell=0,1,2,\ldots.
\]
\item An \(SU(2)\) representation descends to \(SO(3)\) exactly when
\(-1\) acts trivially, equivalently \(j\in\mathbb Z_{\ge0}\).  Thus the
scalar spectrum on \(\mathbb{RP}^3(R)\) retains exactly the even degrees
\[
 \ell=0,2,4,\ldots
\]
with the same eigenvalue and multiplicity formulas.
\end{enumerate}
The \(SO(3)\) scalar doctrine is precisely the center-even part of the
\(S^3\) doctrine.  These are representation-theoretic consequences of the
global form, not a selection of physical spin sectors or a spectral cutoff.
\end{theorem}

\begin{proof}
Pontryagin duality gives the integer characters of the circle.  The highest
weight classification for \(SU(2)\) gives one \(V_j\) for each nonnegative
half-integer \(j\), and Peter--Weyl places
\(V_j\otimes V_j^*\) in scalar \(L^2(SU(2))\).  The quadratic Casimir gives
\(\ell(\ell+2)/R^2\) with \(\ell=2j\), and the summand dimension gives
\((\ell+1)^2\).  The center acts on \(V_j\) by
\((-1)^{2j}\), so descent through \(SU(2)\to SO(3)\) is equivalent to
integral \(j\), or even \(\ell\).  The quotient Laplacian is the restriction
to these invariant functions.
\end{proof}

\begin{corollary}[Exact scalar spectral counting functions]
\label{p2:cor:rank-one-scalar-spectral-counting}
Let \(N_H^{(0)}(\Lambda)\) count scalar Laplace eigenvalues not exceeding
\(\Lambda\ge0\), with multiplicity, for the rank-one metric of radius \(r\)
or \(R\).  Then
\[
 N_{U(1)}^{(0)}(\Lambda)
 =2\left\lfloor r\sqrt{\Lambda}\right\rfloor+1,
\]
\[
 N_{SU(2)}^{(0)}(\Lambda)
 =\frac{(L+1)(L+2)(2L+3)}6,
 \qquad
 L=\left\lfloor\sqrt{1+R^2\Lambda}-1\right\rfloor,
\]
and
\[
 N_{SO(3)}^{(0)}(\Lambda)
 =\frac{(K+1)(2K+1)(2K+3)}3,
 \qquad
 K=\left\lfloor
 \frac{\sqrt{1+R^2\Lambda}-1}{2}
 \right\rfloor.
\]
The first positive scalar eigenvalues and multiplicities are
\[
\begin{array}{c|c|c}
H&\lambda_1&\operatorname{mult}(\lambda_1)\\ \hline
U(1)&r^{-2}&2\\
SU(2)&3R^{-2}&4\\
SO(3)&8R^{-2}&9.
\end{array}
\]
These scalar \(L^2\) counting functions are not the dimension of Paper~I-A's
\(\Spin(4)\)-module \(\mathcal H_m\).
\end{corollary}

\begin{proof}
For \(U(1)\), the condition \(n^2/r^2\le\Lambda\) gives
\(|n|\le\lfloor r\sqrt\Lambda\rfloor\).  For \(SU(2)\),
\(\ell(\ell+2)\le R^2\Lambda\) is equivalent to
\(\ell\le L\), so
\[
 N_{SU(2)}^{(0)}(\Lambda)
 =\sum_{\ell=0}^{L}(\ell+1)^2
 =\frac{(L+1)(L+2)(2L+3)}6.
\]
For \(SO(3)\), put \(\ell=2k\).  The bound is \(k\le K\), and
\[
 N_{SO(3)}^{(0)}(\Lambda)
 =\sum_{k=0}^{K}(2k+1)^2
 =\frac{(K+1)(2K+1)(2K+3)}3.
\]
The first positive rows are the cases
\(|n|=1\), \(\ell=1\), and \(\ell=2\), respectively.
\end{proof}

\begin{theorem}[No canonical maximal torus on a nonabelian rank-one parent]
\label{p2:thm:no-canonical-maximal-torus}
Although \(SU(2)\) and \(SO(3)\) contain circle subgroups, every maximal torus
is conjugate to every other and no maximal torus is invariant under all
operation automorphisms.  Equivalently, choosing one such circle chooses an
axis in \(\mathfrak{su}(2)\) modulo sign.  Therefore the generated
nonabelian rank-one datum supplies neither a preferred scalar subgroup nor a
preferred Hopf axis.  Either choice requires an additional generated
selector.
\end{theorem}

\begin{proof}
The maximal-torus conjugacy theorem gives transitivity of inner
automorphisms on the set of maximal tori.  In the nonabelian rank-one groups
that set is the space of unoriented axes in
\(\mathfrak{su}(2)\cong\mathbb R^3\), hence has more than one point.  A torus
fixed by every automorphism would therefore have to equal every conjugate,
which is impossible.  The equivalence with an axis follows because a maximal
toral algebra is a one-dimensional subspace.
\end{proof}

\begin{theorem}[Rank-one flow periodicity]
\label{p2:thm:rank-one-flow-periodicity}
Let \(H\) be an operational-rank-one parent isomorphic to \(U(1)\), \(SU(2)\),
or \(SO(3)\).  For every \(X\in\mathfrak h\), the one-parameter subgroup
\(t\mapsto\exp(tX)\) is constant when \(X=0\) and periodic when \(X\ne0\).
Thus there exists \(T_X>0\) with \(\exp(T_XX)=e\).
Accordingly neither primitive branch contains an intrinsically unwrapped
additive real flow.  The constitutional chain does not select the numerical period or an
orientation of the parameter.
\end{theorem}

\begin{proof}
The claim is immediate for \(U(1)\).  Every one-parameter subgroup of either
nonabelian rank-one global form is contained in a maximal torus, and every
maximal torus is a circle.  A nontrivial continuous homomorphism
\(\mathbb R\to U(1)\) has a nonzero discrete kernel and is periodic.  The
period depends on the normalization of \(X\), while the constitutional chain leaves the common
comparison scale and parameter orientation unfixed.
\end{proof}

\begin{theorem}[Exhaustion of natural scalar polynomial invariants]
\label{p2:thm:quadratic-invariant-exhaustion}
Let \(q(X)=g_{\rm rk}(X,X)\) be the retained rank-one quadratic comparison.
On the abelian
primitive branch,
\[
 \mathbb R[\mathfrak u(1)]^{O(1)}=\mathbb R[q].
\]
If \(H\) is either nonabelian rank-one global form,
\[
 \mathbb R[\mathfrak{su}(2)]^{\operatorname{Ad}(H)}
 =\mathbb R[q].
\]
Thus every polynomial scalar invariant natural under the full constituted
frame or adjoint symmetry is a polynomial in the quadratic norm.

Equivalently, the algebra of bi-invariant scalar differential operators that
also respects the full constituted rank-one symmetry is generated, up to metric
scale and sign convention, by the Laplace--Casimir operator on either branch.
This does not select a Hamiltonian, mass parameter, physical energy, or
spectral cutoff.
\end{theorem}

\begin{proof}
For \(O(1)=\{\pm1\}\), invariant polynomials are precisely the even
polynomials and hence are generated by \(X^2=q(X)\).  On
\(\mathfrak{su}(2)\cong\mathbb R^3\), the adjoint action is the standard
\(SO(3)\) action for both global forms, whose polynomial invariants are generated by the squared
Euclidean norm.  Bi-invariant differential operators correspond to the center
of the complexified universal enveloping algebra.  For \(\mathfrak{su}(2)\)
that center is generated by the quadratic Casimir.  For the abelian
one-dimensional algebra it is the polynomial algebra in one generator, and
the constituted inversion symmetry retains only the even subalgebra, generated by the
square.  Under \(g_{\rm rk}\) these quadratic generators are the
Laplace--Casimir operators.
\end{proof}

\begin{theorem}[Operational resolution-growth dimension]
\label{p2:thm:resolution-growth-dimension}
Let \(N_H(\varepsilon)\) denote the maximal cardinality of an
\(\varepsilon\)-separated subset of a selected primitive parent \(H\), measured
with the distance of \(g_{\rm rk}\).  There exist constants
\(c,C,\varepsilon_0>0\) such that for
\(0<\varepsilon<\varepsilon_0\),
\[
 c\,\varepsilon^{-d}
 \leq N_H(\varepsilon)
 \leq C\,\varepsilon^{-d},
 \qquad
 d=
 \begin{cases}
 1,&H\cong U(1),\\
 3,&H\cong SU(2)\text{ or }SO(3).
 \end{cases}
\]
Consequently the Lie, manifold, Hausdorff, and Minkowski dimensions coincide
and are operationally recoverable from the small-resolution growth exponent.
The constants depend on the unfixed metric scale; the exponent does not.
\end{theorem}

\begin{proof}
A compact smooth \(d\)-dimensional Riemannian manifold is uniformly locally
bi-Lipschitz to Euclidean \(d\)-balls.  Equivalently, for sufficiently small
radii its metric-ball volumes are bounded above and below by positive
multiples of \(\varepsilon^d\).  Disjoint
\(\varepsilon/2\)-balls around a maximal separated set give the upper bound,
while maximality makes the corresponding \(\varepsilon\)-balls cover the
manifold and gives the lower bound.  The branch classification supplies
\(d=1\) and \(d=3\).  Standard dimension comparison for compact Riemannian
manifolds then identifies the displayed exponent with the Hausdorff and
Minkowski dimensions as well as the manifold dimension.
\end{proof}

\begin{theorem}[Sharp closure frontier of the compact-role chain]
\label{p2:thm:gp-independence-frontier}
The following dependency statements are sharp.
\begin{enumerate}[label=(\roman*), leftmargin=2em]
\item FRC+GP2, even with the canonical full parent metric profiles, do not
determine a least class: higher tori and higher-rank compact Lie groups remain
complete candidate-internally exact realizations in the primitive domain.
Non-Lie compact regular-translation parents are not candidates
for a faithful finite exact-operation role.
\item The strong geometrically primary role condition constitutes effectiveness,
connected local generation, and the intrinsic order-swap diagram.  The
closed-parent FRC filter and candidate-internal GP2 exactification realize this
content as a nontrivial connected exact candidate parent, where the diagram is
classified by its exhaustive commutator branch.  The GP2 relation test then applies
Proposition~\ref{p2:prop:admissible-globalization-operator}.  GP3 then selects
the least local parent-role types---\(\mathfrak u(1)\) for the connected
abelian role and \(\mathfrak{su}(2)\) for the connected noncommutative
role---whose
globalizations are respectively the unique compact class \(U(1)\) and the
source-global class \(SU(2)\).  Thus
operational rank one and the \(SU(2)\) global form are derived at allocation;
finite Lie dimension was already guaranteed by FRC exact-operation closure.
\item A center-odd half-integral sector is a faithful diagnostic of the
selected source form, not an additional selector.
\item The chain selects neither the numerical Killing scale nor equality of
the primitive parent comparison with the derived Killing distance.
\end{enumerate}
The abelian/noncommutative branch is read by faithfully realizing and
classifying the intrinsic
order-swap diagram of \(\operatorname{Geo}_R\); it is not a separate primitive
mark and does not freeze any candidate's full composition law.  GP1--GP2
realize/classify the constituted branch content, and GP3 selects its least
parent-role skeleton.
The frontier statement does not assert uniqueness of an unmarked full
structured lift; it asserts only the common endpoint skeleton, which is placed
below every admissible primitive structured candidate by the whole-candidate
comparison of
Theorem~\ref{p2:thm:primitive-endpoint-lift-reachability}.
\end{theorem}

\begin{proof}
Equip a generic flat \(T^2\) and a bi-invariant \(SU(3)\) with their regular
parent actions and generated distance-profile families, and pair each with a
faithful finite-dimensional unitary realization as in
Lemma~\ref{p2:lem:relation-exact-endpoint-models}.  They are FRC+GP2
candidatewise, and their profiles detect the parent action kernels.  Their
full operation tables and profile algebras are retained exactly inside the
respective models, but are not data in \(P_{<R}\) or independent selection
marks.  They are therefore admissible candidates but not GP3-least classes for
the primitive scalar/order-sensitive roles.  If the role geometry instead
requires a second toral direction or an \(SU(3)\)-specific bracket or test, the
corresponding endpoint candidate fails \(\operatorname{Geo}_R\) and the master
endpoint theorem abstains; the richer candidate is not reduced.  The
parent-only solenoid and
weighted-infinite-torus examples of
Proposition~\ref{p2:prop:gp1-clause-sharpness} do not furnish additional
physical exact-operation candidates because they lack faithful finite FRC
reversible realizations.

Theorems~\ref{p2:thm:master-compact-role-allocation} and
\ref{p2:thm:relation-minimal-globalization} prove (ii), while the central
character formula proves (iii).  Every positive rescaling of the negative
Killing form preserves all three principles, so its numerical scale is not
selected.  To test metric equality, fix \(0<\alpha<1\) on the selected
regular \(SU(2)\) parent and give its primitive parent comparison the
snowflaked metric
\[
 d_\alpha:=d_K^\alpha.
\]
This is a compact bi-invariant parent metric with the same relation-exact
operation table and full generated profile family.  Its induced uniform
operation comparison is
\[
 D_\infty(a,b)=d_K(a,b)^\alpha.
\]
It preserves the allocation and paired finite action, but for \(X\ne0\),
\[
 \frac{D_\infty(e,\exp(tX))}{d_K(e,\exp(tX))}
 =d_K(e,\exp(tX))^{\alpha-1}\longrightarrow\infty
 \qquad(t\to0).
\]
Thus the primitive parent comparison is not the derived Killing distance,
even up to a positive constant factor.  This proves (iv).
\end{proof}

\subparagraph{Tier III: relation-exact source globalization
\([{\rm SU2}]\).}

The preceding tier records both local nonabelian global forms for comparison.
Logically, GP2 removed the quotient before the GP3 leastness comparison; the catalogue is
presented later only to display that comparison.  The remainder of Paper~II
uses the resulting relation-exact source form, as recorded in the following
definition.


\subparagraph{Positive allocation and succession exports.}
The derived momentum-parent allocation, Hopf-provenanced scalar-succession
occurrence, and transported orientation theorem are stated in the positive
main chain before the ownership ledger.  Only the indicated dependency pointer is required here.

\begin{corollary}[Finite spectral realization follows after parent allocation]
\label{p2:cor:dependency-complete-spectral-realization}
After the compact parent has been allocated, suppose the frozen role grammar
constitutes a positive-order elliptic Casimir/Laplacian, the associated spectral
projection/subspace construction, and an FRC-b bounded spectral range as
physical scope.  Then the bounded spectral subspace is the finite-dimensional
carrier of
Theorem~\ref{p2:thm:gp1-bounded-spectrum-finite-carrier}.  The preceding GP3 allocation compares the exact/faithful role realizations by
its structural preorder; this spectral corollary neither resolves any remaining
unmarked lift family nor chooses the numerical spectral bound or physical level.  Thus
boundedness is primitive, while finite-dimensionality of any bounded elliptic
spectral subspace is derived.
\end{corollary}

\begin{proof}
Every compact Lie parent allocated above is a closed smooth manifold and hence
falls in the boundaryless domain of
Theorem~\ref{p2:thm:gp1-bounded-spectrum-finite-carrier}.  Applying that theorem to the constituted nonnegative elliptic Casimir or
Laplacian makes each bounded spectral projection finite rank.  The additional
role-grammar hypothesis in the statement is what promotes such a mathematical
subspace to physical carrier scope.  The allocation and compact-resolvent
steps determine neither \(\Lambda\) nor a tower label, so the final
non-selection sentence follows from Corollary~
\ref{p2:cor:gp1-boundedness-no-scale}.
\end{proof}




\section{Conditional finite-discreteness and local-regularity arguments}
\label{p2f:sec:fd-lr-technical}
\subsection{Packing, rectifiability and local regularity}
\label{p2:app:fd-lr-technical-chain}
This section supplies the packing, local-rectifiability, Finsler-rate, and
John-ellipsoid details behind
Theorem~\ref{p2:thm:fd-lr-regularity-summary}.  The chain is conditional and
is not an antecedent of the primitive allocation theorem.

\begin{definition}[Operation paths and realization length]
\label{p2:def:generated-operation-path}
A \emph{completed operation path} is a continuous map
\(\gamma:[0,1]\to(G,D_\infty)\).  Its length is
\[
 L(\gamma):=\sup\Bigl\{\sum_{i=0}^{m-1}
 D_\infty\bigl(\gamma(t_i),\gamma(t_{i+1})\bigr):
 0=t_0<t_1<\dots<t_m=1\Bigr\}.
\]
Since \(D_\infty\) is bi-invariant
(Theorem~\ref{p2:thm:gp1-bounded-isometry-closure}), left and right
translations preserve the path class and its length.  The parameter interval
carries no physical unit; only metric length enters condition (LR).
\end{definition}

\begin{definition}[Local rectifiable reachability]
\label{p2:def:local-rectifiable-reachability}
Condition (LR) asserts that there exist constants \(\kappa>0\) and
\(C_R<\infty\), \emph{uniform over \(g\)}, such that every \(g\in G\) with
\(D_\infty(g,e)<\kappa D_X\) --- the uniform operation comparison and the
comparison diameter of
Definition~\ref{p2:def:intrinsic-operational-boundedness} --- is joined to
\(e\) by a completed operation
path \(\gamma\) satisfying
\[
 L(\gamma)\le C_R D_\infty(g,e).
\]
It is deliberately imposed on the completed closure.  A bound only on paths
inside the nonclosed forward monoid need not survive completion and is not
used.
\end{definition}

\begin{theorem}[Local rectifiable reachability derives local connectedness]
\label{p2:thm:gp1-reachability-local-connectedness}
Under (LR) the compact operation closure \(G\) is locally
path-connected.  Its identity component \(G^0\) is open and of finite
index.
\end{theorem}

\begin{proof}
If \(D_X=0\), then \(X\) has at most one point, as does its completion, so
\(\operatorname{Isom}(\widehat X)=\{e\}\) and the conclusion is immediate.
Assume henceforth that \(D_X>0\).
For \(0<r<C_R\kappa D_X\), let
\[
 U_r:=\{g\in G:\ \text{some completed operation path from }e
 \text{ to }g\text{ has length }<r\}.
\]
Condition (LR), recorded in
Definition~\ref{p2:def:local-rectifiable-reachability}, gives
\(B_{r/C_R}(e)\subseteq U_r\), while length dominates endpoint displacement,
so \(U_r\subseteq B_r(e)\).  Each witnessing path stays in \(U_r\), because
all of its initial subpaths have length \(<r\); concatenation through \(e\)
therefore makes \(U_r\) path-connected.

The set \(U_r\) is open.  Indeed, if \(g\in U_r\), choose a witnessing path
of length \(L<r\) and then \(\eta>0\) with
\(\eta<\kappa D_X\) and \(L+C_R\eta<r\).  Whenever
\(D_\infty(g,h)<\eta\), condition (LR) gives a path of length
at most \(C_R D_\infty(g,h)\) from \(e\) to \(g^{-1}h\); the hypothesis of
(LR) holds for \(g^{-1}h\) because \(D_\infty(g^{-1}h,e)=D_\infty(g,h)\) by
left invariance, which needs only that \(g\) is an isometry.  Left translation
and concatenation with the path to \(g\) put \(h\) in \(U_r\).
As \(r\downarrow0\), the open path-connected sets \(U_r\) form a
neighborhood basis at \(e\); translation gives such a basis at every point.
Thus \(G\) is locally path-connected and in particular locally connected.
A locally connected group has open identity component; compactness makes
its index finite.
\end{proof}

For the optional packing estimates below, write
\[
 D_G:=\operatorname{diam}_{D_\infty}(G),
 \qquad
 D_{\widehat X}:=\operatorname{diam}(\widehat X)=D_X,
 \qquad
 D_Y:=
 \begin{cases}
 D_X,&Y=X\text{ or }Y=\widehat X,\\
 D_G,&Y=G.
 \end{cases}
\]
These diameters are finite in this compact setting; completion preserves the
diameter of the dense subspace \(X\).

\begin{theorem}[Polynomial packing is a sufficient criterion for (FD)]
\label{p2:thm:gp1-capacity-finite-dimension}
Suppose, independently of GP1, that \(X\) and \(G\) obey polynomial packing
bounds at small scales: for some \(\varepsilon_0>0\),
\[
 N_Y(\varepsilon)
 \le C_Y\max\left\{1,(D_Y/\varepsilon)^{n_Y}\right\},
 \qquad
 Y\in\{X,G\},\quad 0<\varepsilon\le\varepsilon_0,
\]
where \(C_Y<\infty\) and \(n_Y\ge0\).
Then the compact completion \(\widehat X\) and the compact operation closure
\(G\) have finite topological dimension:
\(\dim\widehat X\le n_X\) and \(\dim G\le n_G\), where \(n_X\) and \(n_G\)
are the displayed exponents.  In particular the bound on \(G\) implies (FD).
This stronger sufficient package is not a clause of GP1.  Nor is it needed on
the primary line: for realizations of the derived primitive-role skeletons,
finite Lie dimension follows
unconditionally from FRC exact-operation closure
(Theorem~\ref{p2:thm:frc-automatic-compact-lie-closure}), so the packing
hypothesis is consumed only by the optional (FD)+(LR) route and by no result
that the allocation theorem depends on.
\end{theorem}

\begin{proof}
First take \(Y=G\), which obeys the displayed polynomial packing bound with
exponent \(n_G\).  A maximal \(\varepsilon\)-separated set in \(Y\) is an
\(\varepsilon\)-net, so \(Y\) is covered by \(N_Y(\varepsilon)\) balls of
radius \(\varepsilon\).  For \(s>n_Y\), the sum of the \(s\)th powers of
the covering diameters obeys
\[
 N_Y(\varepsilon)\,(2\varepsilon)^s
 \le C_Y2^s(1+D_Y)^{n_Y}\varepsilon^{\,s-n_Y}
 \longrightarrow0
 \qquad(\varepsilon\to0,\ \varepsilon\le1),
\]
so the \(s\)-dimensional Hausdorff measure of \(Y\) vanishes and the
Hausdorff dimension of \(Y\) is at most \(n_Y\).  The group \(G\) is compact metric, hence separable, so Szpilrajn's comparison
theorem bounds covering dimension by Hausdorff dimension~\cite{Szpilrajn1937}.
Since covering dimension is an integer, this gives the sharper
\(\dim G\le\lfloor n_G\rfloor\), and in particular \(\dim G\le n_G\).

If \(E\subseteq\widehat X\) is \(\varepsilon\)-separated, choose for every
\(z\in E\) a point \(x_z\in X\) with
\(d(x_z,z)<\varepsilon/3\).  Then the \(x_z\) are
\(\varepsilon/3\)-separated and distinct, so
\[
 N_{\widehat X}(\varepsilon)
 \le N_X(\varepsilon/3)
 \le C_X\max\left\{1,(3D_X/\varepsilon)^{n_X}\right\}
\]
for all sufficiently small \(\varepsilon>0\).
Since \(\widehat X\) is compact metric and \(D_{\widehat X}=D_X\), the same
Hausdorff-dimension and Szpilrajn argument therefore gives
\(\dim\widehat X\le n_X\).
\end{proof}

\begin{theorem}[(FD)+(LR) regularize an abstract compact parent comparison]
\label{p2:thm:finite-mode-lie-regularity}
Let \(G\) be the compact parent-operation closure of
Theorem~\ref{p2:thm:gp1-bounded-isometry-closure}, considered without using a
faithful finite-carrier representation.  Under (FD) and (LR), \(G\) is a
finite-dimensional compact Lie group.  Consequently there exists a faithful
finite-dimensional unitary representation
\[
 \rho:G\hookrightarrow U(N)
\]
with closed image.  This is a parent-comparison regularity theorem and a
finite matrix witness.  On the physical exact reversible branch of FRC, the
compact Lie conclusion already follows directly from
Theorem~\ref{p2:thm:frc-automatic-compact-lie-closure}; (FD)+(LR) are retained
only where the primitive parent comparison is used to derive first-order path
regularity or compare it with the Killing metric.
\end{theorem}

\begin{proof}
Compactness is Theorem~\ref{p2:thm:gp1-bounded-isometry-closure}; local
connectedness is
Theorem~\ref{p2:thm:gp1-reachability-local-connectedness}; finite
topological dimension is (FD).  A locally compact, locally connected
topological group of finite dimension is a Lie group
\cite[p.~185]{MontgomeryZippin1955}.  Every compact Lie group admits a faithful
finite-dimensional unitary representation, whose image is compact and hence
closed in \(U(N)\).  The final sentence is the independent direct FRC route
already proved in Theorem~\ref{p2:thm:frc-automatic-compact-lie-closure}.
\end{proof}

\begin{corollary}[No small subgroups is a theorem]
\label{p2:cor:nss-derived}
The operation closure \(G\) has an identity neighborhood containing no
nontrivial subgroup.  No-small-subgroups regularity is therefore a
consequence of the derived Lie structure, not an independent hypothesis.
\end{corollary}

\begin{proof}
Fix a norm on \(V:=T_eG\) and \(r>0\) such that \(\exp\) is injective on
the ball \(B_{2r}\subset V\).  Let \(H\subseteq\exp(B_r)\) be a subgroup
and suppose \(h=\exp(W)\in H\) with \(0\ne W\in B_r\).  Choose the least
integer \(k\ge1\) with \((k+1)W\notin B_r\); then \(kW\in B_r\), so
\((k+1)W=kW+W\in B_{2r}\).  Since \(h^{k+1}\in H\subseteq\exp(B_r)\),
there is \(W'\in B_r\) with \(\exp(W')=\exp((k+1)W)\); injectivity on
\(B_{2r}\) forces \((k+1)W=W'\in B_r\), a contradiction.  Hence
\(H=\{e\}\).
\end{proof}

\begin{lemma}[Structure of the first-order rate]
\label{p2:lem:rate-seminorm}
Let the generated operation closure be the compact Lie group supplied by
Theorem~\ref{p2:thm:finite-mode-lie-regularity}, acting by generated
isometries, so that the uniform comparison \(D_\infty\) is bi-invariant,
and identify continuous one-parameter flows with
\(t\mapsto\exp(tW)\), \(W\in\mathfrak g\).  For
\(f_W(t):=D_\infty(\exp tW,e)\) define
\[
r(W):=\limsup_{t\to0}\frac{f_W(t)}{|t|}\in[0,\infty].
\]
Then \(f_W\) is even and subadditive, the limit superior is a limit and
equals \(\sup_{t>0}f_W(t)/t\); \(r\) is absolutely homogeneous,
\(r(\lambda W)=\lvert\lambda\rvert\,r(W)\) for all \(\lambda\in\mathbb R\)
--- equivalently positively homogeneous and symmetric, since \(f_W\) is even
--- and is subadditive, Ad-invariant, and strictly positive on \(W\ne0\), so
that \(r\) is a genuine norm on its finiteness locus; and its
finiteness locus
\(S:=\{W\in\mathfrak g:r(W)<\infty\}\) is an ideal of \(\mathfrak g\).
\end{lemma}

\begin{proof}
Left invariance alone gives
\(f_W(-t)=D_\infty(\exp(tW)^{-1},e)=D_\infty(e,\exp(tW))=f_W(t)\) and, by
the triangle inequality along left translates,
\(f_W(t+s)\le f_W(t)+f_W(s)\); bi-invariance is used below only for
Ad-invariance.  Subadditivity gives
\(f_W(s)\le n\,f_W(s/n)\), hence \(f_W(s)/s\le f_W(s/n)/(s/n)\) along
\(s/n\to0\), so \(\limsup_{t\to0^+}f_W(t)/t=\sup_{t>0}f_W(t)/t\); and
since \(f_W\) is continuous at \(0\) with \(f_W(0)=0\), the comparison
being compatible and \(\exp\) continuous, writing \(s=mt+\rho\) with
\(m=\lfloor s/t\rfloor\ge1\) and \(0\le\rho<t\) gives
\(f_W(s)/s\le f_W(t)/t+f_W(\rho)/s\), and \(t\to0^+\) yields
\(f_W(s)/s\le\liminf_{t\to0^+}f_W(t)/t\): the limit exists and equals
the supremum.  Homogeneity is
reparametrization.  For subadditivity in \(W\), put
\[
 A_n:=\bigl(\exp(tW/n)\exp(tV/n)\bigr)^n.
\]
Then \(A_n\to\exp(t(W+V))\) by the Lie product formula, which holds
automatically here because the operation group is a finite-dimensional Lie
group.  Continuity of the metric and the bi-invariant triangle inequality give
\[
\begin{aligned}
 f_{W+V}(t)
 &=\lim_{n\to\infty}D_\infty(A_n,e)\\
 &\le\liminf_{n\to\infty}
 n\bigl[f_W(t/n)+f_V(t/n)\bigr]\\
 &\le |t|\bigl(r(W)+r(V)\bigr),
\end{aligned}
\]
where the last step uses \(f_W(s)\le |s|r(W)\) and the corresponding bound
for \(V\).  Dividing by \(|t|\) and taking \(t\to0\) gives
\(r(W+V)\le r(W)+r(V)\).  Bi-invariance gives
\(D_\infty(g\exp(tW)g^{-1},e)=f_W(t)\) and
\(g\exp(tW)g^{-1}=\exp(t\operatorname{Ad}_gW)\), so \(r\) is
Ad-invariant; a subspace that is
invariant under \(\operatorname{Ad}\) of the identity component is an
ideal, and \(S\) is a subspace by homogeneity and subadditivity.
Positivity: for \(W\ne0\) the set \(\{t\in\mathbb R:\exp(tW)=e\}\) is the
kernel of a continuous homomorphism \(\mathbb R\to G\), hence a closed
subgroup of \(\mathbb R\), hence \(\{0\}\), \(T\mathbb Z\) for some
\(T>0\), or all of \(\mathbb R\); the last case would force
\(W=\tfrac{d}{dt}\big|_{0}\exp(tW)=0\).  So \(\exp(t_0W)\ne e\) for some
\(t_0>0\), whence \(f_W(t_0)>0\) and \(r(W)\ge f_W(t_0)/t_0>0\).  No
no-small-subgroups input is required.
\end{proof}

\begin{theorem}[Untimed first-order group rate follows from (FD)+(LR)]
\label{p2:thm:gp1-propagation-derived}
Under GP1, (FD), and (LR), every continuous one-parameter subgroup
\(\varphi:\mathbb R\to G\) has finite first-order displacement rate:
\[
\limsup_{t\to0}\frac{D_\infty(\varphi_t,e)}{|t|}<\infty .
\]
Moreover \(r\) is an Ad-invariant norm on \(\mathfrak g\), and there is a
bi-invariant Finsler metric \(D_F\) on the identity component with
\[
D_\infty\le D_F\le C_R\,D_\infty
\qquad\text{near the identity,}
\]
whose norm is exactly \(r\).  This local Lipschitz-rate
bound for the primitive, freely reparametrized group path is therefore a
consequence of Lie regularity plus (LR), not of carrier-metric total
boundedness alone.  It is not, by itself, the universal propagation ceiling
of Theorem~\ref{p2:thm:gp1-operational-speed-ceiling}: no elapsed scalar cover
has been attached here, and rescaling the auxiliary flow parameter rescales
this rate.  A comparison whose
generated paths are not rectifiable, such as a snowflaked comparison,
violates (LR).
\end{theorem}

\begin{proof}
Let \(G^0\) be the identity component of the compact Lie closure, which
contains every flow value, and let \(D_\ell\) denote the intrinsic length
metric of \(D_\infty\) along all completed operation paths in \(G^0\), with
length as in Definition~\ref{p2:def:generated-operation-path}.

First, \(D_\ell\) is finite, bi-invariant, intrinsic, and compatible.
Condition (LR), recorded in
Definition~\ref{p2:def:local-rectifiable-reachability}, gives
\(D_\ell(e,g)\le C_R\,D_\infty(e,g)\) whenever
\(D_\infty(e,g)<\kappa D_X\), and trivially
\(D_\infty\le D_\ell\); so the two metrics are bi-Lipschitz near the
identity.  For general \(g\in G^0\), connectedness of \(G^0\) in the
compatible topology supplies a finite chain \(e=g_0,\dots,g_N=g\) with
\(D_\infty(g_i,g_{i+1})<\kappa D_X\): the set of points joined to
\(e\) by such a chain and its complement are each unions of open
\(\kappa D_X\)-balls, so the chainable set is clopen and, \(G^0\) being
connected in the topology induced by the compatible metric, exhausts
\(G^0\).  Left invariance transports the local estimate to each step, so
\(D_\ell\) is finite on \(G^0\).
Bi-invariance of \(D_\ell\) is inherited from \(D_\infty\), because left
and right translations preserve completed paths and their lengths.
Length metrics are intrinsic, and the local bi-Lipschitz comparison makes
\(D_\ell\) induce the manifold topology.

By the classification of invariant intrinsic metrics on Lie groups
\cite{Berestovskii1988,agrachev2019comprehensive}, an invariant intrinsic
compatible metric on a Lie group is the Finsler--Carnot--Caratheodory
metric of an invariant bracket-generating distribution equipped with an
invariant norm.  Here the metric is bi-invariant, so every conjugation is an isometry
fixing \(e\), and the horizontal distribution with its norm is
characterized metrically as the directions and speeds of curves of
finite metric dilation; both are therefore preserved by every isometry
fixing \(e\), so the distribution at the identity is an
\(\operatorname{Ad}\)-invariant subspace of \(\mathfrak g\), hence an
ideal; a bracket-generating ideal is all of
\(\mathfrak g\).  Thus \(D_\ell=D_F\) is a genuine bi-invariant Finsler
metric with Ad-invariant norm \(F_0\), and no Carnot--Caratheodory
degeneration can occur.

For any flow \(\varphi_t=\exp(tW)\), the curve
\(s\mapsto\exp(sW)\) has constant Finsler speed, so
\(D_\infty(\varphi_t,e)\le D_F(\varphi_t,e)\le|t|\,F_0(W)\), giving
\(r(W)\le F_0(W)<\infty\).

For the reverse inequality, fix \(t>0\).  On every partition
\(0=s_0<\cdots<s_m=t\), bi-invariance and the one-parameter subgroup law
give
\[
 \sum_{i=0}^{m-1}
 D_\infty\bigl(\exp(s_iW),\exp(s_{i+1}W)\bigr)
 =
 \sum_{i=0}^{m-1}f_W(s_{i+1}-s_i)
 \le t\,r(W),
\]
using \(r(W)=\sup_{u>0}f_W(u)/u\) from
Lemma~\ref{p2:lem:rate-seminorm}.  Equal partitions with mesh \(t/n\)
make the sum \(n f_W(t/n)\to t\,r(W)\), so the
\(D_\infty\)-length of the exponential segment is exactly \(t\,r(W)\).
Since \(D_F=D_\ell\) is the infimum of such path lengths,
\[
 D_F(\exp(tW),e)\le t\,r(W).
\]
In a Finsler metric, distance and norm are first-order tangent at the
identity:
\[
 \lim_{t\to0}\frac{d_F(\exp tW,e)}{t}=F_0(W).
\]
So division by \(t\) and \(t\downarrow0\) give
\(F_0(W)\le r(W)\).  Hence \(F_0=r\) exactly, and
Lemma~\ref{p2:lem:rate-seminorm} supplies all asserted norm properties.
The snowflake statement is the observation,
already recorded in the sharpness discussion, that a snowflaked
comparison admits no rectifiable generated paths, so (LR) fails.
\end{proof}

\begin{lemma}[Uniform local linear operation-rate bound]
\label{p2:lem:second-kind-uniform-propagation}
For any auxiliary norm \(\lVert\cdot\rVert\) on \(\mathfrak g\), there are an
identity neighborhood \(\mathcal W\subset\mathfrak g\) and \(C<\infty\) such
that
\[
 D_\infty(\exp W,e)\le C\lVert W\rVert\qquad(W\in\mathcal W).
\]
The norm is proof-internal and selects no basis or physical scale.
\end{lemma}

\begin{proof}
Theorem~\ref{p2:thm:gp1-propagation-derived} gives
\(D_\infty\le D_F\) near the identity.  In finite dimension the Finsler norm
of \(D_F\) is bounded above by a constant multiple of any auxiliary norm, and
the exponential curve has length at most that norm of \(W\).
\end{proof}

\begin{corollary}[The first-order displacement norm]
\label{p2:cor:gp1-first-order-norm-derived}
For \(X\in\mathfrak g\),
\[
 F(X):=\lim_{t\to0^+}\frac{D_\infty(\exp(tX),e)}{t}
      =\sup_{t>0}\frac{D_\infty(\exp(tX),e)}{t}
\]
exists and is a finite, symmetric, positive-definite, subadditive,
Ad-invariant norm.  It is intrinsic to the generated action and depends on
(FD)+(LR); it is not a consequence of GP1 alone.
\end{corollary}

\begin{proof}
Existence, the supremum formula, homogeneity, subadditivity, positivity, and
Ad-invariance are Lemma~\ref{p2:lem:rate-seminorm}.  Finiteness is
Theorem~\ref{p2:thm:gp1-propagation-derived}.  Thus \(F=r\).
\end{proof}

\begin{corollary}[Sharp small-scale capacity exponent]
\label{p2:cor:gp1-capacity-saturation}
Let \(q:=\dim G\).  The exponent is the topological dimension precisely
because the propagation theorem has already supplied the Lie structure on
which that dimension is defined.  If \(q>0\), there are constants
\(c,C,\varepsilon_0>0\) such that
\[
 c\Bigl(\frac{D_G}{\varepsilon}\Bigr)^q
 \le N_G(\varepsilon)
 \le C\Bigl(\frac{D_G}{\varepsilon}\Bigr)^q
 \qquad(0<\varepsilon<\varepsilon_0).
\]
If \(q=0\), then \(G\) is finite and \(N_G(\varepsilon)=|G|\) for all
sufficiently small \(\varepsilon\).  Hence the infimum of the exponents in
polynomial upper bounds for the completed operation group is \(q\).  The
particular exponent in an optional polynomial-packing hypothesis need not be optimal, and this
corollary does not claim that polynomial capacity is the logically weakest
possible hypothesis for Lie regularity.
\end{corollary}

\begin{proof}
Assume first that \(q>0\).  Theorem~\ref{p2:thm:gp1-propagation-derived}
makes \(D_\infty\) locally bi-Lipschitz to the bi-invariant Finsler distance
\(D_F\) on \(G^0\).  Finsler volume, equivalently Haar measure on the compact
group \(G^0\), assigns to every sufficiently small \(D_F\)-ball a measure
between positive constant multiples of \(\varepsilon^q\).  The same is
therefore true for small \(D_\infty\)-balls.  Disjoint
\(\varepsilon/2\)-balls centred at an \(\varepsilon\)-separated set give the
upper packing bound; the \(\varepsilon\)-balls centred at a maximal
\(\varepsilon\)-separated set cover and give the lower bound.  The identity
component is open and has finite index, so for \(\varepsilon\) below the
minimum distance between its finitely many closed disjoint cosets --- positive
by compactness --- every \(\varepsilon\)-separated set splits across cosets and
the finitely many translates change only the constants.  Replacing \(\varepsilon^{-q}\) by
\((D_G/\varepsilon)^q\) merely normalizes the scale.

If \(q=0\), the compact Lie group \(G\) is discrete and therefore finite.
Its positive minimum separation gives the asserted eventual packing number.
\end{proof}

\begin{theorem}[Operational norm derives a canonical quadratic comparison]
\label{p2:thm:operational-john-metric}
Let \(F\) be the derived first-order displacement norm of
Corollary~\ref{p2:cor:gp1-first-order-norm-derived} on \(V=T_eG\).  Then \(F\) is \(\operatorname{Ad}(G)\)-invariant.  The unique
maximal-volume ellipsoid contained in its unit ball is functorial under every
linear \(F\)-isometry and defines a positive quadratic form \(g_F\), unique
without a chosen basis or volume normalization.  In particular \(g_F\) is
\(\operatorname{Ad}(G)\)-invariant and translates to a bi-invariant
Riemannian proof metric on \(G\).  This is the quadratic tangent comparison
retained alongside, and not in place of, the operational norm \(F\).
\end{theorem}

\begin{proof}
Bi-invariance of \(D_\infty\) gives
\[
 D_\infty(\exp(t\,\operatorname{Ad}_gX),e)
 =D_\infty(g\exp(tX)g^{-1},e)
 =D_\infty(\exp(tX),e),
\]
so taking first-order rates proves Ad-invariance.  The norm unit ball is a
compact convex body with nonempty interior.  John's ellipsoid
theorem~\cite{John1948} gives a unique maximal-volume inscribed ellipsoid;
central symmetry is not needed for that, but it is what centres the ellipsoid
at the origin and so makes it the unit ball of a quadratic form.  Volume
may be computed with any Lebesgue density on \(V\): two such densities differ
by a positive scalar, so they have the same maximizer.  Linear
equivalences carry maximizers to maximizers, so uniqueness makes the
construction functorial and Ad-invariant.  The ellipsoid is the unit ball of a
unique positive quadratic form.  Translation and Ad-invariance make its
Riemannian metric bi-invariant.
\end{proof}


\section{Finite-test normalization and typed descent}
\label{p2f:sec:finite-test-descent}
\subsection{Finite readout normalization and exact descent checks}
\label{p2a:app:repair-normalization-descent}
The main text states the canonical closing word and the domain conditions
needed by the positive chain. This technical section classifies the broader natural
functional family and proves supplementary quotient, compression and
variation identities. All operations used as physical descendants retain
their original constructor licenses; host calculations do not create them.

\begin{lemma}[Natural functional coefficients and licensed sum constraints]
\label{p2a:lem:repair-natural-functional-coefficients}
Let $\lambda_H:\operatorname{End}(H)\to\mathbb C$ be complex-linear and
invariant under the full unitary group of the unmarked carrier. For nonzero $H$,
\[
 \lambda_H=c_H\operatorname{Tr}_H.
\]
Unmarked unitary identifications equate the corresponding coefficients.
Compatibility with a licensed sum $K=H\oplus J$, for independently arbitrary
blocks, implies $c_K=c_H=c_J$ when the summands are nonzero. Normalization on
$\mathbb C$ therefore fixes $c_H=1$ on every coefficient class connected to
that unit by these licensed identifications and sum constraints. No conclusion
about an unconnected coefficient is supplied by these hypotheses alone.
\end{lemma}
\begin{proof}
The trace pairing represents $\lambda_H(T)=\operatorname{Tr}(B_HT)$ uniquely.
Unitary invariance forces $B_H$ into the scalar commutant, so $B_H=c_H1_H$.
In the block-additivity equation, set one input block to zero and vary the
trace of the other. This gives the two coefficient equalities. Propagation
from the normalized unit proves the last assertion. In the absence of the
required sum links, $\lambda_H=\operatorname{Tr}_H/\dim H$ is a counterexample
to unit coefficient on every carrier, even though it is tensor-multiplicative.
\end{proof}

\begin{lemma}[Operator graph descent, restriction, and compression]
\label{p2a:lem:repair-operator-graph-descent}
For a surjective linear map $q:H\to D$ and $A\in\operatorname{End}(H)$,
the projected graph $\{(qv,qAv):v\in H\}$ defines a unique total linear map
$\bar A$ on $D$ if and only if $A\ker q\subseteq\ker q$. Then
$qA=\bar A q$.
For an orthogonal projector $P$ with $Q=I-P$, restriction of $A$ to
$E=\operatorname{ran}P$ requires $QAP=0$, whereas descent through
$P:H\to E$ requires $PAQ=0$. For self-adjoint $A$ these are equivalent to
$[A,P]=0$. Without them, $PAP|_E$ is only the compressed operator; its
physical use requires its own licensed word.
\end{lemma}
\begin{proof}
The value $\bar A(qv)=qAv$ is representative-independent exactly when
$qAv=0$ for every $v\in\ker q$. Surjectivity gives totality, and the
formula gives linearity and uniqueness. The two projector conditions follow
by splitting vectors into $Pv+Qv$. They are adjoints when $A=A^*$.
\end{proof}

\begin{lemma}[Restricted functional calculus and the compression defect]
\label{p2a:lem:repair-restricted-trace}
Let $A=A^*$ on a finite Hilbert space $H$ and let $P=P^*=P^2$.
Put $E=\operatorname{ran}P$, $Q=I-P$, and $B=PAP|_E$.
If $[A,P]=0$, then for every function $F$ on the finite spectrum of $A$,
\[
 F(B)=PF(A)P|_E,\qquad
 \operatorname{Tr}_E F(B)=\operatorname{Tr}_H(PF(A)).
\]
Without the commutation premise,
\[
 PA^2P|_E-B^2=(QA|_E)^*(QA|_E)\ge0.
\]
Its trace is zero if and only if $QAP=0$, equivalently $[A,P]=0$.
Thus compression cannot generally replace restriction before a spectral
process is formed and closed. The trace on $E$ remains its canonical
duality trace; the expression on $H$ retains the insertion $P$.
\end{lemma}
\begin{proof}
Invariance gives $A^ni=iB^n$ by induction for the inclusion $i:E\to H$.
Polynomial interpolation on the finite spectrum gives the functional
calculus identity; rectangular trace cyclicity gives the trace identity.
Insert $I=P+Q$ between the two factors of $A$ to obtain the defect.
Self-adjointness writes it as a Gram matrix. Its trace is the squared
Hilbert--Schmidt norm of $QA|_E$, proving the final equivalence.
\end{proof}

\begin{lemma}[Nonlinear Hessian and stationary Schur transport]
\label{p2a:lem:repair-hessian-schur-transport}
On a specified real $C^2$ coordinate domain, for $f:M\to N$ and $S:N\to\mathbb R$,
\[
 D^2(S\circ f)(u,v)=D^2S(df\,u,df\,v)+dS(D^2f(u,v)).
\]
The last term vanishes at a full target critical point or when $f$ is affine.
If an invariant infinitesimal field is $R(D)=AD$, then
\[
 D^2S_D(v,AD)+dS_D(Av)=0;
\]
thus the symmetry tangent is in the Hessian radical at full stationarity.
For $S(x,y)$, an interior stationary branch with invertible $H_{yy}$ has
\[
 D^2S_{\rm eff}=H_{xx}-H_{xy}H_{yy}^{-1}H_{yx}.
\]
For a symmetric singular quadratic interior block $C$, elimination for every
retained direction requires $B\ker C=0$ for the cross block $B$. The unique
stationary class then lies in the nondegenerate quotient by $\ker C$;
using that quotient physically requires a claim-complete licensed quotient.
\end{lemma}
\begin{proof}
Differentiate $d(S\circ f)=dS\circ df$ and $dS_D(AD)=0$ to get the first
two identities. Differentiate $\partial_yS(x,f(x))=0$ to obtain
$df=-H_{yy}^{-1}H_{yx}$, then differentiate
$dS_{\rm eff}=\partial_xS$ along the branch. For a quadratic form,
completion of the square gives the same Schur expression, with a full null
vector lifting from $x\in\ker(H_{xx}-BC^{-1}B^t)$ as
$(x,-C^{-1}B^tx)$. In the singular case $-B^tx$ must lie in
$\operatorname{im}C=(\ker C)^\circ$, equivalently $B\ker C=0$.
The induced form on the quotient is nondegenerate and its stationary value
is independent of representatives. Rank changes, boundary limits and
physical occurrence remain separate conditions.
\end{proof}




\section{Optional momentum towers, finite supports, and negative selectors}
\label{p2ftech:sec:momentum}

\subsection{Auxiliary pair presentation of the existing geometric coordinates}
\label{p2b:app:v2174-geometric-pair-presentation}
This section evaluates a presentation of the already generated carrier of
Theorem~\ref{p2:thm:momentum-host-H1}. It assumes the canonical smooth
coordinate class of Paper I-B,
Theorem~\ref{p1:thm:intrinsic-tangential-coordinate-class}, and adds no
physical Bose role, occupation boundary, field operator or state.

\begin{proposition}[Exact phase-qualified pair presentation]
\label{p2b:prop:v2174-geometric-pair-presentation}
Write $S=S_L\oplus S_R$ and present the existing geometric carrier as
$\mathcal F_{\rm aux}=\bigoplus_{n=0}^2\operatorname{Sym}^nS$.
Let $b_{Lr},b_{Rc}$ be the four auxiliary compressed symmetric ladder maps,
and $P_1,P_2$ their degree projectors. Use the Hermitian Clifford matrices
\[
 C_j=\begin{pmatrix}0&-i\sigma_j\\i\sigma_j&0\end{pmatrix}
 \ (j=1,2,3),\qquad
 C_4=\begin{pmatrix}0&I_2\\I_2&0\end{pmatrix},\qquad
 \epsilon=\begin{pmatrix}0&1\\-1&0\end{pmatrix}.
\]
Set $A_j=-i\sigma_j$, $A_4=I_2$ and
\[
 B_a^\dagger=\frac1{\sqrt2}\sum_{r,c}(A_a\epsilon)_{rc}
                         b_{Lr}^\dagger b_{Rc}^\dagger.
\]
On the normalized source basis $e_0,u_a,w_{ab}$ ($a<b$), and its spinor
block, a phase-compatible unitary $T$ is specified by
\[
 Te_0=|0_{\rm aux}\rangle,\qquad
 Tu_a=B_a^\dagger|0_{\rm aux}\rangle,\qquad
 Tw_{ab}=\tfrac12d\Gamma(C_a)B_b^\dagger|0_{\rm aux}\rangle,
\]
and by sending the ordered $S_L\oplus S_R$ basis to degree one. All four
native coordinates, with their actual source coefficients, satisfy
\begin{equation}
 TX_a^DT^\dagger=\frac{R_p}{2}(B_a^\dagger+B_a)
       +\frac{\sqrt3R_p}{4}P_1d\Gamma(C_a)P_1
       +\frac{\sqrt2R_p}{6}P_2d\Gamma(C_a)P_2.
 \label{p2b:eq:v2174-geometric-pair-presentation}
\end{equation}
\end{proposition}
\begin{proof}
Use $\{C_a,C_b\}=2\delta_{ab}I_4$ and
$\operatorname{tr}(A_a^\dagger A_b)/2=\delta_{ab}$.
The four pair vectors are orthonormal. The six displayed $Tw_{ab}$ are
orthonormal, perpendicular to the pair vectors, and antisymmetric under
$a\leftrightarrow b$; this follows by substituting the four $2\times2$
matrices in the normalized symmetric basis. Together with the vacuum and
degree-one vectors they form all fifteen columns, so $T$ is unitary.
On the forms, the native phase presentation has
$X_a^De_0=(R_p/2)u_a$ and
$X_a^Du_b=(R_p/2)\delta_{ab}e_0+(\sqrt2R_p/3)w_{ab}$,
with reverse edges given by adjoints. The first term of the formula gives
exactly the scalar/vector edges. On degree one $d\Gamma(C_a)=C_a$,
giving the protected spinor block. On degree two
$d\Gamma(C_a)Tu_b=2Tw_{ab}$, giving the stated vector/bivector edges and
their adjoints. These are every nonzero native block. The exact arithmetic
companion checks all entries and all fifteen columns.
\end{proof}
The formula changes auxiliary degree by zero or two. It is an operator
presentation, not an identification with a physical linear Bose coordinate.
For an already licensed action factor, its appropriate induced unitary
transports the entire action, every product, projector, contact and its full
Gibbs support by conjugation, as in
Theorem~\ref{p3a:thm:v204-action-law-transport}. It does not instantiate a
missing complete joint action member.

\subsubsection{The scope of the proposed linear-occupation identification}
\begin{proposition}[The native geometric grading fails the linear port identity]
\label{p2b:prop:v2174-geometric-linear-port-defect}
Let $N_{\rm geo}$ be zero on the scalar, one on the spinors and two on
the vector and bivectors. For each native coordinate,
\[
 \bigl\|[N_{\rm geo},[N_{\rm geo},X_a^D]]-X_a^D\bigr\|_{\rm HS}^2
       =\frac{79R_p^2}{12}.
\]
Thus at $R_p\ne0$ this marked coordinate/degree pair is not a transported
linear standard Bose port/number pair.
\end{proposition}
\begin{proof}
On a degree-$r$ to degree-$s$ block the double commutator multiplies by
$(s-r)^2$. The scalar/vector gap is two, whereas the vector/bivector and
spinor gaps are zero. Their squared defect contributions per coordinate
are respectively
$9\cdot2(R_p/2)^2$, $2\cdot3(\sqrt2R_p/3)^2$ and
$4(\sqrt3R_p/4)^2$, whose sum is $79R_p^2/12$.
Every standard linear port obeys $[N_B,[N_B,Q]]=Q$, including its
compressed top edge. That identity is preserved under a marked unitary
transport, proving the stated exclusion.
\end{proof}

\begin{proposition}[Restricted full-sublevel classification on this carrier]
\label{p2b:prop:v2174-fifteen-dimensional-port-exclusion}
For $R_p\ne0$, there is no simultaneous identification of the entire native carrier with
a standard complete Bose initial sublevel having a unique degree-zero unit,
a number operator commuting with the inherited $\operatorname{Spin}(4)$
action, and every native $X_a^D$ a linear coordinate.
\end{proposition}
\begin{proof}
Multiplicity one makes number scalar on each of the six irreducibles.
The scalar is degree zero. Its nonzero vector edge forces the vector to
degree one. The two vector/bivector edges force the bivectors to degree
zero or two; the unique unit excludes zero. The nonzero spinor edge forces
the two nontrivial spinor degrees to differ by one, with neither zero.
The dimension equation $\binom{N+d}{d}=15$ has only
$(d,N)=(1,14),(2,4),(4,2),(14,1)$: for $d,N\ge2$, symmetry reduces
to $d=2$, yielding $N=4$, since $\binom63=20>15$.
The first two possibilities cannot hold the degree-one vector of dimension
four; the last lacks the required degree-two bivectors. At $(4,2)$ the
spinor degrees are one and two, forcing degree one to have dimension six,
contrary to $d=4$.
\end{proof}
These exclusions have precisely the marked full-sublevel hypotheses just
stated. They do not exclude every generated compression, alternative first
allocation or constrained quantization. The existing coordinate algebra
$M_{11}(\mathbb C)\oplus M_4(\mathbb C)$ of
Corollary~\ref{p2c:cor:present-coordinate-envelope-condition} remains
unchanged. No geometric unit is assigned to the interacting ground.




\subsection{Algebraic branch non-implication tests}
\label{p2b:app:algebraic-branch-diagnostics}
\begin{theorem}[Kinematic Lie algebra classification]
\label{p2:thm:compact_selection}
Let $\mathfrak{g}$ be a six-dimensional real Lie algebra with basis
$\{J_i,P_i\}_{i=1}^3$.  Suppose its abstract brackets are represented on a
quantum operator domain by a faithful linear map
$J_i\mapsto\hat J_i$, $P_i\mapsto\hat P_i$, so in particular the six
displayed operators are linearly independent, and their operator commutators
satisfy
\[
[\hat J_i,\hat J_j]_{\rm op}= i\hbar\,\epsilon_{ijk}\hat J_k,\qquad
[\hat J_i,\hat P_j]_{\rm op}= i\hbar\,\epsilon_{ijk}\hat P_k,\qquad
[\hat P_i,\hat P_j]_{\rm op}= i\hbar\,\varkappa\,\epsilon_{ijk}\hat J_k.
\]
Equivalently, the underlying real Lie bracket is $(i\hbar)^{-1}[\ ,\ ]_{\rm op}$ in this operator normalization. Then:
\begin{enumerate}
\item If $\varkappa=0$, then $\mathfrak{g}\cong\mathfrak{iso}(3)$.
\item If $\varkappa>0$, then $\mathfrak{g}\cong\mathfrak{so}(4)\cong \mathfrak{su}(2)\oplus\mathfrak{su}(2)$.
\item If $\varkappa<0$, then $\mathfrak{g}\cong\mathfrak{so}(3,1)$.
\end{enumerate}
\end{theorem}

\begin{proof}
Throughout the proof the displayed brackets are operator commutators, or equivalently the real Lie bracket after division by $i\hbar$.  If $\varkappa=0$, then $[\hat P_i,\hat P_j]_{\rm op}=0$, so the $\hat P_i$ span an abelian ideal transforming as a vector under $\mathfrak{so}(3)$; this is precisely $\mathfrak{iso}(3)$.
For $\varkappa\neq 0$, write $\varkappa=s\,\varkappa_0^2$ with $\varkappa_0>0$ and $s\in\{+1,-1\}$. Define $\hat K_i:=\hat P_i/\varkappa_0$. Then
\[
[\hat J_i,\hat K_j]_{\rm op}=i\hbar\,\epsilon_{ijk}\hat K_k,\qquad
[\hat K_i,\hat K_j]_{\rm op}= i\hbar\,s\,\epsilon_{ijk}\hat J_k.
\]
If $s=+1$, define $\hat L_i^\pm:=\frac12(\hat J_i\pm \hat K_i)$. One checks that the two sets $\hat L_i^\pm$ generate commuting $\mathfrak{su}(2)$ algebras, hence $\mathfrak{g}\cong\mathfrak{so}(4)$. For $s=-1$, the real brackets in the basis $(\hat J_i,\hat K_i)$ are
\[
[\hat J_i,\hat J_j]=\epsilon_{ijk}\hat J_k,\qquad
[\hat J_i,\hat K_j]=\epsilon_{ijk}\hat K_k,\qquad
[\hat K_i,\hat K_j]=-\epsilon_{ijk}\hat J_k,
\]
after dividing the operator commutators by $i\hbar$.  These are exactly the rotation--boost commutation relations of $\mathfrak{so}(3,1)$, with $\hat J_i$ rotations and $\hat K_i$ boosts \cite{Hall2015lie}.  The three cases are also separated by a basis-independent invariant: the Killing form has signature $(0,6)$ for $\varkappa>0$, $(3,3)$ for $\varkappa<0$, and is degenerate for $\varkappa=0$, so no two of them are isomorphic.
\end{proof}

\begin{remark}[Generality of the displayed brackets, and hypothesis usage]
\label{p2:rem:bracket-generality}
The displayed \([\hat P_i,\hat P_j]\) is not the most general
\(\mathfrak{so}(3)\)-covariant bracket.  Imposing only the first two relations
and the Jacobi identity leaves a two-parameter family
\([\hat P_i,\hat P_j]_{\rm op}=i\hbar\,\epsilon_{ijk}(a\hat P_k+b\hat J_k)\).
The displayed form is reached without loss of generality by the shift
\(\hat P_i\mapsto\hat P_i-\tfrac a2\hat J_i\), under which
\((a,b)\mapsto(0,\,b+a^2/4)\), so \(\varkappa=b+a^2/4\).  The shift is a
redefinition of the momentum generators that mixes them with the rotation
generators, and is recorded here rather than performed silently.
The hypothesis that \(\mathfrak g\) contain \(\mathfrak{so}(3)\) is exhibited
by the first displayed relation.  The basis and faithfulness hypotheses in
Theorem~\ref{p2:thm:compact_selection} rule out a collapse of the abstract
six-dimensional algebra in its operator image; without them the displayed
commutators classify only that image.
\end{remark}

\begin{remark}[Interpretation of the compact constitution]
\label{p2b:rem:compact-branch-legacy}
Theorem~\ref{p2:thm:compact_selection} is an algebraic trichotomy. The constitution clauses above force the compact spherical carrier, so in the remainder we set
\[
\varkappa=\frac{1}{\Rmom^2}>0,
\qquad
\mathfrak{g}\cong \mathfrak{so}(4)\cong \mathfrak{su}(2)_L\oplus\mathfrak{su}(2)_R.
\]
Here \(\mathfrak g\) is the derived Killing algebra of the one round parent.
The carrier group itself remains \(SU(2)\), with Lie algebra
\(\mathfrak{su}(2)\).
\end{remark}


\subsection{Conditional Clifford completion of the coordinate algebra}
\label{p2b:app:conditional-clifford-package}
This section contains the optional coordinate-completion specialization
advertised in the main text.  It is downstream of the unconditional
GP1--GP3 momentum-carrier theorem and neither selects the carrier nor promotes
\(\mathsf C_{\rm PI}\) to a framework-generated completion.

\subsubsection{Derived level-one structure}
\label{p2:subsec:m1-distinguished}

Theorem~\ref{p2:thm:momentum-host-H1} fixes the generated six-type
momentum support \(\mathcal G^{\rm car}_{\Sigma_p}\cong\mathcal H_1\),
and Corollary~\ref{p2:thm:m1-component} gives \(m_p=1\).
The stronger assertion that \emph{every} GP3-irreducible full
momentum realization has the one-copy carrier \(\mathcal H_p\cong\mathcal H_1\)
requires the additional whole-structured comparison hypotheses of
Theorem~\ref{p2b:thm:momentum-carrier-rigidity}.  The
Paper-I-A representative \(\mathsf C_{\rm PI}\) is needed only for the protected coordinate
Clifford vector below.  We now separate the Hilbert-space count fixed by host
closure from its additional real operator algebra.

\paragraph{Protected even Clifford algebra and the cross-category count}
\label{p2:subsec:hilbert_recon}

\begin{remark}[Residual coordinate-completion selector family]
\label{p2b:rem:coordinate-completion-selector-family}
After quotienting the sector-rephasing presentation freedom already removed in
Paper~I-A, the completion family over the derived carrier remains
\[
 \boxed{
 \mathfrak M_{\rm coord}(\mathcal H_1)/\!\sim
 \ \cong\ \mathbb R_{\ge0}^{3},\qquad
 [\mathsf C]\longleftrightarrow
 (r_{1,0},r_{0,1},\varrho_{\rm sp}).}
\]
No defect, relation, GP3 reduction, factor exchange, or value of \(m_p\)
generated in the momentum-constitution chain selects a point of this cone.
The representative \(\mathsf C_{\rm PI}\) used below is therefore an explicit
conditional completion specialization, not ``the'' coordinate structure
selected by \(\mathcal H_1\).

This residual selector concerns raw ambient coordinate actions and their
completion-specific coefficients and scales.  It must not be promoted to a
claim that every normalized protected Clifford invariant is completion
relative: under the typed geometric input conditions,
Paper~II-C proves on the nondegenerate protected stratum
\(\varrho_{\rm sp}>0\) that the normalized protected even algebra is
completion-independent and represented by \(\Cl(0,3)\).  That result is
not used as an antecedent of the conditional II-B calculation below and does
not select \(\mathsf C_{\rm PI}\).
\end{remark}

\begin{remark}[Host-derived carrier and independent finite-cell imports]
\label{p2:rem:series-selected-minimal-imports}
Theorem~\ref{p2:thm:momentum-host-H1} supplies the generated level-one
support, while the additional whole-candidate certificate in
Theorem~\ref{p2b:thm:momentum-carrier-rigidity} would promote it to the
full GP3-irreducible carrier.  The conditional coordinate input supplies the
\(\mathsf C_{\rm PI}\) protected Euclidean Clifford vector needed in this
subsection.  The algebra and
dimensions below are derived from those relations.  No
\(\mathsf{Supp}_{2T}\) support, Hurwitz-shell identification, or physical
selection of \(m=1\) is used.  The normalized coordinate intertwiners remain
part of \(\mathsf C_{\rm PI}\); they are not derived from the integer label.
\end{remark}

\begin{remark}[Relation to a Hurwitz-link realization]
\label{p2:rem:hurwitz-link-structural-carrier}
Paper~I-A also realizes an algebra abstractly isomorphic to
\(\Cl^0_{4,0}\cong\Cl_{0,3}\) on a Hurwitz-link Walsh space.  Equality of
algebra isomorphism type does not provide a canonical map from that finite
support to \(\mathcal H_1\), and none is used here.
\end{remark}

\begin{theorem}[Protected even Clifford algebra and cross-category count]
\label{p2:thm:minimal-hilbert-structural-15}
For the derived structural momentum-role carrier \(\mathcal H_p\cong\mathcal H_1\),
\[
\mathcal H_1
\cong
V_{0,0}\oplus V_{\frac12,\frac12}\oplus V_{1,0}\oplus V_{0,1}
\oplus V_{\frac12,0}\oplus V_{0,\frac12},
\qquad
\dim_{\mathbb C}\mathcal H_1=15.
\]
If the coordinate completion is chosen to be the Paper-I-A representative
\(\mathsf C_{\rm PI}\), then on the protected complex Dirac module
\[
\mathcal S_{\rm prot}:=V_{\frac12,0}\oplus V_{0,\frac12},
\qquad \dim_{\mathbb C}\mathcal S_{\rm prot}=4,
\]
write \(\mathcal S_{{\rm prot},\mathbb R}\) for its underlying real vector space, so
\(\dim_{\mathbb R}\mathcal S_{{\rm prot},\mathbb R}=8\), and
let \(C_1,\ldots,C_4\) be the imported real Euclidean Clifford vector,
\[
\{C_a,C_b\}=2\delta_{ab}I_{\mathcal S_{\rm prot}}.
\]
The real even operator algebra it generates is
\[
\begin{aligned}
A_{\rm ev}
&:=
\operatorname{span}_{\mathbb R}
\{I,\ C_aC_b\ (1\le a<b\le4),\ C_1C_2C_3C_4\}\\
&\cong\Cl^0_{4,0}\cong\Cl_{0,3},
\end{aligned}
\]
and
\[
\dim_{\mathbb R}A_{\rm ev}=8.
\]
Its matrix-traceless part
\[
\smash{A_{\rm ev}^{0}:=\{A\in A_{\rm ev}:
\Tr_{\mathcal S_{{\rm prot},\mathbb R}}A=0\}}
\]
has real dimension \(7\).  Therefore
\[
\dim_{\mathbb C}\mathcal H_1=15
\quad\text{and}\quad
\dim_{\mathbb R}A_{\rm ev}^{0}+\dim_{\mathbb R}A_{\rm ev}=7+8=15,
\qquad\text{where}\qquad
\dim_{\mathbb R}A_{\rm ev}=8=1+7.
\]
This final equality compares dimensions over different scalar fields and of
different objects.  It asserts no linear, algebraic, or representation-theoretic
isomorphism with \(\mathcal H_1\).  It is also not a decomposition into
complementary pieces: \(A_{\rm ev}^{0}\subset A_{\rm ev}\) is the kernel of
the trace on \(A_{\rm ev}\), and the grade count
\(\dim_{\mathbb R}A_{\rm ev}=8=1+7\), which is the Paper-I-A convention,
displays the \(7\) as a summand of the \(8\).  There is no
fifteen-dimensional object on the Clifford side.  Note also that
\(A_{\rm ev}^{0}\) is defined here by the \emph{real} trace on
\(\mathcal S_{{\rm prot},\mathbb R}\); the Paper-I-A frame identities for the same
operators use the complex trace on \(\mathcal S_{\rm prot}\), under which the
corresponding values are halved.  Both give
\(\dim_{\mathbb R}A_{\rm ev}^{0}=7\).
\end{theorem}

\begin{proof}
The sector dimensions are \(1,4,3,3,2,2\), whose sum is \(15\).
The universal property of \(\Cl_{4,0}\) gives a real representation
\(\Cl_{4,0}\to\End_{\mathbb R}(\mathcal S_{{\rm prot},\mathbb R})\).
The standard real Dirac module is faithful, so restriction identifies the
even image with \(\Cl^0_{4,0}\)
\cite{lawson1989spin}.  Its homogeneous even basis consists of one scalar,
six bivectors, and one volume element, hence has dimension \(1+6+1=8\).
In the Dirac representation every non-scalar displayed basis element is
traceless, while \(I\) has nonzero trace.  The trace therefore has rank one
on \(A_{\rm ev}\), and its kernel has dimension \(7\).
\end{proof}

\begin{remark}[Seven-dimensional object is not a Hilbert complement]
\label{p2:rem:no-seven-hilbert-complement-pointer}
The seven-dimensional traceless even-Clifford space is an operator-space
object, not a seven-real-dimensional Hilbert complement inside
\(\mathcal H_1\).  The elementary dimension check is deferred to
Section~\ref{p2b:app:cross-category-caveat}.
\end{remark}

\begin{corollary}[Seed-carrier and protected Clifford count package]
\label{p2:cor:protected-minimal-package-status}
For the structural momentum-role carrier,
\(\mathcal H_p\cong\mathcal H_1\), \(m_p=1\), and
\(\dim_{\mathbb C}\mathcal H_1=15\).  If the completion is chosen to be
\(\mathsf C_{\rm PI}\), then
\[
\dim_{\mathbb R}A_{\rm ev}=8,\qquad
\dim_{\mathbb R}A_{\rm ev}^{0}=7.
\]
The last two dimensions derive from the protected Clifford vector.  Their sum
equals the first dimension numerically, but no matched carrier, canonical map,
or physical identification follows.
\end{corollary}

\begin{proof}
Combine Theorems~\ref{p2:thm:m1-component} and
\ref{p2:thm:minimal-hilbert-structural-15}.
\end{proof}

\begin{remark}[Noncommutative geometry without points]
\label{p2:rem:noncommutative}
On the derived carrier \(\mathcal H_1\), conditional on
\(\mathsf C_{\rm PI}\), the Toeplitz-compressed coordinate operators generate
a finite-dimensional noncommutative algebra inside
\(\End(\mathcal H_1)\).  Indeed, on the protected block
\(\mathcal S_{\rm prot}\) they are
\[
 \widehat X_{a,\mathcal S_{\rm prot}}=\frac{\rho_{\rm prot}}{2}C_a,
\]
so for \(a\ne b\),
\[
 [\widehat X_{a,\mathcal S_{\rm prot}},\widehat X_{b,\mathcal S_{\rm prot}}]
 =\frac{\rho_{\rm prot}^2}{2}C_aC_b\ne0.
\]
The last inequality follows because each \(C_a^2=I\), so \(C_aC_b\) is
invertible.  This is a proved level-one statement, not a claim for every
integer level: at \(m=0\), \(\mathcal H_0=V_{0,0}\), every compressed
nonconstant coordinate is zero, and the generated unital algebra is only
\(\mathbb C I\), which is commutative.

The protected algebraic scale \(\rho_{\rm prot}\) is independent of the
round-parent radius \(\Rmom\).  The optional convention
\(\rho_{\rm prot}=\Rmom\) may be imposed as a separate scale matching, but it
is not supplied by \(\mathsf C_{\rm PI}\) or by the level label.

The classical carrier remains the continuous round \(S^3\); the notation
introduces no underlying discrete set of points and no second \(S^3\) factor.
The truncation parameter controls the declared module resolution, not the
number of lattice sites.
\end{remark}


\subsection{Finite-support constructions and optional geometric developments}
\label{p2b:app:optional-finite-support}
These results apply only after their finite support, quotient, tower, or Hopf
representative has been separately constituted.  They do not select the
level-one structural momentum-role carrier.

\subsubsection{Finite transitive support means}
\begin{lemma}[Unique invariant mean on a finite transitive support]
\label{p2:lem:finite-transitive-invariant-mean}
Let a finite group \(K\) act transitively on a nonempty finite set \(S\).  The
only \(K\)-invariant probability weight \(\mu:S\to[0,1]\) is
\[
        \mu(s)=\frac1{|S|}\qquad(s\in S).
\]
If \(K_1,K_2\) act transitively on \(S_1,S_2\), their product acts transitively
on \(S_1\times S_2\), and the invariant mean is the product mean.
\end{lemma}

\begin{proof}
Finiteness of \(K\) is not used.  For any \(s,t\in S\), transitivity gives
\(k\in K\) with \(ks=t\);
invariance gives \(\mu(s)=\mu(t)\).  Normalization fixes the common value.
The product assertion follows by acting independently in the two factors.
\end{proof}

\begin{corollary}[Symmetry-derived closure of the Paper-I-A octet descendants]
\label{p2:cor:paper1-octet-descendant-weight-closure}
On an occurring exceptional core-support output
\(\mathsf{Supp}_{2T}\), constituted either by direct core-constructor
occurrence or, from an occurring point-resolution role, after the additional
candidate-independent core constructor and exceptional-domain bridge of
Paper~II-A Corollary~\ref{p2:cor:core-support-inheritance-no-duplicate-witness},
let the octet, scalar-free septet, and oriented rank-three pair supports carry
their constituted transitive permutation actions.  Their \emph{constituted-action-invariant point-probability}
weights are uniquely
\[
    \wp_8=\frac18,
    \qquad
    \wp_7=\frac17,
    \qquad
    \wp_9=\frac19.
\]
Therefore the Paper-I-A finite products
\[
    \wp_8\wp_7=\frac1{56},
    \qquad
    \wp_8\wp_9=\frac1{72}
\]
are fixed as product-orbit means under the product action of the two
constituted transitive groups, although their later
physical, dynamical, or comparison-facing interpretations remain outside
Paper~II.  Generic point resolution alone generates only
\(\mathsf{Supp}_{\varepsilon}(X)\); it does not insert the exceptional
\([2T]\) orbit.  The role occurrence \(\mathsf{Supp}_{2T}\) is likewise not
supplied by \(\mathsf P_{\rm sli}\), a coordinate completion, or \(m=1\).
Once the direct-core or core-constructor/domain-bridge antecedent holds, the
complete unmarked \([2T]\) output co-realizes and is not an additional support
assumption.
No uniqueness is asserted for normalized point weights that are not invariant
under the constituted actions.
\end{corollary}

\begin{proof}
Within \(\mathsf{Supp}_{2T}\), for the octet
\(G=(\mathbb Z_2)^3\), translations act freely and transitively
on the eight cube-link points.  For the septet
\(\widehat G_*=\widehat G\setminus\{\mathbf1\}\), Paper~I-A proves that
\(\operatorname{Aut}(G)\cong GL(3,2)\) acts transitively on the seven
nontrivial characters~\cite{Kugelmann2026a}; the transitivity is also the
standard action of \(GL(3,2)\) on \(\mathbb F_2^3\setminus\{0\}\).  For an oriented trimer \(T\), the constituted cyclic
\(\mathbb Z_3\)-action is free and transitive; hence
\(\mathbb Z_3\times\mathbb Z_3\) is transitive on the pair support
\(T\times T\) of size nine.  Lemma~
\ref{p2:lem:finite-transitive-invariant-mean} gives \(1/8,1/7,1/9\)
directly from these actual symmetries, without enlarging them to a full
orthogonal or unitary group.  These are orbit probabilities, not states on
the reducible matrix algebra \(\End(\mathcal H_1)\).  The same lemma on product supports gives
\(1/56\) and \(1/72\).  These are kinematic invariant-mean identities; later
physical interpretations require their own constituting theorems.
\end{proof}



\subsubsection{Quotient restrictions}
The quotient-restriction result below is independent of the optional
tower machinery.  The separate lower-ideal maximality calculation is deferred
to Section~\ref{p2b:app:lower-ideal-maximality}.

\begin{proposition}[Quotient kinematics are invariant restrictions]
\label{p2:prop:quotient-restriction}
Let \(\Gamma<\mathrm{Isom}^+(S^3)\cong SO(4)\) be a finite freely
acting deck group and let
\(p:S^3\to S^3/\Gamma\) be the quotient covering.  For scalar functions, and
for quotient bundles equipped with the trivial lifted equivariance convention,
the kinematic Hilbert space on \(S^3/\Gamma\) compatible with the lifted round
data is the \(\Gamma\)-invariant subspace of the corresponding simply
connected truncation on \(S^3\).  With unnormalized Riemannian volume on both
spaces, the unitary scalar identification is
\[
 U_\Gamma=|\Gamma|^{-1/2}p^*;
\]
with probability-normalized volume, unscaled pullback \(p^*\) is unitary.
For spinorial bundles the quotient model must supply or generate a
spin-structure member whenever the spin-structure family is non-singleton,
equivalently a homomorphic lift
\[
 \tau\in\operatorname{SpinStruct}(S^3/\Gamma),\qquad
 \tau:\Gamma\longrightarrow\Spin(4),
 \qquad
 \Theta\circ\tau=\iota_\Gamma,
\]
where \(\Theta:\Spin(4)\to SO(4)\) is the covering and
\(\iota_\Gamma\) is the given action.  Such a lift always exists in the
stated free orientation-preserving three-dimensional domain: the preimage
extension
\[
 1\longrightarrow\ZZ_2\longrightarrow
 \Theta^{-1}(\Gamma)\longrightarrow\Gamma\longrightarrow1
\]
therefore splits.  The set of spin structures is a torsor for
\(H^1(S^3/\Gamma;\ZZ_2)\cong\operatorname{Hom}(\Gamma,\ZZ_2)\); in
particular it is unique for odd \(|\Gamma|\), in which case the member is
fixed by uniqueness, while \(\mathbb{RP}^3\) has two.  A twisted associated
bundle may additionally supply an explicit compatible twist member
\[
 \chi\ \text{or}\ \rho_\Gamma\in
 \operatorname{Rep}_{\rm compat}(\Gamma),
\]
and its quotient sections are the corresponding equivariant sections.  These
optional spin/twist members belong to the separately constituted quotient
model and do not
enter the simply connected momentum-parent or carrier derivation.  In the
special left-multiplication cases of Paper~I-A this reduces to a finite subgroup
of \(SU(2)\).  In every case, quotient models are projected equivariant
submodels of the full simply connected \(S^3\) kinematics rather than
independent maximal compact branches.

The state-space identification restricts an operator only when that operator
preserves the chosen deck-equivariant subspace.  More precisely, let
\(\Pi_\Gamma\) be its orthogonal projector on the corresponding finite
truncation \(\mathcal H\).  An operator \(A\in\End(\mathcal H)\) has a
genuine restriction to \(\operatorname{ran}\Pi_\Gamma\) exactly when
\[
 (I-\Pi_\Gamma)A\Pi_\Gamma=0.
\]
Commutation with the prescribed deck action is sufficient.  For self-adjoint
\(A\), the condition is equivalent to \([A,\Pi_\Gamma]=0\), but it need
not imply commutation with every deck transformation on \(\mathcal H\).
Individual \(\hat P_i\) need not descend.  If preservation fails,
\(\Pi_\Gamma A\Pi_\Gamma\) is a compression rather than a genuine
restriction; using it as generated quotient kinematics requires a separately
generated compression constructor.  The orientation-preserving scalar
isometry group is \(N_{SO(4)}(\Gamma)/\Gamma\); a spin/twist model retains
only symmetries admitting compatible lifts that preserve its chosen
equivariance data.
\end{proposition}
\begin{remark}[Optional quotient branch]
A non-singleton spin/twist torsor is live only inside a claim that actually
asserts the corresponding quotient model.  No theorem through Paper~VI-B uses
such a member to constitute the simply connected momentum carrier, so it is
not part of the active structural frontier of that chain.
\end{remark}
\begin{proof}
This follows from the Paper~I-A discussion of free finite-subgroup actions on the selected round \(S^3\) and spherical space forms~\cite{Kugelmann2026a}, together with the standard covering-space identification of quotient functions and sections with invariant or equivariant functions and sections on the cover.  For scalar functions, pullback along the covering map identifies functions on the quotient with \(\Gamma\)-invariant functions on \(S^3\), and finite truncation commutes with taking invariant subspaces.  For a twisted or spinorial associated bundle, the same pullback identifies quotient sections with sections on \(S^3\) satisfying the prescribed lifted \(\Gamma\)-equivariance law; the finite Peter--Weyl truncation again commutes with the corresponding finite-group projector.
The descent and associated-bundle identifications are the standard covering
and bundle results; see \cite{hatcher2002algebraic,steenrod1951topology}.
For the spin clause, \(S^3/\Gamma\) is a closed oriented three-manifold.
Every oriented three-manifold is parallelizable and hence spin
\cite{lawson1989spin}.  Pulling a quotient spin structure back to the unique
spin structure on simply connected \(S^3\) equips that pullback with a genuine
\(\Gamma\)-action.  In the homogeneous frame presentation
\(SO(S^3)\cong SO(4)\), this action is exactly a homomorphism
\(\tau:\Gamma\to\Spin(4)\) lifting \(\iota_\Gamma\), so the displayed
central extension splits.  Conversely every such lift descends the spin
bundle.  The standard affine classification of spin structures gives the
\(H^1(-;\ZZ_2)\) torsor and the two stated examples.
Since \(p\) has degree \(|\Gamma|\),
\[
 \|p^*f\|_{L^2(S^3,d\mathrm{vol})}^2
 =
 |\Gamma|\,\|f\|_{L^2(S^3/\Gamma,d\mathrm{vol})}^2,
\]
which gives the stated Riemannian-volume normalization; dividing both volume
measures by their total volumes gives the probability-volume convention.
The finite-group invariant or character projector commutes with any truncation
projector equivariant for the chosen lifted action.  This is the exact
compatibility required by the phrase ``corresponding simply connected
truncation'' above; no commutation claim is made for an arbitrary
non-equivariant cutoff.

For the operator clause, \(A\) maps \(\operatorname{ran}\Pi_\Gamma\)
into itself if and only if \(A\Pi_\Gamma=\Pi_\Gamma A\Pi_\Gamma\),
which is exactly the displayed condition.  If \(A=A^*\), taking adjoints
also gives \(\Pi_\Gamma A=\Pi_\Gamma A\Pi_\Gamma\); together these
are equivalent to \([A,\Pi_\Gamma]=0\).  Commutation with every deck
operator implies commutation with the corresponding finite-group projector.
Every orientation-preserving quotient isometry lifts to an isometry of the
round universal cover normalizing \(\Gamma\); conversely such normalizing
isometries descend, with kernel \(\Gamma\).  This proves the scalar
normalizer statement, with preservation of the chosen lift/twist data
required in the bundle case.

For a concrete failure of transvection descent, let
\(\Gamma\cong\ZZ_3\) act freely on \(SU(2)\) by left multiplication
by \(\operatorname{diag}(e^{2\pi i/3},e^{-2\pi i/3})\).
In the scalar sector \((1,1)\subset\mathcal H_2^{\rm sc}\), the invariant
left weight is \(m_L=0\).  For the unit weight vector
\(v=|m_L=0,m_R=0\rangle\), the \(L_1^+\) term in
\(\hat P_1=\Punit(L_1^+-L_1^-)\) produces non-invariant left weights
\(\pm1\), while the \(L_1^-\) term remains invariant.  The spin-one
ladder coefficients give
\[
 \|(I-\Pi_\Gamma)\hat P_1v\|=\Punit>0.
\]
The right-factor generators and the scalar Laplacian do preserve this
invariant subspace.
\end{proof}

\subsubsection{Canonical shell filtration}
\label{p2:subsec:shell-filtration}

\begin{definition}[Total-weight shells]
\label{p2:def:total-weight-shells}
For \(s\in \frac12\ZZ_{\ge 0}\) with \(s\le m\), define
\[
\mathcal S_s
:=
\bigoplus_{j_L+j_R=s} V_{j_L}\otimes V_{j_R}.
\]
Then
\[
\mathcal H_m
=
\bigoplus_{\substack{s\in \frac12\ZZ_{\ge 0}\\ s\le m}}
\mathcal S_s.
\]
\end{definition}

\begin{theorem}[Shell filtration relative to the total-weight cutoff of \(\mathcal H_m\)]
\label{p2:thm:shell-filtration}
The family
\[
\mathcal H_{\le s}
:=
\bigoplus_{\substack{r\in \frac12\ZZ_{\ge 0}\\ r\le s}}
\mathcal S_r
\]
defines an increasing \(\Spin(4)\)-invariant filtration of \(\mathcal H_m\), canonical relative to the selected total-weight truncation family. If \(\hat N\) denotes the operator acting by
\[
\hat N|_{\mathcal S_s}=s\,\mathbf 1_{\mathcal S_s},
\]
then \(\hat N\) is self-adjoint and \(\Spin(4)\)-equivariant, and each shell projector \(\Pi_s\) is a polynomial in \(\hat N\).
\end{theorem}

\begin{proof}
The decomposition is just the partition of the defining sum for \(\mathcal H_m\) by total weight \(s=j_L+j_R\). Each shell is a direct sum of \(\Spin(4)\)-irreducibles and is therefore \(\Spin(4)\)-invariant. Since \(\hat N\) acts by distinct scalars on distinct shells, spectral calculus yields the polynomial expression for \(\Pi_s\).
\end{proof}

\begin{proposition}[Shellwise spectral cutoffs]
\label{p2:prop:shellwise-spectral-cutoff}
Let \(T=f(\hat N)\) for some real-valued function \(f\) on the finite shell set. Then every spectral projector of \(T\) is a sum of complete shell projectors \(\Pi_s\). In particular, any cutoff or coarse-graining defined from \(T\) retains or removes whole shells, never partial components inside a shell.  The qualifier is needed: the kinetic operator of Proposition~\ref{p2:prop:component_spectrum} is \emph{not} of the form \(f(\hat N)\), since its eigenvalue is not constant on a shell from total weight one upward, so an energy cutoff does split shells.
\end{proposition}

\begin{proof}
Since
\[
T=\bigoplus_s f(s)\,\mathbf 1_{\mathcal S_s},
\]
its spectral projectors are exactly sums of the shell projectors corresponding to shells with a common eigenvalue \(f(s)\).
\end{proof}


\begin{remark}[Compact branch comparison]
The comparison of excluded spherical quotients, the classical \(SU(2)\)
parent, and optional finite tower members is diagnostic rather than part of
the positive proof.  The full comparison table is deferred to
Section~\ref{p2b:app:compact-branch-comparison}.
\end{remark}

\subsubsection{Hopf phase--direction coupling}

For any representative of the Hopf conjugacy class, the fibration
\[
S^1\hookrightarrow S^3\to S^2
\]
geometrizes a coupling between the quotient direction sphere and its internal
fiber coordinate.  This geometric fiber is not the independently classified
scalar \(S^1_A\) carrier.

\begin{proposition}[Nontrivial Hopf phase-direction coupling]
\label{p2:prop:hopf-phase-direction-coupling}
On the constituted compact momentum parent \(S^3\cong SU(2)\), fix any
maximal torus and the corresponding Hopf representative.  Its fiber coordinate
is not a globally decoupled product coordinate over the quotient direction
sphere.  More precisely, the Hopf projection
\[
\Sone\longrightarrow S^3\overset{\pi}{\longrightarrow}S^2
\]
is the principal \(\mathrm U(1)\)-bundle with first Chern class
\[
c_1(S^3\to S^2)=\pm1\in H^2(S^2;\ZZ),
\]
so it admits no global section and is not isomorphic to the trivial bundle
\(S^2\times \Sone\to S^2\).  Hence the chosen internal fiber coordinate is
globally tied to its quotient-direction base.  In module terms: under the right
\(\mathrm U(1)\), normalize the integer character charge by
\(q:=2m_R\), where \(m_R\) is the usual right spin weight.  The charges in
the \((j,j)\) sector are
\[
 q=-2j,-2j+2,\ldots,2j,
\]
each with multiplicity \(2j+1\).  Thus a charge-\(q\) subspace exists only
when \(j\ge|q|/2\) and \(2j\equiv q\pmod2\), so a fixed phase charge forces
a minimum direction content.  If the half-weight notation \(k=m_R=q/2\) is
used instead, the same condition reads \(j\ge|k|\); the two conventions are
not to be mixed.  With the Paper-I-A orientation convention, the left/common-phase representative
fixed at the start of this paper has \(c_1=+1\).
\end{proposition}

\begin{proof}
The standard Hopf bundle has unit first Chern class up to orientation
\cite{hopf1931abbildung,steenrod1951topology}.
Paper~I-A~\cite{Kugelmann2026a} fixes the left/common-phase representative and
the sign \(c_1=+1\).  A nonzero Chern class forbids a global section, so the
bundle is not a product.  II-D owns the detailed Hopf-class and normalization
theorems; no II-D result is used as a premise here.  No preferred maximal
torus or independent scalar carrier is inferred.

For the representation statement, restrict the right \(V_j\) factor to the
standard maximal torus
\(\theta\mapsto\operatorname{diag}(e^{i\theta},e^{-i\theta})\).  Its spin
weight \(m_R\) has character \(e^{2im_R\theta}\), hence integer charge
\(q=2m_R\) in the displayed range and parity class.  The corresponding right
weight space is one-dimensional, while the left \(V_j\) factor has dimension
\(2j+1\), proving the multiplicity and threshold.
\end{proof}



\subsection{Optional tower and component-ceiling calculus}
\label{p2b:app:optional-tower-machinery}
This section collects optional tower definitions and calculations that are
not used in the constructor-relative proof of the structural momentum-role
seed.  They provide after-the-fact classification of that seed and optional
infrastructure for later constituted roles.

\subsubsection{Word-budget classification}
\label{p2:subsec:spectral-exhaustivity-arbitration}

The following construction is retained as representation-theoretic and
later-role infrastructure.  It classifies the derived momentum seed after the
fact and can serve as an explicit tower grammar for roles whose signature
actually constitutes such a boundary.  It is not a premise of
Theorem~\ref{p2:thm:momentum-host-H1}.

\begin{definition}[Mathematical word stages]
\label{p2:def:spex-fit}
For \(m\in\frac12\mathbb Z_{\ge0}\), define
\[
 \mathcal W(m):=\bigoplus_{0\le r\le2m}(F_L\oplus F_R)^{\otimes r},
 \qquad
 \mathcal C^{\rm word}(m):=\operatorname{IrrSupp}(\mathcal W(m)).
\]
The later-role tags \((\mathrm{TWR}_R)\), \((\mathrm{SPEX}_R)\), and
\((\mathrm{FIT}_R^1)\) retain the series-interface meanings fixed in
Paper~II-A.  Their detailed lower-ideal and ceiling
realizations are supplementary and are collected in
Section~\ref{p2b:app:optional-tower-machinery}.
\end{definition}

\begin{definition}[Complete intrinsic coordinate-type retention]
\label{p2:def:coord-role-antecedent}
Let \(\mathfrak M_R\) be a faithful finite realization of a
\(\Spin(4)\)-typed role \(R\), with carrier \(\mathcal H_R\).  Write
\[
 (\mathrm{COORD}_R)
 \quad\Longleftrightarrow\quad
 V_{\rm coord}:=V_{\frac12,\frac12}
 \preccurlyeq_{\Spin(4)}\mathcal H_R,
\]
where \(\preccurlyeq_{\Spin(4)}\) means an equivariant summand inclusion of
the entire irreducible type, not retention of a selected vector or axis.
The type \(V_{\rm coord}\) is the intrinsic first scalar eigenspace derived in
Theorem~\ref{p2:thm:intrinsic-coordinate-sector}.  Thus
\((\mathrm{COORD}_R)\) is an explicit conditional interface for a later role:
it asserts complete retention of that coordinate type, but it neither selects
a coordinate basis nor follows for every finite realization.  It is not an
antecedent of the momentum-host theorem.
\end{definition}

\begin{theorem}[Word-budget lower-ideal identity]
\label{p2:thm:finite-spectral-exhaustivity}
For every \(m\in\frac12\mathbb Z_{\ge0}\),
\[
 \mathcal C^{\rm word}(m)
 \cong
 \bigoplus_{j_L+j_R\le m}V_{j_L}\otimes V_{j_R}
\]
as a multiplicity-free \(\Spin(4)\)-module.  Thus the mathematical
word-budget prefix is exactly the total-weight lower ideal.  This theorem
classifies a declared word stage; it does not turn unrestricted tensor closure
into physical GP2 generation.
\end{theorem}

\begin{proof}
A word containing \(r_L\) left and \(r_R\) right fundamentals has each
constituent satisfying \(j_L\le r_L/2\) and \(j_R\le r_R/2\), hence
\(j_L+j_R\le m\) for total length at most \(2m\).  Conversely,
\(V_{j_L,j_R}\) occurs in a word with \(2j_L\) left and \(2j_R\) right
fundamentals.  Passing to \(\operatorname{IrrSupp}\) removes repeated copies.
\end{proof}

\begin{remark}[Sharpness and scope of the word grammar]
\label{p2:rem:spex-seed-sharp}
Both primitive chiral types are needed to generate the full representation
ring.  This is a mathematical statement about the frozen parent.  Under the
revised GP2 semantics, however, the physical momentum type-support is fixed by
Theorem~\ref{p2:thm:momentum-host-H1}, while
Theorem~\ref{p2b:thm:momentum-carrier-rigidity} gives the common
one-copy carrier of GP3-irreducible endpoints \emph{when its complete
structured-comparison hypotheses are proved}; the word tower is not
its selector.
Finite-support countermodels to the weaker lower-ideal selector are
compared in Section~\ref{p2b:app:legacy-selector}; this selector is
not part of the GP2 generation theorem.
\end{remark}

\begin{corollary}[Logical boundary of the word-budget theorem]
\label{p2:cor:psli-discharged}
Theorem~\ref{p2:thm:finite-spectral-exhaustivity} proves an identity between
mathematically declared word stages and lower ideals.  It neither asserts nor
needs the universal statement that every finite realized role is a word-stage
prefix.  For the initial momentum role that stronger realized-support
quantifier is unnecessary because constructor-relative host closure already
derives \(\mathcal H_1\).  For later roles, tower membership remains an
independent property to be proved if their signature uses this grammar.
\end{corollary}

\begin{theorem}[Placement of the derived momentum seed in the word tower]
\label{p2:thm:spex-fit-level-one}
The GP2-generated momentum \emph{type-support} of
Theorem~\ref{p2:thm:momentum-host-H1} satisfies
\[
 \mathcal G^{\rm car}_{\Sigma_p}
 \cong\mathcal H_1\cong\mathcal C^{\rm word}(1),
 \qquad m_p=1.
\]
If, in addition, the complete whole-structured comparison hypotheses
of Theorem~\ref{p2b:thm:momentum-carrier-rigidity} hold,
then every GP3-irreducible physical endpoint has full carrier
\(\mathcal H_p\cong\mathcal H_1\).  This is a classification of the
already derived support and, conditionally, its full carrier, not a level
arbitration theorem.
\end{theorem}

\begin{proof}
Combine Theorem~\ref{p2:thm:finite-spectral-exhaustivity} at \(m=1\) with
Theorem~\ref{p2:thm:momentum-host-H1} and
Corollary~\ref{p2:thm:m1-component} for the exact support/word label.
The conditional full-carrier assertion is exactly
Theorem~\ref{p2b:thm:momentum-carrier-rigidity} applied \emph{after}
its hypotheses (i)--(iii) are independently certified.
\end{proof}

\begin{remark}[Weaker selector countermodel]
\label{p2:rem:why-coordinate-hypothesis}
The weaker lower-ideal selector does not force the complete momentum
seed from coordinate content alone.  The sharp countermodels and the reason
the revised host signature excludes them are recorded in
Section~\ref{p2b:app:legacy-selector}; they are not part of the positive
constructor-relative proof.
\end{remark}

\begin{remark}[Status of optional spectral exhaustivity]
\label{p2:rem:spex-proof-status}
Spectral exhaustivity is now a property of an explicitly constituted tower
grammar, not a universal definition of GP2 generation.  The present paper
proves the exact word-stage identity above.  Later roles that use a word tower
must separately prove the corresponding realized-boundary statement; the
initial momentum seed does not consume it.
\end{remark}

\begin{remark}[Optional tower machinery and non-implication]
\label{p2b:rem:optional-tower-machinery-deferred}
The directed-tower construction, the optional spectral-lower-ideal package
\(\mathsf P_{\rm sli}\), registered level ceilings, multiplicity-boundary
lemmas, and component-ceiling arithmetic are not used in the constitutional
momentum-seed proof.  Their precise definitions and proofs are collected in
Section~\ref{p2b:app:optional-tower-machinery}.  The only main-text fact
needed below is that such tower structures remain available as separately
constituted later-role or recovery interfaces and do not select
\(\mathcal H_p\) or \(m_p\).
\end{remark}
These optional structures remain available for later roles whose frozen
signature explicitly constitutes
a word/lower-ideal grammar or a component-spectrum ceiling.  In particular,
none of the statements below is an antecedent of
Theorem~\ref{p2:thm:momentum-host-H1} or
Theorem~\ref{p2b:thm:momentum-carrier-rigidity}.

\begin{definition}[Complete half-integral lower-ideal envelopes]
\label{p2:def:half-integral-complete-envelope}
For \(s\in\tfrac12\mathbb Z_{\ge0}\), set
\[
 \mathcal H_{\le s}^{\rm comp}
 :=\bigoplus_{\substack{j_L,j_R\in\frac12\mathbb Z_{\ge0}\\
                        j_L+j_R\le s}}
       V_{j_L}\otimes V_{j_R}.
\]
By Theorem~\ref{p2:thm:finite-spectral-exhaustivity}, this is the
multiplicity-free module underlying \(\mathcal C^{\rm word}(s)\).  At an
integer boundary \(s=m\), it agrees as a \(\Spin(4)\)-module with
\(\mathcal H_m\); the latter notation additionally refers to the separately
adopted Paper-I-A integer coordinate-stable package.
\end{definition}

\begin{proposition}[Directed finite tower]
\label{p2:prop:directed-finite-tower}
For \(0\le s\le t\) in \(\tfrac12\mathbb Z_{\ge0}\), the complete
half-integral envelopes have the equivariant inclusion and orthogonal
projection, canonical because each \(\mathcal H_{\le t}^{\rm comp}\) is
multiplicity-free and its isotypic
components are determined,
\[
 J^{\rm comp}_{s,t}:\mathcal H_{\le s}^{\rm comp}
   \hookrightarrow\mathcal H_{\le t}^{\rm comp},
 \qquad
 \Pi^{\rm comp}_{t,s}:\mathcal H_{\le t}^{\rm comp}
   \twoheadrightarrow\mathcal H_{\le s}^{\rm comp},
 \qquad
 \Pi^{\rm comp}_{t,s}J^{\rm comp}_{s,t}
   =I_{\mathcal H_{\le s}^{\rm comp}}.
\]
For \(s\le t\le u\), they satisfy
\[
 J^{\rm comp}_{t,u}J^{\rm comp}_{s,t}=J^{\rm comp}_{s,u},
 \qquad
 \Pi^{\rm comp}_{t,s}\Pi^{\rm comp}_{u,t}=\Pi^{\rm comp}_{u,s}.
\]
Hence \(\{\mathcal H_{\le s}^{\rm comp}\}_{s<\infty}\) is a directed
finite tower.  Restriction to integer indices gives the corresponding
directed underlying-module tower \(\{\mathcal H_m\}_{m\in\mathbb Z_{\ge0}}\)
for the integer coordinate-stable subfamily.
Its algebraic union and Hilbert completion are presentation objects used to
type parent/finite maps; neither is an infinite realized carrier.
\end{proposition}

\begin{proof}
The label set at level \(s\) is a subset of the label set at level \(t\).
Use the direct-sum inclusion and the orthogonal projection onto the retained
summands.  Equivariance and the retraction identity are immediate, and both
composition identities follow from nested label sets.  The integer statement
is the restriction \(s=m\), \(t=n\).
\end{proof}

\begin{definition}[Spectral-lower-ideal finite-role package]
\label{p2:def:spectral-lower-ideal-finite-role-package}
Let
\[
 \Lambda_{\rm fin}
 :=
 \tfrac12\mathbb Z_{\ge0}\times\tfrac12\mathbb Z_{\ge0},
 \qquad
 w(j_L,j_R):=j_L+j_R,
\]
and define the total-weight preorder
\[
 \lambda\preceq_w\mu
 \quad\Longleftrightarrow\quad
 w(\lambda)\le w(\mu).
\]
The generic component-bounded equivariant domain permits any finite
\(\Spin(4)\)-stable support whose labels lie below the component ceiling.
The optional hypothesis \(\mathsf P_{\rm sli}\) is an explicit total-weight
lower-ideal tower grammar for later roles and comparative calculations.  It
records:
\begin{enumerate}[label=(\roman*),leftmargin=2em]
\item the complete unmarked coefficient universe
\(\{V_{j_L}\otimes V_{j_R}\}_{(j_L,j_R)\in\Lambda_{\rm fin}}\);
\item the scalar unit and the total-weight function \(w\);
\item the rule that every nonempty finite realized label support is a
\(\preceq_w\)-lower ideal; and
\item no label-specific deletion relation, chirality selector, multiplicity
mark, or candidate-specific compressed-product relation.
\end{enumerate}
Thus, if a support contains a label of weight \(s\), it contains every label
of weight at most \(s\), including every label of the same weight.  The
notation \(\mathsf P_{\rm sli}\) is retained for compatibility with the
Paper-I-A total-left-right-weight family.  It is not a momentum-seed antecedent.  Round \(S^3\), its \(\Spin(4)\) symmetry,
and a component bound do not alone imply this optional lower-ideal grammar.



When this optional tower is used, GP3 still acts through the universal
host-faithful reduction semantics of Paper~II-A.  In particular, unmarked multiplicity may be removed under the additional
generated-comparison hypothesis of
Proposition~\ref{p2:prop:gp3-multiplicity-bound}; the tower itself
supplies neither that comparison nor a physical upper level.
\end{definition}

\begin{definition}[Level-ceiling hypothesis \(\mathsf P_{\rm lvl}(R,s)\)]
\label{p2:def:ceiling-hypothesis}
Let \(R\) be a finite \(\Spin(4)\)-typed role with labels in
\(\Lambda_{\rm fin}\), and let \(s\in\tfrac12\mathbb Z_{\ge0}\).  The
registered hypothesis \(\mathsf P_{\rm lvl}(R,s)\) asserts
\[
 \operatorname{supp}\mathfrak F_R
 \subseteq\{\lambda\in\Lambda_{\rm fin}:w(\lambda)\le s\}.
\]
It is a pure upper bound.  It licenses no reduction, names no preorder, and is
not a GP3 statement.  Under \(\mathsf P_{\rm sli}\), the support already has
shape \(\{w\le s_R\}\) for some \(s_R\); the level hypothesis says only
\(s_R\le s\).  No value of \(s\) is selected by \(\mathsf P_{\rm sli}\), by
GP1--GP3, or by the existence of a generated grading.
\end{definition}

\begin{lemma}[Equivariant truncation with explicit multiplicity boundary]
\label{p2:lem:equivariant-truncation-no-insertion}
Let a finite unitary \(\Spin(4)\)-module have isotypic decomposition
\[
 \mathcal E
 =
 \bigoplus_\lambda V_\lambda\otimes M_\lambda ,
\]
where the \(V_\lambda\) are pairwise inequivalent irreducibles.  If an
orthogonal projector \(Q\) preserves the constituted \(\Spin(4)\) symmetry,
then
\[
 Q
 =
 \bigoplus_\lambda I_{V_\lambda}\otimes Q_\lambda
\]
for orthogonal projectors \(Q_\lambda\) on the multiplicity spaces
\(M_\lambda\).  Thus equivariance retains whole irreducible
\emph{isotypic types} but may retain a proper multiplicity subspace.  If the
ambient envelope is multiplicity-free, then
\[
 Q|_{V_\lambda}\in\{0,I_{V_\lambda}\},
\]
so the cutoff is a sum of whole irreducible-sector projectors.  The
component-spectrum bound may therefore be imposed sectorwise after this
multiplicity boundary has been made explicit.
\end{lemma}

\begin{proof}
The compact-group commutant theorem gives
\[
 \End_{\Spin(4)}(\mathcal E)
 =
 \bigoplus_\lambda
 I_{V_\lambda}\otimes\End(M_\lambda).
\]
Hence an equivariant orthogonal projector has the displayed block form.  The
isotypic components of a unitary module are mutually orthogonal, which is what
makes each \(Q_\lambda\) an orthogonal projector rather than merely idempotent.  When
\(\dim M_\lambda=1\), idempotence makes the scalar block \(0\) or \(1\);
see, for example, \cite{fulton1991representation,Hall2015lie}.
\end{proof}

\begin{proposition}[Conditional GP3 reduction of an unmarked multiplicity]
\label{p2:prop:gp3-multiplicity-bound}
Write a finite equivariant carrier as
\[
 \mathcal E=\bigoplus_\lambda V_\lambda\otimes M_\lambda.
\]
Suppose the GP2-exact host signature is unmarked on a multiplicity
space \(M_\lambda\)---its generated transformations act trivially on that
factor and no generated test distinguishes its directions.  Suppose
\emph{in addition} that whenever \(\dim M_\lambda>1\), there is a
complete GP1-admissible, GP2-exact/host-faithful candidate with one
fewer redundant multiplicity direction, together with a \emph{generated,
role-faithful, proper whole-candidate comparison} in the sense of
Postulate~\ref{p2:post:GP3} preserving all frozen objects, operations,
relations and tests.  Then any GP3-irreducible realization in this
comparison class has
\[
 \dim M_\lambda\le1.
\]
The unmarked hypothesis by itself does \emph{not} provide the generated
comparison: in particular, selecting a one-dimensional copy by a
non-generated projector is forbidden.  The mathematical obstruction to
canonical copy selection is Proposition~\ref{p2b:prop:v279-no-canonical-copy-selector}.
Theorem~\ref{p2b:thm:v293-operational-multiplicity-quotient}
provides a canonical multiplicity-blind \emph{state/effect quotient} when
all generated tests and processes descend; unlike a non-generated copy
projection, it preserves every test in that algebra.  It becomes a GP3
whole-candidate witness only after its GP1/GP2 admissibility, generated
operation status, and preservation of the complete frozen host signature
have been independently certified.
Paper~II-C, Theorem~\ref{p2c:thm:v294-native-multiplicity-algebra},
now discharges multiplicity blindness for the \emph{specified native
rotation--transvection algebra} on every isotypic implementation. Its
coordinate-component equality test remains conditional on representing
the full intrinsic coordinate algebra.
Proposition~\ref{p2c:prop:v294-finite-process-descent} supplies a finite
basis test for any specified process grammar. Neither result certifies
the absence of other frozen multiplicity-sensitive ports nor manufactures
a whole-candidate comparison.
If independently generated data act nontrivially on or separate
\(M_\lambda\), an alleged comparison must preserve them and may fail to
exist.

There are now two independent, exactly checkable output classes.  For
additional frozen linear test/effect ports, the conditional-expectation
defect of Theorem~\ref{p2c:thm:v295-complete-port-twirl} vanishes
exactly when the port factors through the operational quotient; its
process analogue tests every specified Heisenberg generator.  For
\emph{scalar closed-loop traces}, Theorem~\ref{p3b:thm:v295-trace-owner-multiplicity-firewall}
shows that an ordinary unnormalized same-owner trace distinguishes
replicated implementations even if the operator algebra is
multiplicity blind.  Any such \emph{occurring} scalar readout belongs to
the frozen whole-candidate signature and cannot be discarded by a
state/effect quotient; if absent, its hypothetical availability is not
an obstruction.  Neither test alone creates a generated GP3 comparison.
\end{proposition}

\begin{proof}
Under the additional comparison hypothesis, a candidate with
\(\dim M_\lambda>1\) admits a proper generated role-faithful comparison
to another complete GP1/GP2-admissible candidate, so it cannot be
GP3-irreducible by Postulate~\ref{p2:post:GP3}.  Iteration is unnecessary:
GP3-irreducibility already rules out the first proper comparison.
Unmarkedness alone does not prove its existence; neither arbitrary
projection nor compression can replace the stipulated witness.
\end{proof}

\begin{remark}[Physical matter multiplicity, scalar owners, and GP3 comparison]
\label{p2:rem:v329-matter-gp3-whole-owner}
For the finite-matter construction of Paper~V-A, classification of the
specified abstract Clifford--color--family algebra, unit-solder rank
matching, and anomaly-compatible common-copy reduction concern
mathematical representations or conditional physical candidate classes.
Even exact descent of each \emph{specified} state/effect/process port
to a multiplicity-blind algebra does not itself establish the proper
GP3 whole-candidate comparison required by
Proposition~\ref{p2:prop:gp3-multiplicity-bound}.
Where the ordinary unnormalized common-heat trace is an occurring,
frozen scalar action owner, its trace value is part of the complete
physical signature and must be retained. In particular,
Theorem~\ref{p3b:thm:v295-trace-owner-multiplicity-firewall}
prevents an unmodified same-profile state/effect quotient from being
used as a whole-candidate comparison for a replicated ordinary-trace
action. The final physical one-copy conclusion therefore needs an
actual complete comparison or a generated non-occurrence theorem;
abstract algebra closure and the absence of a marked copy selector
alone do not supply either result.
\end{remark}





\begin{theorem}[Positive source-response descent: kernel and Gram certificates]
\label{p2f:thm:v296-source-quadratic-quotient}
Let $J$ be a finite-dimensional real \emph{specified generated source
span}, $L:J\to J_0$ a surjective proposed linear source quotient, and
$G:J\to J^*$ a symmetric positive-semidefinite same-owner response
form. Write $q(j)=\langle j,Gj\rangle$. The following are equivalent:
\begin{enumerate}[label=(\roman*),leftmargin=2em]
\item The complete quadratic response $q(j)$ depends only on $Lj$
for every $j\in J$.
\item $\ker L\subseteq\ker G$.
\item With $P_{\ker L}$ the orthogonal projector for any declared
positive source-space metric,
$\operatorname{Tr}(P_{\ker L}GP_{\ker L})=0$
when $G$ is represented using that same metric.
\end{enumerate}
If these hold, a \emph{unique} positive-semidefinite bilinear form
$G_0$ on $J_0$ satisfies $G=L^*G_0L$. In particular, if $G$ is
positive definite on $J$, \emph{no proper noninjective} linear source
quotient preserves the full response. For a positive source inverse
$T^{-1}$ and an insertion map $I:J\to W$, the criterion applies to
$G=I^*T^{-1}I$; strict positivity of $T^{-1}$ makes
$\ker G=\ker I$. Thus a proper source quotient is possible only if
every discarded direction is already zero under the actual insertion.

This theorem tests a \emph{complete linear source span}. Agreement at
one occurring source section does not imply the criterion, and a
source quotient need not equal the state/effect restriction $Q$ of
Theorem~\ref{p2b:thm:v293-operational-multiplicity-quotient}.
The physical GP3 use remains conditional on the occurrence and typing
of the source family, its full quadratic response and a licensed
whole-candidate comparison; see
Proposition~\ref{p2:prop:gp3-multiplicity-bound}.
\end{theorem}
\begin{proof}
If $q$ factors through $L$, then $q(k)=q(0)=0$ for every
$k\in\ker L$. Positivity gives $Gk=0$, proving (ii). If (ii)
holds, the polarization of $q(j+k)=q(j)$ shows that the form is
constant on $L$-fibres and defines a unique bilinear form on the
surjective quotient $J_0$. It is positive and satisfies
$G=L^*G_0L$. For any positive-semidefinite operator $G$, its
restriction $P_{\ker L}GP_{\ker L}$ is positive, so vanishing of
its trace is equivalent to vanishing of the quadratic form on
$\ker L$, hence to (ii). Finally
$\langle j,I^*T^{-1}Ij\rangle=\|T^{-1/2}Ij\|^2$,
which vanishes exactly when $Ij=0$ for positive definite $T^{-1}$.
The final physical qualification is an independent typing claim, not
part of the finite linear-algebra equivalence.
\end{proof}
\begin{theorem}[Exact component-spectrum truncation theorem]
\label{p2:thm:truncation}
Let momentum space be the round sphere $S^3(\Rmom)$, with momentum unit
$\Punit=\hbar/\Rmom$.  For an auxiliary component-spectrum admissibility
ceiling \(P_{\rm ceil}>0\), set
\[
\rho:=\frac{P_{\rm ceil}}{\Punit}.
\]
If admissible finite momentum sectors are required to have component spectra
bounded by this ceiling, $\|\hat P_i\|\le P_{\rm ceil}$, then an irreducible
$\Spin(4)$ sector can occur only if
\begin{equation}
j_L+j_R\le \rho.
\end{equation}
Since $j_L+j_R\in\frac12\ZZ_{\ge0}$, the exact component-spectrum ceiling is equivalently
\begin{equation}
j_L+j_R\le s_\rho,
\qquad
s_\rho:=\frac12\left\lfloor 2\rho\right\rfloor.
\end{equation}
This determines the allowed label set.  It does not require every allowed
label to occur in a realized carrier.  The maximal label-complete unmarked
component envelope associated with the same ceiling, after adopting
\(\mathsf P_{\rm sli}\), is the envelope of
Definition~\ref{p2:def:half-integral-complete-envelope},
\begin{equation}
\mathcal H_{\le s_\rho}^{\mathrm{comp}}
=\bigoplus_{\substack{j_L,j_R\in\frac12\ZZ_{\ge0}\\ j_L+j_R\le s_\rho}}\bigl(V_{j_L}\otimes V_{j_R}\bigr).
\end{equation}
It contains one representative of every allowed label;
Proposition~\ref{p2:prop:gp3-multiplicity-bound} excludes surplus
unmarked copies from a GP3-irreducible endpoint only when its additional
generated whole-structure comparison has been certified.  This is an
optional label-complete envelope inside \(\mathsf P_{\rm sli}\), not the
constitutional momentum carrier.
If one restricts, as in the Paper~I-A finite truncation and finite-coordinate construction~\cite{Kugelmann2026a}, to the integer coordinate-stable family
\begin{equation}
\mathcal H_m:=\bigoplus_{\substack{j_L,j_R\in\frac12\ZZ_{\ge0}\\ j_L+j_R\le m}}\bigl(V_{j_L}\otimes V_{j_R}\bigr),
\qquad m\in\ZZ_{\ge0},
\end{equation}
then the maximal integer member contained below the same component ceiling has label
\begin{equation}
m:=\left\lfloor \rho\right\rfloor
=\left\lfloor \frac{P_{\rm ceil}}{\Punit}\right\rfloor.
\end{equation}
Throughout this theorem, \(P_{\rm ceil}\) is only an auxiliary
component-spectrum ceiling.  The two
labels \(s_\rho=\frac12\lfloor2\rho\rfloor\) and
\(m=\lfloor\rho\rfloor\) therefore have different status: \(s_\rho\) is the
exact component-spectrum ceiling, while \(m\) is the integer coordinate-stable
family label.  For a fixed ratio \(\rho\), evaluating the Paper-I-A integer
coordinate-stable family gives these labels by the displayed spectral rule.
The optional family \(\{\mathcal H_{\le s}^{\rm comp}\}\) is not itself the physical
GP2 generated closure of the momentum seed.  Paper~II does not fix a numerical
value of \(\rho\), and GP1 boundedness does not select a tower member.  The
physical momentum carrier is instead fixed independently by
Theorem~\ref{p2:thm:momentum-host-H1}.  A bound on a separately defined total
momentum operator would be a different condition and is not used here.
\end{theorem}

\begin{proof}
By Proposition~\ref{p2:prop:component_spectrum}, the component norm on
\((j_L,j_R)\) is \(\Punit(j_L+j_R)\).  The finite component-spectrum
condition therefore excludes exactly the labels with
\(j_L+j_R>\rho\), and half-integrality gives
\[
 s_\rho=\frac12\lfloor2\rho\rfloor.
\]
Lemma~\ref{p2:lem:equivariant-truncation-no-insertion} gives the exact
isotypic-projector form, while
Proposition~\ref{p2:prop:gp3-multiplicity-bound} conditionally excludes
unmarked duplicate copies only under a generated-comparison certificate.
The one-copy maximal envelope here is the declared representative of the
optional \(\mathsf P_{\rm sli}\) module family, not an unconditional GP3
outcome.  The largest \(\preceq_w\)-lower ideal permitted by the ceiling in
\(\mathsf P_{\rm sli}\) contains every label with \(w\le s_\rho\), giving
\(\mathcal H_{\le s_\rho}^{\rm comp}\).  Evaluating the separately declared
Paper-I-A integer coordinate-stable family~\cite{Kugelmann2026a}, the maximal
integer label satisfying \(m\le\rho\) is \(m=\lfloor\rho\rfloor\).
\end{proof}

\begin{remark}[Envelope versus realized carrier]
\label{p2:rem:maximal-envelope-versus-realized-carrier}
Equivariance and the component bound determine which irreducible labels are
allowed; by themselves they do not select a Paper-I-A integer level or require
the label-complete envelope.  The latter is supplied by
\(\mathsf P_{\rm sli}\).  Proposition~\ref{p2:prop:gp3-multiplicity-bound}
makes its unmarked GP3-irreducible endpoint multiplicity-free only if
the required proper generated whole-candidate comparisons exist; a missing
copy mark alone does not supply those comparisons.  Once an integer Paper-I-A level
\(m\) is fixed inside \(\mathsf P_{\rm sli}\), its carrier is the complete
direct sum \(\mathcal H_m\).  A proper submodule is a different finite-role
model, not another realization of that same lower-ideal-complete level.  A
marked multiplicity space likewise requires a distinct chronological origin
and pairing.
\end{remark}

\begin{remark}[Auxiliary ceiling versus sharp realized maximum]
\label{p2:rem:pmax-selected-carrier-reading}
The \(P_{\rm ceil}\) used in Theorem~\ref{p2:thm:truncation} is an arbitrary
admissibility ceiling and need not be attained.  Once an integer
coordinate-stable carrier \(\mathcal H_m\) has been selected, its sharp
realized component maximum is instead
\[
\Pmax^{(m)}:=m\Punit.
\]
Any larger \(P_{\rm ceil}\) that leaves the realized carrier unchanged is
merely a non-sharp inequality and is not a further physical datum.  Thus
\(\Pmax\) in the realized theory denotes the actual spectral maximum, while
\(P_{\rm ceil}\) remains auxiliary cutoff notation.  This sharpness does not
fix an absolute scale: simultaneous rescaling of \(\Punit\) and \(\Pmax\)
leaves the carrier and their ratio unchanged.
The physical momentum carrier is fixed independently of this auxiliary ceiling:
Theorem~\ref{p2:thm:momentum-host-H1} fixes its type-support and
Theorem~\ref{p2b:thm:momentum-carrier-rigidity} identifies
\(\mathcal H_p\cong\mathcal H_1\) at every GP3-irreducible
endpoint \emph{if its whole-structured comparison hypotheses hold};
Corollary~\ref{p2:thm:m1-component} derives \(m_p=1\)
independently from the generated support.  The optional tower
therefore classifies the derived carrier after the fact; it is not the source
of the physical level.  No proper submodule may replace \(\mathcal H_1\)
without violating the generated host signature.
\end{remark}

\begin{remark}[Mathematical Nature of $\mathcal{H}_m$]
\label{p2:rem:extended-hilbert}
The Hilbert space $\mathcal H_m$ is not a truncation of scalar $L^2(S^3)$. By Peter--Weyl, scalar $L^2(S^3)$ contains only diagonal sectors, whereas $\mathcal H_m$ contains all compact Spin(4)-equivariant sectors allowed by the integer coordinate-stable ceiling $j_L+j_R\le m$.  The exact component ceiling may lie at a half-integral value $s_\rho$ before this integer-family restriction is imposed.
\end{remark}

\begin{remark}[Terminology: coordinate-stable]
\label{p2:rem:coordinate-stable-terminology}
The adjective \emph{coordinate-stable} is used in Paper~II for the Paper~I-A
integer truncation family equipped with its coordinate action and Toeplitz
compression.  It is not an additional Paper~I-A theorem name.  The distinction
matters because the component-spectrum ceiling naturally admits half-integral
cutoff values, whereas the Paper-I-A scalar-seeded coordinate-degree subfamily
is labeled by integer \(m\).  The following proposition records this optional tower status.  It does not
participate in the momentum-seed proof, whose physical level is already derived from generated support by
Theorem~\ref{p2:thm:momentum-host-H1} and
Corollary~\ref{p2:thm:m1-component}, without assuming the optional
whole-structured comparison of
Theorem~\ref{p2b:thm:momentum-carrier-rigidity}.
\end{remark}


\begin{corollary}[Level-one consistency inside an optional lower-ideal tower]
\label{p2:cor:coord-plus-ceiling-m1}
If a later role independently adopts
\(\mathsf P_{\rm sli}+(\mathrm{COORD}_R)+\mathsf P_{\rm lvl}(R,1)\), then its
complete lower-ideal carrier is \(\mathcal H_1\).  This is a consistency
statement about the optional tower grammar, not the constitutional proof of
the momentum seed.
\end{corollary}

\begin{proof}
By Definition~\ref{p2:def:coord-role-antecedent},
\((\mathrm{COORD}_R)\) retains the full label
\((\tfrac12,\tfrac12)\), whose total weight is one.  The lower-ideal rule in
\(\mathsf P_{\rm sli}\) therefore retains every label of weight at most one,
while \(\mathsf P_{\rm lvl}(R,1)\) excludes every label of larger weight.
The complete unmarked lower-ideal carrier is consequently exactly
\(\mathcal H_{\le1}^{\rm comp}=\mathcal H_1\).
\end{proof}

\begin{proposition}[Support form of an optional level-one ceiling]
\label{p2:prop:level-one-support-ceiling-bridge}
Within \(\mathsf P_{\rm sli}+(\mathrm{COORD}_R)\), the condition
\(\mathsf P_{\rm lvl}(R,1)\) is equivalent to absence of all retained labels
of total weight greater than one; in particular the first exterior-spinor
labels \((1,\tfrac12)\) and \((\tfrac12,1)\) are absent.  This proposition is
not used to derive \(m_p=1\).
\end{proposition}

\begin{proof}
Definition~\ref{p2:def:ceiling-hypothesis} says precisely that
\(\mathsf P_{\rm lvl}(R,1)\) is the inclusion of the retained support in
\(\{\lambda:w(\lambda)\le1\}\).  This is equivalent to the absence of every
retained label with weight greater than one.  The two displayed exterior-spinor
labels each have weight \(3/2\), so they are excluded.  Conversely, if no
retained label has weight greater than one, the defining support inclusion
holds.  When either equivalent ceiling condition holds, its conjunction with
\((\mathrm{COORD}_R)\) and \(\mathsf P_{\rm sli}\) lets the same argument as
in Corollary~\ref{p2:cor:coord-plus-ceiling-m1} identify the complete support
as \(\mathcal H_1\).
\end{proof}


\subsection{Integer coordinate-stable tower calculations}
\label{p2b:app:integer-coordinate-stable-details}
\begin{proposition}[Status of the integer coordinate-stable family]
\label{p2:prop:integer-coordinate-stable-family}
After the compact parent has been fixed as the round
\(S^3\cong SU(2)\), Theorem~\ref{p2:thm:intrinsic-coordinate-sector}
derives the real coordinate sector as
\(V_{\rm coord}=(\tfrac12,\tfrac12)\).  After adopting
\(\mathsf P_{\rm sli}\), the maximal label-complete half-integral envelope at
terminal weight \(s\) is
\[
\mathcal H_{\le s}^{\rm comp}
=
\bigoplus_{j_L+j_R\le s}V_{j_L}\otimes V_{j_R},
\qquad s\in\tfrac12\ZZ_{\ge0},
\]
whose scalar-seeded integer boundaries give the family
\[
\mathcal H_m
=
\bigoplus_{j_L+j_R\le m}V_{j_L}\otimes V_{j_R},
\qquad m\in\ZZ_{\ge0}.
\]
Thus the restriction from the exact component ceiling to the integer
coordinate-stable label is not forced by the component bound alone.  It is the
boundary convention of the Paper-I-A scalar-seeded coordinate package:
a half-integral terminal shell may be kept as part of the exact component
carrier, while the Paper-I-A coordinate-degree filtration records its unital
boundaries at integer degree.
\end{proposition}

\begin{proof}
Proposition~\ref{p2:prop:component_spectrum} derives the allowed
half-integral labels, while
Lemma~\ref{p2:lem:equivariant-truncation-no-insertion} gives the exact
isotypic form of an equivariant projector.  Neither result requires every
allowed label.  The lower-ideal clause of
Definition~\ref{p2:def:spectral-lower-ideal-finite-role-package} supplies the
complete initial segment \(\mathcal H_{\le s}^{\rm comp}\).

The extra word \emph{coordinate-stable}, however, refers here to the finite unital
package that also carries the Paper-I-A coordinate-sector rule.  The coordinate
symbol has total weight one and preserves \(j_L+j_R\bmod1\).  The integer parity
class contains the scalar unit \((0,0)\) and the order-one coordinate generator
\((\tfrac12,\tfrac12)\); it is therefore the scalar-seeded unital coordinate
filtration.  The half-integral parity class contains no scalar unit and is a
spinorial module over that filtration, not a second coordinate algebra.  Hence
the primitive coordinate boundary is counted by integer coordinate degree.  The Paper-I-A truncations
\(\mathcal H_m\), \(m\in\ZZ_{\ge0}\), are precisely the full total-weight initial
segments at such integer coordinate boundaries, equipped with the canonical
projection used in the Toeplitz coordinate quadruple.  Choosing instead to
stop with a terminal half-integral shell retains the next spinorial-module
shell without the next integer coordinate boundary.  That selection is a valid
exact component cutoff when imposed by \(P_{\rm ceil}\); it is simply not a member of
the Paper-I-A integer coordinate-degree subfamily.  Accordingly, Paper~II
derives the broader half-integral allowed-label set,
\(\mathsf P_{\rm sli}\) supplies its complete lower-ideal envelopes, and the
Paper-I-A coordinate package supplies the integer coordinate-degree boundaries
--- \(\mathsf P_{\rm sli}\) contains no integer clause, while \(\mathsf C_{\rm PI}\) equips
those boundaries with the normalized Paper-I-A coordinate realization.  None
of these statements selects the later physical level.
\end{proof}

\subsection{Cross-category dimensional distinctions}
\label{p2b:app:cross-category-caveat}
\begin{lemma}[No literal seven-dimensional Hilbert complement]
\label{p2:lem:no-literal-seven-hilbert-complement}
The protected spinor module $\mathcal S_{\rm prot}\subset\mathcal H_1$ does not have a seven-real-dimensional Hilbert-space complement inside $\mathcal H_1$.
\end{lemma}

\begin{proof}
Any Hilbert-space complement of the complex four-dimensional
\(\mathcal S_{\rm prot}\) in the complex fifteen-dimensional \(\mathcal H_1\) has
complex dimension \(11\), hence real dimension \(22\).  The
seven-dimensional object is instead \(A_{\rm ev}^{0}\).
\end{proof}

\subsection{Maximality in the optional lower-ideal tower}
\label{p2b:app:lower-ideal-maximality}
\begin{theorem}[Maximal label-complete envelope within \(\mathsf P_{\rm sli}\)]
\label{p2:thm:uniqueness}
Fix the selected global momentum parent manifold
\[
M_p\cong S^3(\Rmom),
\]
with composition law transported from $SU(2)\cong S^3(1)$, and the finite
auxiliary component-spectrum ceiling \(P_{\rm ceil}\).  Assume
\(\mathsf P_{\rm sli}\).
The label-complete unmarked equivariant envelope compatible with that
ceiling is the half-integral component module
\[
\mathcal H_{\le s_\rho}^{\mathrm{comp}}
=\bigoplus_{j_L+j_R\le s_\rho}V_{j_L}\otimes V_{j_R},
\qquad
s_\rho=\frac12\left\lfloor2\frac{P_{\rm ceil}}{\Punit}\right\rfloor.
\]
Within the Paper~I-A integer coordinate-stable family, the maximal member contained below the same ceiling is
\[
\mathcal H_m=\bigoplus_{j_L+j_R\le m}V_{j_L}\otimes V_{j_R},
\qquad
m=\left\lfloor\frac{P_{\rm ceil}}{\Punit}\right\rfloor.
\]
Within this fixed unmarked module family, any other equivariant subcarrier
satisfying the same integer-family bound is obtained by restricting to a
$\Spin(4)$-invariant subrepresentation of \(\mathcal H_m\); without the
integer-family restriction, it is a subrepresentation of
\(\mathcal H_{\le s_\rho}^{\mathrm{comp}}\). This is a label-completeness
statement inside the unmarked compact equivariant construction, not a
classification of all compact quantum models and not an exclusion of
constituted multiplicity spaces.  A proper realized submodule is a different
finite-role model when it violates the declared lower-ideal rule; the explicit
models of Proposition~\ref{p2:prop:spectral-lower-ideal-sharpness} show that
round geometry, equivariance, faithfulness, and left--right symmetry do not
exclude all such alternatives.  Theorem~\ref{p2:thm:no-cutoff-from-finiteness}
also forbids selecting any one of them merely from the existence of a finite
bound.
\end{theorem}
\begin{proof}
Theorem~\ref{p2:thm:truncation} gives the component-spectrum ceiling for every
allowed irreducible label.  Definition~
\ref{p2:def:spectral-lower-ideal-finite-role-package} makes the maximal
permitted support the complete lower ideal \(w\le s_\rho\), while
The notation \(\mathcal H_{\le s_\rho}^{\mathrm{comp}}\) chooses one abstract
copy per retained label, as part of the mathematical envelope definition;
Proposition~\ref{p2:prop:gp3-multiplicity-bound} identifies when that
representative can additionally be justified as a GP3-irreducible endpoint.
This gives
\(\mathcal H_{\le s_\rho}^{\mathrm{comp}}\).  Restriction to the
\(\mathsf P_{\rm sli}\) integer boundary gives \(\mathcal H_m\);
\(\mathsf C_{\rm PI}\) is needed only when its normalized coordinate maps
are invoked.
Lemma~\ref{p2:lem:equivariant-truncation-no-insertion} shows that any
equivariant subcarrier of these multiplicity-free envelopes is obtained by
deleting whole irreducible summands.  This argument is
representation-theoretic and must not be cited as a scalar Peter--Weyl
completeness theorem: scalar functions occupy only the diagonal subpackage
\(\mathcal H_m^{\rm sc}\).
\end{proof}

\subsection{Conditional compact spherical branch comparison}
\label{p2b:app:compact-branch-comparison}
\begin{table}[ht]
\centering
\small
\begingroup
\setlength{\tabcolsep}{2.5pt}
\renewcommand{\arraystretch}{1.18}

\begin{tabular}{|
  L{3.2cm}|
  C{1.4cm}|
  C{2.0cm}|
  C{2.5cm}|
  L{\dimexpr\linewidth-3.2cm-1.4cm-2.0cm-2.5cm-10\tabcolsep-6\arrayrulewidth\relax}|
}
\hline
\textbf{Compact spherical branch}
& \textbf{Local axioms}
& \textbf{Global group law}
& \textbf{Coordinate-stable package}
& \textbf{Status} \\
\hline

Exceptional quotient \(S^3/\Gamma\), \(\Gamma\neq \{1\},\{\pm1\}\)
& \(\checkmark\)
& \(\times\)
& not intrinsic
& Not a group quotient compatible with the selected operation law \\
\hline

\(\mathbb{RP}^3\cong \SO(3)\)
& \(\checkmark\)
& \(\checkmark\)
& \(\times\)
& Excluded by relation-exact globalization; the untwisted coordinate package
also does not descend \\
\hline

Classical \(S^3\cong \SU(2)\)
& \(\checkmark\)
& \(\checkmark\)
& \(\checkmark\)
& Constituted classical carrier \\
\hline

Finite Toeplitz \(S^3_m\)
& \(\checkmark\)
& \(\checkmark\)
& \(\checkmark\)
& Mathematical finite tower member; the physical seed is \(\mathcal H_1\),
with \(\mathsf C_{\rm PI}\) coordinate values only when that representative is chosen \\
\hline

\end{tabular}

\caption[Compact spherical carrier selection]{Compact spherical carrier selection.}
\TableNotes{%
Candidatewise actual-profile closure gives effective realizations.  At the unmarked primitive-role stage, GP2 fixes the prior-generated exact and
host-faithful class, and GP3 selects the irreducible least parent-role
skeleton: this gives \(U(1)\) for the connected abelian role and the \(A_1\)
germ for the connected noncommutative role.  GP2 exact/host-faithful
relation globalization then selects the compact source form \(SU(2)\); the
\(SO(3)\) quotient would add \(-1=1\).  Thus operational rank one and
center-odd sensitivity are consequences or diagnostics, not selection
inputs.
Perfectness and conjugacy of
maximal tori exclude a hidden canonical scalar branch inside \(SU(2)\).  The
finite coordinate package independently verifies why the chosen untwisted
package does not descend to
\(SO(3)\cong\mathbb{RP}^3\).  The resulting classical carrier is
\(S^3_A\cong SU(2)\).  Its physical finite momentum seed is
\(\mathcal H_1\) by Theorem~\ref{p2:thm:momentum-host-H1}; the larger
\(S^3_m\) family is optional mathematical/tower infrastructure, and
\(\mathsf C_{\rm PI}\) supplies normalized Toeplitz coordinate values only
when chosen as the completion representative.
}
\endgroup
\end{table}

\subsection{Selector-independence and weaker-signature countermodels}
\label{p2b:app:legacy-selector}
The previous II-B architecture attempted to select the momentum seed by
restricting realized supports to total-weight word-stage lower ideals and then
combining tower membership, spectral exhaustivity, coordinate content, and a
level-one condition.  That route remains useful as a diagnostic because it
shows why the revised host signature is stronger.  In the present
constructor-relative proof the formal nondependency statement is
\[
 \boxed{
 \mathcal H_p\cong\mathcal H_1
 \quad\text{is independent of}\quad
 \mathsf P_{\rm sli},\ \mathrm{TWR},\ \mathrm{SPEX},\ \mathrm{FIT},\
 \mathsf P_{\rm lvl},\ P_{\rm ceil},\ \mathsf C_{\rm PI}.}
\]
The first six are unnecessary for carrier selection; \(\mathsf C_{\rm PI}\)
enters only later completion-dependent coordinate specializations.

In particular, under the weaker signature
\[
 \mathcal E_9
 =\mathbf1\oplus S_+\oplus S_-\oplus V_{\frac12,\frac12}
\]
is finite, multiplicity-free, left--right symmetric, and faithful, but omits
\((1,0)\oplus(0,1)\).  Hence symmetry, coordinate content, faithfulness, and a
plain support ceiling do not by themselves force \(\mathcal H_1\).  Under the
revised GP2 signature this countermodel is excluded because
\((Z_\beta)_{\mathbb C}=(\mathcal K_{\rm Kill})_{\mathbb C}\), with
\(Z_\beta=\operatorname{span}_{\mathbb R}\operatorname{im}\beta\), inherits
the same-role section-object typing and the same-role no-splitting lemma places
its finite-dimensional object copy in the momentum carrier, so no-omission
forbids deleting it.  Conversely, higher tower sectors
are not inserted unless a later constituted constructor generates them.

The optional package \(\mathsf P_{\rm sli}\), together with TWR/SPEX/FIT and
level-ceiling hypotheses, may still be used to analyze explicitly declared
word-stage roles.  Its use as an independent constitutional explanation of the
initial momentum carrier is superseded by
Theorem~\ref{p2:thm:momentum-host-H1}.

\subsubsection{Optional lower-ideal classification and weaker-signature sharpness}
\label{p2b:app:legacy-lower-ideal}
The following results are retained for comparison with the superseded
selector architecture.  They are statements inside the optional
\(\mathsf P_{\rm sli}\) grammar or about the weaker lower-ideal signature; they are
not premises of the constructor-relative momentum theorem.

\begin{theorem}[Classification inside the spectral-lower-ideal finite-role package]
\label{p2:thm:gp2-exact-label-complete-level-class}
Assume the declared package \(\mathsf P_{\rm sli}\) of
Definition~\ref{p2:def:spectral-lower-ideal-finite-role-package}.  Then the \emph{type-support} of every finite realized model is a
complete initial segment, whose canonical one-copy representative is
\[
 \mathcal H_{\le s}^{\rm comp}
 =
 \bigoplus_{j_L+j_R\le s}V_{j_L}\otimes V_{j_R},
 \qquad
 s\in\tfrac12\mathbb Z_{\ge0}.
\]
The displayed direct sums are type-support representatives and do not assert
multiplicity one for every actual GP2-exact carrier; extra multiplicity can
be excluded by GP3 only with the additional generated-comparison certificate
of Proposition~\ref{p2:prop:gp3-multiplicity-bound}.
At the scalar-seeded unital coordinate boundaries \(s=m\in\mathbb
Z_{\ge0}\), these are exactly the Paper-I-A carriers
\[
        \mathcal H_m
        =\bigoplus_{j_L+j_R\le m}V_{j_L}\otimes V_{j_R}.
\]
In particular, a proper equivariant submodule such as
\[
        (0,0)\oplus(\tfrac12,\tfrac12)
\]
is not another realization of this same \(\mathsf P_{\rm sli}\)-role because
its support is not a lower ideal for \(w\).  It can nevertheless be a valid
realization of a different, narrower finite-role signature.

The comparison is made before compressed products belonging only to a
higher candidate are frozen.  When \(\mathsf C_{\rm PI}\) is also chosen
and complete top shells are deleted, the endpoint receives its own
canonically compressed coordinate doctrine.  No claim is made that
compression is a homomorphism from the higher candidate's full Toeplitz
algebra.  The arguments presently proved here do not yet derive
\(\mathsf P_{\rm sli}\) from GP2 alone.  In particular, the word-budget
identity computes the support after a word boundary has been selected; it
does not prove that every realized finite boundary is selected by word degree
rather than by another generated spectral variable.
\end{theorem}

\begin{proof}
Let \(I\subset\Lambda_{\rm fin}\) be a nonempty finite realized support.
Choose \(\mu\in I\) of maximal weight \(s=w(\mu)\).  If
\(\lambda\in\Lambda_{\rm fin}\) satisfies \(w(\lambda)\le s\), then
\(\lambda\preceq_w\mu\), so lower-ideal completeness gives
\(\lambda\in I\).  Conversely maximality of \(s\) excludes every label of
larger weight.  Hence
\[
 I=\{\lambda\in\Lambda_{\rm fin}:w(\lambda)\le s\},
\]
which is precisely the support of
\(\mathcal H_{\le s}^{\rm comp}\).
The one-copy envelope is the canonical representation of the type-support;
Proposition~\ref{p2:prop:gp3-multiplicity-bound} removes unmarked duplicate
copies from physical GP3 endpoints only under its additional comparison
hypothesis.  Proposition~\ref{p2:prop:integer-coordinate-stable-family}
identifies the scalar-seeded integer boundaries with \(\mathcal H_m\).

The maximal weight \(s\) used in the proof is a property of the given finite
support, not a value selected by this theorem.  Any later physical upper level
must therefore enter through an explicit premise such as
Definition~\ref{p2:def:ceiling-hypothesis}.
\end{proof}

\begin{proposition}[Sharpness of spectral lower-ideal completeness]
\label{p2:prop:spectral-lower-ideal-sharpness}
At component ceiling \(s=1\), define the multiplicity-free
\(\Spin(4)\)-modules
\[
\begin{aligned}
 \mathcal E_5
 &:=
 V_{0,0}\oplus V_{\frac12,\frac12},\\
 \mathcal E_7^L
 &:=
 \mathcal E_5\oplus V_{\frac12,0},\\
 \mathcal E_7^R
 &:=
 \mathcal E_5\oplus V_{0,\frac12},\\
 \mathcal E_9
 &:=
 \mathcal E_5\oplus
 V_{\frac12,0}\oplus V_{0,\frac12}.
\end{aligned}
\]
They have complex dimensions \(5,7,7,9\), respectively.  Each is finite,
unital, \(\Spin(4)\)-equivariant, and contains the complete coordinate sector
\(V_{\frac12,\frac12}\).  The two modules
\(\mathcal E_7^{L,R}\) are faithful \(\Spin(4)\)-modules, while
\(\mathcal E_9\) is faithful, invariant under interchange of the left and
right factors, and contains both protected chiral spinor sectors.  Nevertheless
\[
 \mathcal H_1
 =
 \mathcal E_9\oplus V_{1,0}\oplus V_{0,1}.
\]
Thus the scalar unit and coordinate role require only \(\mathcal E_5\);
adding \(\Spin(4)\)-faithfulness permits the incomparable pair
\(\mathcal E_7^L,\mathcal E_7^R\); adding left--right symmetry gives
\(\mathcal E_9\).  None of these requirements forces the two adjoint sectors.
More strongly, their conjunction---finiteness, unitality,
multiplicity-freedom, \(\Spin(4)\)-equivariance, component ceiling \(s=1\),
retention of the complete coordinate sector and both chiral spinor sectors,
faithfulness, and invariance under interchange of the left and right
factors---still does not force the two adjoint sectors or
\(\mathcal H_1\).  The module \(\mathcal E_9\) is an explicit countermodel.
Any theorem concluding \(\mathcal H_1\) from these data must therefore add a
condition that excludes \(\mathcal E_9\); the lower-ideal completeness clause
of \(\mathsf P_{\rm sli}\) is one such sufficient condition.

All four modules are built inside finite-dimensional complex \(\Spin(4)\)
representation theory, and that is the ambient theory in which the
countermodels are to be read.

A fifth module is needed to see what \(\mathsf P_{\rm sli}\) does
\emph{not} settle by itself.  Each of
\(\mathcal E_5,\mathcal E_7^{L,R},\mathcal E_9\) already violates
lower-ideal completeness, so none of them bears on the question.  The
complete lower ideal at half-integral terminal weight,
\[
 \mathcal H_{\le1/2}
 =
 V_{0,0}\oplus V_{\frac12,0}\oplus V_{0,\frac12},
 \qquad
 \dim_{\mathbb C}\mathcal H_{\le1/2}=5,
\]
does satisfy \(\mathsf P_{\rm sli}\) in full, and is finite, unital,
multiplicity-free, \(\Spin(4)\)-equivariant, faithful, invariant under
interchange of the factors, and in possession of both chiral spinor sectors.
It fails exactly one of the listed properties: it does not contain the
coordinate sector \(V_{\frac12,\frac12}\).  Hence
\(\mathsf P_{\rm sli}\) together with faithfulness, left--right symmetry and
both chiral sectors still does not force \(\mathcal H_1\); the coordinate
requirement is what forces it, and conversely the least
\(\preceq_w\)-lower ideal containing \(V_{\frac12,\frac12}\) is exactly
\(\mathcal H_1\).

Under the Casimir normalization of
Proposition~\ref{p2:prop:component_spectrum}, \(\mathcal H_{\le1/2}\) and
\(\mathcal E_9\) are the two indicated nontrivial proper sublevel sets above
the scalar vacuum inside \(\mathcal H_1\), at
\(\lambda\le\tfrac32\Rmom^{-2}\) and \(\lambda\le3\Rmom^{-2}\)
respectively.  So \(\mathcal E_9\) is not an
arbitrarily punctured module: it is the energy truncation of
\(\mathcal H_1\), and the lower-ideal clause of \(\mathsf P_{\rm sli}\) is
the choice of the component norm \(\|\hat P\|\), rather than the Casimir, as
the variable whose sublevel sets are admissible supports.
\end{proposition}

\begin{proof}
The dimension formula
\(\dim_{\mathbb C}V_{j_L,j_R}=(2j_L+1)(2j_R+1)\) gives the displayed
dimensions.  On \(V_{\frac12,\frac12}\), the kernel of the
\(\Spin(4)=SU(2)_L\times SU(2)_R\) action is
\(\Delta\mathbb Z_2=\{(1,1),(-1,-1)\}\).  Intersecting it with the kernel of
either chiral defining representation leaves only \((1,1)\), proving
faithfulness of \(\mathcal E_7^L\) and \(\mathcal E_7^R\), hence also of
\(\mathcal E_9\).  The factor-swap exchanges the two chiral summands of
\(\mathcal E_9\).  The displayed decomposition of \(\mathcal H_1\) follows
by listing the six labels with \(j_L+j_R\le1\).

Each \(\mathcal E\)-module is a proper equivariant submodule allowed by the
component ceiling.  It fails \(\mathsf P_{\rm sli}\) precisely because
retaining weight-one coordinate content while deleting other labels of
weight at most one violates lower-ideal completeness.  Under the superseded weaker signature, the missing labels were not generated
carrier distinctions and these modules therefore survived as countermodels.
Under the revised host signature they fail GP2 no-omission because
\((Z_\beta)_{\mathbb C}=(\mathcal K_{\rm Kill})_{\mathbb C}
\cong(1,0)\oplus(0,1)\) is a canonically generated finite-dimensional linear
section-object subspace whose same-role carrier typing is inherited.  These modules are unrelated to the seven-point finite
support denoted a septet elsewhere in the Paper-I-A finite datum.

In particular, \(\mathcal E_9\) is finite, unital, multiplicity-free, and
\(\Spin(4)\)-equivariant; it obeys the component ceiling \(s=1\), contains
\((\tfrac12,\tfrac12)\), \((\tfrac12,0)\), and
\((0,\tfrac12)\), is faithful, and is invariant under factor interchange.
Its displayed decomposition omits \(V_{1,0}\) and \(V_{0,1}\), so it satisfies
the full conjunction stated in the proposition while remaining a proper
submodule of \(\mathcal H_1\).  Lower-ideal completeness excludes it because
its support contains weight one but not every label of weight at most one.
\end{proof}





\setlength{\emergencystretch}{2em}

\section{Reciprocal operator matrices and ambient-coordinate constructions}
\label{p2ftech:sec:operators}

\subsection{\texorpdfstring{Explicit operator realization on $\mathcal{H}_1$}{Explicit operator realization on H\_1}}
\label{p2:app:explicit-operators}

This section records an explicit basis for the $m=1$ Hilbert space, gives blockwise matrix representations of the momentum operators $\hat{P}_i$ and of the closed spinor coordinate block, and states the basis-independent checks that define an admissible geometric-sector coordinate table.

\subsubsection{Spin(4) decomposition and basis}
The $m=1$ Hilbert space decomposes into irreducible $\Spin(4)\cong \SU(2)_L\times \SU(2)_R$ subspaces as
\begin{equation}
\label{p2:eq:H1_decomp}
\mathcal{H}_1
=
V_{0,0}\oplus V_{1/2,0}\oplus V_{0,1/2}\oplus V_{1,0}\oplus V_{0,1}\oplus V_{1/2,1/2},
\end{equation}
with dimensions $1,2,2,3,3,4$ respectively.

We order the orthonormal basis states $\ket{j_L,m_L;j_R,m_R}$ by sector:
\begin{center}
\renewcommand{\arraystretch}{1.15}
\begin{tabular}{clccc}
\toprule
Index & State & $(j_L,j_R)$ & Sector & dim \\
\midrule
0  & $\ket{0,0;0,0}$ & $(0,0)$ & scalar & 1 \\
1  & $\ket{\tfrac12,\tfrac12;0,0}$ & $(\tfrac12,0)$ & left spinor & 2 \\
2  & $\ket{\tfrac12,-\tfrac12;0,0}$ & $(\tfrac12,0)$ & left spinor & 2 \\
3  & $\ket{0,0;\tfrac12,\tfrac12}$ & $(0,\tfrac12)$ & right spinor & 2 \\
4  & $\ket{0,0;\tfrac12,-\tfrac12}$ & $(0,\tfrac12)$ & right spinor & 2 \\
5  & $\ket{1,1;0,0}$ & $(1,0)$ & self-dual 2-form & 3 \\
6  & $\ket{1,0;0,0}$ & $(1,0)$ & self-dual 2-form & 3 \\
7  & $\ket{1,-1;0,0}$ & $(1,0)$ & self-dual 2-form & 3 \\
8  & $\ket{0,0;1,1}$ & $(0,1)$ & anti-self-dual 2-form & 3 \\
9  & $\ket{0,0;1,0}$ & $(0,1)$ & anti-self-dual 2-form & 3 \\
10 & $\ket{0,0;1,-1}$ & $(0,1)$ & anti-self-dual 2-form & 3 \\
11 & $\ket{\tfrac12,\tfrac12;\tfrac12,\tfrac12}$ & $(\tfrac12,\tfrac12)$ & vector & 4 \\
12 & $\ket{\tfrac12,\tfrac12;\tfrac12,-\tfrac12}$ & $(\tfrac12,\tfrac12)$ & vector & 4 \\
13 & $\ket{\tfrac12,-\tfrac12;\tfrac12,\tfrac12}$ & $(\tfrac12,\tfrac12)$ & vector & 4 \\
14 & $\ket{\tfrac12,-\tfrac12;\tfrac12,-\tfrac12}$ & $(\tfrac12,\tfrac12)$ & vector & 4 \\
\bottomrule
\end{tabular}
\end{center}

\subsubsection{Momentum operators}
\label{p2:app:momentum_matrices}

Let $L_i^+$ and $L_i^-$ denote the chiral $\mathfrak{su}(2)$ generators acting on the left and right
quantum numbers, respectively. The momentum operators are
\begin{equation}
\label{p2:eq:P_def_app}
\hat{P}_i=\frac{\hbar}{\Rmom}\,(L_i^+ - L_i^-).
\end{equation}
Since $L_i^\pm$ preserve $(j_L,j_R)$, the matrices of $\hat{P}_i$ on $\mathcal{H}_1$ are block-diagonal
with respect to the decomposition \eqref{p2:eq:H1_decomp}.

\begin{proposition}[Diagonal form of $\hat{P}_3$]
\label{p2:prop:P3_diag}
In the basis above,
\[
\hat{P}_3\,\ket{j_L,m_L;j_R,m_R}=\frac{\hbar}{\Rmom}(m_L-m_R)\,\ket{j_L,m_L;j_R,m_R}.
\]
Hence
\[
\hat{P}_3=\frac{\hbar}{\Rmom}\,\mathrm{diag}\Big(0,\tfrac12,-\tfrac12,-\tfrac12,\tfrac12,1,0,-1,-1,0,1,0,1,-1,0\Big).
\]
\end{proposition}

\begin{proof}
\leavevmode\par\noindent
By definition,
\[
  \hat{P}_3=\frac{\hbar}{\Rmom}(L_3^{+}-L_3^{-}).
\]
On the weight basis,
\[
  L_3^{+}\ket{j_L,m_L}=m_L\ket{j_L,m_L},
  \qquad
  L_3^{-}\ket{j_R,m_R}=m_R\ket{j_R,m_R}.
\]
Hence the tensor product state is an eigenvector with eigenvalue
\[
  \frac{\hbar}{\Rmom}(m_L-m_R).
\]
The explicit diagonal matrix follows by ordering the basis states as specified
in the preceding construction.
\end{proof}

\begin{proposition}[Block matrices for $\hat{P}_1$ and $\hat{P}_2$]
\label{p2:prop:P12_blocks}
Let $J_i^{(j)}$ be the standard $(2j+1)$-dimensional $\mathfrak{su}(2)$ generators. Then the blocks of
$\hat{P}_i$ are:
\begin{enumerate}[label=(\roman*), leftmargin=2.1em]
\item On $(0,0)$: $\hat{P}_i=0$.
\item On $(\tfrac12,0)$: $\hat{P}_i=\frac{\hbar}{\Rmom}\,J_i^{(1/2)}=\frac{\hbar}{2\Rmom}\,\sigma_i$.
\item On $(0,\tfrac12)$: $\hat{P}_i=-\frac{\hbar}{\Rmom}\,J_i^{(1/2)}=-\frac{\hbar}{2\Rmom}\,\sigma_i$.
\item On $(1,0)$: $\hat{P}_i=\frac{\hbar}{\Rmom}\,J_i^{(1)}$.
\item On $(0,1)$: $\hat{P}_i=-\frac{\hbar}{\Rmom}\,J_i^{(1)}$.
\item On $(\tfrac12,\tfrac12)$:
\[
\hat{P}_i=\frac{\hbar}{\Rmom}\Big(J_i^{(1/2)}\otimes \mathbf{1}_2-\mathbf{1}_2\otimes J_i^{(1/2)}\Big)
=\frac{\hbar}{2\Rmom}\Big(\sigma_i\otimes \mathbf{1}_2-\mathbf{1}_2\otimes \sigma_i\Big).
\]
\end{enumerate}
\end{proposition}

\begin{proof}
On \(V_{j_L}\otimes V_{j_R}\) one has \(L_i^{+}=J_i^{(j_L)}\otimes \mathbf{1}\) and \(L_i^{-}=\mathbf{1}\otimes J_i^{(j_R)}\).
Substituting into \(\hat{P}_i=\frac{\hbar}{\Rmom}(L_i^{+}-L_i^{-})\) gives
\[
\hat{P}_i=\frac{\hbar}{\Rmom}\bigl(J_i^{(j_L)}\otimes \mathbf{1}-\mathbf{1}\otimes J_i^{(j_R)}\bigr),
\]
from which the stated block forms follow. For $j=\tfrac12$ one uses $J_i^{(1/2)}=\sigma_i/2$.
\end{proof}


\subsubsection{Protected block and phase-gauge representatives}
\label{p2:app:X_spinor_block}

\begin{remark}[Worked factorized-channel calibration]
\label{p2:rem:factorized-channel-calibration-status}
The chiral factorization of
Theorem~\ref{p2:thm:ambient-symmetric-carrier} motivates the
factorized-channel terminology, but no global factorized ambient operator or
unsigned product ansatz is defined in this paper.  The numerical choice
\[
\gamma=\varrho_{\rm sp}=\frac{\Rmom}{2\sqrt2}
\]
is therefore a worked appendix calibration only.  Proposition~
\ref{p2:prop:ambient-completion-existence-nonuniqueness} guarantees that some
bounded ambient completion can realize this protected value, but it neither
selects a unique global coefficient table nor coincides with the Paper-I-A
representative \(\mathsf C_{\rm PI}\) unless all reduced coefficients agree.
\end{remark}

The general Hermitian \(\Spin(4)\)-equivariant protected block is
\[
\widehat X_a^{\rm sp}
=\alpha C_a+\beta\,i\Gamma_{\rm sp}C_a.
\]
The matrices below display its relative chiral phase gauge
\(\widehat X_a^{\rm sp}=\gamma C_a\), where
\(\gamma=\varrho_{\rm sp}=(\alpha^2+\beta^2)^{1/2}\ge0\).  Substitution
\(\gamma=\Rmom/(2\sqrt2)\) gives the factorized-channel
calibration; \(\gamma=\rho_{\rm prot}/2\) gives the Paper-I-A protected
representative under \(\mathsf C_{\rm PI}\).  It becomes
\(\gamma=\Rmom/2\) only under the optional scale match
\(\rho_{\rm prot}=\Rmom\).  None of these substitutions is forced by the
scalar symbol or by \(m=1\).

Tensoring the Killing-spinor module
$(\tfrac12,0)\oplus (0,\tfrac12)$ with the vector irrep $(\tfrac12,\tfrac12)$ does produce ambient Clebsch--Gordan channels with $j_L+j_R>1$, for example $(1,\tfrac12)$ and $(\tfrac12,1)$.
These are precisely the representation-theoretically allowed exterior
channels from the protected block listed among the omitted targets in the
proof of Proposition~\ref{p2:prop:h1-coordinate-coupling-graph}.  They
contribute to protected-block leakage only when the corresponding reduced
coefficients of the chosen ambient completion are nonzero; representation
theory alone does not force that nonvanishing.  They lie outside the \(m=1\)
retained carrier and are projected out by the unique finite orthogonal
compression relative to the chosen ambient action.  Hence \(\hat X_\mu\)
restricts to a \emph{closed} \(4\times 4\) block on
\[
\mathcal{H}_{\mathrm{sp}}:=V_{1/2,0}\oplus V_{0,1/2}.
\]

\begin{proposition}[Explicit \(4\times4\) Clifford block]
\label{p2:prop:X_spinor_block}
Let $\sigma_i$ be the Pauli matrices and set
\(\gamma=\varrho_{\rm sp}\ge0\) after the relative chiral phase choice.  In the ordered basis
$\mathcal{H}_{\mathrm{sp}}=(V_{1/2,0}\oplus V_{0,1/2})$,
define
\begin{equation}
\label{p2:eq:X_spinor_explicit}
\hat X_4^{\mathrm{sp}}=\gamma
\begin{pmatrix}
0 & \mathbf{1}_2\\
\mathbf{1}_2 & 0
\end{pmatrix},
\qquad
\hat X_i^{\mathrm{sp}}=\gamma
\begin{pmatrix}
0 & -i\sigma_i\\
i\sigma_i & 0
\end{pmatrix}\quad (i=1,2,3).
\end{equation}
Then
\begin{equation}
\label{p2:eq:clifford_relation_correct}
\{\hat X_\mu^{\mathrm{sp}},\hat X_\nu^{\mathrm{sp}}\}
=2\gamma^2\,\delta_{\mu\nu}\,\mathbf{1}_4,
\qquad
\sum_{\mu=1}^4 (\hat X_\mu^{\mathrm{sp}})^2
=4\gamma^2\,\mathbf{1}_4.
\end{equation}
The phase-gauge value
\(\gamma=\Rmom/(2\sqrt2)\) is the factorized-channel calibration.
\end{proposition}

\begin{proof}
The displayed matrices are the standard Euclidean chiral Clifford
representation, and the Pauli identity
\(\sigma_i\sigma_j=\delta_{ij}\mathbf{1}_2+i\epsilon_{ijk}\sigma_k\)
gives \eqref{p2:eq:clifford_relation_correct}.  The classification before
phase fixing and the constitutional normalization rows for \(\varrho_{\rm sp}\)
are listed in Corollary~\ref{p2:cor:protected-phase-normalization-freedom};
the worked factorized calibration is stated in the appendix remark above.
\end{proof}

\begin{remark}[Noncommutativity on the spinor block]
For \(\mu\ne\nu\), the exact Clifford formula is
\[
 [\hat X_\mu^{\mathrm{sp}},\hat X_\nu^{\mathrm{sp}}]
 =2\gamma^2 C_\mu C_\nu.
\]
Thus distinct protected coordinates are noncommutative exactly when the
protected coordinate action is nonzero.  At the allowed boundary
\(\gamma=0\), all four protected coordinates vanish and commute.  Under
\(\mathsf C_{\rm PI}\), where \(\gamma=\rho_{\rm prot}/2\) with
\(\rho_{\rm prot}>0\), and under the worked calibration
\(\gamma=\Rmom/(2\sqrt2)\), noncommutativity follows from the already supplied
positive scale.  This exact case distinction is compatible with the
sector-dependent uncertainty relations discussed in the main text.
\end{remark}


\subsubsection{\texorpdfstring{Coordinate operators on the geometric sector}{Coordinate operators on the geometric sector}}
\label{p2:app:X_geo_block}

The geometric-sector coordinate block is the unique finite compression of a
\emph{chosen} ambient action.  Let
\[
\mathcal H_{\mathrm{geo}}:=V_{0,0}\oplus V_{1,0}\oplus V_{0,1}\oplus V_{1/2,1/2},
\]
with the ordered basis
\begin{multline*}
\ket{0,0;0,0} \oplus
\{\ket{1,1;0,0},\ket{1,0;0,0},\ket{1,-1;0,0}\} \\
\oplus \{\ket{0,0;1,1},\ket{0,0;1,0},\ket{0,0;1,-1}\} \\
\oplus \{\ket{\tfrac12,\tfrac12;\tfrac12,\tfrac12},\ket{\tfrac12,\tfrac12;\tfrac12,-\tfrac12},
\ket{\tfrac12,-\tfrac12;\tfrac12,\tfrac12},\ket{\tfrac12,-\tfrac12;\tfrac12,-\tfrac12}\}.
\end{multline*}
Write
\[
\hat X_a^{\mathrm{geo}}:=\Rmom\,\widetilde X_a^{\mathrm{geo}},
\qquad
\widetilde X_a^{\mathrm{geo}}:=\Rmom^{-1}\Pi_{\mathrm{geo}}\mathsf X_a\Pi_{\mathrm{geo}},
\]
where \(\Pi_{\mathrm{geo}}\) is the orthogonal projection onto
\(\mathcal H_{\mathrm{geo}}\).  Clebsch--Gordan coefficients fix the tensorial
shape, while the ambient completion supplies the reduced matrix elements.

\begin{proposition}[Exact geometric-coordinate radius entries]
\label{p2:prop:geo-coordinate-admissibility-checks}
For every compatible Hermitian equivariant coordinate vector on $\mathcal H_1$,
the geometric block is Hermitian and its quadratic radius is positive and
$\Spin(4)$-invariant. With $r_+=r_{1,0}$ and $r_-=r_{0,1}$, its entries in
scalar, vector, positive-adjoint and negative-adjoint order are
\[
 \sum_a(\widehat X_a^{\rm geo})^2
 =\Rmom^2 I_{V_{0,0}}
 \oplus\left[\frac{\Rmom^2}{4}
             +\frac32(r_+^2+r_-^2)\right]I_{V_{1/2,1/2}}
 \oplus2r_+^2 I_{V_{1,0}}
 \oplus2r_-^2 I_{V_{0,1}}.
\]
The protected value is $4\varrho_{\rm sp}^2I_4$.
Thus $\sigma_{0,0}=1$ is fixed by scalar matching, while the remaining
geometric entries depend on the two geometric reduced magnitudes.
\end{proposition}
\begin{proof}
In an orthonormal exterior realization let
$W_av=e_a\wedge v$ and $P_\pm=(I\pm *)/2$. The vector-to-two-form edge is
$T_a=(r_+P_++r_-P_-)W_a$. The exterior identities give
\[
 \sum_aT_aT_a^*=2r_+^2P_++2r_-^2P_-,
 \qquad
 \sum_aT_a^*T_a=\frac32(r_+^2+r_-^2)I_V.
\]
For volume $V_{\rm vol}$, the normalized scalar and first harmonics are
$e_0=V_{\rm vol}^{-1/2}$ and $u_a=2x_a/(\Rmom\sqrt{V_{\rm vol}})$;
the scalar unit $1$ is distinct from $e_0$. Thus the scalar edge is
$\Rmom/2$, contributing $\Rmom^2$ on the scalar and
$\Rmom^2I_V/4$ on the vector. The scalar-to-two-form cross term vanishes
because $W_ae_a=0$. Adding the protected Clifford value proves the formula. In scalar, vector,
positive-adjoint, negative-adjoint, positive-spinor and negative-spinor
order, the intrinsic point gives the dimensionless entries
$(1,11/12,4/9,4/9,3/4,3/4)$.
\end{proof}
\begin{remark}[Status of coordinate tables]
A numerical \(11\times11\) representative is admissible only after its ambient
completion and reduced coefficients have been declared and it passes
Proposition~\ref{p2:prop:geo-coordinate-admissibility-checks}.  In the
\(m=1\) compression, the protected and geometric blocks are invariant, so
after the unitary basis permutation that groups those two invariant
subspaces,
\[
\hat X_a=\hat X_a^{\mathrm{sp}}\oplus \hat X_a^{\mathrm{geo}}
\]
In the declared global sector ordering above, the scalar state precedes the
protected block and the remaining geometric sectors follow it, so this is an
invariant direct-sum statement rather than a claim of literal contiguous
block order before that permutation.  Relative to one fixed ambient action,
different bases give unitary-conjugate tables and an \(O(4)\) frame rotation.
Different allowed ambient coefficients or nonzero cycle phases/signs need not
be equivalent.
\end{remark}



\subsection{Subleading correction and carrier-density diagnostics}
\label{p2:app:subleading-correction-diagnostics}

This section collects finite correction data used only as kinematic diagnostics.  These formulae sharpen the leading-order statements in the body of the paper; they do not introduce additional dynamics, renormalization assumptions, or later-paper phenomenological inputs.

\subsubsection{Scalar Toeplitz leakage identities}
\label{p2:app:toeplitz-leakage-identities}

Let \(\Pi=\Pi_m^{\rm sc}\) be the orthogonal projection onto a finite scalar diagonal sector, let \(Q=1-\Pi\), and write \(T_f:=\Pi M_f\Pi\) for the Toeplitz compression of the multiplication operator by a classical scalar symbol \(f\).  For coordinate symbols \(x_a\) on the round embedding \(S^3(\Rmom)\hookrightarrow\RR^4\), set \(X_a:=T_{x_a}\).  This subsection concerns the scalar multiplication sector only; the full non-diagonal coordinate action obeys Proposition~\ref{p2:prop:extended-coordinate-compression-identity}.

\begin{proposition}[Exact Toeplitz product and leakage identities]
\label{p2:prop:toeplitz-leakage-identities}
For any two classical symbols \(f,g\),
\begin{equation}
T_fT_g=T_{fg}-\Pi M_fQM_g\Pi .
\end{equation}
In particular,
\begin{equation}
X_aX_b=T_{x_ax_b}-\Pi M_{x_a}QM_{x_b}\Pi,
\end{equation}
\begin{equation}
[X_a,X_b]
=\Pi\bigl(M_{x_b}QM_{x_a}-M_{x_a}QM_{x_b}\bigr)\Pi,
\end{equation}
and
\begin{equation}
\sum_{a=1}^4X_a^2=\Rmom^2\Pi-E_m,
\qquad
E_m:=\sum_{a=1}^4\Pi M_{x_a}QM_{x_a}\Pi\ge0.
\end{equation}
For every \(\psi\) in the scalar Hilbert space,
\begin{equation}
\langle\psi,E_m\psi\rangle
=
\sum_{a=1}^4\|QM_{x_a}\Pi\psi\|^2.
\end{equation}
In particular, when \(\psi\in\operatorname{Ran}\Pi\), the final \(\Pi\) may
be omitted.  Moreover,
\begin{equation}
        0\le E_m\le\Rmom^2\Pi.
\end{equation}
If \(P_m^{\rm term}\) denotes the projection onto the terminal scalar shell
\[
V_{m/2}\otimes V_{m/2}^*
\subset\mathcal H_m^{\rm sc},
\]
then
\begin{equation}
        E_m=P_m^{\rm term}E_mP_m^{\rm term}.
\end{equation}
\end{proposition}

\begin{proof}
Insert \(1=\Pi+Q\) between the two multiplication operators:
\[
T_fT_g=\Pi M_f\Pi M_g\Pi
=\Pi M_f(1-Q)M_g\Pi
=\Pi M_fM_g\Pi-\Pi M_fQM_g\Pi.
\]
The coordinate product, commutator, and radius identities follow by taking \(f=x_a\), \(g=x_b\), using \(M_{x_a}M_{x_b}=M_{x_b}M_{x_a}\), and using \(\sum_a x_a^2=\Rmom^2\) on the round sphere.  Positivity follows from
\[
\langle\psi,\Pi M_{x_a}QM_{x_a}\Pi\psi\rangle=\|QM_{x_a}\Pi\psi\|^2.
\]
The radius identity and positivity of
\(\sum_aX_a^2\) give \(E_m\le\Rmom^2\Pi\).  Finally, on the scalar
Peter--Weyl carrier, pointwise coordinate multiplication sends a diagonal
sector \((j,j)\) only to the diagonal targets
\((j\pm\tfrac12,j\pm\tfrac12)\).  If \(j<m/2\), every nonnegative target
lies inside \(\mathcal H_m^{\rm sc}\), so
\[
        QM_{x_a}\Pi(1-P_m^{\rm term})=0.
\]
Thus every leakage factor, and hence \(E_m\), is supported on the terminal
scalar shell.
\end{proof}

\begin{remark}[Linear compact Heisenberg relation is projection-stable]
\leavevmode\par\noindent
On the finite scalar sector, \(\Pi_m^{\rm sc}\) is
\(\Spin(4)\)-equivariant, so the transvection generators commute with it and
the linear transvection-coordinate commutators are compressed exactly:
\begin{equation}
\begin{aligned}
  \relax[P_i,X_j]
  &= \Pi_m^{\rm sc}[P_i,M_{x_j}]\Pi_m^{\rm sc}
   = -\frac{i\hbar}{\Rmom}\delta_{ij}X_4,
  \\
  [P_i,X_4]
  &= \frac{i\hbar}{\Rmom}X_i.
\end{aligned}
\end{equation}
The same linear relation holds on the full extended finite carrier by vector covariance
of \(\mathsf X_a\) and equivariance of \(\Pi_m\).  Scalar Toeplitz leakage,
and more generally extended-coordinate boundary leakage, affects nonlinear
coordinate products, coordinate-coordinate commutators, and the corresponding
quadratic-radius operator, but not the linear compact Heisenberg relation.
\end{remark}




\subsection{Fixed-action coordinate-closure comparison}
\label{p2c:app:e9-diagnostic}
This section retains the fixed-action \(\mathcal E_9\)-closure calculation
as a sharp counterexample to the weaker carrier-selection signature.  It is not used
to select the structural momentum-role carrier in the revised architecture.
Formally,
\[
 \boxed{
 \mathcal H_p\cong\mathcal H_1
 \quad\text{is independent of}\quad
 \mathcal E_9,\ r_{1,0},\ r_{0,1},\ \mathsf C,
 \ \text{and the no-recompression closure condition}.}
\]
Using any of those completion-dependent data as provenance for the already
frozen momentum carrier would reverse the chronology and is not licensed by
the current GP2/GP3 construction.

\begin{proposition}[Exact coordinate closure of the Casimir sublevel]
\label{p2c:prop:e9-coordinate-closure}
\label{p2c:thm:joint-coordinate-closure-locus}
Fix a compatible Hermitian equivariant coordinate vector
\(\widehat X=(\widehat X_1,\ldots,\widehat X_4)\) on \(\mathcal H_1\), and
set
\[
\mathcal E_9
:=
V_{0,0}\oplus V_{\frac12,\frac12}
\oplus V_{\frac12,0}\oplus V_{0,\frac12}.
\]
Let \(\langle\mathcal E_9\rangle_{\widehat X}\) be the smallest subspace of
\(\mathcal H_1\) that contains \(\mathcal E_9\) and is invariant under all
four operators \(\widehat X_a\).  In terms of the two geometric reduced
magnitudes of
Corollary~\ref{p2:cor:h1-coordinate-classification},
\[
\langle\mathcal E_9\rangle_{\widehat X}
=
\begin{cases}
\mathcal E_9,
&r_{1,0}=r_{0,1}=0,\\
\mathcal E_9\oplus V_{1,0},
&r_{1,0}\ne0,\ r_{0,1}=0,\\
\mathcal E_9\oplus V_{0,1},
&r_{1,0}=0,\ r_{0,1}\ne0,\\
\mathcal H_1,
&r_{1,0}r_{0,1}\ne0.
\end{cases}
\]
Equivalently, set
\[
 \mathfrak A_{\rm clos}
 :=\left\{
 \mathcal E_9,
 \mathcal E_9\oplus V_{1,0},
 \mathcal E_9\oplus V_{0,1},
 \mathcal H_1
 \right\},
 \qquad
 \mathfrak A_{\rm coord}:=\mathbb R_{\ge0}^3,
\]
and let \(\kappa(c)\) denote the displayed closure for
\(c=(r_{1,0},r_{0,1},\varrho_{\rm sp})\).  When the licensed relation says
that its first entry is the \emph{least} subspace containing
\(\mathcal E_9\) and closed under the same coordinate action represented by
its second entry, its complete zero-defect locus is
\[
 \boxed{
 \mathfrak Z_{\rm clos}
 =\{(\kappa(c),c):c\in\mathfrak A_{\rm coord}\}
 \subset\mathfrak A_{\rm clos}\times\mathfrak A_{\rm coord}.}
\]
The projection
\(\mathfrak Z_{\rm clos}\to\mathfrak A_{\rm coord}\) is a bijection, so the
common closure constructor propagates every coordinate class and selects
none.  An independently generated equation
\(\kappa(c)=\mathcal E_9\) cuts out exactly
\[
 \{(0,0,\varrho_{\rm sp}):\varrho_{\rm sp}\ge0\},
\]
whereas \(\kappa(c)=\mathcal H_1\) cuts out exactly
\[
 \{(r_{1,0},r_{0,1},\varrho_{\rm sp}):
 r_{1,0}r_{0,1}>0,\ \varrho_{\rm sp}\ge0\}.
\]
If the fixed coordinate action is invariant under interchange of the left and
right factors, then \(r_{1,0}=r_{0,1}\), so only the first and last cases
remain.

Consequently, a realization that contains \(\mathcal E_9\), is required to
be closed under this fixed action without recompression, and has both adjoint
channels nonzero must contain the full \(\mathcal H_1\).  Neither the
nonvanishing of those channels nor the no-recompression closure requirement
is derived from GP1--GP3 or from the scalar coordinate symbol.
\end{proposition}

\begin{proof}
Proposition~\ref{p2:prop:h1-coordinate-coupling-graph} lists every retained
edge.  All edges internal to \(\mathcal E_9\) remain there.  The only edges
from \(\mathcal E_9\) to its complement are
\[
V_{\frac12,\frac12}\longleftrightarrow V_{1,0},
\qquad
V_{\frac12,\frac12}\longleftrightarrow V_{0,1},
\]
with reduced magnitudes \(r_{1,0}\) and \(r_{0,1}\), respectively.  By the
multiplicity-one Clebsch--Gordan rule, a nonzero channel map onto either
irreducible target is surjective: its image is a nonzero invariant subspace
of that target.  Thus each nonzero magnitude adjoins exactly the indicated
adjoint sector.  Once adjoined, that sector has no other retained edge outside
the displayed graph, so no further sector is generated.  This proves the
four cases.  Omitting a generated target gives a closure/omission defect,
while adjoining a target whose channel vanishes gives an insertion defect for
the declared least-closure role; hence the zero-defect set is the displayed
graph.  Its projection has inverse \(c\mapsto(\kappa(c),c)\), proving
nonselection.  Intersecting the graph with
\(\kappa(c)=\mathcal E_9\) or \(\kappa(c)=\mathcal H_1\) gives the two stated
fibres.  Finally, factor interchange exchanges the two adjoint sectors and
hence their reduced magnitudes, proving the symmetry clause.
\end{proof}




\subsection{Ambient-host constructions and conditional representatives}
\label{p2c:app:deferred-host-infrastructure}
This section collects supplementary dictionary, optional kinetic, and
symmetric-host details.  They classify available mathematics and do not
constitute a physical seed, apparatus, or dynamics.

\subsubsection{Supplementary fibration/contact/spin dictionary}
\begin{table}[!htbp]
\centering
\footnotesize
\begingroup
\setlength{\tabcolsep}{2.5pt}
\renewcommand{\arraystretch}{1.16}

\begin{tabular}{|
  L{3.6cm}|
  L{4.3cm}|
  L{\dimexpr\linewidth-3.6cm-4.3cm-6\tabcolsep-4\arrayrulewidth\relax}|
}
\hline
\textbf{Structure}
& \textbf{Geometric property}
& \textbf{Physical representative} \\
\hline

\multicolumn{3}{|c|}{\textbf{\textit{Fibration and Contact Structure}}} \\
\hline

Hopf fibration \(\Sone \to S^3 \to S^2\)
& Chosen representative of the canonical conjugacy class
& Paper-I-A left/common-phase Hopf/Reeb generator; not the scalar \(S^1_A\) parent \\
\hline

Chern class \(|c_1|=1\)
& Representative-independent magnitude; \(c_1=+1\) after the Paper-I-A orientation convention
& Topological charge carrier; Paper~IV gives exact or effective
  superselection only after the accessible algebra preserves the charge \\
\hline

Contact form \(\eta\)
& \(\eta\wedge d\eta\)
  \newline \(=(2/\Rmom)\vol_{S^3(\Rmom)}\)
& Scale structure \\
\hline

Reeb field \(\xi=(2/\Rmom)\partial_\psi\)
& Determined after the Hopf representative, orientation, and normalization
  \(\eta(\xi)=1\) are fixed; Killing
& Chosen Hopf \(\mathrm{U}(1)\) generator; Paper~IV's ray-phase theorem is a
  separate projective-state realization \\
\hline

Charge sectors \(\mathcal V_q\)
& Eigenspaces of \(\xi\)
& Topological sector labels; exact or effective superselection in Paper~IV
  is conditional on charge preservation by the accessible algebra \\
\hline

\multicolumn{3}{|c|}{\textbf{\textit{Spin and Symmetry}}} \\
\hline

Parallelizability
& \(TS^3\cong S^3\times\RR^3\)
& Global frame / trivial Clifford bundle \\
\hline

Unique spin structure
& \(H^1(S^3;\ZZ_2)=0\)
& Unique spin-structure class; a fermion model still requires its declared
  field/module and dynamics \\
\hline

Killing spinors
& Maximal round-sphere Killing-spinor space.  Under the stated convention, the
  complex dimension is \(2\) per Killing-constant sign, hence \(4\) over both
  signs.
& Maximal spinorial symmetry data \\
\hline

\(\SO(4)=(\SU(2)_L\times\SU(2)_R)/\ZZ_2\)
& Orientation-preserving isometry group; lifted carrier symmetry is
  \(\Spin(4)\)
& Full kinematic symmetry \\
\hline

\end{tabular}

\caption[Extended physics dictionary]{Extended physics dictionary.}
\label{p2:tab:physics-dictionary2}
\TableNotes{%
This table extends the physics dictionary by recording constitutive mappings from
vector-momentum geometry to physics.
}
\endgroup
\end{table}



\paragraph{Casimir/Laplacian representative}
After a nonzero real scale \(M\in\mathbb R\setminus\{0\}\) has been chosen,
the \(SO(4)\)-invariant scaled Casimir/Laplacian representative has quadratic
operator \(\mathcal C\):
\[
\mathcal{C}:=\frac{1}{\hbar^2}\left(\sum_{i=1}^3 \hat{J}_i^2 + \Rmom^2\sum_{i=1}^3 \hat{P}_i^2\right).
\]
The scaled operator, denoted \(\hat H_{\rm kin}\) for compatibility with the
series notation, is:
\begin{equation}
\hat H_{\rm kin} = \frac{\hbar^2}{2M\Rmom^2} \mathcal{C}= \frac{1}{2M}\left(\sum \hat{P}_i^2 + \frac{1}{\Rmom^2}\sum \hat J_i^2\right).
\end{equation}
For \(M>0\) this operator is nonnegative and may be used as a mathematical
kinetic-energy representative.  Neither \(M\) nor the assertion that this
operator is the physical kinetic Hamiltonian is generated by the II-C
kinematic constitution.  No positivity assumption on \(M\) is made here, so
for general \(M\ne0\) only the scaled-Casimir statement is asserted.  On the
diagonal sector ($j_L = j_R = k/2$), the Casimir factor reproduces the
standard Laplacian eigenvalues $\lambda_k = k(k+2)/\Rmom^2$. Using the sum of
squared momenta $\sum \hat{P}_i^2$ alone would fail to represent the full
Laplacian on $S^3$.


\begin{theorem}[Symmetric-module factorization]
\label{p2:thm:ambient-symmetric-carrier}
For \(S_L=(\tfrac12,0)\), \(S_R=(0,\tfrac12)\), and
\(W=S_L\oplus S_R\), define the Hilbert-completed symmetric algebra by
\[
\widehat{\operatorname{Sym}}^{\bullet}W
:=\widehat{\bigoplus}_{n\ge0}\operatorname{Sym}^{n}W.
\]
Then
\begin{equation}
\label{p2:eq:ambient-symmetric-factorization}
\mathcal H_\infty
\cong
\widehat{\operatorname{Sym}}^{\bullet}S_L\widehat\otimes
\widehat{\operatorname{Sym}}^{\bullet}S_R
\cong
\widehat{\operatorname{Sym}}^{\bullet}W,
\end{equation}
and
\begin{equation}
\label{p2:eq:finite-symmetric-initial-segment}
\mathcal H_m
\cong
\bigoplus_{n=0}^{2m}\operatorname{Sym}^{n}W.
\end{equation}
The coordinate type is the multiplicity-one mixed summand
\begin{equation}
\label{p2:eq:coordinate-mixed-quadratic}
V_{\rm coord}
\cong
S_L\otimes S_R
\subset\operatorname{Sym}^2W.
\end{equation}
\end{theorem}

\begin{proof}
Since \(V_j\cong\operatorname{Sym}^{2j}\mathbb C^2\), the bidegree
\((2j_L,2j_R)\) part is \(V_{j_L}\otimes V_{j_R}\), and total degree is
\(2(j_L+j_R)\).  Finally,
\[
\operatorname{Sym}^2(S_L\oplus S_R)
\cong
(1,0)\oplus(\tfrac12,\tfrac12)\oplus(0,1),
\]
which proves the last assertion.
\end{proof}

\begin{remark}[Ambient symmetric powers are classification, not GP2 closure]
\label{p2c:rem:ambient-symmetric-not-gp2}
The factorization above is an ambient representation-theoretic identity.  It
does not reinstate unrestricted symmetric or tensor closure as physical GP2
generation.  Only constructor outputs present in the frozen host signature
enter the physical seed.
\end{remark}


\subsection{Supplementary geometric and group-valued interfaces}
\label{p2c:app:supplementary-geometric-interfaces}
The following standard or supplementary constructions support later series
interfaces but are not part of the finite-coordinate proof chain.

\subsubsection{Cone, trace, and Hopf/contact interfaces}
\begin{proposition}[Dilations on the metric cone]
\label{p2:prop:cone-dilations}
The metric cone $C(S^3)=(0,\infty)\times S^3$ equipped with the cone metric
\[
\tilde g = dr^2 + r^2 g_{S^3}
\]
admits a canonical homothetic vector field $r\partial_r$ generating dilations
$r\mapsto e^t r$. This flow commutes with the $\SO(4)$ isometries of $S^3$ and
provides the intrinsic geometric notion of scale transformation used below.
\end{proposition}

\begin{proof}
For the cone metric $\tilde g=dr^2+r^2 g_{S^3}$ one checks directly that
$\mathcal{L}_{r\partial_r}\tilde g = 2\tilde g$, hence $r\partial_r$ is homothetic.
Since the dilation acts only on the radial factor, it commutes with the lifted
action of $\SO(4)$ on the $S^3$ factor.
\end{proof}



\paragraph{Finite kinematic traces at fixed truncation}
Standard compact elliptic theory
\cite{berline1992heat} implies that the nonnegative scalar Laplacian
\(\Delta_{S^3}\) has compact resolvent and discrete spectrum with finite
multiplicities. Consequently, every fixed spectral cutoff contains finitely
many modes, and heat traces \(\mathrm{Tr}(e^{-t\Delta_{S^3}})\) are finite
for \(t>0\). Compactness by itself does not make arbitrary unregularized
traces such as \(\mathrm{Tr}(1)\) or vacuum sums finite; those require either
a convergence domain, analytic regularization, or the finite Toeplitz
projection. On the finite module \(\mathcal H_m\), all operator traces are
finite because \(\dim\mathcal H_m<\infty\).

\paragraph{Scale transformations}
The natural extension of the momentum space $S^3$ to include a scale degree of freedom is the metric cone $C(S^3) = \RR^+ \times S^3$ equipped with the cone metric $\tilde{g} = dr^2 + r^2 g_E$, where $g_E=g_{S^3(\Rmom)}/\Rmom^2$ is the unit-radius round metric on the base. The radial vector field $\rho = r \partial_r$ is the canonical $\SO(4)$-invariant radial homothetic vector field satisfying $\mathcal{L}_\rho \tilde{g} = 2\tilde{g}$. This identifies $\rho$ as the canonical geometric representative of infinitesimal scale transformations. No renormalization-group interpretation is required at this kinematic stage.

\paragraph{Phase isometry}
Fix the standard contact \emph{form} $\eta$ on $S^3$ (normalized by $\eta(\xi)=1$). Then the Reeb vector field $\xi$ is uniquely determined by $\eta$ via $\iota_\xi d\eta=0$ and $\eta(\xi)=1$.
In Hopf coordinates with $\psi\sim\psi+4\pi$ one has $\xi=\frac{2}{\Rmom}\partial_\psi$, whose closed orbits have length $2\pi\Rmom$ in the round metric.


\paragraph{Scaling weights}
Let \(Y_k\) be a scalar harmonic of degree
\(k\in\mathbb Z_{\ge0}\) on the unit-radius base, with nonnegative Laplacian
eigenvalue \(k(k+2)\).  A homogeneous function $r^\alpha Y_k$ on the cone has dilation weight $\alpha$, since $\rho(r^\alpha Y_k)=\alpha r^\alpha Y_k$.  The Laplacian eigenvalue fixes this weight only after an additional extension rule is imposed.  For the harmonic cone extension one has
\[
\Delta_C(r^\alpha Y_k)=r^{\alpha-2}\bigl(k(k+2)-\alpha(\alpha+2)\bigr)Y_k,
\]
so harmonicity gives exactly the two weights $\alpha=k$ and $\alpha=-k-2$.  Without imposing harmonicity, or some other specified cone extension rule, the same spherical harmonic can be extended with arbitrary homogeneous weight.  Thus Paper~II uses the cone dilation as a canonical kinematic scale generator, not as an automatic derivation of scaling dimensions from spherical Laplacian eigenvalues alone.

\paragraph{Topological invariants}
For any Hopf representative on the simply connected carrier, the principal
$\mathrm{U}(1)$-bundle has $|c_1|=1$; after the Paper-I-A orientation convention
is fixed, this is written $c_1=+1$.  More generally, $c_1\in\ZZ$ classifies
principal circle bundles and associated complex line bundles over $S^2$.
Thus the bundle class supplies an intrinsic integer topological label, whereas
its sign and its realization along a displayed fiber axis are
choice-relative.  Interpreting that label as a conserved physical charge
still requires dynamical or superselection closure.


\paragraph{\texorpdfstring{Contact form $\eta$}{Contact form}}
For a fixed round Hopf representative, orientation, and normalization
\(\eta(\xi)=1\), the standard Hopf contact form fixes the corresponding
horizontal/vertical splitting of \(TS^3\).  Its associated Reeb field \(\xi\)
generates that representative's phase isometry, while the curvature
\(d\eta_0\) pulls back from the monopole form on the base \(S^2\) (equivalently
\(d\eta = \Rmom\,\pi^*\omega\) with \(\eta=\Rmom\eta_0\)).  The normalisation
\(\eta \wedge d\eta = \frac{2}{\Rmom}\vol_{S^3(\Rmom)}\) fixes the contact
form relative to those choices and its symplectic properties on horizontal
subspaces.  The underlying tight contact structure is unique up to isotopy,
but arbitrary contact forms defining it are not unique merely up to sign and
scale.
\begin{remark}[Paper~I-A rigidity of the Hopf/contact data]
\label{p2:rem:hopf-rigidity}
The intrinsic statement imported from Paper~I-A is rigidity of the
\emph{conjugacy class}, not selection of an axis.  On any chosen simply
connected round Hopf realization, the principal \(\mathrm U(1)\)-bundle with
\(|c_1|=1\) is unique up to bundle isomorphism, the tight contact structure is
unique up to isotopy, and after fixing the representative, orientation, and
normalization \(\eta(\xi)=1\), its Reeb field is determined.  Changing the
maximal torus conjugates this whole package.  Hence the unmarked round parent
canonically determines the Hopf class and Chern magnitude, but neither a
preferred embedded \(U(1)\) nor an oriented charge generator.  Its later
identification with a physical charge is a separate closure problem.
\end{remark}





\subsubsection{Group-valued momentum summary}
\label{p2c:app:group-valued-momentum}
\label{p2:sec:group-valued-momentum}

With the intrinsic \(\SU(2)\) composition law derived for the momentum parent, momentum
composition is group multiplication and is therefore non-abelian.  The standard
group, exponential-map, and representation-theory facts used in the summary are
reviewed in \cite{Kirillov2004lectures,Hall2015lie,fulton1991representation}.
Table~\ref{p2:tab:kinematic-consequences} records the resulting kinematic
consequences.  Here the infinitesimal generators are one regular chiral
\(\mathfrak{su}(2)\) action.  They are distinct from the symmetric-pair
transvections \(\hat P_i\) of
Definition~\ref{p2:def:composition-versus-transvection}; the latter do not
close as a subgroup because their commutator produces \(\hat J_i\).

\begin{table}[htbp]
\centering
\small
\begingroup
\setlength{\tabcolsep}{2.5pt}
\renewcommand{\arraystretch}{1.18}

\begin{tabular}{|
  L{3.0cm}|
  L{\dimexpr\linewidth-3.0cm-4.5cm-6\tabcolsep-4\arrayrulewidth\relax}|
  L{4.5cm}|
}
\hline
\textbf{Structure}
& \textbf{Mathematical Statement}
& \textbf{Derivation / reference} \\
\hline

Lie group momentum space
& Momentum composition is given by group multiplication:
  \[
    p_1\oplus p_2
    \leftrightarrow
    g_1g_2.
  \]
& Definition of \(\SU(2)\). \\
\hline

Non-abelian addition
& In general,
  \[
    p_1\oplus p_2\neq p_2\oplus p_1.
  \]
& Non-commutativity of \(\SU(2)\); this work. \\
\hline

BCH expansion
& In dimensionless Lie-algebra coordinates \(\xi_a\), for
  \(\xi_1,\xi_2\) in a sufficiently small neighborhood of \(0\) on which
  \(\log(\exp\xi_1\exp\xi_2)\) is defined, put \(g_a=\exp\xi_a\).  Then
  \[
  \begin{aligned}
    \xi_1\oplus\xi_2
    ={}& \xi_1+\xi_2+\frac12[\xi_1,\xi_2] \\
       &+\frac1{12}[\xi_1,[\xi_1,\xi_2]] \\
       &+\frac1{12}[\xi_2,[\xi_2,\xi_1]]
        +R_4(\xi_1,\xi_2).
  \end{aligned}
  \]
  For any fixed norm on the finite-dimensional Lie algebra, the neighborhood
  can be chosen with a constant \(C\) such that
  \[
  \|R_4(\xi_1,\xi_2)\|
  \le C(\|\xi_1\|+\|\xi_2\|)^4.
  \]
  Dimensionful momentum coordinates carry the corresponding powers of the
  momentum scale.
& Local BCH theorem for the Lie algebra exponential map. \\
\hline

Fusion rules
& Allowed momentum couplings obey Clebsch--Gordan rules.
& Clebsch--Gordan representation theory. \\
\hline

Intrinsic momentum parent
& The group-valued momentum parent \(SU(2)\cong S^3\) is compact and its
  geodesic distance is bounded.  Exponential/logarithm coordinates are only
  local and require boundary identifications; they are not one global
  coordinate chart.  Bounded component-generator spectra arise on finite projected
  modules; the structural momentum-role seed is fixed independently by II-B.
& Compact group geometry plus finite module projection; this work. \\
\hline

\end{tabular}

\caption[Group-valued momentum consequences]{Group-valued momentum consequences.}
\label{p2:tab:kinematic-consequences}
\endgroup
\end{table}







\section{Covariant POVMs, uncertainty, and exact moment-radius calculations}
\label{p2ftech:sec:localization}

\subsection{Extended finite-coordinate uncertainty and covariant POVMs}
\label{p2c:app:extended-finite-coordinate-structure}
These verified constructions are supplementary finite-model results.  They
form a mathematical POVM realization atlas, not a selector for a physical
localization apparatus.  A claim that one member is the physical localization
POVM requires a particular-POVM G0 occurrence together with an explicit
label-to-localization map; where coordinate first moments are identified with
a chosen \(\widehat X_a\), a separate completion-dependent moment-matching
bridge is required.  Covariance and identity resolution alone do not select
any of these interfaces.

\subsubsection{Sector-dependent finite-coordinate Robertson bounds}
\label{p2:subsec:fuzzy-gup}
\label{p2:subsec:finite-coordinate-robertson}
The non-commutative finite-coordinate structure leads to Robertson bounds that depend explicitly on the geometric sector.  Unlike standard deformations with a constant parameter $\theta$, the finite-coordinate commutator is sector-dependent.

\begin{theorem}[Sector-dependent finite-coordinate Robertson bound]
\label{p2:thm:spin-dependent-gup}
\label{p2:thm:sector-dependent-coordinate-robertson}
Define the dimensionless ``spin-density'' operators by
\[
\hat\Sigma_{\mu\nu}:=\frac{1}{i\Rmom^2}\,[\hat X_\mu,\hat X_\nu],\qquad \mu,\nu\in\{1,2,3,4\}.
\]
On a finite module, the Robertson uncertainty relation implies the sector-dependent bounds
\begin{align}
\Delta X_\mu\,\Delta X_\nu &\ge \frac{1}{2}\bigl|\langle[\hat X_\mu,\hat X_\nu]\rangle\bigr|
= \frac{\Rmom^2}{2}\,|\langle \hat{\Sigma}_{\mu\nu} \rangle|, \\
\Delta P_i\,\Delta P_j &\ge \frac{1}{2}\bigl|\langle[\hat{P}_i,\hat{P}_j]\rangle\bigr|
= \frac{\hbar}{2\Rmom^2}\,|\epsilon_{ijk}\langle \hat J_k\rangle|,
\end{align}
where \(\hat J_k=\hbar(L^+_k+L^-_k)\).
\end{theorem}

\begin{proof}
On the finite module all displayed operators are bounded.  For any observables $A,B$, Robertson gives $\Delta A\,\Delta B\ge \frac12|\langle[A,B]\rangle|$. Apply this with $(A,B)=(\hat X_\mu,\hat X_\nu)$ and use the definition of $\hat\Sigma_{\mu\nu}$. For momenta, Definition~\ref{p2:def:compact-momentum-generators} gives $[\hat{P}_i,\hat{P}_j]=\frac{i\hbar}{\Rmom^2}\epsilon_{ijk}\hat J_k$, which yields the second bound.
\end{proof}


\begin{corollary}[Protected spinor coordinate-uncertainty bound]
\label{p2:cor:fermionic-precision}
For the normalized invariant scalar vector
\( |0\rangle\in V_{0,0}\) (often called the scalar vacuum vector in the
finite-module notation, but with no dynamical vacuum interpretation here),
\(\langle0|\hat\Sigma_{\mu\nu}|0\rangle=0\) for all \(\mu,\nu\) by
\(\Spin(4)\)-invariance.  On the protected block
\[
\widehat X_\mu^{\rm sp}
=\alpha C_\mu+\beta\,i\Gamma_{\rm sp}C_\mu,
\qquad
\varrho_{\rm sp}=(\alpha^2+\beta^2)^{1/2},
\]
and every distinct Clifford pair obeys
\[
\|\hat\Sigma_{\mu\nu}\|=\frac{2\varrho_{\rm sp}^2}{\Rmom^2}.
\]
For every normalized eigenstate belonging to an extremal eigenvalue of
\(\hat\Sigma_{\mu\nu}\),
\[
\langle\hat X_\mu\rangle=\langle\hat X_\nu\rangle=0,
\qquad
\Delta X_\mu=\Delta X_\nu=\varrho_{\rm sp},
\]
and hence
\[
\Delta X_\mu\,\Delta X_\nu=\varrho_{\rm sp}^2.
\]
The Robertson bound is therefore saturated exactly in every extremal
eigenspace.  Conversely, each Hermitian coordinate
\(\hat X_\mu^{\rm sp}\) has eigenstates with \(\Delta X_\mu=0\), so the
coordinate-product uncertainty can be zero in other states.  There is no
state-independent positive minimum.  In the extremal states the Paper-I-A
\(\mathsf C_{\rm PI}\) construction gives the saturated product
\(\rho_{\rm prot}^2/4\) and individual uncertainties
\(\rho_{\rm prot}/2\).  Under the optional scale match
\(\rho_{\rm prot}=\Rmom\), these become \(\Rmom^2/4\) and \(\Rmom/2\).
\end{corollary}

\begin{proof}
The first claim follows because \(\hat\Sigma_{\mu\nu}\) transforms as an
antisymmetric tensor under \(\Spin(4)\), while \(|0\rangle\) is invariant,
so the expectation value must vanish.  Conjugate the protected family to the
relative chiral phase gauge; simultaneous unitary conjugation of operators
and states preserves spectra, expectations, and variances.  For
\(\mu\ne\nu\), in that gauge,
\[
\hat\Sigma_{\mu\nu}
=\frac{\varrho_{\rm sp}^2}{i\Rmom^2}[C_\mu,C_\nu]
\]
and, since \([C_\mu,C_\nu]=2C_\mu C_\nu\), has eigenvalues
\(\pm2\varrho_{\rm sp}^2/\Rmom^2\).  For
\(\varrho_{\rm sp}>0\), define the Hermitian involution
\[
S_{\mu\nu}
:=
\frac{\Rmom^2}{2\varrho_{\rm sp}^2}\hat\Sigma_{\mu\nu}
=-iC_\mu C_\nu.
\]
It anticommutes with both \(C_\mu\) and \(C_\nu\).  In an eigenstate of
\(S_{\mu\nu}\), this anticommutation forces the two coordinate expectations
to vanish.  Since
\((\hat X_\mu^{\rm sp})^2=(\hat X_\nu^{\rm sp})^2
=\varrho_{\rm sp}^2I\), both variances equal
\(\varrho_{\rm sp}^2\), proving exact saturation.  The case
\(\varrho_{\rm sp}=0\) is immediate.  Finally, finite-dimensional
self-adjointness supplies coordinate eigenstates, whose corresponding
coordinate variance is zero.  The numerical specializations follow from
Corollary~\ref{p2:cor:protected-phase-normalization-freedom}.
\end{proof}


\subsubsection{Diagonal-\texorpdfstring{$SU(2)$}{SU(2)}-Covariant POVMs}
Since the finite coordinate operators do not commute, they admit no common
joint spectral PVM or orthogonal basis of simultaneous embedding-coordinate
eigenstates.  Individual coordinate operators do have eigenstates.  A
joint localization observable may instead be modeled by a Positive
Operator-Valued Measure (POVM).  A single diagonal-$SU(2)$ fiducial vector
does not resolve the identity on all of $\mathcal H_1$, because the diagonal
restriction has multiplicities.  The normalized full-module construction
below is multi-fiducial.  It is a diagonal-symmetry covariant POVM labeled by
the group manifold $SU(2)$.  Its interpretation as momentum-space
localization requires the chosen identification between diagonal-group
labels and localization labels; it is not a point-localization PVM and is not
the unique possible covariant POVM on the compact momentum parent.

\subparagraph{Center convention for the POVM label action.}
Here and below, ``center'' means the center of the label group
\(SU(2)_{\rm diag}\), represented inside
\(\Spin(4)=SU(2)_L\times SU(2)_R\) by \((-1,-1)\).  It is distinct from the
center of the source composition carrier acting through a single chiral
regular action, represented for left composition by \((-1,1)\).  Antipodal
evenness of a diagonal-label density is therefore a statement about the
former action only.

\begin{definition}[Diagonal decomposition of the derived momentum carrier]
\label{p2:def:diag-decomp-H1}
Under the diagonal subgroup $SU(2)_{\rm diag}\subset \Spin(4)$, the derived structural momentum-role carrier $\mathcal H_p\cong\mathcal H_1$ decomposes as
\begin{equation}
\mathcal H_1\downarrow_{SU(2)_{\rm diag}}
\cong
V_0\otimes M_0
\oplus V_{1/2}\otimes M_{1/2}
\oplus V_1\otimes M_1,
\end{equation}
with multiplicity dimensions
\begin{equation}
\dim M_0=2,
\qquad
\dim M_{1/2}=2,
\qquad
\dim M_1=3.
\end{equation}
Indeed,
\[
(0,0)\oplus(\tfrac12,0)\oplus(0,\tfrac12)\oplus(1,0)\oplus(0,1)\oplus(\tfrac12,\tfrac12)
\downarrow_{\rm diag}
\cong
V_0^{\oplus2}\oplus V_{1/2}^{\oplus2}\oplus V_1^{\oplus3}.
\]
\end{definition}

\begin{proposition}[No single-fiducial full-module frame]
\label{p2:prop:no-single-fiducial-povm}
No single diagonal-$SU(2)$ orbit $\{U(g)\Omega:g\in SU(2)\}$ gives an identity resolution on all of $\mathcal H_1$.
\end{proposition}

\begin{proof}
The multiplicity obstruction is deferred to
Section~\ref{p2c:app:povm-obstruction-proofs}.
\end{proof}

\begin{remark}[Scope of the single-orbit obstruction]
Proposition~\ref{p2:prop:no-single-fiducial-povm} is only an obstruction to a Haar-uniform rank-one orbit for the diagonal \(SU(2)\) action on the full \(m=1\) carrier.  It does not rule out different group actions, non-uniform constructions, or the multi-orbit/multi-fiducial POVMs used below.
\end{remark}

\begin{definition}[Multi-fiducial finite-coordinate coherent POVM]
\label{p2:def:multi-fiducial-povm}
\leavevmode\par\noindent
For each
\[
  \ell\in\{0,\tfrac12,1\},
\]
choose a unit vector \(v_\ell\in V_\ell\) and an orthonormal basis
\[
  \{e_{\ell,a}\}_{a=1}^{\dim M_\ell}
  \qquad\text{of}\qquad
  M_\ell .
\]
Define
\begin{equation}
  |g;\ell,a\rangle:=U(g)v_\ell\otimes e_{\ell,a}.
\end{equation}
Let \(d\mu\) denote the unnormalized Haar measure
\cite{haar1933mass}, equivalently the Riemannian
volume measure of the round metric, normalized by
\[
  \mu(S^3)=\mathrm{Vol}(S^3).
\]
For a Borel set \(\Omega\subset S^3\simeq SU(2)\), set
\begin{equation}
  E(\Omega):=
  \sum_{\ell,a}
  \frac{\dim V_\ell}{\mathrm{Vol}(S^3)}
  \int_\Omega |g;\ell,a\rangle\langle g;\ell,a|\,d\mu(g).
\end{equation}
\end{definition}

\begin{theorem}[Diagonal-$SU(2)$-covariant multi-fiducial identity resolution]
\label{p2:thm:multi-fiducial-resolution}
The operators $E(\Omega)$ define a normalized POVM on $S^3$:
\begin{equation}
E(S^3)=\Id_{\mathcal H_1}.
\end{equation}
\end{theorem}

\begin{proof}
Schur averaging on the irreducible module \(V_\ell\)
\cite{Hall2015lie,fulton1991representation} gives
\[
\int_{SU(2)} |U(g)v_\ell\rangle\langle U(g)v_\ell|\,d\mu(g)
=\frac{\mathrm{Vol}(S^3)}{\dim V_\ell}\Id_{V_\ell}.
\]
Tensoring with $|e_{\ell,a}\rangle\langle e_{\ell,a}|$ and multiplying by $\dim V_\ell/\mathrm{Vol}(S^3)$ yields
\[
\Id_{V_\ell}\otimes |e_{\ell,a}\rangle\langle e_{\ell,a}|.
\]
Summing over $a$ gives $\Id_{V_\ell}\otimes\Id_{M_\ell}$, and summing over $\ell$ gives $\Id_{\mathcal H_1}$. Countable additivity and positivity follow from the construction as a finite positive sum of covariant rank-one effects.
\end{proof}

\begin{corollary}[Diagonal-label center blindness and vanishing odd moments of the sectorwise frame]
\label{p2:cor:sectorwise-povm-center-blind}
Let \(F(g)\) denote the operator density of
Definition~\ref{p2:def:multi-fiducial-povm}, so that
\(E(\Omega)=\int_\Omega F(g)\,d\mu(g)\).  Then
\[
        F(-g)=F(g),
        \qquad
        E(-\Omega)=E(\Omega),
\]
where \(-\Omega=\{-g:g\in\Omega\}\).  Hence this diagonal-label POVM descends through
\(SU(2)\to SO(3)\).  Every odd embedding-coordinate moment vanishes; in
particular,
\[
        \int_{SU(2)}x_a(g)\,E(dg)=0,
        \qquad a=1,\ldots,4.
\]
Thus the sectorwise frame is normalized and covariant but cannot realize a
nonzero selected embedding-coordinate operator as its first moment.
\end{corollary}

\begin{proof}
The antipodal-cancellation proof is deferred to
Section~\ref{p2c:app:povm-obstruction-proofs}.
\end{proof}

\begin{remark}[Faithful full-\(S_B^3\) localization firewall]
\label{p2c:rem:center-blind-localization-firewall}
If a later physical localization role is explicitly required to retain the
antipodal distinction already separated by the generated \(E_1^x\) position
tests on \(S_B^3\), the sectorwise center-blind frame above cannot realize
that faithful role: it identifies \(g\) and \(-g\) and naturally factors
through \(S_B^3/\{\pm1\}\cong SO(3)\).  This is a conditional typing
consequence of the existing corollary, not a declaration that every physical
localization apparatus must be faithful or point-resolving.  The separate
label-to-position map and any first-moment matching to a chosen
\(\widehat X_a\) remain later bridge data.
\end{remark}

\begin{proposition}[Minimal diagonal-label-center-sensitive two-orbit frame]
\label{p2:prop:center-sensitive-two-orbit-frame}
A normalized diagonal-\(SU(2)\)-covariant frame on \(\mathcal H_1\) can be
diagonal-label-center-sensitive, and two weighted rank-one orbits are sufficient and
minimal.

Let \(d\nu=d\mu/\mathrm{Vol}(S^3)\), choose orthonormal bases
\(\{u_{\ell,r}\}_{r=1}^{2\ell+1}\) of \(V_\ell\), and retain the
multiplicity bases of
Definition~\ref{p2:def:multi-fiducial-povm}.  Put
\[
\Phi_{1/2}:=\sum_{r=1}^{2}
u_{1/2,r}\otimes e_{1/2,r},
\qquad
\Phi_1:=\sum_{r=1}^{3}u_{1,r}\otimes e_{1,r},
\]
and, for \(s=1,2\),
\[
\Omega_s
:=
u_{0,1}\otimes e_{0,s}
+\Phi_{1/2}
+\sqrt{\frac32}\,\Phi_1.
\]
Then
\[
\widetilde E(\Omega)
:=
\sum_{s=1}^{2}
\int_\Omega
|U(g)\Omega_s\rangle\langle U(g)\Omega_s|\,d\nu(g)
\]
is a normalized covariant POVM:
\[
        \widetilde E(S^3)=I_{\mathcal H_1}.
\]
Equivalently, each \(\Omega_s\) may be normalized and its orbit assigned
weight \(\|\Omega_s\|^2=15/2\).  The density of
\(\widetilde E\) is not antipodally even.
\end{proposition}

\begin{proof}
For a vector component
\(\Omega_{\ell}\in V_\ell\otimes M_\ell\), Schur averaging gives
\[
\int_{SU(2)}
|U(g)\Omega_{\ell}\rangle\langle U(g)\Omega_{\ell}|\,d\nu(g)
=
\frac{I_{V_\ell}}{\dim V_\ell}
\otimes
\operatorname{Tr}_{V_\ell}
|\Omega_{\ell}\rangle\langle\Omega_{\ell}|.
\]
Cross terms between inequivalent \(\ell\) vanish.  For the two displayed
fiducials, the sums of the multiplicity-space partial traces are
\[
\begin{array}{c|c}
\ell & \displaystyle
\sum_{s=1}^{2}\operatorname{Tr}_{V_\ell}
|\Omega_{s,\ell}\rangle\langle\Omega_{s,\ell}|\\
\hline
0 & I_{M_0}\\
\tfrac12 & 2I_{M_{1/2}}\\
1 & 3I_{M_1}.
\end{array}
\]
These are exactly
\((\dim V_\ell)I_{M_\ell}\), proving the identity resolution.

The diagonal-label center acts with sign \(+1\) on the \(\ell=0,1\) components and sign
\(-1\) on the \(\ell=\tfrac12\) component.  Each \(\Omega_s\) contains
nonzero components of both signs, and the sum of the corresponding
even--odd cross terms is nonzero.  Its rank-one density therefore changes
under \(g\mapsto-g\).  Finally, on the \(\ell=0\) isotypic component a
single fiducial contributes a multiplicity-space operator of rank at most
one, whereas \(\dim M_0=2\).  At least two orbits are necessary, so the
construction is minimal.  No claim is made that its first moments equal a
chosen \(\widehat X_a\); that moment-matching condition is a separate bridge.
\end{proof}

\begin{remark}[Protected diagonal-irrep coherent projections]
\label{p2:rem:protected-sector-coherent-projections}
The obstruction above concerns a single fiducial vector resolving the identity
on the full multiplicity-bearing module \(\mathcal H_1\).  On either protected
chiral summand,
\[
V_{\frac12,0}\downarrow_{\rm diag}\cong V_{\frac12},
\qquad
V_{0,\frac12}\downarrow_{\rm diag}\cong V_{\frac12},
\]
the usual single-irrep Schur coherent-state normalization remains valid.
These diagonal irreducible modules are not the Hopf/Reeb charge sectors
\(\mathcal V_q\); no corresponding single-orbit statement for a general
\(\mathcal V_q\) is asserted here.  Later protected-diagonal-irrep overlap
calculations should cite this sector construction, not the invalid
single-fiducial full-module frame.
\end{remark}




\subsection{Finite POVM obstruction proofs}
\label{p2c:app:povm-obstruction-proofs}
The main text retains the obstruction statements because they motivate the
successful multi-fiducial and center-sensitive constructions.  Their routine
Schur-averaging and antipodal-cancellation proofs are collected here.

\subsubsection{Single-fiducial obstruction}
\begin{proof}
For a normalized vector $\Omega$, the frame operator
\[
F_\Omega=\int_{SU(2)} |U(g)\Omega\rangle\langle U(g)\Omega|\,d\mu(g)
\]
commutes with the diagonal $SU(2)$ action. By Schur's lemma
\cite{fulton1991representation} on each isotypic component,
\[
F_\Omega=\bigoplus_{\ell\in\{0,1/2,1\}}\Id_{V_\ell}\otimes A_\ell,
\]
where $A_\ell$ acts on $M_\ell$. A single vector contributes at most rank $\dim V_\ell$ to $A_\ell$. For $\ell=0$, one has $\dim V_0=1$ but $\dim M_0=2$, so $A_0$ cannot be full rank. Hence $F_\Omega$ cannot be a nonzero scalar multiple of the identity on $\mathcal H_1$.
\end{proof}

\subsubsection{Center blindness and odd moments}
\begin{proof}
On \(V_\ell\), the central element \(-1\) of the diagonal label group acts as
\((-1)^{2\ell}I\).  Each fiducial in
Definition~\ref{p2:def:multi-fiducial-povm} lies in a single
\(\ell\)-isotypic component, so replacing \(g\) by \(-g\) multiplies its
orbit vector by a scalar sign and leaves the rank-one projector unchanged.
This proves \(F(-g)=F(g)\), and Haar invariance gives the set identity.
Because \(x_a(-g)=-x_a(g)\), the first moment cancels under the antipodal
change of variables.  The same argument applies to every odd embedding
polynomial.
\end{proof}

\subsection{Exact classical moment-saturation obstruction}
\label{p2c:app:exact-moment-saturation-obstruction}
This section proves the scope-protecting no-go statement in
Proposition~\ref{p2c:prop:exact-intrinsic-classical-moment-no-go}.  It is not
part of the positive unsharp-POVM construction.

\begin{lemma}[Ball-valued POVM saturation forces commutativity on the saturated vector]
\label{p2c:lem:povm-ball-saturation-commutativity}
Let \(E\) be a normalized POVM on
\(B_R^d=\{x\in\mathbb R^d:|x|\le R\}\), and let
\(X_a=\int x_aE(dx)\).  Then
\[
 \sum_aX_a^2\le R^2I.
\]
If \(\psi\) satisfies
\[
 \left(R^2I-\sum_aX_a^2\right)\psi=0,
\]
then \([X_a,X_b]\psi=0\) for all \(a,b\).
\end{lemma}

\begin{proof}
Choose a Naimark dilation \(E(\Omega)=V^*P(\Omega)V\), with \(V\) an
isometry and \(P\) a projection-valued measure.  Put
\(Z_a=\int x_aP(dx)\) and \(P_H=VV^*\).  The \(Z_a\) commute, and
\[
 \begin{aligned}
 R^2I-\sum_aX_a^2
 &=V^*\left(R^2I-\sum_aZ_a^2\right)V\\
 &\quad+\sum_aV^*Z_a(I-P_H)Z_aV.
 \end{aligned}
\]
Both terms are positive.  Saturation on \(\psi\) therefore gives
\((I-P_H)Z_aV\psi=0\) for every \(a\), so
\(Z_aV\psi=VX_a\psi\).  Consequently
\[
 V[X_a,X_b]\psi
 =P_H(Z_aZ_b-Z_bZ_a)V\psi=0.
\]
Injectivity of \(V\) proves the commutator claim, while the displayed positive
decomposition proves the moment inequality.
\end{proof}

\begin{proof}[Detailed application to \(X^D\)]
The Paper-I-A quadratic fingerprint gives
\(\sum_a(X_a^D)^2e_0=\Rmom^2e_0\).  Its explicit edges are
\[
 X_a^De_0=\frac{\Rmom}{2}u_a,
 \qquad
 X_a^Du_b=\frac{\Rmom}{2}\delta_{ab}e_0
 +\frac{\sqrt2\Rmom}{3}w_{ab},
 \qquad w_{ba}=-w_{ab}.
\]
For \(a\ne b\), therefore,
\[
 [X_a^D,X_b^D]e_0
 =\frac{\sqrt2\Rmom^2}{3}w_{ab}\ne0.
\]
Lemma~\ref{p2c:lem:povm-ball-saturation-commutativity} would force this
commutator to vanish if \(X^D\) were the exact first moment of a normalized
classical POVM supported in \(B_{\Rmom}^4\).  This contradiction proves the
no-go theorem, including the sphere-supported special case.
\end{proof}

\subsection{Exact classical moment radius on the finite coordinate atlas}
\label{p2c:app:optimal-classical-moment-radius}
This section supplies mathematical POVMs and a radius-optimality certificate.
It adds no apparatus, preparation, physical metric or occurrence premise.
Throughout, $\Rmom>0$, $q\in S^3(1)$, and $C(q)=\sum_aq_aC_a$.
Use probability Haar measure $d\sigma$ and the sector-phase gauge
with nonnegative magnitudes $r_+,r_-,\rho$.

\begin{theorem}[Classical moment radius for the full coordinate atlas]
\label{p2c:thm:full-atlas-povm-radius}
For a compatible quadruple with scalar edge $\Rmom/2$ and reduced magnitudes
$r_+,r_-,\rho$, a normalized POVM with exact first moments exists on the
closed ball of radius $r>0$ if and only if
\[
 r\ge\max\left\{
        \sqrt{\Rmom^2+6\max(r_+^2,r_-^2)},\,4\rho\right\}.
\]
A covariant POVM on the boundary sphere attains every admissible radius.
\end{theorem}
\begin{proof}
Any ball-valued POVM can be moved to the sphere while preserving its first
moments: replace $x\ne0$ by $\pm r x/|x|$ with weights
$(1\pm|x|/r)/2$, and replace $0$ by uniform sphere measure.
Average over rotations and matching unitary conjugations. Target covariance
preserves the first moments. Haar disintegration on
$\Spin(4)/\Spin(3)_{\rm diag}$ gives a positive stabilizer-invariant seed without assuming absolute
continuity of the original POVM. Explicitly, let $G=\Spin(4)$,
$H=\Spin(3)_{\rm diag}$, and choose a Borel section $s(q)e_4=q$.
Regard the sphericalized measure $E(dq)$ as a measure on unit directions,
with outcome vector $rq$. Its averaged seed is
\[
 K_0=\int_{S^3}\int_H U_{h s(q)^{-1}}E(dq)U_{h s(q)^{-1}}^*\,dh.
\]
It is positive and commutes with $U(H)$. Writing
$g=s(p)h s(q)^{-1}$ in the group average gives the density
$K(p)=U_{s(p)}K_0U_{s(p)}^*$ against $d\sigma(p)$.
Commutation with $H$ makes this expression independent of the section.
The stabilizer decompositions are
$\mathcal H_{\rm geo}|_H=2V_0\oplus3V_1$ and
$\mathcal H_{\rm sp}|_H=2V_{1/2}$; their inequivalent integer and half-integer
types exclude cross blocks in the seed.

On the two singlet copies and three triplet copies of the geometric
stabilizer representation, normalization and the first moments give seeds
\[
 S_0=\begin{pmatrix}1&b\\ b&a\end{pmatrix},\qquad
 T_0=\begin{pmatrix}B&d_+&d_-\\d_+&1&c\\d_-&\bar c&1\end{pmatrix},
\quad
 a+3B=4,
\]
where the triplet order is tangent vector, positive adjoint, negative
adjoint. Haar normalization fixes each irreducible diagonal trace;
the real sector-phase gauge and prescribed first moments fix
\[
 b=\frac{2\Rmom}{r},\qquad d_\pm=\frac{2\sqrt2\,r_\pm}{r}.
\]
Indeed, the vector-to-adjoint seed field is
$\sqrt2(d_+P_++d_-P_-)W_q$, whose first moment is
$r\sqrt2(d_+P_++d_-P_-)W_a/4$; the scalar-vector moment is $rb/4$.
These identities give the stated coefficients directly.
Positivity implies $a\ge b^2$ and $B\ge d_\pm^2$, hence
$r^2\ge\Rmom^2+6\max(r_+^2,r_-^2)$.
On the two chiral spinor copies, normalization fixes both diagonal
entries to one, and the first moment fixes the off-diagonal entry to
$4\rho/r$. Thus the protected seed is $I_4+(4\rho/r)C_4$, so positivity also requires
$r\ge4\rho$.

Conversely, choose $a=b^2$, $B=(4-b^2)/3$, and
$c=d_+d_-/B$ when $B>0$. The Schur complement of $B$ in $T_0$ is
$\operatorname{diag}(1-d_+^2/B,1-d_-^2/B)\ge0$.
If $B=0$, the inequalities force $d_\pm=0$; use $c=0$.
Thus both geometric seeds and the protected seed are positive.

For an explicit covariant extension on
$\mathbb C\oplus\mathbb C^4\oplus\Lambda^2\mathbb C^4$ set
\[
 W_qv=q\wedge v,\quad P_\pm=(I\pm *)/2,\quad
 H_q=2W_qW_q^*-I,\quad
 D(q)=\sqrt2(d_+P_++d_-P_-)W_q.
\]
Use
\[
 K_{\rm geo}(q)=
 \begin{pmatrix}
 1&bq^*&0\\
 bq&B I+(a-B)qq^*&D(q)^*\\
 0&D(q)&I+cH_q
 \end{pmatrix},
 \qquad
 K_{\rm sp}(q)=I_4+(4\rho/r)C(q).
\]
Set $K(q)=K_{\rm geo}(q)\oplus K_{\rm sp}(q)$ and define
$E(B)=\int\mathbf1_B(rq)K(q)\,d\sigma(q)$.
The seed at $q=e_4$ is the positive seed above.
The identities
$\int q_aq_b\,d\sigma=\delta_{ab}/4$ and
$\int W_qW_q^*\,d\sigma=I/2$
prove normalization and every prescribed first moment for outcomes $rq$.
Covariance proves positivity at every $q$. Zero and unequal adjoint edges
are included, and the $B=0$ branch covers the scalar-only endpoint $r=\Rmom$.
In terms of the radius entries $Q_\pm=2r_\pm^2$ and $Q_{\rm sp}=4\rho^2$,
$r_{\min}^2=\max\{\Rmom^2+3\max(Q_+,Q_-),4Q_{\rm sp}\}$.
This is a moment-realizability invariant, not a physical metric.
\end{proof}

\begin{proof}[Proof of the optimal intrinsic radius and an explicit endpoint]
On the intrinsic point $r_+=r_-=\sqrt2\Rmom/3$ and
$\rho=\sqrt3\Rmom/4$, the two bounds are
$\sqrt{7/3}\Rmom$ and $\sqrt3\Rmom$.
For an explicit endpoint density one may instead choose
\[
 b=2/\sqrt3,\quad a=4/3,\quad B=8/9,\quad
 c=1/3,\quad D(q)=(4\sqrt6/9)W_q,\quad K_{\rm sp}=I_4+C(q).
\]
The scalar/radial block has eigenvalues $0,7/3$.
Each tangent/two-form block has eigenvalues $0,20/9$; the remaining
three two-form eigenvalues are $2/3$. The protected eigenvalues are $0,2$,
each twice. Thus this density is positive.
With outcomes $x=\sqrt3\Rmom q$, its first-moment coefficients are
$\Rmom/2,\sqrt2\Rmom/3,\sqrt3\Rmom/4$, exactly those of $X^D$.

Independently, compress any ball-valued POVM to the protected block:
$E_{\rm sp}=P_{\rm sp}EP_{\rm sp}$ and $\int E_{\rm sp}=I_4$.
Its prescribed moments and Clifford traces give
\[
 4\sqrt3\Rmom
 =\int\operatorname{Tr}_4((x\cdot C)E_{\rm sp}(dx))
 \le4r,
\]
because $x\cdot C\le rI_4$. This proves the universal lower bound directly.
Using outcomes $\Rmom q$ in the endpoint density, mixing with the identity
density and reflecting $q$ gives every $|t|\le1/\sqrt3$; the same trace
bound proves necessity. The endpoint is not unique: with the same intrinsic $a,b,B,D$, every
$1/3\le c\le1$ also gives a positive normalized density with the same
moments. No physical instrument is selected.
\end{proof}

\begin{corollary}[Finite Hurwitz-support attainment]
\label{p2c:cor:optimal-hurwitz-povm}
For the intrinsic endpoint density, the 24 effects $E_h=K(h)/24$ on the
unit Hurwitz support
$H_{24}=\{\pm e_a:1\le a\le4\}\cup
\{(\pm1,\pm1,\pm1,\pm1)/2\}$, with outcomes $x_h=\sqrt3\Rmom h$, form a normalized
POVM with exact moments $X_a^D$. More generally, the polynomial density in
Theorem~\ref{p2c:thm:full-atlas-povm-radius} gives a 24-outcome POVM at every
admissible radius by the same quadrature.
\end{corollary}
\begin{proof}
The 24-point support has the sphere's second moments $1/4$ and fourth
moments $1/8$ and $1/24$, with odd-coordinate moments zero by reflection.
It is therefore exact through degree five. Positivity holds pointwise, while
normalization and first moments have polynomial degree at most three, so
their exact Haar identities persist under the finite sum.
The fixed finite support is not invariant under the entire continuous
$\Spin(4)$ group; full covariance is that of the transported support family.
Higher harmonic sampling can alias and is not covered by the continuous
all-band kernel statement below. The finite map agrees with its continuous
counterpart through degree two. Its degree-five cubature does not replace
the $2I$ degree-eight exactness used for quartic degree-two comparisons,
and no claim of a minimal number of outcomes is made.
\end{proof}

\begin{corollary}[Positive covariant map with the sharp harmonic kernel]
\label{p2c:cor:optimal-povm-harmonic-kernel}
For the explicit intrinsic endpoint density $K(q)$ above, define
$\Phi(f)=\int f(q)K(q)\,d\sigma(q)$.
On $C(S^3(1))$ this is a unital completely positive covariant map.
Writing $\mathcal P(m)=\bigoplus_{k=0}^mE_k$ for unit-sphere scalar
harmonics, for $m\ge2$,
\[
 \operatorname{rank}\Phi|_{\mathcal P(m)}=14,\qquad
 \ker\Phi|_{\mathcal P(m)}=\bigoplus_{k=3}^m E_k.
\]
\end{corollary}
\begin{proof}
The density is polynomial of degree at most two, so harmonic orthogonality
annihilates all higher degrees. The constant maps to $I$, first harmonics
map nontrivially to the intrinsic coordinates, and
\[
 \Phi(q_1^2-q_2^2)|_V
       =\frac1{27}\operatorname{diag}(1,-1,0,0)\ne0.
\]
Here $V=V_{1/2,1/2}$ is the four-dimensional scalar-vector sector.
The three source types are irreducible and inequivalent, so their nonzero
images are injective with direct sum of dimensions $1+4+9$.
Positivity, complete positivity and normalization follow from integration
against a positive density. The map is not multiplicative: its first moments are noncommuting
intrinsic coordinates. Existence does not actualize a physical assignment.
This exact all-band kernel statement concerns the continuous density;
finite Hurwitz sampling can alias higher harmonics.
\end{proof}



\subsubsection{Eight outcomes and the full covariant intrinsic endpoint domain}
\label{p2c:app:v233-eight-outcomes}

\begin{corollary}[Quaternion-support attainment and its subgroup domain]
\label{p2c:cor:v233-optimal-quaternion-povm}
Let \(K_c(q)\) be the intrinsic endpoint density below.
For \(q\in Q_8=\{\pm e_a:1\le a\le4\}\), the eight effects and outcomes
\[
 E_q=\tfrac18K_c(q),\qquad x_q=\sqrt3\Rmom q
\]
form a normalized POVM with exact first moments \(X_a^D\).
Every noncyclic finite subgroup supplies such an attaining rule;
no cyclic subgroup can support exact first moments proportional
to this nondegenerate quadruple at any nonzero scale.
Among subgroup supports the minimum cardinality is eight, attained
by \(Q_8\), but there is no subgroup-least class.
This makes no minimum-cardinality claim for arbitrary finite POVMs.
The density of Theorem~\ref{p2c:thm:full-atlas-povm-radius}
likewise admits \(Q_8\) quadrature at every admissible radius.
\end{corollary}
\begin{proof}
The densities have degree at most two, so normalization and
first moments require cubature only through degree three.
All noncyclic finite subgroups are 3-designs by
Theorem~\ref{p1b:thm:v233-subgroup-sampler-classification}
and its proof. Positivity holds at each node.
A cyclic group's directions lie in a two-plane, so any vector
normal to it gives a zero linear combination of first moments.
On the protected block, however,
\((\sum v_a C_a)^2=|v|^2I\), so the four native coordinates
are linearly independent. This is a contradiction for \(v\ne0\).
The finite-subgroup classification gives the cardinality claim.
The incompatible \(Q_8\) and \(2D_3\) inclusions exclude a least
class as in the cited theorem. A fixed embedded support is not
invariant under all of \(\Spin(4)\); covariance applies to its
transported family.
\end{proof}

\begin{theorem}[Complete covariant seeds on the optimal intrinsic sphere]
\label{p2c:thm:complete-intrinsic-endpoint-seeds}
Fix the intrinsic tuple, its real sector-phase gauge, and the outcome sphere
of radius $r=\sqrt3\Rmom$.  All normalized $\Spin(4)$-covariant POVMs on
this sphere with first moments $X_a^D$ are obtained from stabilizer seeds
\[
 S_0=\begin{pmatrix}1&b\\b&a\end{pmatrix},\qquad
 T_0=\begin{pmatrix}B&d&d\\d&1&\zeta\\d&\bar\zeta&1\end{pmatrix},
 \qquad K_{{\rm sp},0}=I_4+C_4,
\]
where $b=2/\sqrt3$, $d=4/(3\sqrt3)$, $B=(4-a)/3$, and precisely
\[
 \boxed{\frac43\le a\le\frac{20}{9},\qquad
         |\zeta-u|\le1-u,\qquad
         u=\frac{16}{9(4-a)},\quad\zeta\in\mathbb C.}
\]
Here $S_0$ acts on the two stabilizer singlets and $T_0$ on the multiplicity
space of the three stabilizer triplets.  The domain has three real
parameters in its interior.  It classifies covariant POVMs on the boundary
sphere; it makes no assertion of uniqueness for instruments or for arbitrary
ball-valued measures with additional interior support.
\end{theorem}
\begin{proof}
Use the full-atlas seed decomposition in
Theorem~\ref{p2c:thm:full-atlas-povm-radius}.  Haar normalization fixes
$a+3B=4$, and the intrinsic first moments fix $b,d$ and the protected seed.
Integer and half-integer stabilizer types exclude geometric/protected cross
blocks.  These exhaust the commuting stabilizer blocks, so no seed parameter
has been discarded.
Positivity of $S_0$ is $a\ge b^2=4/3$.
Positivity of $T_0$ forces $B>0$ because $d\ne0$.
Its Schur complement at $B$ is
\[
 \begin{pmatrix}1-u&\zeta-u\\\bar\zeta-u&1-u\end{pmatrix},
 \qquad u=d^2/B.
\]
It is positive if and only if $u\le1$ and
$|\zeta-u|\le1-u$; $u\le1$ is $a\le20/9$.
Thus the displayed region is necessary and sufficient.
Transport by a section $s(q)e_4=q$ gives
$K(q)=U_{s(q)}K_0U_{s(q)}^*$, independently of the section, because the
seed commutes with the stabilizer.  The covariant-measure disintegration
used in the atlas theorem proves that every covariant POVM has this form.

For an explicit polynomial representative set $W_qv=q\wedge v$,
$H_q=2W_qW_q^*-I_6$, and use the oriented Hodge star $*$.  Then
\[
 K_{a,\zeta}(q)=
 \begin{pmatrix}
 1&bq^*&0\\
 bq&B I_4+(a-B)qq^*&(4\sqrt6/9)W_q^*\\
 0&(4\sqrt6/9)W_q&I_6+(\operatorname{Re}\zeta)H_q
                  +i(\operatorname{Im}\zeta)*H_q
 \end{pmatrix}\oplus(I_4+C(q)).
\]
The relative self-dual/anti-self-dual basis is chosen so its north-pole
off-diagonal seed entry is $\zeta$; reversing that basis convention
conjugates $\zeta$.  Since $*H_q=-H_q*$, the displayed matrix is Hermitian.
At $q=e_4$ its stabilizer blocks are exactly $S_0,T_0$ and $I_4+C_4$.
The Haar identities in the atlas theorem prove normalization and the exact
first moments; covariance then proves positivity everywhere.
\end{proof}
The previously displayed real family below is the slice
$a=4/3$, $\zeta=c\in[1/3,1]$.
Even requiring a real seed leaves the two-dimensional domain
$4/3\le a\le20/9$, $2u-1\le\zeta\le1$.
The quaternion and Hurwitz quadratures apply to the full domain, because
the density still has degree at most two.  This freedom changes the
classical readout map, rather than the generated coordinate class.

\begin{proposition}[A real endpoint slice with fixed native first moments]
\label{p2c:prop:v233-endpoint-map-family}
In the intrinsic real sector-phase gauge, put
\(W_qv=q\wedge v\), \(H_q=2W_qW_q^*-I_6\),
\(C(q)=\sum_aq_aC_a\), and
\[
 K_c(q)=
 \begin{pmatrix}
 1&(2/\sqrt3)q^*&0\\
 (2/\sqrt3)q&(8/9)I_4+(4/9)qq^*&(4\sqrt6/9)W_q^*\\
 0&(4\sqrt6/9)W_q&I_6+cH_q
 \end{pmatrix}\oplus(I_4+C(q)).
\]
For real \(c\), this is positive at every point if and only if
\(1/3\le c\le1\). Throughout that interval,
\[
 \int K_c\,d\sigma=I_{15},\qquad
 \sqrt3\Rmom\int q_aK_c(q)\,d\sigma=X_a^D .
\]
The maps \(\Phi_c(f)=\int fK_c\,d\sigma\) have all the properties
in Proposition~\ref{p2c:prop:v233-unit-and-positive-maps}.
For \(f=q_1^2-q_2^2\), their vector and two-form blocks are
\[
 \Phi_c(f)|_V=\frac1{27}\operatorname{diag}(1,-1,0,0),\qquad
 \Phi_c(f)|_{\Lambda^2V}
 =\frac c6\operatorname{diag}(0,1,1,-1,-1,0),
\]
in the two-form order \(12,13,14,23,24,34\).
\end{proposition}
\begin{proof}
Rotate \(q\) to one coordinate axis. The scalar/radial block has
eigenvalues \(0,7/3\). Each of the three tangent/two-form blocks is
\[
 \begin{pmatrix}8/9&4\sqrt6/9\\4\sqrt6/9&1+c\end{pmatrix},
 \qquad \det=\frac{8(3c-1)}{27}.
\]
The remaining two-form eigenvalues are \(1-c\); the protected
eigenvalues are \(0,2\), each twice. These conditions give exactly
the stated interval. The Haar identities
\(\int qq^*\,d\sigma=I_4/4\) and
\(\int W_qW_q^*\,d\sigma=I_6/2\), with odd moments zero,
give normalization and all three native edge coefficients
\(\Rmom/2,\sqrt2\Rmom/3,\sqrt3\Rmom/4\).
Equivariance of the wedge and Clifford operations proves covariance.
Integration against a positive matrix density is completely
positive and unital.

The fourth moments \(1/8\) and \(1/24\) give the two displayed
quadratic images; in particular changing \(c\) changes the map.
The density has degree at most two, so it annihilates every
\(E_k\), \(k\ge3\). The scalar image is nonzero, the first
harmonics have nonzero native images, and the displayed vector
quadratic image is nonzero for every \(c\). These are three
inequivalent irreducible types. Their equivariant images are
injective and have direct sum of dimensions \(1+4+9=14\).
Their first images do not commute, so the maps are not
multiplicative.
\end{proof}

\begin{corollary}[Invariant-state scalar readout is Haar and map-independent]
\label{p2c:cor:v340-s3lit-invariant-povm}
For every positive unital $\operatorname{Spin}(4)$-covariant map
$\Phi_c:C(S^3)\to\operatorname{End}(\mathcal H_1)$ in
Proposition~\ref{p2c:prop:v233-endpoint-map-family} and any
$\operatorname{Spin}(4)$-invariant density operator $\rho$,
\[
 \operatorname{Tr}\bigl(\rho\Phi_c(f)\bigr)
 =\int_{S^3}f(q)\,d\mu_{\rm Haar}(q),
 \qquad f\in C(S^3).
\]
For the maximally mixed $\rho=I_{15}/15$, the explicit kernel of
that proposition has $\operatorname{Tr}K_c(q)=15$ at each $q$.
This does not choose a member of the localization-map family for
anisotropic charged states, or identify any such map with a
position-space physical electromagnetic source distribution.
\end{corollary}
\begin{proof}
Positivity and unitality give a normalized Borel probability measure
for $f\mapsto\operatorname{Tr}(\rho\Phi_c(f))$.
Covariance and invariance of $\rho$ make this measure invariant
under the transitive $\operatorname{Spin}(4)$ action on $S^3$;
uniqueness of invariant probability gives normalized Haar measure.
The displayed $15$-trace follows directly from the block diagonal
traces in the source-owned density $K_c(q)$.
\end{proof}

\begin{proposition}[Ambient chiral momentum--position comparison algebra]
\label{p2c:prop:v340-s3lit-chiral-algebra}
On the existing smooth $\mathrm{SU}(2)\cong S^3_p$ momentum parent,
let $p_i=M(q_i)$ and $p_0=M(q_0)$ be commuting quaternion-coordinate
multipliers for $q_0^2+\sum_iq_i^2=\Rmom^2$.
In the regular left/right generator convention with
$J_i=\hbar(L_i^++L_i^-)$ and
$P_i=(\hbar/\Rmom)(L_i^+-L_i^-)$, put
$\ell=\hbar/\Rmom$ and
$X_i^+=(\ell/\hbar)(J_i+\Rmom P_i)=2\ell L_i^+$.
Choose the orientation of the left regular action consistently with
the quaternion-product convention.  On its smooth common core,
\[
 [X_i^+,X_j^+]=2i\ell\epsilon_{ijk}X_k^+,
 \quad[p_i,p_j]=0,
 \quad[X_i^+,p_j]=i\ell(\epsilon_{ijk}p_k-\delta_{ij}p_0),
\]
with $p_0^2+\sum_ip_i^2=\Rmom^2I$.  This is an ambient
algebraic comparison with the $\mathrm{SU}(2)$ phase-space bracket in
Kupriyanov--Kurkov--Sharapov, \emph{not} an action-selected physical
Hamiltonian, the intrinsic 15-state coordinates $X_a^D$, or the
noncommuting transvection operator $P_i$ itself.  In particular, no
nonzero exactly commuting momentum triple transforming as a full
continuous $\mathrm{SO}(3)$ vector exists \emph{within} a finite
Hilbert carrier.
\end{proposition}
\begin{proof}
The existing $\mathrm{Spin}(4)$ regular chiral commutators and their
action on quaternion coordinates give the displayed algebra by direct
substitution; $p_0^2+\sum_i p_i^2=\Rmom^2I$ holds by multiplication.
For the finite-carrier claim, three commuting Hermitian matrices have
a finite joint spectrum.  If they are covariant under all continuous
rotations, this spectrum is rotationally invariant.  A nonzero joint
eigenvalue has an infinite rotation orbit, impossible for a finite
spectrum; hence all joint eigenvalues vanish.
\end{proof}


The all-band kernel assertion concerns the continuous maps.
Eight-node quadrature preserves normalization and first moments,
but need not preserve \(\Phi_c\) on quadratic tests, whose
integrands can have degree four. The original Hurwitz 5-design
does preserve that quadratic comparison. Neither finite rule
replaces the \(2I\) degree-eight requirement for quartic
degree-two fields.

\begin{corollary}[Map-independent total readout noise at the optimal sphere]
\label{p2c:cor:optimal-endpoint-total-noise}
For any normalized POVM on $S^3(\sqrt3\Rmom)$ with exact first moments
$X_a^D$, let $\Phi$ denote its integration map on classical outcome
functions.  Its total noise is
\[
 \boxed{N:=\sum_a[\Phi(x_a^2)-\Phi(x_a)^2]
        =3\Rmom^2I-Q_D\succeq2\Rmom^2I.}
\]
In the ordered sectors of
Theorem~\ref{p2c:thm:intrinsic-tangential-coordinate-class}, the eigenvalues
of $N/\Rmom^2$ are
\[
 2,\quad \frac{25}{12},\quad\frac{23}{9},\quad\frac{23}{9},
 \quad\frac94,\quad\frac94.
\]
\end{corollary}
\begin{proof}
Integration against a POVM is unital completely positive, so each
$\Phi(x_a^2)-\Phi(x_a)^2$ is positive.  The sphere identity
$\sum_a x_a^2=3\Rmom^2$ and the fixed first moments give the equality.
Subtract the native quadratic fingerprint to obtain the entries and bound.
\end{proof}
This readout noise is distinct from the operator variance of
Theorem~\ref{p2c:thm:intrinsic-total-variance}.

\begin{proposition}[Fixed-map ranks and family-independent readout]
\label{p2c:prop:endpoint-fixed-map-qualification}
For every seed in
Theorem~\ref{p2c:thm:complete-intrinsic-endpoint-seeds}, its continuous
integration map has rank fourteen on $\mathcal P(m)$, $m\ge2$, with kernel
$\bigoplus_{k=3}^mE_k$.  This does not make the maps equal.  Across the full
family, the scalar symbols on which all maps agree are exactly
\[
 E_0\oplus E_1\oplus\bigoplus_{k=3}^mE_k;
\]
their common nonzero readout has dimension five.
Native operator words in $X_a^D$ are fixed throughout.  A classical
quadratic image $\Phi(x_ax_b)$ must retain its chosen map, or a proof of
independence from that map.
\end{proposition}
\begin{proof}
The polynomial density annihilates $E_k$ for $k\ge3$.  Constants and first
harmonics have their nonzero fixed images.  The vector block gives
\[
 \Phi_{a,\zeta}(q_1^2-q_2^2)|_V
 =\frac{a-B}{12}\operatorname{diag}(1,-1,0,0)\ne0,
\]
because $a-B=(4a-4)/3>0$.  Equivariance and the three inequivalent
irreducible types give rank $1+4+9=14$.  On the existing real slice,
the difference of two distinct maps on $E_2$ is a nonzero equivariant
map, by the two-form image in
Proposition~\ref{p2c:prop:v233-endpoint-map-family}; irreducibility makes
it injective.  Thus no nonzero degree-two symbol has a common value
across the family.  Higher degrees vanish and the first two types agree,
giving exactly the displayed common-symbol space.
\end{proof}

\begin{theorem}[Full quadratic endpoint-POVM exactness if and only if fourth-degree cubature]
\label{p2c:thm:v268-endpoint-quadratic-iff-design4}
Fix any covariant positive endpoint density $K_{a,\zeta}(q)$ from
Theorem~\ref{p2c:thm:complete-intrinsic-endpoint-seeds} on the unit
outcome sphere and any finitely supported normalized measure
$\nu=\sum_{z\in Z}w_z\delta_z$ ($w_z\ge0$,
$\sum_zw_z=1$).  Write
$\Phi_{a,\zeta}^{\nu}(f)=\int f(q)K_{a,\zeta}(q)\,d\nu(q)$ and
$\Phi_{a,\zeta}^{H}(f)=\int f(q)K_{a,\zeta}(q)\,d\sigma(q)$.
Then the following are equivalent:
\begin{enumerate}[label=(\roman*),leftmargin=2em]
\item $\Phi_{a,\zeta}^{\nu}(f)=\Phi_{a,\zeta}^{H}(f)$ for
      \emph{every} $f\in\mathcal P_2$;
\item $\nu$ integrates every restricted polynomial of degree at most
      four exactly, i.e. $\nu$ is a spherical $4$-design in the
      positive weighted sense;
\item evaluation with weights $w_z$ is an isometry on
      $\mathcal P_2$.
\end{enumerate}
The equivalence holds for \emph{every} admissible endpoint seed,
not merely the one-parameter real slice.  Thus exact first moments
on a three-design do not imply exact quadratic-symbol readout.
\end{theorem}
\begin{proof}
The density is polynomial of degree at most two.  Consequently
(ii) implies (i).  Conversely, its scalar block is identically
$K_{a,\zeta}(q)_{00}=1$, so equality in (i) implies
$\nu(f)=\sigma(f)$ for every $f\in\mathcal P_2$.
Its vector block is
\[
 K_{a,\zeta}(q)|_V=B I_4+(a-B)qq^*,
 \qquad B=\frac{4-a}{3},\qquad
 a-B=\frac{4(a-1)}3>0
 \quad(4/3\le a\le20/9).
\]
Equality of the vector-block matrix entries for every $f\in\mathcal P_2$,
after subtracting the already equal scalar $B\nu(f)$ terms, gives
$\nu(fq_iq_j)=\sigma(fq_iq_j)$ for all $i,j$ and all $f\in\mathcal P_2$.
Such products, together with $\mathcal P_2$, span $\mathcal P_4$
by splitting coordinate monomials into a degree-at-most-two factor
and a quadratic factor.  This proves (ii).  Equivalence of (ii) and
(iii) is Lemma~\ref{p1b:lem:v233-weighted-band-products}.
\end{proof}

\begin{corollary}[Universal factor-three quaternionic quadratic-readout defect]
\label{p2c:cor:v268-q8-endpoint-factor-three}
For $Q_8=\{\pm e_1,\ldots,\pm e_4\}$ with uniform weights and
$u(q)=q_1^2-q_2^2$, let $D=\operatorname{diag}(1,-1,0,0)$.
For every admissible $K_{a,\zeta}$,
\[
 \Phi_{a,\zeta}^{Q_8}(u)|_V=\frac{a-B}{4}D
   =3\frac{a-B}{12}D=3\Phi_{a,\zeta}^{H}(u)|_V.
\]
On the real slice $a=4/3$ this is $D/9\ne D/27$.
In contrast, every $2T$ uniform sample reproduces the continuous
quadratic readout exactly for all $f\in\mathcal P_2$ and all endpoint
seeds.  This obstruction does not invalidate the $Q_8$ eight-fibre
sampling isometry or its normalized first-moment POVM.
\end{corollary}
\begin{proof}
Use $\sigma(u)=\tau_{Q_8}(u)=0$,
$\sigma(uqq^*)=D/12$, and
$\tau_{Q_8}(uqq^*)=D/4$.  The scalar $B I_4$ term therefore vanishes
in both integrals.  The $2T$ assertion follows from its strength-five
design and the preceding theorem.
\end{proof}



\subsection{Exact exceptional frames, POVM barycentres and projector firewalls}
\label{p2f:sec:v273-exceptional-superposition}

\begin{theorem}[Exact finite-fibre barycentre of a continuous endpoint POVM]
\label{p2f:thm:v273-finite-fibre-povm-barycentre}
Let $\sigma$ be normalized Haar measure on $G=SU(2)$ and let
$K(q)=K_{a,\zeta}(q)\ge0$ be an admissible degree-two
native endpoint-POVM density satisfying $\int_GK\,d\sigma=I_{15}$.
For a finite subgroup $\Gamma\le G$ of spherical-design strength
$t_\Gamma\ge2$, and a coset $[g]\in G/\Gamma$, set
\[
 E^{[g]}_\Gamma(B)=\frac1{|\Gamma|}\sum_{\gamma\in\Gamma}
              \mathbf1_B(g\gamma)K(g\gamma),\qquad B\subset G
              \ \text{Borel}.
\]
Each $E^{[g]}_\Gamma$ is a positive, normalized finite-outcome POVM.
For the invariant quotient measure $\bar\sigma$ one has
\[
 \boxed{\int_{G/\Gamma}E^{[g]}_\Gamma(B)
              \,d\bar\sigma([g])
       =\int_BK(q)\,d\sigma(q)=E(B)}
\]
for \emph{all} Borel outcome events. On each individual coset,
all polynomial endpoint moments $\int f\,dE^{[g]}_\Gamma$ agree with
the continuous moments whenever $\deg f\le t_\Gamma-2$.
The barycentre is a convex/classical mixture of mathematical POVMs,
not a coherent superposition of finite quantum states or a newly
constituted physical measurement instrument.
\end{theorem}
\begin{proof}
The formula is independent of the representative of $[g]$ because
$g\mapsto g\gamma_0$ permutes the support. Positivity is termwise.
The density $K$ has polynomial degree two; the coset design identity
then gives $|\Gamma|^{-1}\sum_\gamma K(g\gamma)=I_{15}$,
and similarly gives the moment assertion for $fK$.
Haar disintegration into left cosets reads
\[
 \int_GF(q)\,d\sigma(q)=\int_{G/\Gamma}
  \frac1{|\Gamma|}\sum_{\gamma\in\Gamma}F(g\gamma)
  \,d\bar\sigma([g]).
\]
Apply this entrywise with $F=\mathbf1_BK$. The infinite-dimensional
$L^2(G)$ disintegration should be understood as the associated
rank-$|\Gamma|$ fibre bundle, since a change of coset representative
permutes its fibre labels; it is not an identification with the native
15-dimensional $\mathcal H_1$.
\end{proof}

\begin{theorem}[Twin-icosahedral Reynolds mean and exact non-idempotence]
\label{p2f:thm:v273-twin-projector-spectral-identity}
Fix one embedded $H\cong2T$, put $O=N_G(H)\cong2O$, and let
$I_+,I_-$ be the two icosahedral completions of
Theorem~\ref{p1:thm:v273-two-icosahedral-completions}.
On $V=\operatorname{Sym}^{12}(\mathbb C^2)$, of dimension thirteen,
write $P_\Gamma=|\Gamma|^{-1}\sum_{\gamma\in\Gamma}\rho_{12}(\gamma)$.
Then
\[
 \dim V^H=2,\quad\dim V^O=1,\quad \dim V^{I_+}=\dim V^{I_-}=1,
\]
\[
 \boxed{A_H:=\frac12(P_{I_+}+P_{I_-})
 =\frac{11}{32}P_O+\frac{21}{32}(P_H-P_O)
 =\frac{21P_H-10P_O}{32},}
\]
\[
 \boxed{A_H^2-A_H=-\frac{231}{1024}P_H.}
\]
The nonzero eigenvalues are $11/32$ and $21/32$ on the
$O$-even and $O$-odd one-dimensional lines inside $V^H$.
On the full scalar harmonic $E_{12}=V\otimes V^*$ each nonzero
eigenvalue has multiplicity $13$; zero has multiplicity $143$.
This is a canonical operator \emph{relative to a given $H$},
not a projection or a physical branch-selection device.
\end{theorem}
\begin{proof}
The character averages for the Hurwitz, octahedral, and icosahedral
groups in the degree-twelve $SU(2)$ irrep give the stated invariant
ranks $2,1,1$. Let $v$ be a unit $I_+$-fixed vector and let
$a=(1+\mathbf i)/\sqrt2\in O\setminus H$, so that
$v_-=\rho_{12}(a)v$ spans the $I_-$-fixed line.
The exact 600-cell quaternionic character sum is
\[
 \langle v,v_-\rangle=
 \frac1{120}\sum_{\gamma\in I_+}
 U_{12}(\operatorname{Re}(a\gamma))=-\frac5{16}.
\]
Since $a^2\in H$, its action on the $H$-fixed plane is an
involution, with $O$-even line and the orthogonal $O$-odd line.
The rank-one mean of the projectors onto $v$ and $v_-$ has eigenvalues
$(1-5/16)/2=11/32$ and $(1+5/16)/2=21/32$ on these lines, and
vanishes on $(V^H)^\perp$.
Squaring gives $\lambda(1-\lambda)=231/1024$ on either line.
The scalar $E_{12}$ has an independent right factor of dimension $13$.
\end{proof}

\begin{corollary}[Unmarked twirl gives a scalar trace fraction, not a sharp projector]
\label{p2f:cor:v273-icosahedral-projector-twirl}
On $E_k=V_{k/2}\otimes V_{k/2}^*$ let $P_I$ be the one-sided
binary-icosahedral Reynolds projector and
$\mu_k^I=\dim (V_{k/2})^I$. With $U(g)$ the left $SU(2)$ action,
\[
 \boxed{\int_GU(g)P_IU(g)^*\,d\sigma(g)
           =\frac{\mu_k^I}{k+1}I_{E_k}.}
\]
In particular, at $k=12$, the unmarked twirl is $I_{E_{12}}/13$
(rank $169$), whereas every individual $P_I|_{E_{12}}$ has rank $13$.
For scalar covariance $C_k=c_kI$, the trace fraction
$\operatorname{Tr}(P_IC_kP_I)/\operatorname{Tr}C_k
=\mu_k^I/(k+1)$ is independent of the orientation of $I$.
This identity does not perform or generate the rank-thirteen
projection itself.
\end{corollary}
\begin{proof}
The Haar twirl commutes with the irreducible left $SU(2)$ action;
Schur's lemma makes it scalar on $V_{k/2}$. Its trace there is
$\mu_k^I$, fixing the scalar $\mu_k^I/(k+1)$.
The unaffected right multiplicity factor gives the full-shell formula.
The covariance fraction follows from $P_I^2=P_I$.
\end{proof}

\begin{remark}[Sharp distinction between cubature and invariant-sector readouts]
\label{p2f:rem:v273-finite-readout-versus-projection}
The positive nineteen-design of
Theorem~\ref{p1b:thm:v273-triple-icosahedral-nineteen-design}
integrates polynomials through degree nineteen exactly, but is not
a substitute for any embedded $2I$ Reynolds projector. In fact
$P_I|_{E_{12}}$ has a nonzero rank-thirteen image, whereas Haar
averaging annihilates $E_{12}$. A physical sparse-harmonic comb,
frame permutation, ten-axis Clifford incidence, or repeated finite
convolution requires the corresponding generated operator \emph{and}
its typed action-consumer; it does not follow from finite cubature
or a projector's mere mathematical availability. Conversely,
continuous endpoint-POVM statistics, scalar Dirichlet forms and
bounded-band native $\mathcal W_2^D$ symbols need no actual finite
support to have their exact mathematical finite-frame presentations.
In particular simultaneous sharp $I_+$ and $I_-$ invariance is
excluded for nonconstant scalar modes by
Corollary~\ref{p1:cor:v273-twins-sharp-common-invariants}.
\end{remark}




\section{Independent radius, spectral and finite-observable certificates}
\label{p2f:sec:v260-three-gap-closure}

These statements merge the independent POVM-radius certificate and exact
observable readout frontier from the earlier parallel branch. The main
signed-spectrum theorem is already present as
Theorem~\ref{p2f:thm:v260-signed-rs-s3}; the Casimir-difference
statement below is retained as a second, equivalent first-order proof.
Neither statement promotes an
available mathematical POVM to an occurring detector or change the constituted
level-one momentum carrier. All operator identities concern the declared
coordinate representative or the standard round spin connection on their
stated domains.

\begin{theorem}[Unrestricted positive-operator dual certificate for the coordinate moment radius]
\label{p2f:thm:v260-full-atlas-dual-radius}
Let \(R=\Rmom>0\) and fix a compatible self-adjoint coordinate tuple
\(X_a\) on
\(\CC e_0\oplus V\oplus\Lambda^2V\oplus\mathcal H_{\rm sp}\), where
\(V=\CC^4\), with the geometric scalar--vector edge \(R/2\),
vector--two-form edges
\(T_a=(r_+P_++r_-P_-)W_{e_a}\), and protected Clifford edge
\(\rho C_a\), all in the real sector-phase gauge. The scalar inputs
\(r_+,r_-,\rho\) are nonnegative and
\(P_\pm=(I\pm*)/2\) on \(\Lambda^2 V\). Any normalized POVM
\(E\) on \(\{x\in\RR^4:|x|\le r\}\) with exact first moments
\(\int x_aE(dx)=X_a\) satisfies
\[
 r\ge\max\{\sqrt{R^2+6r_+^2},\sqrt{R^2+6r_-^2},4\rho\}.
\]
The inequality applies to arbitrary POVMs, without covariance, absolute
continuity, or a prescribed number of outcomes. In view of the constructive
positive sphere density in
Theorem~\ref{p2c:thm:full-atlas-povm-radius}, it is sharp and proves the
same necessary-and-sufficient radius formula by an independent dual argument.
\end{theorem}
\begin{proof}
Fix one chirality \(\varepsilon\in\{+,-\}\), abbreviate
\(s=r_\varepsilon\), \(P=P_\varepsilon\), and set
\(S=(R^2+6s^2)^{1/2}\), \(t=S/2>0\).
In the geometric ordered blocks \(\CC\oplus V\oplus\Lambda^2V\), define
\[
 Y_a=\begin{pmatrix}
 0&Re_a^*&0\\
 Re_a&0&2sW_{e_a}^*P\\
 0&2sPW_{e_a}&0
 \end{pmatrix},\qquad
 Z=\begin{pmatrix}
 R^2/t&0&0\\0&tI_V&0\\0&0&(2s^2/t)P
 \end{pmatrix}.
\]
On \(V\), \(W_qv=q\wedge v\). For a unit \(q\), the chiral map
\(PW_q:q^\perp\to\Lambda^2_\varepsilon V\) has three singular
values \(1/\sqrt2\), and vanishes on \(q\). This follows from
\(W_q^\dagger W_q=I-qq^*\) and
\(W_q^\dagger\circ\star\circ W_q=0\), where \(\star\) is the
Hodge-star operator. Equivalently,
\(W_q^\dagger PW_q=(I-qq^*)/2\).
Thus \(Z-\sum_aq_aY_a\) decomposes into the scalar/radial block
\[
 \begin{pmatrix}R^2/t&-R\\-R&t\end{pmatrix}
\]
and three tangent/chiral blocks
\[
 \begin{pmatrix}t&-\sqrt2s\\-\sqrt2s&2s^2/t\end{pmatrix},
\]
plus zero blocks on the opposite chirality. Each displayed block is
positive semidefinite of determinant zero. The same holds for every
\(|q|\le1\), by writing an interior \(q\) as a convex combination of
opposite unit vectors; hence
\(Z-\sum_a(x_a/r)Y_a\succeq0\) throughout the outcome ball.
Compress the POVM to the geometric part and integrate the positive trace:
\[
 0\le\int\operatorname{Tr}_{\mathcal H_{\rm geo}}
  \left[\left(Z-\frac1r\sum_ax_aY_a\right)E_{\rm geo}(dx)\right]
 =\operatorname{Tr}_{\mathcal H_{\rm geo}}Z
 -\frac1r\sum_a\operatorname{Tr}_{\mathcal H_{\rm geo}}(Y_aX_a).
\]
The wedge identity
\(\operatorname{Tr}_V(W_{e_a}^*PW_{e_a})=3/2\)
now gives, by direct blockwise calculation,
\[
 \sum_a\operatorname{Tr}_{\mathcal H_{\rm geo}}(Y_aX_a)
 =4R^2+24s^2=4S^2,
 \qquad
 \operatorname{Tr}_{\mathcal H_{\rm geo}}Z
 =\frac{R^2+6s^2}{t}+4t=4S.
\]
Thus \(4S^2\le4rS\), proving \(r\ge S\) for either chirality.
For the protected compression \(E_{\rm sp}\), use the Clifford identity
\(C(x)^2=|x|^2I_4\), so \(C(x)\preceq rI_4\). Then
\[
 16\rho=\sum_a\operatorname{Tr}_{\CC^4}(C_aX_a|_{\rm sp})
 =\int\operatorname{Tr}_{\CC^4}(C(x)E_{\rm sp}(dx))
 \le r\operatorname{Tr}_{\CC^4}I_4=4r.
\]
This proves \(r\ge4\rho\). The positive density already constructed
in the cited theorem attains the maximum of these three sharp bounds,
including the vanishing-edge boundary cases.
\end{proof}

\begin{theorem}[First-order homogeneous Casimir-difference formula for the round Rarita--Schwinger operator]
\label{p2f:thm:v260-rs-signed-casimir}
Let \(S^3(R)=\Spin(4)/\Spin(3)_{\rm diag}\) have its round metric
of radius \(R>0\), and let \(\slashed D_{3/2}\) be the standard
self-adjoint, projected twisted Dirac operator on the gamma-traceless
vector-spinor bundle. Choose the left-invariant frame
\(e_a=(2/R)\mathscr R_a\) with \(\mathscr R_a\) generated by
\(t_a=-i\sigma_a/2\), and Clifford multiplication
\(\gamma_a=i\sigma_a\). Denote the positive quadratic Casimirs of
\(SU(2)_L,SU(2)_R\) by \(\mathcal C_L,\mathcal C_R\). On smooth
sections,
\[
 \boxed{\slashed D_{3/2}=-\frac{2}{3R}(\mathcal C_L-\mathcal C_R).}
\]
The two signed spectral series are therefore
\[
 \begin{array}{ll}
 \ker P^*:\quad \lambda_{k,\pm}^{(A)}=\pm(k+5/2)/R,
 &\operatorname{mult}_{k,\pm}^{(A)}=(k+1)(k+4),\\[2pt]
 \operatorname{Im}P:\quad \lambda_{k,\pm}^{(B)}=\pm(k+5/2)/(3R),
 &\operatorname{mult}_{k,\pm}^{(B)}=(k+2)(k+3),
 \end{array}\qquad k\ge0.
\]
With the stated convention, \(j_L>j_R\) yields the negative sign and
swapping the two factors yields the positive sign. The overall sign
reverses if the global Clifford/frame convention reverses; the signed
spectral multiset is unchanged.
\end{theorem}
\begin{proof}
The stabilizer fiber is the spin-\(3/2\) summand
\(V_1\otimes V_{1/2}=V_{3/2}\oplus V_{1/2}\). Use
\emph{skew-Hermitian} Lie generators \(T_a^{(1)}\) and
\(T_a^{(1/2)}\), normalized by
\(T_a^{(1/2)}=-i\sigma_a/2\), and put
\(T_a=T_a^{(1)}+T_a^{(1/2)}\) on the highest-spin summand.
The positive quadratic Casimir is \(-\sum_aT_a^2\), so its
spin-\(j\) eigenvalue is \(j(j+1)\). Equivariance forces the
projected spin-\(1/2\) generator to be proportional to the
spin-\(3/2\) generator. Contracting with \(T\) fixes the coefficient:
\[
 \Pi_{3/2}\sum_aT_a^{(1/2)}T_a\Pi_{3/2}
 =\tfrac12\bigl[-\tfrac{15}{4}-\tfrac34+2\bigr]\Pi_{3/2}
 =-\tfrac54\Pi_{3/2}.
\]
Since \(\sum_aT_a^2=-15/4\) on this summand, the ratio is \(1/3\):
\[
 \Pi_{3/2}T_a^{(1/2)}\Pi_{3/2}=\tfrac13T_a,
 \qquad \Pi_{3/2}\gamma_a\Pi_{3/2}=-\tfrac{2}{3}T_a.
\]
Here \(\gamma_a=i\sigma_a=-2T_a^{(1/2)}\), fixing
Clifford and connection signs consistently. This is also an exact
Clebsch--Gordan projector identity in the six-dimensional fiber,
independently testable by matrix multiplication using Hermitian
angular momenta multiplied by \(-i\).
For the bi-invariant round connection, the left-frame Levi--Civita
formula \(\nabla_{e_a}e_b=\tfrac12[e_a,e_b]\) induces
\[
 \nabla_{e_a}f=(2/R)\mathscr R_af+(1/R)T_af
\]
on spin-\(3/2\)-valued functions. Parallelity of the gamma-trace
projection then gives
\[
 \slashed D_{3/2}
 =-\frac23\sum_aT_a\bigl((2/R)\mathscr R_a+T_a/R\bigr).
\]
In the same trivialization, a homogeneous section is represented by
\(f(x)\in V_{3/2}\) with \(\Spin(4)\) action
\((U(a,b)f)(x)=\rho_{3/2}(b)f(a^{-1}xb)\). The left Casimir is
\(-\sum_a\mathscr L_a^2\) and the right Casimir is
\(-\sum_a(\mathscr R_a+T_a)^2\). Equality of the left and right
scalar Laplacians implies
\[
 \mathcal C_L-\mathcal C_R
 =2\sum_a\mathscr R_aT_a+\sum_aT_a^2.
\]
Substitution proves the displayed \emph{first-order} identity.
Frobenius reciprocity and the diagonal Clebsch--Gordan rule put the
multiplicity-one spin-\(3/2\) bundle sectors at
\(\left|j_L-j_R\right|\in\{1/2,3/2\}\), with
\(j_L+j_R\ge3/2\). Since
\(\mathcal C_{L/R}\) acts by \(j_{L/R}(j_{L/R}+1)\),
its scalar on each sector is
\[
 -\frac{2}{3R}\bigl[j_L(j_L+1)-j_R(j_R+1)\bigr].
\]
For \((j_L,j_R)=((k+3)/2,k/2)\) it equals
\(-(k+5/2)/R\), with dimension \((k+4)(k+1)\).
For \(((m+1)/2,m/2)\), \(m=k+1\ge1\), it equals
\(-(k+5/2)/(3R)\), with dimension \((k+3)(k+2)\).
The mirror modules reverse the sign. The former are \(\ker P^*\),
while the latter are \(\operatorname{Im}P\), by the already verified
twistor/representation allocation. Their union exhausts the smooth
harmonic decomposition. Standard compact ellipticity makes the
operator's self-adjoint closure have precisely these eigenvalues.
To see directly that the two first-order expressions are equivalent,
write $J_a=i\mathscr R_a$ and $S_a=iT_a$ in the Hermitian
angular-momentum notation of
Theorem~\ref{p2f:thm:v260-signed-rs-s3}. Then
\[
  \mathcal C_L-\mathcal C_R=-2\mathbf S\cdot\mathbf J-\tfrac{15}{4}I,
  \qquad
  -\frac{2}{3R}(\mathcal C_L-\mathcal C_R)
    =\frac4{3R}\mathbf S\cdot\mathbf J+\frac5{2R}I.
\]
This is exactly Equation~\ref{p2f:eq:v260-rs-round-operator},
so the two derivations use the same round Levi--Civita operator, not
competing conventions or a second physical selection.
\end{proof}

\begin{theorem}[Quadratic-moment inverse of the complete intrinsic covariant seed family]
\label{p2f:thm:v260-quadratic-moment-inverse}
Fix the intrinsic radius \(r_0=\sqrt3R\) and the covariant family
\(K_{a,\zeta}(q)\) of
Theorem~\ref{p2c:thm:complete-intrinsic-endpoint-seeds}, with
\(4/3\le a\le20/9\) and
\(|\zeta-u|\le1-u\), \(u=16/[9(4-a)]\).
Put \(f(q)=q_1^2-q_2^2\) and
\(F=\Phi_{a,\zeta}(f)=\int f(q)K_{a,\zeta}(q)\,d\sigma(q)\).
In the oriented basis \(e_{ij}=e_i\wedge e_j\), the parameters are
recovered from just three real matrix coefficients:
\[
 \boxed{a=1+9\langle e_1,Fe_1\rangle,\qquad
 \operatorname{Re}\zeta=6\langle e_{13},Fe_{13}\rangle,\qquad
 \operatorname{Im}\zeta=-6\operatorname{Im}
                  \langle e_{24},Fe_{13}\rangle.}
\]
For the physical outcome variable \(x=r_0q\), replace \(F\) in these
formulas by \(r_0^{-2}\Phi(x_1^2-x_2^2)\). Thus a complete action-generated
quadratic operator readout would remove this family's mathematical
parameter ambiguity by inversion. Three continuously reported real
coefficients are locally minimal for complete reconstruction on the
nonempty three-dimensional interior of the seed domain. The theorem
does not generate that quadratic readout, identify a physical coherent
preparation protocol for its off-diagonal entry, or assert that a
detector performs it.
\end{theorem}
\begin{proof}
The unit-sphere fourth moments are
\(\int q_i^4d\sigma=1/8\) and
\(\int q_i^2q_j^2d\sigma=1/24\) for \(i\ne j\).
The vector block \(BI+(a-B)qq^*\) gives
\[
 F|_V=\frac{a-B}{12}\operatorname{diag}(1,-1,0,0)
 =\frac{a-1}{9}\operatorname{diag}(1,-1,0,0),
 \qquad B=(4-a)/3.
\]
In the two-form order \((12,13,14,23,24,34)\), put
\(D=\operatorname{diag}(0,1,1,-1,-1,0)\). Direct integration of
\(H_q=2W_qW_q^*-I_6\) gives
\(\int f(q)H_qd\sigma=D/6\). Consequently
\[
 F|_{\Lambda^2V}
 =\frac16\bigl[(\operatorname{Re}\zeta)D
                   +i(\operatorname{Im}\zeta)*D\bigr].
\]
The chosen orientation obeys \(*e_{13}=-e_{24}\), so the \(13,13\)
coefficient is \(\operatorname{Re}\zeta/6\), and the \(24,13\)
coefficient is \(-i\operatorname{Im}\zeta/6\). These yield the inverse identities. Minimality follows from invariance
of domain: if an open three-dimensional parameter neighborhood admitted
a continuous injective descriptor into \(\RR^m\) for \(m<3\), adjoining
\(3-m\) zero coordinates would give a continuous injection into
\(\RR^3\) whose image lies in a proper linear subspace, contradicting
invariance of domain. This is injectivity within the stated covariant
family, not full tomography on the native \(137\)-real-dimensional algebra.
\end{proof}

\begin{corollary}[Maximal instrument-independent harmonic readout inside the covariant endpoint family]
\label{p2f:cor:v260-common-harmonic-readout}
For \(f\in C(S^3)\), all admitted continuous covariant endpoint seeds
\(K_{a,\zeta}\) give the same \(\Phi_{a,\zeta}(f)\) if and only if the
\(E_2\) harmonic projection of \(f\) vanishes. On this exact common domain,
\[
 \Phi(f)=\left(\int f\,d\sigma\right)I_{15}
 +\frac4{r_0}\sum_{a=1}^4
   \left(\int f(q)q_a\,d\sigma(q)\right)X_a^D.
\]
In particular, every affine coordinate observable is uniquely represented,
and all harmonics of degree \(\ge3\) lie in the common kernel of these
quadratic-density maps. Also, for any sphere-supported endpoint POVM
with the same moments (even if noncovariant),
\[
 \sum_a\int x_a^2E(dx)=r_0^2I_{15},\qquad
 \mathcal N_{\rm tot}=r_0^2I_{15}-\sum_a(X_a^D)^2
\]
is a unique total added-noise operator. These are conditional
mathematical readout identities, not an occurrence claim.
\end{corollary}
\begin{proof}
The family densities contain only harmonics of degrees zero, one and
two. Their degree-zero and degree-one coefficients are fixed by
normalization and the first moments, giving the displayed formula
whenever the degree-two component of \(f\) is zero.
Conversely, the variable vector-block coefficient
\((a-B)(qq^*-I/4)\) detects every nonzero traceless quadratic harmonic:
for \(f_2(q)=q^THq\) with \(H\) traceless symmetric,
\[
 \int f_2(q)\left(qq^*-\tfrac14I\right)d\sigma=\tfrac1{12}H.
\]
Since the admissible interval for \(a\) has nonempty interior, varying
\(a\) changes \(\Phi(f)\) for every nonzero \(E_2\) projection.
This proves maximality relative to the \emph{declared} family.
The total-noise identity uses only support on \(|x|=r_0\) and
\(\int E(dx)=I\); it requires no covariance.
\end{proof}

\begin{proposition}[Exact nonselection of a physical POVM or instrument from first moments]
\label{p2f:prop:v260-readout-nonselection}
The structural data consisting of the intrinsic coordinate tuple,
normalization, full \(\Spin(4)\) covariance, radius \(r_0\), and exact
first moments do not determine one POVM, its higher moments, or an
occurring apparatus. Even a fixed POVM does not determine a unique
post-measurement completely positive instrument. This conclusion
concerns the stated structural data, and does not exclude selection by
some \emph{additional, actually generated} action--detector incidence.
\end{proposition}
\begin{proof}
On the real endpoint slice \(a=4/3\), choose \(c=1/3\) and \(c=1\).
Both covariant positive normalized densities have identical first
moments, but the two-form coefficient of \(f(q)=q_1^2-q_2^2\) on
\(e_{13}\) is \(\langle e_{13},\Phi_c(f)e_{13}\rangle=c/6\).
For outcomes \(x=r_0q\) their quadratic expectations in the same
normalized state differ by
\[
 \frac{r_0^2}{6}\left(1-\frac13\right)=\frac{R^2}{3}>0.
\]
They therefore cannot be the same POVM despite identical structural
first-moment data. Separately, for any supplied finite POVM \(\{E_j\}\),
\(\mathcal I_j(\tau)=E_j^{1/2}\tau E_j^{1/2}\) and
\(\mathcal J_j(\tau)=\operatorname{Tr}(\tau E_j)\sigma_j\), with fixed
normalized output states \(\sigma_j\), are completely positive
operations with the same event probabilities and, in general, distinct
post-measurement outputs. Neither structural positivity nor GP3 licenses
minimizing an unconstituted effect rank as an apparatus selector.
Paper~II-A's action-consumer/observable-quotient principles give common
observable values when their actual action-origin hypotheses hold;
they do not create an occurrence or a detector coupling.
\end{proof}

\begin{corollary}[Individual-coordinate sharp readout without a joint-POVM selector]
\label{p2f:cor:v260-sharp-coordinate-pvm}
For the intrinsic tuple, every individual Hermitian \(X_a^D\) has the
canonical spectral PVM given by polynomial functional calculus. Its
spectrum, with multiplicities on \(\mathcal H_1\), is
\[
 \begin{array}{c|c}
 \lambda & \mathrm{mult}\\\hline
 +R/2,-R/2&1\ \text{each}\\
 +\sqrt2R/3,-\sqrt2R/3&3\ \text{each}\\
 +\sqrt3R/4,-\sqrt3R/4&2\ \text{each}\\
 0&3
 \end{array}
\]
(where each signed pair has the indicated per-sign multiplicity).
An \emph{actual ideal sharp test} of a named coordinate therefore
has mathematically fixed Born probabilities given its preparation,
without selecting a joint localization POVM. The performance of such
a test and its post-measurement instrument remain occurrence/process
claims.
\end{corollary}
\begin{proof}
Rotate to one unit direction. The geometric coordinate has one
scalar/radial \(2\times2\) block with edge \(R/2\), three tangent/bivector
blocks with edge \(\sqrt2R/3\), and three zero bivectors. The protected
block is \((\sqrt3R/4)C(q)\), which has eigenvalues \(\pm\sqrt3R/4\)
with multiplicity two each. The spectral theorem on this finite native
algebra yields the unique projection
\(P_\lambda(X_a^D)=\prod_{\mu\ne\lambda}(X_a^D-\mu I)/(\lambda-\mu)\)
for every distinct eigenvalue \(\lambda\). The projections and their
Born pairing are generated algebraic content; no performed test is inferred.
\end{proof}



\section{Spectral, Hopf, Berry, and higher-spin technical calculations}
\label{p2ftech:sec:spectra}

\subsection{Alternative momentum-selector non-implications}
\label{p2d:app:legacy-selector-diagnostics}
The following material records the superseded momentum-selector diagnostics
for comparison with earlier versions.  Under the current Paper-II
constitution they are not premises for the structural momentum-role seed.  In
particular,
\[
 \boxed{\begin{gathered}
 \mathcal H_p\cong\mathcal H_1,\qquad m_p=1\\
 \text{are independent of}\\
 \mathrm{TWR},\ \mathrm{SPEX},\ \mathcal E_9,\
 \mathsf C_{\rm PI},\\
 \text{spectral gaps, and Hopf charge}.
 \end{gathered}}
\]
The listed structures remain valid conditional mathematics in their declared
domains, but none may be used to re-select the already constituted momentum
carrier.

\begin{remark}[Displayed background classes of this paper]
In the dependency convention of Paper~II-A, this paper consumes three
background classes beyond the parent-allocation chain.  First, standard elliptic
spectral theory is applied only after a smooth compact parent and a specified
elliptic operator have been supplied.  Second, the Paper-I-A harmonic,
Hopf--contact, geodesic, and finite-module theorem packages are applied with
their stated normalizations.  Under one optional lower-ideal comparison, \(\mathrm{TWR}_p\) is
an independent support-stage condition and \((\mathrm{LVL}_p^1)\)
is a further level-one condition; a separate fixed-coordinate comparison
uses a full-carrier premise \(\mathfrak F_p=\mathcal H_1\).
Neither comparison is an antecedent of the GP2 result: Paper~II-B derives
\(\mathcal H_p\cong\mathcal H_1\) and \(m_p=1\) from the revised host
constitution, without either selector package.  Paper~II-C's exact coordinate-closure
proposition narrows one possible independent route by proving, for a fixed
coordinate action on the supplied comparison module \(\mathcal H_1\),
\[
 \langle\mathcal E_9\rangle_{\widehat X}=\mathcal H_1
 \quad\Longleftrightarrow\quad r_{1,0}r_{0,1}\ne0.
\]
In the fixed-action comparison, a level-one ceiling together with closure under the fixed
action without recompression yielded the full carrier.  The ceiling,
nondegeneracy, and no-recompression hypotheses remained explicit; the
coordinate-closure theorem alone was not an unrestricted physical selector.

The optional condition \((\mathrm{SPEX}_p)\) computes the contents of each
declared word stage but did not prove that the realized support was one of
those stages.
Thus
\[
   (\mathrm{SPEX}_p)\not\Longrightarrow(\mathrm{TWR}_p),
\]
with the Casimir sublevel \(\mathcal E_9\) providing the sharp countermodel.
This conditional implication does not derive
\(\mathsf P_{\rm sli}\), did not identify the complete-shell Casimir on
non-diagonal sectors with a scalar \(L^2(S^3)\) Laplacian, and did not supply
or select the Paper-I-A coordinate representative \(\mathsf C_{\rm PI}\) from
the completion family \(\mathfrak M_{\rm coord}\).  Third,
convention-sensitive higher-spin operators and
their sphere spectra are used in the theorems that require them.  For the
Rarita--Schwinger operator,
Proposition~\ref{p2:prop:rs-spectrum-cross-check} independently derives the
sector inventory, all multiplicities, and the twistor-image series, but leaves
the transverse A-series eigenvalue anchored to the cited spectrum theorem.
\end{remark}

\subsection{Conditional spectral-energy representatives}
\label{p2d:app:oscillator-examples}
This section collects worked vacuum/oscillator representatives that require
additional dispersion, operator-identification, or renormalization data and
are not used to prove the kinematic chain.  For a
physical continuum renormalized value retains
\[
 \mathrm{REN}=(\text{subtraction scheme},\ \mu_{\rm ref},\
 \text{allowed local counterterms}),
\]
while a physical spectral-energy interpretation retains a typed interface
\[
 \mathrm{DISP}=(c_{\rm field},\ A\leftrightarrow\Delta,\
 \text{field role},\ \text{clock/calibration}).
\]
Neither tuple is generated by the kinematic spectrum itself.

\begin{remark}[Downstream dispersion domain and distinct speed symbols]
\label{p2d:rem:dispersion-speed-details}
Write \(c_{\rm op}\) for the operational ceiling obtained in the timed
propagation domain of Paper~II-A, \(c_{\rm field}\) for the principal-symbol
speed of a specified free field, and \(c_{\rm light}^{\rm th}\) for the front
speed of a later-identified theoretical vacuum-light channel.  Paper~II-D
does not itself prove the principal ancestry of any specified field.
If the characteristic histories of that specified field are proved to
belong to the full timed propagation family used to define \(c_{\rm op}\),
and \(c_{\rm field}\) is identified in the same comparison and clock
normalization with the supremal invariant rate of those histories, then the
definition gives only
\[
        c_{\rm field}\le c_{\rm op}.
\]
This containment argument alone does not prove equality.  There is now a
separate stronger route: if a typed principal constructor proves that the
field is an undeformed same-owner descendant of the operational host
principal relation, Paper~II-A,
Theorem~\ref{p2a:thm:undeformed-host-principal-inheritance}, gives
\[
        c_{\rm field}=c_{\rm op}.
\]
Unique neutral support without that principal-owner hypothesis is
insufficient.  A later vacuum photon satisfying the same hypothesis has
\(c_{\rm light}^{\rm th}=c_{\rm op}\); empirical channel identification is
an additional measurement bridge.

If a later finite-dispersion theorem denotes its principal-symbol coefficient
by \(c_{\rm field}\), its minimal finite free-sector result proves, only on
an intrinsic finite block
in the minimal uninserted second-order free scalar-cover domain,
\[
 H_1^2
 =
 \hbar^2c_{\rm field}^2A
 +\mu^2c_{\rm field}^4I,
 \qquad
 E_\lambda^2
 =
 \hbar^2c_{\rm field}^2\lambda
 +\mu^2c_{\rm field}^4.
\]
For \(\mu=0\),
\[
        \omega_\lambda=c_{\rm field}\sqrt{\lambda}.
\]
This is a forward interface and is not used to prove any Paper-II kinematic
theorem.  It does not by itself identify \(A\) with the momentum-carrier
Laplacian used below and does not extend the finite theorem to the smooth parent
continuum.  Thus it removes a bare dispersion postulate only for a block for
which the later finite-dispersion hypotheses and the required operator identification have
both been independently proved.  Every oscillator formula below therefore
states explicitly the massless dispersion and operator identification it uses.
\end{remark}

\subparagraph{\texorpdfstring{Smooth parent $S^3$ (requires a stated regularization):}{Smooth parent S3 (requires a stated regularization):}}
\begin{proposition}[Classical scalar spectral finite-part prescription conditional on free dispersion]
\label{p2:prop:casimir-classical}
Fix \(c_{\rm field}>0\) and assume a massless minimally coupled scalar whose
nonzero modes obey
\[
        \omega_k=c_{\rm field}\sqrt{\lambda_k}.
\]
This continuum dispersion is a stated free-field hypothesis; a finite-block
dispersion theorem does not by itself prove its parent-continuum extension.  For the
minimally coupled scalar Laplacian on classical \(S^3(\Rmom)\), the formal
zero-point series is \cite{Dowker1976,Camporesi1991}
\begin{equation}
E_{\mathrm{vac}}^{\mathrm{formal}}
=
\frac{\hbar c_{\rm field}}{2\Rmom}
\sum_{k=1}^\infty (k+1)^2[k(k+2)]^{1/2},
\end{equation}
where the zero mode is excluded.  To avoid confusing the dimensionful operator zeta function of Section~\ref{p2:sec:spectral-zeta} with the dimensionless series appearing in the frequency sum, introduce
\[
\bar\zeta_\Delta(s)=\sum_{k=1}^\infty (k+1)^2 [k(k+2)]^{-s},
\qquad
\zeta_\Delta(s)=\Rmom^{2s}\bar\zeta_\Delta(s).
\]
In this minimally coupled normalization, Proposition~\ref{p2:prop:spectral-zeta-continuation} places \(s=-\tfrac12\) among the heat-kernel pole locations, and on the round \(S^3(\Rmom)\) that pole is \emph{present}: the residue is
\[
  \operatorname*{Res}_{s=-1/2}\zeta_\Delta(s)
  =\frac{a_2}{(4\pi)^{3/2}\,\Gamma(-\tfrac12)},
  \qquad
  a_2=\frac{\Vol}{2\Rmom^{4}}=\frac{\pi^{2}}{\Rmom},
\]
which is nonzero, and equals \(-\tfrac1{16}\) at \(\Rmom=1\).  Thus
\(\bar\zeta_\Delta(-\tfrac12)\), equivalently \(\zeta_\Delta(-\tfrac12)\), does not exist as an ordinary value, and a regularity or subtraction convention is \emph{required} rather than merely advisable.  The classical vacuum energy is defined only after a renormalization prescription, for example by the finite part
\begin{equation}
E_{\mathrm{vac}}^{\mathrm{ren}}
:=
\frac{\hbar c_{\rm field}}{2\Rmom}\,
\operatorname*{FP}_{s=-1/2}\bar\zeta_\Delta(s),
\end{equation}
Equivalently one may formulate the prescription using the dimensionful operator zeta function, but then a reference scale must be included; at a pole the finite part is shifted by a residue-dependent logarithmic term.  This is a mathematical finite-part prescription; a physical renormalized vacuum energy would require specifying the subtraction scheme, reference scale, and allowed local counterterms.  The finite Toeplitz sums below require no such analytic prescription.
\end{proposition}

\begin{proof}
Under the displayed free-field hypothesis, the frequency and degeneracy are
\[
 \omega_k
 =
 c_{\rm field}\sqrt{\lambda_k}
 =
 \frac{c_{\rm field}}{\Rmom}\sqrt{k(k+2)},
 \qquad
 d_k=(k+1)^2.
\]
Formally,
\[
E_{\mathrm{vac}}
=\frac{\hbar}{2}\sum_{k\ge1}(k+1)^2\,\omega_k
=\frac{\hbar c_{\rm field}}{2\Rmom}
\sum_{k\ge1}(k+1)^2[k(k+2)]^{1/2}.
\]
The exponent \(+\tfrac12\) in the formal series corresponds to the spectral-zeta parameter \(s=-\tfrac12\).  The meromorphic continuation of \(\bar\zeta_\Delta(s)\), or equivalently of the dimensionful operator zeta \(\zeta_\Delta(s)=\Rmom^{2s}\bar\zeta_\Delta(s)\), therefore represents the regularized series only after specifying how the possible pole at \(s=-\tfrac12\) is handled.  Since the residue displayed in the statement is nonzero, that point is a genuine pole for the minimally coupled scalar Laplacian, and the mathematically well-defined classical quantity is the finite part, or another explicitly chosen renormalized value, rather than the unqualified evaluation of \(\bar\zeta_\Delta(-\tfrac12)\).
\end{proof}


\subparagraph{\texorpdfstring{Finite Toeplitz $S^3_m$ (no regularization needed):}{Finite Toeplitz S^3_m (no regularization needed):}}
\begin{proposition}[Exact finite spectral oscillator sums under a typed dispersion]
\label{p2:prop:casimir-fuzzy}
Fix \(c_{\rm field}>0\).  On the scalar diagonal, assume
\(A=\Delta_{S^3}\), \(\mu=0\), and
\(\omega_\lambda=c_{\rm field}\sqrt{\lambda}\).
A later finite-dispersion theorem may discharge the affine massless dispersion
on an intrinsic finite block in its minimal uninserted second-order free
domain.  It does not, by
itself, identify that paper's positive spatial operator with the present
momentum-carrier scalar Laplacian.  Accordingly the complete hypothesis in
the preceding sentence is discharged only if an independent bridge supplies
that operator identification; absent such a bridge, the following scalar
sum remains a typed conditional representative.  The scalar diagonal finite
zero-point sum and the full-module
kinematic oscillator sum are distinct.  For the latter, separately
\emph{define} the representative sector frequencies by
\[
\omega_{j_L,j_R}^{\rm rep}
:=
\frac{c_{\rm field}}{\Rmom}
\sqrt{2\bigl[j_L(j_L+1)+j_R(j_R+1)\bigr]}.
\]
This definition on non-diagonal sectors is not a consequence of the scalar
free-field hypothesis or of a later finite-dispersion theorem.

For the scalar diagonal sector,
\[
\boxed{
E_{\mathrm{vac,diag}}^{(m)}
=
\frac{\hbar c_{\rm field}}{2\Rmom}
\sum_{k=1}^{m}(k+1)^2\sqrt{k(k+2)}.
}
\]
For the full \(\Spin(4)\)-equivariant module package, using the separately
defined frequencies \(\omega_{j_L,j_R}^{\rm rep}\) as a finite kinematic
oscillator representative rather than a scalar field mode sum,
\[
\boxed{
E_{\mathrm{vac,full}}^{(m)}
=
\frac{\hbar c_{\rm field}}{2\Rmom}
\sum_{\substack{j_L,j_R\in\frac12\ZZ_{\ge0}\\j_L+j_R\le m\\(j_L,j_R)\ne(0,0)}}
(2j_L+1)(2j_R+1)
\sqrt{2\bigl[j_L(j_L+1)+j_R(j_R+1)\bigr]}.
}
\]
No analytic continuation is required in either case.

For \(m=1\), the scalar diagonal contribution is
\[
E_{\mathrm{vac,diag}}^{(1)}
=
\frac{\hbar c_{\rm field}}{2\Rmom}\,4\sqrt3.
\]
The non-trivial full-module data at \(m=1\) are
\[
\begin{aligned}
(\tfrac12,0),(0,\tfrac12)&:\quad d=2,\quad \lambda\Rmom^2=\tfrac32,\\
(1,0),(0,1)&:\quad d=3,\quad \lambda\Rmom^2=4,\\
(\tfrac12,\tfrac12)&:\quad d=4,\quad \lambda\Rmom^2=3.
\end{aligned}
\]
Thus
\begin{equation}
E_{\mathrm{vac,full}}^{(1)}
=
\frac{\hbar c_{\rm field}}{2\Rmom}
\left[4\sqrt{3/2}+12+4\sqrt3\right].
\end{equation}
\end{proposition}

\begin{proof}
For the scalar diagonal sector of a free minimally coupled scalar field, the formal scalar zero-point representative is
\(E_{\mathrm{vac}}=\frac12\sum_n\hbar\omega_n\) with
\(\omega_n=c_{\rm field}\sqrt{\lambda_n}\), excluding the zero mode.  On the scalar
diagonal sector of \(S^3(\Rmom)\), the finite truncation gives
\(1\le k\le m\), hence the first displayed sum.  The non-diagonal summands
of the full \(\Spin(4)\)-equivariant module package are not scalar
\(L^2(S^3)\) modes; there the same finite oscillator formula is used only as
the paper's declared kinematic full-module representative, not as an
application of the later finite-dispersion scalar-field theorem, with eigenvalues labeled by
\((j_L,j_R)\) and \(j_L+j_R\le m\).  This gives the second displayed sum,
and the displayed \(m=1\) data give its explicit evaluation.
\end{proof}







\subsection{Discrete-symmetry supplements}
\label{p2d:app:discrete-supplements}
This section supplies the extended proof of the finite cycle criterion
stated in the main text, then records the remaining discrete-symmetry
clarifications and independently supplied constructions.

\begin{lemma}[Real intertwiner frames and the time-reversal cycle criterion]
\label{p2d:lem:time-reversal-cycle-criterion}
Consider the finite multiplicity-free coordinate-completion problem of
Proposition~\ref{p2d:prop:finite-completion-discrete-symmetry-scope},
with a sector-preserving time-reversal lift.
Normalized time-reversal-real frames exist on every allowed
one-dimensional coordinate-intertwiner space.
For arbitrary normalized reference frames, the class is time-reversal fixed
if and only if \(H_C=\kappa_C\overline{H_C}\) on every fundamental cycle of
the nonzero-edge graph.
When this criterion holds, a sectorwise modification of the lift implements
time reversal on the chosen representative and preserves its square.
\end{lemma}

\begin{proof}
On \(V_{j_L,j_R}\), write
\(\Theta^2=\varepsilon I\), where
\(\varepsilon=(-1)^{2(j_L+j_R)}\).
Every allowed coordinate edge changes both spin labels by
\(\pm\tfrac12\), so its two endpoint signs agree.
The coordinate-vector representation has real structure with square \(+1\).
Consequently the anti-linear map on the intertwiner line,
\[
 \mathfrak t_e(T)_a=\Theta_vT_a\Theta_u^{-1},
\]
has square \(1\).  Since this line is one-dimensional and the map preserves
norm, \(\mathfrak t_e(T_e)=\kappa_eT_e\) with \(|\kappa_e|=1\).
Under \(T'_e=e^{i\phi_e}T_e\),
\[
 c'_e=e^{-i\phi_e}c_e,\qquad
 \kappa'_e=e^{-2i\phi_e}\kappa_e.
\]
Choose \(e^{2i\phi_e}=\kappa_e\) to obtain a real frame.
The reverse-adjoint convention is maintained by the opposite reverse phase.

On a cycle, sector phases telescope, whereas time reversal sends
\(H_C\) to \(\kappa_C\overline{H_C}\); this proves necessity.
For sufficiency set \(r_e=\kappa_e\overline{c_e}/c_e\).
Choose a spanning forest, put \(z=1\) at one root in each component, and
solve \(z_v=r_ez_u\) recursively along its edges.
For a remaining edge the consistency equation is
\(\prod_{e\in C}r_e=1\) on its fundamental cycle, exactly the stated
condition.  This proves sufficiency.
The geometric scalar lift fixes the scalar backbone, where \(r_e=1\);
including that chain in the forest preserves the scalar symbol.

For a frame change of total cycle phase \(\Phi_C\),
\[
 H'_C=e^{-i\Phi_C}H_C,\qquad
 \kappa'_C=e^{-2i\Phi_C}\kappa_C.
\]
The criterion is therefore frame covariant.  Equivalently,
\(H_C/(\kappa_C\overline{H_C})=1\) is frame invariant.
Rephasing the sector lifts multiplies the edge \(\kappa_e\)'s by a
vertex coboundary and leaves \(\kappa_C\) unchanged.

Finally, if
\(\widehat X^{\mathcal T}=U_z\widehat XU_z^{-1}\), set
\(\widetilde\Theta=U_z^{-1}\Theta\).
Then \(\widetilde\Theta\widehat X_a\widetilde\Theta^{-1}
=\widehat X_a\).
As \(\Theta U_z^{-1}\Theta^{-1}=U_z\), one has
\(\widetilde\Theta^2=\Theta^2\).
All products and phase divisions above are on the nonzero-edge graph;
no phase is assigned to a vanished coefficient.
\end{proof}

\begin{corollary}[Antipodal Eigenstates]
\label{p2:cor:antipodal-eigenstates}
Let $\mathcal A:S^3\to S^3$ be the antipodal map $\mathcal A(x)=-x$. The hyperspherical harmonics $Y_{klm}$ on $S^3$ obey \cite{Camporesi1991}
\begin{equation}
\mathcal{A} Y_{klm} = (-1)^k Y_{klm}.
\end{equation}
Thus scalar modes with even $k$ are antipodally even and scalar modes with odd $k$ are antipodally odd. This is an antipodal $\ZZ_2$ grading, not the spatial parity operator of Definition~\ref{p2:def:parity}.
\end{corollary}

\begin{proof}
Hyperspherical harmonics $Y_{klm}$ are restrictions of homogeneous harmonic polynomials of degree $k$ on $\RR^4$. Under the antipodal map $x\mapsto -x$, such a polynomial acquires the factor $(-1)^k$, yielding the stated formula. The spatial parity operator $\mathcal P$ fixes $x_4$ and flips only $(x_1,x_2,x_3)$, so it acts within a fixed $k$-eigenspace but does not give the universal eigenvalue $(-1)^k$ on all scalar harmonics.  Other spatial-parity choices are conjugate by $SO(4)$ unless additional structures select this embedding axis.
\end{proof}
\begin{lemma}[An anti-unitary square is invariant under scalar rephasing]
\label{p2:lem:antiunitary-phase-square}
If \(\mathcal T\) is anti-unitary and \(z\in U(1)\), then
\[
        (z\mathcal T)^2=\mathcal T^2.
\]
Consequently a central phase cannot change a time-reversal square.
\end{lemma}

\begin{proof}
Anti-linearity gives \(\mathcal T z=\bar z\,\mathcal T\).  Hence
\[
 (z\mathcal T)^2
 =z\mathcal T z\mathcal T
 =z\bar z\,\mathcal T^2
 =\mathcal T^2.
\]
\end{proof}
\paragraph{Optional independent internal doublet}

\begin{proposition}[Externally supplied two-state exchange]
\label{p2:prop:independent-internal-doublet}
Let \(\mathcal H_{\mathrm{geom}}\) carry the geometric coordinate,
transvection, rotation, and Hopf operators, and enlarge it by an independently
supplied factor:
\[
        \mathcal H_{\mathrm{ext}}
        :=\mathcal H_{\mathrm{geom}}\otimes\mathbb C^2.
\]
For \(q\in\mathbb R\), define
\[
        Q_{\mathrm{int}}
        :=qI_{\mathrm{geom}}\otimes\sigma_3,
        \qquad
        \mathcal E_{\mathrm{int}}
        :=I_{\mathrm{geom}}\otimes\sigma_1.
\]
Then \(\mathcal E_{\mathrm{int}}\) is a unitary involution which fixes every
geometric operator \(O\otimes I_2\) and reverses the internal label:
\[
 \mathcal E_{\mathrm{int}}Q_{\mathrm{int}}
 \mathcal E_{\mathrm{int}}^{-1}=-Q_{\mathrm{int}}.
\]
This construction is external representation data.  In particular,
\(Q_{\mathrm{int}}\neq\hat{\mathfrak h}\otimes I_2\); it is not Hopf
momentum, and neither the factor \(\mathbb C^2\) nor its exchange is derived
from the geometric kinematic algebra.
\end{proposition}

\begin{proof}
The Pauli identities
\(\sigma_1^*=\sigma_1\),
\(\sigma_1^2=I_2\), and
\(\sigma_1\sigma_3\sigma_1=-\sigma_3\)
give the exchange statement.  The tensor factors show that
\(\mathcal E_{\mathrm{int}}\) commutes with all geometric operators,
including \(\hat{\mathfrak h}\otimes I_2\).
\end{proof}


\subsection{Detailed diagonal-weight isolation}
\label{p2d:app:diagonal-weight-isolation}
The main text retains the weight table and protected-block outcome.  This
appendix records the Fourier projector and its detailed application.

\subsubsection{Isolation of the Protected Diagonal-Weight Carrier}
\label{p2:subsec:auxiliary-decoupling}

The 15-dimensional Hilbert space \(\mathcal H_1\) decomposes into finite
diagonal spin-weight sectors.  The protected carrier isolated below is
\[
  \mathcal S_{\rm prot}=\mathcal H_1^{(-1/2)}\oplus\mathcal H_1^{(+1/2)},
  \qquad
  \dim_\CC\mathcal S_{\rm prot}=4,
\]
equivalently an eight-real-dimensional spinor module.  The remaining
diagonal-weight sectors
\[
  \mathfrak s=0,\qquad \mathfrak s=\pm1
\]
are not, by themselves, the structural Clifford complement
\[
  7=\dim_\RR\Cl(0,3)-1.
\]
That seven-dimensional complement belongs to the real operator algebra on the
protected block.  Thus the statement proved here is the finite diagonal-weight
projection onto the protected carrier, not a Hilbert-space decomposition
\(15=8\oplus7\), and not an identification with the geometric
left-fiber Hopf
charge.

\begin{proposition}[Diagonal-weight projection]
\label{p2:prop:diagonal-weight-projection}
Let \(\hat{\mathfrak s}\) denote the Hermitian finite diagonal spin-weight generator, so that
\[
\mathcal H_1=\bigoplus_{\mathfrak s}\mathcal H_1^{(\mathfrak s)},
\qquad
\hat{\mathfrak s}\psi=\mathfrak s\psi.
\]
The corresponding action is \(U(\theta)=e^{i\theta\hat{\mathfrak s}}\).  Since \(\mathfrak s\) may be half-integral, the Fourier projector is
\begin{equation}
\Pi_{\mathfrak s_0}
:=
\frac{1}{4\pi}
\int_0^{4\pi}
e^{-i\mathfrak s_0\theta}\,U(\theta)\,d\theta.
\end{equation}
It is the orthogonal projector onto \(\mathcal H_1^{(\mathfrak s_0)}\). In particular, choosing \(\mathfrak s_0=\pm\tfrac12\) projects onto the spinor diagonal-weight sectors and removes all modes with \(\mathfrak s\neq\pm\tfrac12\).
\end{proposition}

\begin{proof}
Using the orthogonal decomposition into \(\hat{\mathfrak s}\)-eigenspaces, any \(\psi=\sum_{\mathfrak s}\psi_{\mathfrak s}\) satisfies
\[
U(\theta)\psi=\sum_{\mathfrak s}e^{i\mathfrak s\theta}\psi_{\mathfrak s}.
\]
Hence
\[
\Pi_{\mathfrak s_0}\psi
=
\frac{1}{4\pi}\int_0^{4\pi}
\sum_{\mathfrak s}e^{i(\mathfrak s-\mathfrak s_0)\theta}\psi_{\mathfrak s}\,d\theta
=
\psi_{\mathfrak s_0},
\]
because \(\mathfrak s-\mathfrak s_0\in\frac12\ZZ\).
\end{proof}

\begin{corollary}[Diagonal-weight isolation of the protected carrier at \(m=1\)]
\label{p2:cor:protected-diagonal-weight-carrier}
At \(m=1\),
\[
\mathcal S_{\rm prot}
=
\mathcal H_1^{(-1/2)}\oplus \mathcal H_1^{(+1/2)},
\qquad
\dim_\CC \mathcal H_1^{(\pm 1/2)}=2.
\]
Hence the protected complex \(4\)-dimensional spinor module, equivalently the \(8\) real low-grade modes, is isolated by the finite diagonal spin-weight grading.  The complementary sectors are the diagonal-weight sectors \(\mathfrak s=0,\pm 1\).
\end{corollary}

\begin{proof}
The diagonal-weight table above gives
\[
\dim_\CC \mathcal H_1^{(-1/2)}=\dim_\CC \mathcal H_1^{(+1/2)}=2.
\]
The preceding proposition shows that the Fourier-averaging projectors onto \(\mathfrak s=\pm\tfrac12\) isolate exactly those two sectors. Their direct sum is therefore the protected carrier.  Although the projector formula uses a chosen Cartan axis, the resulting direct sum is the invariant spinor submodule \(V_{1/2,0}\oplus V_{0,1/2}\) inside \(\mathcal H_1\).
\end{proof}

The diagonal-weight grading therefore isolates the protected eight-real-dimensional spinor module as a canonical-relative kinematic block. Later operator or coupling bases may realize this block, but any physical low-energy identification requires finite-field and dynamical closure in later papers.



\subsection{Berry and Chern--Simons technical lemmas}
\label{p2d:app:berry-technical}
The following lemmas supply the technical global-lift, integrality,
normalization, and homogeneous-uniqueness facts cited in the main Hopf
normalization argument.

\begin{definition}[Projected Reeb-weight doubled-orbit data]
\label{p2:def:berry-connection}
Let \(X_H\) be the principal Hopf generator normalized by
\(\eta_0(X_H)=1\), and let \(r\in\frac12\ZZ\) be a projected Reeb weight.
Suppose \(f_r:S^3(\Rmom)\to\mathbb P(\mathcal V)\) is a smooth ray family
constant on every principal orbit:
\[
        f_r(h_\alpha g)=f_r(g).
\]
For each orbit map \(o_g(\alpha)=h_\alpha g\), pull its ray line back along
the double covering
\[
 p:\mathbb R/(4\pi\mathbb Z)\longrightarrow
   \mathbb R/(2\pi\mathbb Z),
 \qquad
 p(\widetilde\alpha)=\widetilde\alpha\bmod2\pi.
\]
Choose a normalized lift
\(\widetilde\psi_{r,g}\) of that pullback line satisfying
\begin{equation}
\label{p2:eq:psi-q0}
\bigl[\widetilde\psi_{r,g}(\widetilde\alpha)\bigr]
=f_r(o_g(p(\widetilde\alpha))),
\qquad
\widetilde\psi_{r,g}(\widetilde\alpha)
=e^{-ir\widetilde\alpha}\widetilde\psi_{r,g}(0),
\qquad
\widetilde\alpha\sim\widetilde\alpha+4\pi .
\end{equation}
The projected ray is single-valued after \(2\pi\), whereas the normalized
double-cover lift may change sign.  Half-integrality is exactly the
single-valuedness condition \(e^{-4\pi i r}=1\) on the doubled cover.

Throughout this subsection Hermitian inner products are conjugate-linear in
the first slot and linear in the second.  For any local normalized lift
\(\psi_r\) of the ray line on \(S^3\), define the local Berry form
\begin{equation}
\label{p2:eq:berry-A}
A_r^{\rm loc}
\;:=\;
i\,\langle \psi_r(g),\,d\psi_r(g)\rangle.
\end{equation}
Under the local phase change
\(\psi_r'=e^{i\chi}\psi_r\), this convention gives
\[
        (A_r^{\rm loc})'=A_r^{\rm loc}-d\chi.
\]
For \(r\in\frac12+\mathbb Z\), the doubled-orbit lift is not a global
normalized lift over the \(2\pi\)-periodic principal orbit.  The underlying
ray line is globally trivial by the next lemma, but no global lift preserves
this half-integral vertical normalization.  Thus this definition supplies
local projective spin-lift data; it does not supply a global Berry one-form of
weight \(r\) on \(S^3\).

This projected weight is independent of the integer Chern class or classical
Hopf charge \(q\).  The definition does not assert that the finite
diagonal-weight projector \(\Pi_{\mathfrak s_0}\) is a classical Hopf-charge
projector.
\end{definition}








\begin{proposition}[Arbitrary-real Reeb-form diagnostic]
\label{p2d:prop:arbitrary-real-reeb-form-diagnostic}
For any real \(r\), the global real one-form
\[
        B_r:=r\eta_0
\]
satisfies \(Q_{\CS}[B_r]=r^2\).  Thus
\(Q_{\CS}[B_{\pm1/2}]=\tfrac14\), but this is only a quadratic one-form
diagnostic and is not the global Berry invariant of the half-integral
doubled-orbit data in Definition~\ref{p2:def:berry-connection}.  The
independent associated-character generator \(q=\pm1\) has unit value.
\end{proposition}

\begin{proof}
Quadratic scaling in
\(Q_{\CS}[A]=(4\pi^2)^{-1}\int A\wedge dA\), together with
\(Q_{\CS}[\eta_0]=1\), gives the formula.
\end{proof}


\begin{lemma}[Global Berry trivialization on \(S^3\)]
\label{p2:lem:global-berry-trivialization}
Let \(f:S^3\to\mathbb P(\mathcal V)\) be a smooth ray family and let
\[
        L_f:=f^*\mathcal O(-1).
\]
Then \(L_f\) is trivial and admits a global normalized lift
\(\psi:S^3\to\mathcal V\).  The form
\[
        A:=i\langle\psi,d\psi\rangle
\]
is a globally defined real one-form.  If
\(\psi'=e^{i\chi}\psi\) is another global normalized lift, then
\[
        A'=A-d\chi
\]
for a globally defined real function \(\chi\).  Consequently \(dA\) and
\[
        \frac1{4\pi^2}\int_{S^3}A\wedge dA
\]
are independent of the chosen global normalized lift.
\end{lemma}

\begin{proof}
Complex line bundles are classified by their first Chern class.  Since
\(H^2(S^3;\mathbb Z)=0\), one has \(c_1(L_f)=0\), so \(L_f\) is trivial;
normalizing a nowhere-zero global section gives \(\psi\).  Differentiating
\(\langle\psi,\psi\rangle=1\) shows that
\(\langle\psi,d\psi\rangle\) is imaginary, hence \(A\) is real.

Any two global normalized lifts differ by a smooth map
\(u:S^3\to U(1)\).  Since \(H^1(S^3;\mathbb Z)=0\), equivalently since
\(S^3\) is simply connected, \(u\) lifts through
\(\mathbb R\to U(1)\) to \(u=e^{i\chi}\) with global real \(\chi\).  Direct
calculation gives \(A'=A-d\chi\).  Finally,
\[
 (A-d\chi)\wedge d(A-d\chi)-A\wedge dA
 =
 -d(\chi\,dA),
\]
whose integral vanishes on closed \(S^3\).
Equivalently,
\([S^3,U(1)]\cong H^1(S^3;\mathbb Z)=0\): there is no separate large
\(U(1)\) gauge class omitted by the global-\(\chi\) calculation.  Together
with \(H^2(S^3;\mathbb Z)=0\), this proves both global trivializability of
the ray line and invariance of the displayed quadratic integral under every
global phase gauge transformation.
\end{proof}

\begin{corollary}[Integrality on a constant-ray principal orbit]
\label{p2:cor:global-berry-weight-integrality}
Suppose a genuine smooth ray family satisfies
\[
        f(h_\alpha g)=f(g)
\]
on each \(2\pi\)-periodic principal Hopf orbit.  If its global Berry form is
globally gauge-equivalent to \(r\eta_0\), namely
\[
        A=r\eta_0-d\chi
\]
with global \(\chi\), then
\[
        r\in\mathbb Z.
\]
In particular, the half-integral doubled-orbit data of
Definition~\ref{p2:def:berry-connection} do not define such a global Berry
connection when \(r\in\frac12+\mathbb Z\).
\end{corollary}

\begin{proof}
Choose the global normalized lift supplied by
Lemma~\ref{p2:lem:global-berry-trivialization}.  On a principal orbit the ray
is constant, so the lift has the form
\[
        \psi(h_\alpha g)=u_g(\alpha)\psi(g)
\]
for a loop \(u_g:S^1\to U(1)\).  Therefore
\[
        \frac1{2\pi}\oint_{S^1}A
\]
is, up to the fixed sign convention, the integer winding number of \(u_g\).
On the other hand,
\(\oint_{S^1}\eta_0=2\pi\) and
\(\oint_{S^1}d\chi=0\), so the assumed global form gives
\(\oint A=2\pi r\).  Hence \(r\) is an integer.
\end{proof}

\begin{lemma}[Hopf normalization]
\label{p2:lem:hopf-normalization}
The Hopf connection satisfies
\begin{equation}
\label{p2:eq:hopf-volume}
\int_{S^3}\eta_0\wedge d\eta_0 \;=\; 4\pi^2.
\end{equation}
\end{lemma}

\begin{proof}
This is the Hopf/contact normalization of Paper~I-A \cite{Kugelmann2026a}: \(\oint_{S^1}\eta_0=2\pi\), \(d\eta_0=\pi^*\omega\), and \(\int_{S^2}\omega=2\pi\).  Fiber integration gives \(\int_{S^3}\eta_0\wedge d\eta_0=4\pi^2\).
\end{proof}

\begin{lemma}[Homogeneous uniqueness of the Reeb-normalized connection]
\label{p2:rem:berry-connection-status}
Let \(f:S^3\to\mathbb P(\mathcal V)\) be a genuine ray family satisfying
\(f(h_\alpha g)=f(g)\), and let \(A_r\) be a global Berry form supplied by
Lemma~\ref{p2:lem:global-berry-trivialization}.  Suppose the corresponding
Berry connection, as a connection rather than merely its underlying ray
family, admits a homogeneous gauge under the commuting right
\(SU(2)\)-action and principal left \(U(1)\)-action.  If in that homogeneous
gauge its vertical normalization is \(A_r(X_H)=r\) for
\(\eta_0(X_H)=1\), then it equals \(r\eta_0\) in the homogeneous gauge and
\[
        A_r=r\eta_0-d\chi
\]
in every global phase gauge, for a global real \(\chi\), and necessarily
\(r\in\mathbb Z\).  No independent horizontal invariant one-form is
available.  This statement does not apply to a half-integral doubled-orbit
lift that fails to descend to a global lift on the principal orbit.
\end{lemma}

\begin{proof}
In a homogeneous gauge the difference \(A_r-r\eta_0\) is horizontal.  At the
identity coset, the isotropy \(U(1)\) rotates the two-dimensional horizontal
plane and fixes no nonzero horizontal covector.  Homogeneity therefore forces
the invariant horizontal difference to vanish.  Every other global normalized
lift differs by \(e^{i\chi}\), and
Lemma~\ref{p2:lem:global-berry-trivialization} gives the global correction
\(-d\chi\).  The integrality statement is
Corollary~\ref{p2:cor:global-berry-weight-integrality}.
\end{proof}

\begin{corollary}[Conditional global homogeneous Berry normalization]
\label{p2d:cor:global-homogeneous-berry-normalization}
Under the exact hypotheses of
Lemma~\ref{p2:rem:berry-connection-status}, the global Berry form satisfies
\begin{equation}
\label{p2:eq:cs-q0}
 A_r=r\eta_0-d\chi,
 \qquad r\in\mathbb Z,
 \qquad Q_{\CS}[A_r]=r^2.
\end{equation}
The displayed connection identity is asserted on this globally trivialized,
principal-orbit-constant, homogeneous branch.  It
does not apply to a half-integral doubled-orbit lift that fails to descend to
the principal \(2\pi\)-orbit.
\end{corollary}

\begin{proof}
The form and the integrality of \(r\) are exactly the conclusions of
Lemma~\ref{p2:rem:berry-connection-status}.  Since
\(dA_r=r\,d\eta_0\),
\[
 A_r\wedge dA_r
 =r^2\eta_0\wedge d\eta_0-r\,d(\chi\,d\eta_0).
\]
The exact term integrates to zero on closed \(S^3\), while
Lemma~\ref{p2:lem:hopf-normalization} gives
\(\int_{S^3}\eta_0\wedge d\eta_0=4\pi^2\).  Definition~\ref{p2:def:cs-charge}
therefore yields \(Q_{\CS}[A_r]=r^2\).
\end{proof}

\begin{corollary}[Fibre-basic Berry integral without homogeneity]
\label{p2d:cor:fibre-basic-berry-cs}
Let \(\pi:S^3\to S^2\) be the normalized Hopf projection, and suppose the
smooth ray map is genuinely fibre-basic:
\[
 f=\bar f\circ\pi:S^3\longrightarrow\mathbb P(\mathcal V),\qquad
 \bar f:S^2\longrightarrow\mathbb P(\mathcal V).
\]
Give the descended ray line \(\bar L:=\bar f^*\mathcal O(-1)\) its Berry
connection and set
\[
 n:=\langle c_1(\bar L),[S^2]\rangle\in\mathbb Z.
\]
For every global normalized lift of \(f\), its real Berry form has
\[
 \boxed{A=n\eta_0+\pi^*\beta-d\chi,\qquad
 \beta\in\Omega^1(S^2;\mathbb R),\quad
 \chi\in C^\infty(S^3;\mathbb R),\qquad
 Q_{\CS}[A]=n^2.}
\]
No homogeneous gauge is required for this integral.  The horizontal
correction \(\beta\) need not vanish or be exact, so the stronger
connection identity \(A=n\eta_0-d\chi\) does not follow.  Nor does this
corollary quantize the integral of an arbitrary ray map that does not
descend through \(\pi\).
\end{corollary}

\begin{proof}
Complex line bundles over \(S^2\) are classified by their first Chern
number.  A unitary isomorphism therefore identifies \(\bar L\) with the
degree-\(n\) associated line of
Proposition~\ref{p2:prop:hopfion-charge}.  On this fixed line any two
unitary connections differ by a global real one-form \(\beta\) in the
Berry convention.  Pulling back to \(S^3\) and using the associated-line
trivialization gives \(n\eta_0+\pi^*\beta\).  Every other global
normalized lift differs by \(e^{i\chi}\) with global real \(\chi\), by
Lemma~\ref{p2:lem:global-berry-trivialization}; hence the displayed form of
\(A\), including the sign of its gauge correction.

The exact gauge correction does not change the integral.  Since
\(d\eta_0=\pi^*\omega\), all terms containing three base differential
degrees vanish, leaving
\[
 (n\eta_0+\pi^*\beta)\wedge
 d(n\eta_0+\pi^*\beta)
 =n^2\eta_0\wedge d\eta_0+n\eta_0\wedge\pi^*d\beta.
\]
The second term integrates to
\(2\pi n\int_{S^2}d\beta=0\) by fibre integration and Stokes' theorem.
The first integrates to \(4\pi^2n^2\) by
Lemma~\ref{p2:lem:hopf-normalization}, giving the result.
\end{proof}

\subsection{Supplementary quantization, geodesic, and monopole dictionaries}
\label{p2d:app:supplementary-geometry}
These standard positive consequences are useful interpretive dictionaries but
are not required by the main Paper-II-D proof chain.

\begin{corollary}[\texorpdfstring{$\mathbb{CP}^1$}{CP1} geometric quantization]
\label{p2:cor:cp1-geometric-quantization}
For every integer \(k\ge1\), geometric quantization of the integral
symplectic manifold \((S^2,k\omega)\), equivalently its holomorphic
quantization by the tensor power \(L^{\otimes k}\), yields the irreducible
representation of \(\SU(2)\) of dimension \(k+1\) and spin \(j=k/2\)
\cite{Kirillov2004lectures}.  At \(k=0\), the form \(0\omega\) is
degenerate and therefore is not a symplectic form on \(S^2\); the
degree-zero member is instead included through the Borel--Weil line-bundle
dictionary,
\[
H^0(\mathbb{CP}^1,\mathcal O)\cong\mathbb C,
\]
which is the trivial spin-zero representation.  Thus the holomorphic
line-bundle family yields the spin-\(k/2\) irreducibles for every
\(k\ge0\), while the symplectic geometric-quantization formulation on
\(S^2\) applies only for \(k\ge1\).
\end{corollary}

\begin{proof}
For \(k\ge1\), this is the \(\mathbb{CP}^1\)
geometric-quantization dictionary recorded in
Paper~I-A~\cite{Kugelmann2026a} and in the cited standard source.  For
\(k=0\), Borel--Weil gives
\(H^0(\mathbb{CP}^1,\mathcal O)=\mathbb C\) directly; this case is not
described as geometric quantization of the nonsymplectic pair
\((S^2,0)\).  Paper~II retains both branches only as a
representation-theoretic dictionary for its compact kinematic sectors.
\end{proof}



\begin{remark}
Since the spin is \(j=k/2\), integer spin occurs for even \(k\) and half-integer spin occurs for odd \(k\).  The half-integer sectors therefore arise naturally from the quantization of higher monopole charges or associated bundles, consistent with the unique spin structure on \(S^3\).
\end{remark}


\subsubsection{Momentum Space Geodesics and Periodicity}
\label{p2:sec:geodesic-periodicity}

The round-sphere geodesic facts used here are imported from the Paper~I-A geometry package \cite{Kugelmann2026a}.  Their role in Paper~II is kinematic: they identify the periodicity and finite intrinsic range of momentum-space trajectories.

\begin{proposition}[Geodesic Closure]
\label{p2:prop:geodesic-closure}
Every nonconstant geodesic on $S^3(\Rmom)$ is a great circle.  When parametrized by arclength, equivalently at unit speed, it closes after the minimal arclength period $2\pi \Rmom$.
The injectivity radius is $\mathrm{inj}(S^3) = \pi \Rmom$, with cut locus consisting 
of the antipodal point.
\end{proposition}

\begin{proof}
This is the Paper~I-A round \(S^3(\Rmom)\) geodesic package \cite{Kugelmann2026a}, with the standard cut-locus and injectivity-radius facts as in \cite{do1992riemannian}.  Paper~I-A proves that unit-speed geodesics are great circles with arclength period \(2\pi\Rmom\) and first cut point at the antipode, at distance \(\pi\Rmom\).  The present statement records the same facts in momentum-space language.
\end{proof}

\begin{corollary}[Geodesic-parameter periodicity in momentum space]
\label{p2:cor:momentum-space-geodesic-periodicity}
For a unit-speed momentum-space geodesic, the arclength parameter is periodic
modulo the universal geometric interval
\[
2\pi \Rmom = 2\pi\hbar/\Punit.
\]
Every such trajectory is a closed great circle.  This interval is an
arclength-parameter period on the momentum-space parent, not a physical-time
period; identifying it with a physical evolution time requires a separately
declared dynamics and dispersion interface.
\end{corollary}

\begin{proof}
By Proposition~\ref{p2:prop:geodesic-closure}, a unit-speed geodesic follows a
great circle and returns after arclength parameter \(2\pi\Rmom\).  Using
\(\Punit:=\hbar/\Rmom\) yields
\(2\pi\Rmom=2\pi\hbar/\Punit\).  No map from this arclength parameter to
physical time is used.
\end{proof}





\paragraph{Monopole-harmonic interpretation of Hopf charge sectors}

\begin{proposition}[Magnetic Monopole Analogy]
\label{p2:prop:monopole}
The charge-$q$ sector of $L^2(S^3)$ is isomorphic to the space of \textbf{monopole harmonics} on $S^2$ with line-bundle degree $q$ in the present Hopf-charge normalization:
\begin{equation}
\mathcal V_q \cong L^2(S^2, L^q).
\end{equation}
The finite scalar-diagonal truncation of this charge sector is
\begin{equation}
\label{p2:eq:scalar-diagonal-charge-truncation}
\mathcal V^{\rm diag}_{q,m}
:=
\mathcal V_q\cap
\bigoplus_{0\le k\le m}\left(V_{k/2}\otimes V_{k/2}^*\right)
=
\bigoplus_{\substack{|q|\le k\le m\\ k-|q|\in2\ZZ}}
\operatorname{span}\{D^{k/2}_{q/2,b}:-k/2\le b\le k/2\}.
\end{equation}
Here $L^q$ is the line bundle of degree $q$ over $S^2$.  The notation
\(\mathcal V^{\rm diag}_{q,m}\) deliberately refers only to the scalar diagonal
Peter--Weyl truncation; the full finite module $\mathcal H_m$ is the non-diagonal
finite \(\Spin(4)\)-equivariant package and is not itself a single scalar charge
sector. The spectral gap arises from the same mechanism as the angular momentum quantization of a charged particle in a monopole field.
\end{proposition}

\begin{proof}
Paper~I-A proves the continuum identification
\(\mathcal V_q\cong L^2(S^2,L^q)\), the allowed levels, and their
multiplicities~\cite{Kugelmann2026a}.  The only Paper-II addition is the
intersection with the already declared auxiliary finite scalar-diagonal truncation.
\end{proof}





\subsection{Exterior-parity-preserving same-symbol counterfamily}
\label{p2d:app:form-zero-order-counterfamily}

Let \(P_k\) denote the pointwise projection to differential-form degree
\(k\), and define the local bundle map
\[
 A:=P_2\circ\star\circ P_1:\Omega^1(S^3)\longrightarrow\Omega^2(S^3),
 \qquad
 Z_\tau:=\frac{i\tau}{\Rmom}(A-A^\dagger),
 \quad \tau\in\mathbb R.
\]

\begin{proposition}[Constant-annihilating form counterfamily]
\label{p2d:prop:form-zero-order-counterfamily}
The operator \(Z_\tau\) is local, self-adjoint,
\(\Spin(4)\)-equivariant, odd in exterior degree, zero order, and zero on
constants.  Consequently
\[
 \mathscr D_\tau^{\rm form}:=i(d-d^\dagger)+Z_\tau
\]
has the same principal symbol and scalar coordinate commutators as the
canonical form operator.  Nevertheless its square differs on the retained
Hodge-dual Killing sector whenever \(\tau\ne0\) outside the isolated
coincidence values.  Thus even constant annihilation and exterior parity do
not restore same-symbol rigidity.
\end{proposition}

\begin{proof}
The degree projections and Hodge star are pointwise natural bundle maps.
Hence \(A-A^\dagger\) is skew-adjoint and equivariant, so \(Z_\tau\) has all
the stated properties and commutes with scalar multiplication.

For a normalized coexact curl eigen-one-form \(\alpha\), write
\[
 \star d\alpha=\lambda\alpha,
 \qquad \beta:=\star\alpha.
\]
Then \(d\alpha=\lambda\beta\),
\(d^\dagger\beta=\lambda\alpha\), and on
\(\operatorname{span}\{\alpha,\beta\}\),
\[
 i(d-d^\dagger)\alpha=i\lambda\beta,
 \qquad
 i(d-d^\dagger)\beta=-i\lambda\alpha,
\]
while
\[
 Z_\tau\alpha=\frac{i\tau}{\Rmom}\beta,
 \qquad
 Z_\tau\beta=-\frac{i\tau}{\Rmom}\alpha.
\]
Therefore the squared eigenvalue becomes
\[
 \left(\lambda+\frac{\tau}{\Rmom}\right)^2
\]
instead of \(\lambda^2\).  On the retained Killing branches
\(\lambda=\pm2/\Rmom\), giving
\((\pm2+\tau)^2/\Rmom^2\).  This proves the claimed spectral variation under
a zero-order term satisfying the strengthened conditions.
\end{proof}

\subsection{Extended Rarita--Schwinger spectrum cross-check}
\label{p2d:app:rs-crosscheck}
The imported spectrum theorem used in the main text is independently checked
as far as representation theory and the twistor image determine it; the
transverse A-series eigenvalue is derived directly in
Theorem~\ref{p2f:thm:v260-signed-rs-s3} below.

\begin{proposition}[Representation-theoretic cross-check of the Rarita--Schwinger import]
\label{p2:prop:rs-spectrum-cross-check}
The imported theorem admits the following independent cross-check.
\begin{enumerate}[label=(\roman*), leftmargin=2em]
\item Sections of the \(\gamma\)-traceless vector-spinor bundle occupy
precisely the \(\Spin(4)\) sectors satisfying
\[
 |j_L-j_R|\in\{\tfrac12,\tfrac32\},
 \qquad
 j_L+j_R\ge\tfrac32,
\]
each with multiplicity one.
\item The chirality-difference-\(\tfrac32\) sectors
\[
\left(\frac{k+3}{2},\frac{k}{2}\right)
\quad\text{and their mirrors}
\]
have dimension \((k+1)(k+4)\), exactly the imported A-series
branch multiplicity per sign.  Representation theory identifies these as the
transverse sectors, but it does not by itself determine the scalar by which
the Rarita--Schwinger operator acts on them.  The A-series eigenvalue is not fixed by representation theory alone;
Theorem~\ref{p2f:thm:v260-signed-rs-s3} supplies its independent
signed round-operator calculation. The imported standard account remains
\cite{HommaTomihisa2020RaritaSchwinger}.
\item The twistor-image B-series is rederived.  We use as explicit external
twistor/Rarita--Schwinger inputs the intertwining identity, the
Killing-spinor kernel, and injectivity on higher Dirac levels
\cite{Bures1999HigherSpinDirac,HommaTomihisa2020RaritaSchwinger}.  In
dimension \(n=3\), the twistor operator \(P\) obeys
\[
\slashed D_{3/2}P
=\frac{n-2}{n}PD_{S^3}
=\frac13PD_{S^3}.
\]
Its kernel is the four-complex-dimensional Killing-spinor space, namely the
\(m=0\) Dirac level.  Hence Dirac levels \(m\ge1\), with eigenvalues
\(\pm(m+\tfrac32)/\Rmom\) and multiplicities \((m+1)(m+2)\) per sign,
produce
\[
\pm\frac{m+\tfrac32}{3\Rmom}.
\]
Putting \(k=m-1\) gives exactly
\[
\lambda_k^{(B)}
=\pm\frac{k+\tfrac52}{3\Rmom},
\qquad
\operatorname{mult}_{B,\pm}(k)=(k+2)(k+3).
\]
\item The full-bundle gap \(5/(6\Rmom)\) follows from the rederived B-series.
The transverse gap \(5/(2\Rmom)\) follows from the independently derived A-series
eigenvalue together with the independently verified transverse-sector
inventory.
\end{enumerate}
\end{proposition}

\begin{proof}
Write \(G=\Spin(4)\) and \(K=\Spin(3)_{\rm diag}\), so
\(S^3\cong G/K\).  The isotropy representation of a
\(\gamma\)-traceless vector-spinor is the spin-\(\tfrac32\) representation
\([\tfrac32]\) of \(K\).  Frobenius reciprocity gives the multiplicity of
\(V_{j_L}\otimes V_{j_R}\) in the section space as
\[
\dim\operatorname{Hom}_{K}
\left(
[\tfrac32],
(V_{j_L}\otimes V_{j_R})|_K
\right).
\]
The diagonal restriction is the multiplicity-free Clebsch--Gordan sum
\[
(V_{j_L}\otimes V_{j_R})|_K
\cong
\bigoplus_{s=|j_L-j_R|}^{j_L+j_R}[s].
\]
It contains \([\tfrac32]\) precisely when
\(|j_L-j_R|\le\tfrac32\le j_L+j_R\) and the parity agrees.  The allowed
nonnegative half-integral differences are therefore
\(\tfrac12\) and \(\tfrac32\), proving (i).

For difference \(\tfrac32\), the smallest parametrization is the pair in
(ii), whose dimension is
\[
\left(2\frac{k+3}{2}+1\right)
\left(2\frac{k}{2}+1\right)
=(k+4)(k+1).
\]
For difference \(\tfrac12\), the equivariant twistor operator maps the
Dirac sectors into the vector-spinor section space.  It is injective away
from the Killing-spinor kernel, and the displayed intertwining identity
therefore transports every Dirac level \(m\ge1\) with its multiplicity.
These images exhaust the one-dimensional multiplicity spaces of the
difference-\(\tfrac12\) sectors.  Their orthogonal complement is consequently
the difference-\(\tfrac32\) sector family, which is the transverse summand.
This proves the B-series, the sector allocation, and all branch-multiplicity
claims.
On each transverse irreducible the invariant first-order operator is scalar
by Schur's lemma, but representation theory alone does not determine that
scalar; this is the stated boundary of the cross-check.  Taking the minima
gives (iv).
\end{proof}



\begin{theorem}[Direct signed Rarita--Schwinger spectrum from the round parent]
\label{p2f:thm:v260-signed-rs-s3}
On the oriented round $S^3(R)$ with its inherited spin structure, let
$D_{3/2}=\mathsf P_{3/2}\,\gamma^a\nabla_{e_a}\mathsf P_{3/2}$
act on the gamma-traceless summand of $S\otimes T^*S^3$; the
Clifford convention is $\gamma_a=i\sigma_a$, and
$\mathsf P_{3/2}$ is the fibrewise orthogonal spin-$3/2$ projector.
With the Levi--Civita connection (and no zero-order deformation),
its complete signed spectrum is
\[
 \begin{aligned}
 \lambda^{A}_{k,\pm}&=\pm\frac{k+5/2}{R}, &
 \operatorname{mult}_{A,\pm}(k)&=(k+1)(k+4),\\
 \lambda^{B}_{k,\pm}&=\pm\frac{k+5/2}{3R}, &
 \operatorname{mult}_{B,\pm}(k)&=(k+2)(k+3),
 \qquad k\ge0.
 \end{aligned}
\]
The A series lies in $\ker P^*$ and the B series in $\operatorname{Im}P$.
These are complex multiplicities per sign; a numerical A--B collision
has additive full-bundle multiplicity.
\end{theorem}
\begin{proof}
Identify $S^3(R)$ with $SU(2)$ and take the positively oriented
left-invariant orthonormal frame $e_a$ normalized by
$[e_a,e_b]=(2/R)\varepsilon_{abc}e_c$.
In this global frame, set $J_a=iRe_a/2$; then
$[J_a,J_b]=i\varepsilon_{abc}J_c$.
The gamma-traceless fibre is precisely the highest-spin summand in
$V_{1/2}\otimes V_1=V_{3/2}\oplus V_{1/2}$.
Let $s_a$ act on the spin-$1/2$ factor, $\ell_a$ on the vector factor,
and $S_a=s_a+\ell_a$.
The round Levi--Civita spin/vector connection in the chosen frame is
$\nabla_{e_a}=e_a-iS_a/R$; it follows either from
$\nabla_{e_a}e_b=\varepsilon_{abc}e_c/R$ and spin compatibility, or
from the homogeneous spin connection.
Since $\gamma_a=2is_a$, the projected operator is
\[
 D_{3/2}=\frac4R\mathsf P_{3/2}(\mathbf s\cdot\mathbf J)
                         \mathsf P_{3/2}
          +\frac2R\mathsf P_{3/2}(\mathbf s\cdot\mathbf S)
                         \mathsf P_{3/2}.
\]
The angular-momentum projection identity on an irreducible total-spin
space of spin $h$ is
\[
 \mathsf P_h s_a\mathsf P_h
 =\frac{h(h+1)+s(s+1)-\ell(\ell+1)}{2h(h+1)}
   \mathsf P_h S_a\mathsf P_h.
\]
For $(h,s,\ell)=(3/2,1/2,1)$ this coefficient is $1/3$.
As $\mathbf S^2=15/4$ on the projected fibre, one obtains the
\emph{signed} operator identity
\begin{equation}
\label{p2f:eq:v260-rs-round-operator}
\boxed{D_{3/2}=\frac{4}{3R}\,\mathbf S\cdot\mathbf J
                   +\frac{5}{2R}\,I.}
\end{equation}
This determines the otherwise imported connection shift directly;
no eigenvalue from the target theorem is assumed.
For a right-regular spin $j=0,1/2,1,\ldots$, couple $V_{3/2}\otimes V_j$.
On the total-spin-$h$ summand,
\[
 2\mathbf S\cdot\mathbf J
  =h(h+1)-j(j+1)-15/4,
 \qquad h\in\{j+3/2,j+1/2,j-1/2,j-3/2\}
\]
with only nonnegative values satisfying the triangle inequalities retained.
The signed eigenvalues, in this order, are
\[
 \frac{2j+5/2}{R},\quad
 \frac{2j/3+1/2}{R},\quad
 -\frac{2j/3+1/6}{R},\quad
 -\frac{2j-1/2}{R}.
\]
For the first and fourth branches put respectively $2j=k$ and
$2j=k+3$; for the second and third put respectively
$2j=k+1$ and $2j=k+2$. Their multiplicities are
$(2j+1)(2h+1)$, giving $(k+1)(k+4)$ and $(k+2)(k+3)$
for the appropriate two sign pairs. The complete Peter--Weyl
section decomposition and its multiplicity-one diagonal-$SU(2)$
branching are the ones proved in
Proposition~\ref{p2:prop:rs-spectrum-cross-check}.
Its twistor intertwining identifies the $|j_L-j_R|=1/2$ family as
$\operatorname{Im}P$ and its orthogonal $|j_L-j_R|=3/2$ family as
$\ker P^*$. All signed eigenvalues and multiplicities therefore follow.
\end{proof}
\begin{remark}[Domain and provenance of the signed derivation]
The homogeneous computation acts on smooth $L^2$ spin-$3/2$ sections
and extends to the canonical self-adjoint elliptic realization.
It proves the round operator spectrum, not occurrence of a physical
spin-$3/2$ field, its source term, or a particle mass. A generic
same-principal-symbol zeroth-order deformation is not covered.
\end{remark}


\subsection{Coordinate-derived symmetry and finite spectral error certificates}
\label{p2d:sec:v226-spectral-propagation}
The six skew-Hermitian round generators are the explicit coordinate
polynomials of Theorem~\ref{p1:thm:v226-rotations}. Their Casimirs and the
six irreducible sector projectors are therefore native operation outputs;
no separate rotation-generator admission is needed after the intrinsic
class. This does not identify the coordinate square $Q$ with an independently
chosen kinetic operator or derive a physical clock from an algebra.

\begin{proposition}[Coefficient errors on a supplied kinetic pair]
\label{p2d:prop:v226-coefficient-spectral-error}
Let $K,\widetilde K$ be actual Hermitian native kinetic matrices in the
same fixed frame and let their common normalization satisfy $Z\succeq zI>0$.
Put $\delta q_j=\operatorname{Tr}(E_j(\widetilde K-K))$ and
$\epsilon=(\sqrt{137\cdot136}/\eta)\|\delta q\|_2$, $\eta=1/2$.
Then corresponding ordered generalized eigenvalues obey
\[
 |\widetilde\lambda_k-\lambda_k|\le\epsilon/z.
\]
If both normalized spectra lie above $g>0$, the corresponding positive
square roots differ by at most $\epsilon/(2z\sqrt g)$.
\end{proposition}
\begin{proof}
Theorem~\ref{p1b:thm:v226-dual-coefficients} gives
$\|\widetilde K-K\|\le\epsilon$; congruence by $Z^{-1/2}$ gives
$\|Z^{-1/2}(\widetilde K-K)Z^{-1/2}\|\le\epsilon/z$.
The Hermitian min--max principle bounds the ordered eigenvalues.
Use $|\sqrt a-\sqrt b|=|a-b|/(\sqrt a+\sqrt b)$ for the last assertion.
\end{proof}
The frequencies in the last statement are physical only when the already
supplied action and clock identify these roots with frequencies. If $Z$
also varies, its actual variation must be included before applying a
same-normalization bound. Growing position bands and physical continuum
comparisons retain their separate recovery estimates.


\subsection{Finite frequency tests and real-zero retention}
\label{p2d:sec:v227-frequency-tests}
\begin{corollary}[Complete frequency comparison has finite tests]
\label{p2d:cor:v227-frequency-tests}
For two actual Hermitian elimination words on a common target and kinetic
frame, complementary dimensions $m_1,m_2$ give the complete comparison
certificate of Theorem~\ref{p1b:thm:v227-finite-response}. Source jets of
order at most $k$ use
$\lfloor(k+1)(m_1+m_2)/2\rfloor+1$ upper-half-plane energies.
No strict positivity or inverse of the complementary block at real zero
is an antecedent to these nonreal tests.
\end{corollary}
\begin{proof}
Hermitian complementary spectra are real. Apply the finite numerator and
jet degree bounds, then Corollary~\ref{p1b:cor:v227-zero-mode-comparison}.
\end{proof}
At a real complementary zero, the actual coefficient $BP_0B^*$ vanishes
exactly when $BP_0=0$. Otherwise the mode must be retained in a stationary
zero-energy treatment. This propagates the existing zero-mode criterion;
it does not license deletion of a coupled critical mode. Frequency samples
at fixed curvature do not certify a global curvature profile.



\section{Projective calibrations, trace calculations, and equivariance obstructions}
\label{p2ftech:sec:exports}

\subsection{Fixed-action carrier-selection comparison}
\label{p2e:app:legacy-selector}
The following result characterizes a weaker conditional selector family.
It is not used in the generated momentum-seed proof,
which is owned by Paper~II-B.  In particular,
\[
 \boxed{\begin{gathered}
 \mathcal H_p\cong\mathcal H_1,\quad m_p=1\quad\text{are independent of}\\
 \mathrm{TWR},\ \mathrm{SPEX},\ \mathcal E_9,\ \mathsf C_{\rm PI},\
 n_{\rm H},\ q_{\rm prot},\ \text{and spectral-gap data}.
 \end{gathered}}
\]
The first three belong to optional selector comparisons, while the remaining
objects are completion, charge-atlas, or downstream spectral data and cannot
re-enter the upstream carrier derivation by citation.

\begin{proposition}[Conditional carrier-selection diagnostic and fixed-action closure alternative]
\label{p2e:prop:carrier-selection-boundary}
Specialize Paper~II-A,
Definition~\ref{p2:def:budgeted-finite-closure}, and
Proposition~\ref{p2:prop:word-budget-lower-ideal}, to the momentum-role word
family.  Put
\[
 \Lambda_{\rm fin}:=\left(\tfrac12\mathbb Z_{\ge0}\right)^2,
 \qquad
 w(j_L,j_R):=j_L+j_R,
\]
and, for \(m\in\mathbb Z_{\ge0}\), define
\[
 \mathcal C_p^{\rm word}(m)
 :=\mathcal H_m
 :=\bigoplus_{\substack{(j_L,j_R)\in\Lambda_{\rm fin}\\
                        w(j_L,j_R)\le m}}
        V_{j_L,j_R}.
\]
In this conditional comparison, \((\mathrm{TWR}_p)\) means that the unrestricted
realized carrier \(\mathfrak F_p\) is isomorphic to one declared member
\(\mathcal H_m\), for some \(m\in\mathbb Z_{\ge0}\).  The statement
\((\mathrm{SPEX}_p)\) means that every declared word stage has the exact
spectral support
\[
 \operatorname{supp}\mathcal C_p^{\rm word}(m)
 =\{(j_L,j_R)\in\Lambda_{\rm fin}:w(j_L,j_R)\le m\},
 \qquad m\in\mathbb Z_{\ge0}.
\]
Then
\[
        (\mathrm{SPEX}_p)\not\Longrightarrow(\mathrm{TWR}_p).
\]
Indeed, in the unrestricted finite equivariant domain the Casimir sublevel
\[
 \mathcal E_9
 :=V_{0,0}\oplus V_{\frac12,\frac12}
   \oplus V_{\frac12,0}\oplus V_{0,\frac12}
\]
is a nine-dimensional support compatible with the same declared-stage
schema, but it is not any complete word stage \(\mathcal H_m\).

There is a narrower sufficient alternative.  Fix a compatible Hermitian
equivariant coordinate vector
\(\widehat X=(\widehat X_1,\ldots,\widehat X_4)\) on \(\mathcal H_1\), with
adjoint-channel magnitudes \(r_{1,0}\) and \(r_{0,1}\).  If a realized
carrier \(\mathfrak F_p\) satisfies
\[
 \mathcal E_9\subseteq\mathfrak F_p\subseteq\mathcal H_1,
 \qquad
 \widehat X_a\mathfrak F_p\subseteq\mathfrak F_p
 \quad(a=1,\ldots,4),
 \qquad
 r_{1,0}r_{0,1}\ne0,
\]
where the middle condition uses the same fixed action without recompression,
then
\[
        \mathfrak F_p=\mathcal H_1.
\]
If that fixed action is invariant under interchange of the left and right
factors, then \(r_{1,0}=r_{0,1}\), so nonvanishing of either adjoint-channel
magnitude implies nonvanishing of both.
\end{proposition}

\begin{proof}
The statement \((\mathrm{SPEX}_p)\) quantifies over the declared stages and
contains no quantifier requiring an unrestricted realized support to equal
one of them.  The displayed \(\mathcal E_9\) has dimension
\(1+4+2+2=9\) and maximum total weight one.  It is not
\(\mathcal H_0\), which has dimension one, not \(\mathcal H_1\), which has
dimension fifteen, and not any \(\mathcal H_m\) with
\(m\in\mathbb Z_{\ge2}\), whose
dimension is \(\binom{2m+4}{4}>15\).  It is therefore the sharp finite
equivariant countermodel to the implication: retain the declared tower
\(\{\mathcal H_m\}_{m\in\mathbb Z_{\ge0}}\), for which
\((\mathrm{SPEX}_p)\) holds, and
choose the unrestricted realized carrier to be \(\mathcal E_9\).  Since the
antecedent does not constrain that separate choice, it remains true while
\((\mathrm{TWR}_p)\) is false.

For the second assertion, Paper~II-C,
Proposition~\ref{p2c:prop:e9-coordinate-closure}, proves that the only
coordinate channels from
\(\mathcal E_9\) to its complement in \(\mathcal H_1\) are
\[
 V_{\frac12,\frac12}\longleftrightarrow V_{1,0},
 \qquad
 V_{\frac12,\frac12}\longleftrightarrow V_{0,1}.
\]
Each channel has multiplicity one.  A nonzero equivariant map onto either
irreducible target is surjective, so
\(r_{1,0}r_{0,1}\ne0\) gives
\[
        \langle\mathcal E_9\rangle_{\widehat X}=\mathcal H_1.
\]
Since \(\mathfrak F_p\) contains \(\mathcal E_9\) and is invariant under the
same operators, it contains this generated closure.  The level-one ceiling
then gives
\(\mathcal H_1\subseteq\mathfrak F_p\subseteq\mathcal H_1\), hence equality.
Factor interchange swaps the two adjoint targets and therefore their reduced
magnitudes, proving the last clause.
\end{proof}

\subparagraph{Joint realizability under a fixed coordinate action.}
The four separately admissible least-closure values
\[
 \mathcal E_9,\quad
 \mathcal E_9\oplus V_{1,0},\quad
 \mathcal E_9\oplus V_{0,1},\quad
 \mathcal H_1
\]
are not freely co-realizable with every coordinate class once the same-action
closure relation is imposed.  Their joint locus is the graph determined by the
zero/nonzero pattern of \((r_{1,0},r_{0,1})\): respectively
\((0,0)\), \((>0,0)\), \((0,>0)\), and \((>0,>0)\), with
\(\varrho_{\rm sp}\) unrestricted.  Hence the graph alone selects no
completion, while an independent demand that the least closure equal
\(\mathcal E_9\) or \(\mathcal H_1\) cuts out the corresponding fibre product.
This is a diagnostic of joint realizability and does not re-enter the current
carrier proof.

\subsection{Optional protected-scale and factorized coordinate calibrations}
\label{p2e:app:optional-calibration}
The main text proves the parameterized identity
\(\sum_a(X_a^{\rm sp})^2=4\varrho_{\rm sp}^{\,2}I\) on the raw atlas,
and fixes \(\varrho_{\rm sp}=\sqrt3\Rmom/4\) on the generated
\(\mathsf C_D\) branch.  The specializations below impose additional
representative or calibration data.  The optional Paper-I-A family is
distinct from \(\mathsf C_D\).

\begin{corollary}[Exact protected-scale specializations]
\label{p2e:cor:protected-scale-branches}
The Paper-I-A representative \(\mathsf C_{\rm PI}\) contains an independent
scale \(\rho_{\rm prot}>0\) and imposes
\[
        \sum_{a=1}^4(X_a^{\rm sp})^2
        =\rho_{\rm prot}^{\,2}I_{\mathcal S_{\rm prot}}.
\]
Consequently its protected coordinate scale is
\[
        \varrho_{\rm sp}=\frac{\rho_{\rm prot}}2,
        \qquad
        \sum_a(X_a^{\rm sp})^2
        =\rho_{\rm prot}^{\,2}I_{\mathcal S_{\rm prot}}.
\]
Only after the additional scale-matching convention
\(\rho_{\rm prot}=\Rmom\) does this become
\[
        \varrho_{\rm sp}=\frac{\Rmom}{2},
        \qquad
        \sum_a(X_a^{\rm sp})^2=\Rmom^2I_{\mathcal S_{\rm prot}}.
\]
The two displayed rows are distinct logical branches; neither the derived
momentum seed nor \(\mathsf C_{\rm PI}\) alone identifies
\(\rho_{\rm prot}\) with \(\Rmom\).  A further factorized-channel
calibration used in worked representatives is supplementary and is recorded
below in this technical section.
\end{corollary}

\begin{proof}
Comparing the parameterized identity
\(4\varrho_{\rm sp}^{\,2}I_{\mathcal S_{\rm prot}}\) with the Paper-I-A condition gives
\(4\varrho_{\rm sp}^{\,2}=\rho_{\rm prot}^{\,2}\).  Since
\(\varrho_{\rm sp}\ge0\) and \(\rho_{\rm prot}>0\), this fixes
\(\varrho_{\rm sp}=\rho_{\rm prot}/2\).  Substituting the optional equality
\(\rho_{\rm prot}=\Rmom\) gives the second row.
\end{proof}

A separately declared factorized-channel representative sometimes used for
worked matrices instead imposes
\[
        \varrho_{\rm sp}=\frac{\Rmom}{2\sqrt2},
        \qquad
        \sum_a(X_a^{\rm sp})^2
        =\frac{\Rmom^2}{2}I_{\mathcal S_{\rm prot}}.
\]
This is a noncanonical calibration.  It is not implied by
\(\mathcal H_p\cong\mathcal H_1\), \(m_p=1\), or \(\mathsf C_{\rm PI}\), and
it is not used in the positive constitutional chain.  If physically realized,
it is a G0 calibration certificate whose apparatus/standard map must be fixed
independently of the measured result.  It changes representation only if that
standard and map were already constituted; otherwise their result-independent
signature opens a new downstream frontier.  In neither case is the carrier
changed.  Proposition~\ref{p2e:prop:joint-protected-scale-loci} gives the exact
same-completion intersections: combined with
\(\mathsf C_{\rm PI}\) alone this factorized row forces
\(\rho_{\rm prot}=\Rmom/\sqrt2\), whereas combined also with
\(\rho_{\rm prot}=\Rmom\) its joint locus is empty for \(\Rmom>0\).
If the two rows refer to distinct representatives and no same-object relation
is supplied, they may co-realize without contradiction.

\begin{proposition}[Joint protected-scale loci and incompatible
same-completion rows]
\label{p2e:prop:joint-protected-scale-loci}
Let \(\Rmom>0\).  The three optional normalization relations are
\[
 \mathsf C_{\rm PI}:\quad
   \varrho_{\rm sp}=\frac{\rho_{\rm prot}}2,\quad
   \rho_{\rm prot}>0,
 \qquad
 \mathsf M_R:\quad \rho_{\rm prot}=\Rmom,
 \qquad
 \mathsf F_R:\quad
   \varrho_{\rm sp}=\frac{\Rmom}{2\sqrt2}.
\]
Their exact joint loci, when imposed on the \emph{same} completion, are:
\begin{align*}
 \mathsf C_{\rm PI}\cap\mathsf M_R
 &:\quad
 (\rho_{\rm prot},\varrho_{\rm sp})
   =(\Rmom,\Rmom/2),\\
 \mathsf C_{\rm PI}\cap\mathsf F_R
 &:\quad
 (\rho_{\rm prot},\varrho_{\rm sp})
   =(\Rmom/\sqrt2,\Rmom/(2\sqrt2)),\\
 \mathsf C_{\rm PI}\cap\mathsf M_R\cap\mathsf F_R
 &=\varnothing,\\
 \{\varrho_{\rm sp}=\Rmom/2\}\cap\mathsf F_R
 &=\varnothing.
\end{align*}
Equivalently, the radius-matched representative row and the factorized row
cannot describe one completion, because they would require
\(\Rmom/2=\Rmom/(2\sqrt2)\).

This empty-locus conclusion does not forbid two differently typed or
differently marked completion representatives from occurring side by side,
one on each branch; absent a same-object relation their composition is
relation-free.  Nor does generated-owner normalization identify
\(\rho_{\rm prot}\) with \(\Rmom\): Paper~I-A,
Definition~\ref{p1:def:s3one-finite-geometry-datum} and
Theorem~\ref{p1:thm:s3one-cliffordization}, assign them independent owner
slots and show that the optional finite protected scale is not inherited from
scalar radius-\(\Rmom\) multiplication.  Thus \(\mathsf M_R\) remains an
independent cross-owner calibration, whereas
\(\varrho_{\rm sp}=\rho_{\rm prot}/2\) is a graph internal to
\(\mathsf C_{\rm PI}\).
\end{proposition}

\begin{proof}
The first two loci follow by substitution into
\(\varrho_{\rm sp}=\rho_{\rm prot}/2\).  The triple intersection would give
\[
 \frac{\Rmom}{2}=\frac{\Rmom}{2\sqrt2}.
\]
Multiplication by \(2/\Rmom\), permitted because \(\Rmom>0\), gives
\(1=1/\sqrt2\), a contradiction.  Hence both displayed same-completion
intersections are empty.  Without a common-object relation no equality is
imposed, and independent owner slots are not identified by numerical
comparability.
\end{proof}

\subsection{Supplementary trace-normalization diagnostic}
\label{p2e:app:supplementary-normalization}
The following negative comparison is useful for preventing a normalization
mix-up, but it is not used in the positive trace-kernel chain.

\begin{proposition}[Empty nonzero joint rescaling locus for spinor and species
normalizations]
\label{p2e:prop:empty-spinor-species-rescaling-locus}
On the traceless line \(Q=q_{\rm prot}\Gamma_{\rm sp}\), let the
species-normalized linear and cubic expressions be
\[
 \mathcal A^{\rm spc}_1(Q)=2q_{\rm prot},
 \qquad
 \mathcal A^{\rm spc}_3(Q)=2q_{\rm prot}^{\,3}.
\]
There is no single positive charge rescaling \(\lambda\) that, for
\(q_{\rm prot}\neq0\), simultaneously converts the spinor traces of
\(\lambda Q\) into the species-normalized traces of \(Q\):
\[
 \left.
 \begin{aligned}
  \mathcal A_1(\lambda Q)&=\mathcal A^{\rm spc}_1(Q),\\
  \mathcal A_3(\lambda Q)&=\mathcal A^{\rm spc}_3(Q)
 \end{aligned}
 \right\}
 \quad\Longrightarrow\quad
 \lambda=\frac12
 \quad\text{and}\quad
 \lambda=2^{-1/3}.
\]
Hence the exact joint locus in
\(\mathbb R\times\mathbb R_{>0}\) is
\[
 \boxed{
 \mathfrak Z_{\rm norm}
 =\{(q_{\rm prot},\lambda):
       q_{\rm prot}=0,\ \lambda>0\}.}
\]
The two trace conventions have the same zero locus but are distinct
normalization owners away from zero.
\end{proposition}

\begin{proof}
Homogeneity gives
\(\mathcal A_1(\lambda Q)=4\lambda q_{\rm prot}\) and
\(\mathcal A_3(\lambda Q)=4\lambda^3q_{\rm prot}^{\,3}\).
For \(q_{\rm prot}\neq0\), equality with the two displayed species values
forces respectively \(4\lambda=2\) and \(4\lambda^3=2\), which are
incompatible because \(1/2\neq2^{-1/3}\).  For \(q_{\rm prot}=0\), all four
expressions vanish for every \(\lambda>0\), proving the locus.
\end{proof}

\subsection{Natural left-\texorpdfstring{\(2T\)}{2T} equivariance obstruction}
\label{p2e:app:natural-action-no-go}
For the natural left actions, Paper~I-B,
Proposition~\ref{p1:prop:h1-e15-left-2t-obstruction}, computes
\[
 \dim\mathcal H_1^{2T}=6,
 \qquad
 \dim(\mathcal E_{15}^{\rm BM})^{2T}=1.
\]
An equivariant isomorphism
\(\mathcal H_1\cong\mathcal E_{15}^{\rm BM}\) would restrict to an
isomorphism of fixed-point subspaces and therefore preserve their dimensions,
a contradiction.  This proves the natural-action boundary stated in
Proposition~\ref{p2e:prop:bm-supported-trace-boundary}; it does not exclude
arbitrary vector-space isomorphisms or actions transported through a separately
chosen isomorphism.

\begin{proposition}[Sharp rank of the natural Hurwitz comparison]
\label{p2e:prop:sharp-natural-2t-rank}
For the natural left-\(2T\) actions on \(\mathcal H_1\) and
\(\mathcal E_{15}^{\rm BM}\), every complex-linear intertwiner satisfies
\[
 \operatorname{rank}F\le5,\qquad \dim\ker F\ge10.
\]
The rank bound is attained by
\[
 F_{\rm sc}=\operatorname{ev}_{2T}\Pi_{\rm sc},
 \qquad
 \Pi_{\rm sc}:\mathcal H_1\longrightarrow
 V_{0,0}\oplus V_{1/2,1/2},
\]
where the projection is orthogonal and the scalar evaluation map is that of
Paper~I-B, Proposition~\ref{p1:prop:h1-scalar-hurwitz-evaluation-bridge}.
The space of such intertwiners has complex dimension \(12\).  For the full
regular target \(\mathbb C^{2T}\), the distinct sharp rank bound is \(8\).
\end{proposition}
\begin{proof}
Write \(V_2,V_3\) for the defining and adjoint irreducibles of \(2T\).
Restriction to the left factor gives
\[
 \mathcal H_1|_{2T}=6\mathbf1\oplus3V_2\oplus V_3.
\]
For a unit quaternion \(q\), their characters are
\(\chi_2(q)=2\operatorname{Re}q\) and
\(\chi_3(q)=4(\operatorname{Re}q)^2-1\).
The real parts \(1,-1,0,1/2,-1/2\) occur with multiplicities
\(1,1,6,8,8\), respectively. Character averaging verifies that
\(\mathbf1,V_2,V_3\) are distinct irreducibles.

The adjacency \(A_+\) is central convolution by the eight elements with
\(\operatorname{Re}q=1/2\). Its eigenvalue on \(V_3\) is
\(3^{-1}\sum\chi_3(q)=0\). The regular \(V_3\)-isotypic component has
dimension nine, equal to \(\dim\ker A_+\). Thus
\(\mathcal E_9^{\rm BM}=3V_3\), and its orthogonal complement has
multiplicities \(1,2,0\) for \(\mathbf1,V_2,V_3\).
Schur's lemma gives
\[
 \operatorname{Hom}_{2T}(\mathcal H_1,\mathcal E_{15}^{\rm BM})
 \cong\operatorname{Mat}_{1\times6}(\mathbb C)
 \oplus\operatorname{Mat}_{2\times3}(\mathbb C).
\]
For the two multiplicity matrices \(B_1,B_2\),
\(\operatorname{rank}F=\operatorname{rank}B_1+2\operatorname{rank}B_2\le5\).
Rank--nullity gives the kernel bound; the two matrix spaces have dimensions
six and six.

The projection \(\Pi_{\rm sc}\) commutes with the left action, and the
cited scalar evaluation is an equivariant isometry onto the adjacency
eigenspaces of eigenvalues eight and four. Their sum has dimension five
and lies in \(\mathcal E_{15}^{\rm BM}\), so \(F_{\rm sc}\) attains the
bound. Its kernel is the ten-dimensional nonscalar complement.
The full regular target restores multiplicity three for \(V_3\); the same
calculation gives the attained maximum \(1+4+3=8\).
\end{proof}

\begin{corollary}[Natural scalar-Hurwitz blindness to the left chiral adjoint]
\label{p2e:cor:v261-natural-hurwitz-v3-blindness}
For the natural left-\(2T\) actions of
Proposition~\ref{p2e:prop:sharp-natural-2t-rank}, every intertwiner
\[
 F:\mathcal H_1\longrightarrow\mathcal E_{15}^{\rm BM}
\]
annihilates the full source \(V_3\)-isotypic component:
\[
 \boxed{V_3\subseteq\ker F.}
\]
Under the parent \(\Spin(4)\) decomposition this \(V_3\) is the restriction of
the left chiral adjoint block \(V_{1,0}\).  Thus the natural scalar Hurwitz
target is structurally blind to an entire chiral-adjoint sector, not merely
rank deficient in aggregate.
\end{corollary}

\begin{proof}
The proof of Proposition~\ref{p2e:prop:sharp-natural-2t-rank} gives
\(\mathcal H_1|_{2T}=6\mathbf1\oplus3V_2\oplus V_3\), whereas
\(\mathcal E_{15}^{\rm BM}\) has multiplicities \((1,2,0)\) for
\((\mathbf1,V_2,V_3)\).  Schur's lemma therefore forces every natural
left-\(2T\) intertwiner to vanish on the source \(V_3\) summand.
\end{proof}


This is a scalar vertex-function comparison with the stated natural action.
It is not a full-carrier isomorphism or a physical Born exchange. Changing
the actions or using a fiber-valued sampling target changes the theorem.





\section{Independent reverse certificates for the II-B carrier and IV-A population}
\label{p2f:sec:v264-two-bracket-action-audit}
The original II-B constitution is logically earlier than the intrinsic
II-C tuple, the III-F correlated connection and the V-A fermion action.
The following checks are retrospective implications on those
\emph{already declared} mathematical domains; they neither enlarge the
frozen GP2 constructor grammar nor create a G0 event, second carrier,
process or test from mathematical representation theory.

\begin{proposition}[Sharp positive scalar multiplication defect]
\label{p2f:prop:v264-sharp-scalar-compression-defect}
Let $P_{\le1}$ project $L^2(S^3(R))$ onto the scalar harmonic space
$E_0\oplus E_1=\operatorname{span}\{1,x_1,x_2,x_3,x_4\}$.
For multiplication by the ambient coordinate $x_1$,
\[
\begin{split}
\Delta_{x_1}&=P_{\le1}M_{x_1^2}P_{\le1}
   -(P_{\le1}M_{x_1}P_{\le1})^2\\
 &=P_{\le1}M_{x_1}(I-P_{\le1})M_{x_1}P_{\le1}\ge0.
\end{split}
\]
Its spectrum on $E_0\oplus E_1$ is
\[
 \boxed{\{0,\ R^2/4,\ R^2/6,\ R^2/6,\ R^2/6\}},
 \qquad\boxed{\|\Delta_{x_1}\|=R^2/4.}
\]
This is a mathematical retained-scalar no-go for exact pointwise
multiplication; it is not a physical error bar for an unconsumed
coordinate word.
\end{proposition}
\begin{proof}
The first equality follows by inserting $I=P_{\le1}+(I-P_{\le1})$
between the two self-adjoint multipliers. On the round three-sphere
with uniform normalized volume,
$\mathbb E[x_i^2]=R^2/4$,
$\mathbb E[x_i^4]=R^4/8$ and
$\mathbb E[x_i^2x_j^2]=R^4/24$ for $i\ne j$.
On the constant vector $M_{x_1}$ is retained exactly, so the defect
vanishes. On $x_1$ the full squared norm ratio is $R^2/2$ and the
retained squared norm ratio is $R^2/4$, leaving $R^2/4$.
On each orthogonal $x_j$ with $j\ne1$, the retained product
$P_{\le1}x_1x_j$ vanishes; its full squared norm ratio is
$(R^4/24)/(R^2/4)=R^2/6$. Orthogonality under the sign flips of
independent ambient coordinates removes all off-diagonal entries.
\end{proof}

\begin{proposition}[Exact native indistinguishability counterexample]
\label{p2f:prop:v264-native-state-nonconverse}
Let $u\in\mathcal H_g$ and $v\in\mathcal H_s$ be unit vectors and
$\psi=(u+v)/\sqrt2$. The pure density
$\rho_{\rm coh}=|\psi\rangle\langle\psi|$ and the mixed density
$\rho_{\rm mix}=(|u\rangle\langle u|+|v\rangle\langle v|)/2$
agree on every native observable $A\in M_{11}(\CC)\oplus M_4(\CC)$
but are not equal as full $15\times15$ matrices.
In particular even full native tomography and the exact coordinate
fingerprint do not reconstruct off-block coherences or physical
construction provenance.
\end{proposition}
\begin{proof}
The difference is
$\rho_{\rm coh}-\rho_{\rm mix}
=(|u\rangle\langle v|+|v\rangle\langle u|)/2$.
Its trace against any block-diagonal operator vanishes, while the
difference is nonzero on $M_{15}(\CC)$.
\end{proof}

\begin{proposition}[Zero-fermion first jet cannot create the missing odd source]
\label{p2f:prop:v264-action-jet-obstruction}
On the printed III-F uninserted local bilinear action domain
$S(B,\bar\Psi,\Psi)=S_B(B)+c_F\bar\Psi D(B;Y)\Psi$,
its first fermionic variation at $\Psi=\bar\Psi=0$ vanishes,
whereas its mixed second variation is the paired-index operator
$\delta_{\bar\Psi}\delta_\Psi S=c_FD(B;Y)$.
Without a separately typed owner map, this even two-jet neither
constructs a nonzero \emph{one-spinor} source nor defines the
$\mathcal H_1\otimes\mathcal H_1$ mixed exchange Hessian required
by the IV-A population theorem.
\end{proposition}
\begin{proof}
The derivative with respect to $\bar\Psi$ is
$c_FD(B;Y)\Psi$, which is zero at $\Psi=0$, and taking the
remaining derivative yields $c_FD(B;Y)$ with both fermion indices.
Under the simultaneous central sign change of $\Psi,\bar\Psi$ the
bilinear is invariant. Moreover Paper~V-A's rooted chiral operator
acts on a separate protected/internal field owner, whereas the II-B
central-odd mixed block has source $\mathcal H_s$ and target
$\mathcal H_g$. Equality of coefficients with different domains
requires a bridge map that the printed source has not supplied.
\end{proof}

\begin{remark}[Scope of the paired-exchange contraction]
\label{p2f:rem:v264-action-pair-cert-status}
The canonical algebraic $K$ from
Theorem~\ref{p2b:thm:v264-canonical-paired-exchange} is invariant
under simultaneous $\Spin(4)$, but is not fixed by
$P_g\otimes I$ and cannot be expressed in the native even test
algebra. A proposed action Hessian $H_{\rm act}$ must therefore
be supplied with its own two-system owner and normalization before
one evaluates
$\|(P_s\otimes P_g)H_{\rm act}(P_g\otimes P_s)\|_{\rm HS}^2$.
A nonzero value proves a cross block only; it does not imply the
canonical exchange law, bath preparation, collision statistics,
full field-envelope $\mathsf{FENV}$, or an actually observed
stationary population.
\end{remark}

\begin{remark}[Homogeneous covariant readout does not decide a detector or instrument]
\label{p2f:rem:v270-homogeneous-readout-firewall}
The newly constructed continuous six-channel $G/U(1)_{\rm diag}$
POVM of Paper~II-C is native and informationally complete on
$M_{11}\oplus M_4$ at every radius $r>\sqrt3\Rmom$; it is not
unique and does not select a physical effect law or CP instrument.
The $r=\sqrt3\Rmom$ blindness and existing equal-first-moment,
unequal-second-moment countermodels remain unchanged. The marked
noncovariant finite $137$-outcome construction remains distinct from
continuous $\Spin(4)$-covariance.
\end{remark}



\begin{corollary}[Arbitrary finite polynomial endpoint moments admit exact positive finite presentations]
\label{p2f:cor:v274-finite-endpoint-moment-presentations}
Let $K:S^3\to M_{15}(\mathbb C)$ be an admissible continuous
positive endpoint-POVM density of polynomial degree at most two,
with $\int K\,d\sigma=I_{15}$. For each integer $r\ge0$ there is
a positive finite probability $\nu_r=\sum_iw_i\delta_{q_i}$ with
at most
\[
 d_{r+2}=\dim\mathcal P_{r+2}
   =\frac{(r+3)(r+4)(2r+7)}6
\]
points such that $E_r(B)=\sum_{q_i\in B}w_iK(q_i)$ is a normalized
finite-outcome POVM satisfying
\[
 \boxed{\int f\,dE_r=\int_{S^3}f(q)K(q)\,d\sigma(q)
       \qquad\text{for all } f\in\mathcal P_r.}
\]
For $r\ge1$ it preserves the native exact coordinate first moments.
For $r=17$, one can instead use the explicit positive degree-nineteen
design of Theorem~\ref{p1b:thm:v273-triple-icosahedral-nineteen-design}
or the at-most-fourteen $2I$-coset mixture of
Corollary~\ref{p1b:cor:v273-positive-coset-mixture-no-ceiling},
which has at most $1680$ weighted points.

There is also a \emph{branch-independent}, quantitatively convergent
finite presentation relative to a fixed tetrahedral core $H$: if
$E_{H,n}(B)=\int_BK\,d\xi_{H,n}$ for the finite twin measure of
Theorem~\ref{p1b:thm:v274-branch-symmetric-twin-convolution}, then
$E_{H,n}(S^3)=I_{15}$ exactly and, for every $f\in\mathcal P_{17}$,
\[
 \boxed{\left\|\int f\,dE_{H,n}-\int fK\,d\sigma\right\|_{\rm op}
 \le \sqrt{13(a^{2n}+b^{2n})}\,\|f\|_\infty\,
          \sup_{q\in S^3}\|K(q)\|_{\rm op},}
\]
where $a=275/8192$ and $b=525/8192$.
All these are exact or controlled \emph{mathematical endpoint readout}
representations on the pre-existing carrier. They do not constitute
a physical finite-outcome instrument, retain full continuous covariance
for a single finite support, or establish finite-subgroup occurrence.
\end{corollary}
\begin{proof}
The normalized Haar moment vector for a real basis of the
$d_{r+2}-1$ nonconstant elements of $\mathcal P_{r+2}$ belongs to
the convex hull of their compact point-evaluation image.
Carath\'eodory's theorem supplies at most $d_{r+2}$ positive
nodes reproducing the Haar integral on that space.
Because $K$ has degree at most two, each entry of $fK$ has degree
at most $r+2$, so matrix integration is exact, and $f=1$ gives
normalization. The degree-nineteen assertions follow from the cited
positive designs and the same degree count.

For $E_{H,n}$, the twin measure already integrates $K$ exactly,
because it agrees with Haar through degree eleven. For unit vectors
$u,v\in\mathbb C^{15}$, the scalar polynomial
$f(q)\langle u,K(q)v\rangle$ has degree at most nineteen and
$L^2$ norm at most
$\|f\|_\infty\sup_q\|K(q)\|_{\rm op}$. Apply the sharp scalar
moment-functional norm from the twin theorem and then take the
supremum over $u,v$ to prove the operator-norm bound.
\end{proof}


\begin{corollary}[Exact positive finite endpoint-readout tower and its sharp topology]
\label{p2f:cor:v275-endpoint-finite-tower}
Let $K:S^3\to M_{15}(\mathbb C)$ be the continuous positive
polynomial density of degree at most two, with
$\int K\,d\sigma=I_{15}$. For the nested positive finite word
measures $\nu_m$ of
Theorem~\ref{p1b:thm:v275-nested-positive-resolution-tower},
put $E_m(B)=\int_B K(q)\,d\nu_m(q)$ for $m\ge1$. Then
$E_m$ is a normalized finite-outcome POVM and
\[
 \int f\,dE_m=\int fK\,d\sigma
 \quad(f\in\mathcal P_{2m-2}).
\]
For each fixed continuous scalar $f$, the matrix moments converge
in operator norm:
\[
 \lim_{m\to\infty}
 \left\|\int f\,dE_m-\int fK\,d\sigma\right\|_{\rm op}=0.
\]
But no individual $E_m$ is the same POVM on all Borel events as
$E(B)=\int_BK\,d\sigma$: for the finite event
$B_m=\operatorname{supp}\nu_m$,
\[
 \boxed{E_m(B_m)=I_{15},\quad E(B_m)=0,\quad
 \sup_B\|E_m(B)-E(B)\|_{\rm op}\ge1.}
\]
Thus exact finite moments and weak-continuous readout convergence
do not imply equality of sharp event statistics or a physical
finite-outcome instrument.
\end{corollary}
\begin{proof}
Degree-$2m$ exactness includes the entries of $K$ for $m\ge1$,
so normalization follows. The product $fK$ has degree at most
$2m$ whenever $\deg f\le2m-2$, proving exact moments.
For continuous $f$, each entry of $fK$ is a continuous scalar
function on compact $S^3$ and is uniformly approximable by
restricted polynomials. The positive probability rules have
norm one on continuous scalars and are eventually exact on
any fixed polynomial, so their entrywise integrals converge;
finite matrix dimension upgrades this to operator norm.
Haar assigns zero measure to each finite $B_m$; the finite
measure $\nu_m$ is concentrated there. The displayed event
values follow directly, completing the obstruction.
\end{proof}


\begin{corollary}[Finite observable-history certificate with a whole-candidate firewall]
\label{p2f:cor:v297-outcome-history-gp3-firewall}
Fix an actually registered finite outcome/effect subgrammar of a proposed
whole-candidate momentum signature, on a finite carrier $H$. For its
represented multiplicity-blind test algebra $\mathcal A$ and the
state restriction $Q$ of
Theorem~\ref{p2b:thm:v293-operational-multiplicity-quotient},
all finite probabilities and terminal readouts in \emph{that}
subgrammar are quotient-faithful if and only if the finite closure
$\mathcal O_{N^2-d}$ of
Theorem~\ref{p4a:thm:v297-finite-outcome-word-saturation}
is contained in $\mathcal A_{\rm sa}$. A failure supplies a finite
history and two quotient-equivalent states with unequal probabilities.

This criterion neither verifies the registration of the subgrammar
nor certifies preservation of additional frozen operations, complete
state updates, source response forms, or same-owner scalar trace ports.
In particular, the weaker observable-history condition cannot replace
the complete generated GP1/GP2-faithful proper comparison required by
Proposition~\ref{p2:prop:gp3-multiplicity-bound}. Thus it can
\emph{exclude} a quotient for a specified observable without selecting
or excluding every full physical momentum candidate.
\end{corollary}
\begin{proof}
Theorem~\ref{p4a:thm:v297-finite-outcome-word-saturation} gives the
finite if-and-only-if and the witness; the quotient state/effect duality
is exactly Theorem~\ref{p2b:thm:v293-operational-multiplicity-quotient}.
The whole-candidate requirement is independent by
Proposition~\ref{p2:prop:gp3-multiplicity-bound}, which retains every
other frozen port. The distinction from full-process descent is
witnessed by Proposition~\ref{p4a:prop:v297-specific-versus-full-process}.
\end{proof}


\begin{corollary}[Informational completeness compresses the outcome-port audit]
\label{p2f:cor:v298-ic-one-step-gp3-firewall}
If a proposed physical momentum candidate registers a complete
informationally complete POVM for its represented exact test algebra
$\mathcal A$ together with $b$ finite outcome maps, then the outcome
history portion of a candidate operational quotient is faithful exactly
when the finitely many one-step port pullbacks pass the test of
Theorem~\ref{p4a:thm:v298-ic-one-step-outcome-descent}.
The native kinematic algebra admits a mathematical $43$-effect frame by
Corollary~\ref{p4a:cor:v298-kinematic-43-effect-frame}; on a faithful
coordinate implementation the mathematical $137$-effect position frame
is supplied by
Corollary~\ref{p2c:cor:v298-amplified-coordinate-ic-one-step}.
These yield $43b$ or $137b$ finite checks respectively, \emph{conditional}
on actual frame and instrument registration. The complete frozen
GP3 comparison still additionally requires preserved sources,
unnormalized scalar trace owners, all other occurring ports, and an
admissible proper generated whole-candidate comparison as in
Proposition~\ref{p2:prop:gp3-multiplicity-bound}.
\end{corollary}
\begin{proof}
The equivalence and one-step witness follow from
Theorem~\ref{p4a:thm:v298-ic-one-step-outcome-descent}; the two finite
readout-frame capacities follow from the indicated corollaries.
None of those purely observable-algebra tests establishes source or
scalar-trace preservation, nor promotes a mathematical POVM to an actual
instrument. Therefore the full GP3 comparison must retain its independent
whole-candidate premises.
\end{proof}


\begin{corollary}[Exact covariance selects the continuous IC audit route]
\label{p2f:cor:v299-covariant-ic-gp3-firewall}
For the native six-type $\operatorname{Spin}(4)$ test signature,
an unmarked fixed finite outcome set with exact permutation
covariance cannot give an informationally complete readout:
Theorem~\ref{p2c:thm:v299-finite-covariant-six-ceiling} bounds its
Hermitian effect span by six. The available noncovariant $43$- and
$137$-effect mathematical frames therefore cannot silently be
promoted to exact continuously covariant fixed-outcome instruments.

On the separate directed-tangent branch at
$r>\sqrt3\Rmom$, Corollary~\ref{p2c:cor:v299-directed-tangent-137-event-basis}
provides $137$ measurable events in the existing continuous
covariant six-channel POVM. If \emph{that} readout and $b$ specific
outcome maps are physically registered on one owner, the exact
observable-history quotient test is the $137b$ finite membership
criterion of Theorem~\ref{p4a:thm:v299-continuous-ic-event-test}.
Passing it preserves the specified terminal-event grammar and its
finite histories; it does not establish a generated proper
whole-candidate GP3 comparison or preservation of additional
source, normalized/unnormalized trace, action or other registered
ports as required by Proposition~\ref{p2:prop:gp3-multiplicity-bound}.
A continuously rotating \emph{family} of finite detectors with an
explicit frame mark is not a single finite-outcome POVM with a
permutation action and remains logically distinct from this no-go.
\end{corollary}
\begin{proof}
The first conclusion is the exact-covariance ceiling. The existing
continuous directed-tangent construction and its finite event basis
supply all Hermitian coordinate-test directions without a finite
covariant outcome set. The event-basis theorem then reduces every
registered finite outcome history to the indicated finite tests.
The GP3 proposition additionally requires a complete generated
comparison preserving all frozen roles; no part of this mathematical
readout argument manufactures those remaining inputs.
\end{proof}

\begin{corollary}[Mathematical instrument availability does not register a GP3 comparison]
\label{p2f:cor:v300-instrument-existence-gp3-firewall}
The continuously $\Spin(4)$-covariant, full-span, strict-radius
mathematical directed-tangent POVM of Paper II-C admits a canonical
covariant square-root instrument that preserves every
multiplicity-blind intrinsic coordinate test and every finite
outcome history by
Corollary~\ref{p2c:cor:v300-directed-tangent-luders}.
For post-unitary updates, the exact compatibility class is the
normalizer of Proposition~\ref{p2c:prop:v300-coordinate-normalizer}.
An alternative instrument can share the same POVM effects yet fail
quotient preservation by
Proposition~\ref{p4a:prop:v300-swap-instrument-witness}.
Thus mathematical availability of an exact symmetry-compatible
measurement/update \emph{closes the existence and outcome-descent
compatibility question}, not the physical selection question.

In particular, the complete proper GP3 comparison of
Proposition~\ref{p2:prop:gp3-multiplicity-bound} still requires
an occurring, generated same-owner instrument/measurement port,
all source and scalar-action trace outputs, and preservation of every
other frozen physical input. No conclusion of universal momentum
multiplicity-one follows merely from the new square-root instrument.
\end{corollary}
\begin{proof}
The cited construction provides all instrument and observable algebra
identities. The normalizer result distinguishes all compatible
post-unitary updates inside the specified full-rank class, while the
swap witness shows why effects alone cannot select an update rule.
The GP3 proposition independently demands a proper, generated,
whole-candidate comparison preserving all frozen ports; no
mathematical CP instrument construction supplies its G0 occurrence.
\end{proof}

\begin{corollary}[Conditional naturality certificate for all momentum measurement histories]
\label{p2f:cor:v301-gauge-naturality-gp3-certificate}
For a faithful intrinsic-coordinate momentum implementation of
Theorem~\ref{p2c:thm:v301-intrinsic-gauge-fixedpoints}, assume the $U(k_g)\times U(k_s)$ implementation gauge
belongs to $\operatorname{Aut}_0(P;w_B)$ for the \emph{complete
frozen input presentation} of every actually registered, GP2-licensed
outcome constructor word $w_B$, with its outcome label fixed by that
gauge. In particular, the gauge must preserve all retained sources,
marks, boundary data, instrument schemas and settings; unmarkedness
of the carrier alone does not certify this premise. The GP2
same-owner equivariance theorem
\ref{p2a:thm:v368-same-owner-unmarked-equivariance}
then forces equivariance of each registered outcome map, and
Theorem~\ref{p4a:thm:v301-gauge-natural-instrument-descent}
proves quotient preservation for every finite outcome history and
every native coordinate terminal effect; no further $137$-event
membership test or particular postmeasurement-update selection
is necessary for that specified grammar.

The premise is materially stronger than geometric
$\Spin(4)$ covariance: the same full-rank continuous POVM can
be followed by a $\Spin(4)$-covariant update that distinguishes
multiplicity, by
Proposition~\ref{p4a:prop:v301-spin4-same-povm-leak}.
Showing that the full gauge is indeed a complete-presentation
unmarked automorphism for every \emph{occurring} outcome word remains a
separate source-level certificate; it is not an additional physical
symmetry postulate. Once this antecedent holds, GP2 supplies the
naturality automatically. Even when it holds, the GP3 whole-candidate
comparison of Proposition~\ref{p2:prop:gp3-multiplicity-bound}
still requires same-owner source ports, scalar-action/trace outputs,
other frozen observable roles, and an actually generated proper
whole-candidate comparison. None of those requirements follows
from measurement-history blindness alone.
\end{corollary}
\begin{proof}
The fixed-point identity identifies the exact observables of the
unmarked multiplicity gauge. The Paper II-A same-owner theorem
transports each licensed output under every certified complete-input
automorphism, giving the required equivariance of its CP map.
Apply the outcome-wise intertwining
criterion and its word-stability conclusion to every registered
primitive and hence every generated outcome word. The countermodel
proves strict separation from $\Spin(4)$ covariance. The final
GP3 statement is the complete-comparison firewall of the cited
proposition; no mathematical gauge redundancy automatically
licenses a proper physical competitor.
\end{proof}


\section{Petersen axis readout, group actions, and exact finite certificates}
\label{p2f:sec:v318-petersen-comparison-certificate}

\begin{proposition}[Matrix and incidence certificates for the ten-axis frame]
\label{p2f:prop:v318-ten-axis-matrix-certificate}
Let $D\in\{0,1\}^{10\times5}$ be the unsigned incidence matrix of
all unordered core pairs, $D_{ij,m}=\mathbf1_{m\in\{i,j\}}$.
Then $D^{\mathsf T}D=3I_5+J_5$ and the orthogonal projector onto
the trace-free quadrupole coefficient sector is
\[
 P_Q=I_{10}-D(D^{\mathsf T}D)^{-1}D^{\mathsf T}
     =\frac16(4I_{10}+2A_P-J_{10}).
\]
For $Q_{ij}=u_{ij}u_{ij}^{\mathsf T}-I_3/3$,
\[
 \operatorname{Tr}(Q_{ij}Q_{kl})=\frac43(P_Q)_{ij,kl},
 \qquad \sum_{i<j}u_{ij}u_{ij}^{\mathsf T}=\frac{10}{3}I_3.
\]
Together these are exact certificates of the frame reconstruction
and auxiliary positive POVM of Paper~II-C.
\end{proposition}
\begin{proof}
Every core label belongs to four pairs and each two distinct
labels share one pair, giving $D^{\mathsf T}D=3I+J$.
The projector identity and Gram identity are supplied by the
Paper-I-B Petersen/Hopf angular theorem.  The tight-vector identity
also follows from the invariant real three-dimensional $A_5$ action,
whose commutant is scalar, and the total trace ten.
\end{proof}

\begin{lemma}[Five-probe rank and exact combinatorial count]
\label{p2f:lem:v318-five-probe-count}
For a chosen five-set $S$ of pair labels, let
$T=\binom{[5]}2\setminus S$.  Write $D_T$ for the restriction of $D$
to rows in $T$ and $\Phi_S$ for the restriction of the quadrupole
synthesis $\Phi(a)=\sum_{ij}a_{ij}Q_{ij}$ to $S$.
Then
\[
 \boxed{\operatorname{rank}\Phi_S=\operatorname{rank}D_T.}
\]
There are $162$ rank-five choices and $90$ rank-four choices.
\end{lemma}
\begin{proof}
$\ker\Phi=\operatorname{im}D$ and $D$ is injective.
Thus $\ker\Phi_S\cong\ker D_T$, and with $|S|=5$,
$\operatorname{rank}\Phi_S=5-\dim\ker D_T=\operatorname{rank}D_T$.
For an edge $ij\in T$, the equation $f_i+f_j=0$ alternates signs
along the edge.  A connected component has nonzero solutions
exactly when it is bipartite.  Since $T$ has five vertices and five
edges, no bipartite component requires one connected odd cycle.
The permitted graphs have cycle length five ($12$ cases) or a
triangle with two attached labelled vertices ($10\cdot15=150$
cases).  The remaining $252-162=90$ cases have rank four;
one can confirm the final rank count by enumerating all $252$
integer $5\times5$ submatrices $D_T$.
\end{proof}

\begin{proposition}[Finite reciprocal product readout versus real bilocal correlations]
\label{p2f:prop:v318-reciprocal-real-effect-certificate}
For real three-dimensional vector spaces $V_x,V_p$,
$\operatorname{End}(V_x)_{\rm sym}$ and
$\operatorname{End}(V_p)_{\rm sym}$ each have dimension six.
The $10\times10$ local Hopf-axis product effects span their
36-dimensional tensor product, while
\[
 \operatorname{End}(V_x\otimes V_p)_{\rm sym}
 \cong\big(\operatorname{Sym}^2V_x\otimes
           \operatorname{Sym}^2V_p\big)
  \oplus\big(\Lambda^2V_x\otimes\Lambda^2V_p\big)
\]
has dimension $36+9=45$.
No product of the stated local real-symmetric effects detects
its nine-dimensional skew--skew complement.  This is a finite
mathematical blind sector, not a physical Born/position theorem.
\end{proposition}
\begin{proof}
Transpose acts on elementary tensors by
$(A\otimes B)^{\mathsf T}=A^{\mathsf T}\otimes B^{\mathsf T}$.
Hence a symmetric total tensor has either symmetric--symmetric
or skew--skew parity.  The local effects linearly generate all
symmetric operators by the ten-axis frame theorem, proving the
36-dimensional span and the nine-dimensional complement.
\end{proof}

\begin{remark}[Reproducibility, spectral types, and typed limits]
\label{p2f:rem:v318-certificate-status}
The accompanying exact-arithmetic verifier checks all ten-axis Gram
entries, the positive frame identity, state reconstruction, the full
252-subset rank classification, the $S_3$ vertex stabilizer, and the
odd-$S_5$ exchange of the two edge-cycle rank-three projectors.
These are mathematical consequences of an embedded $2I$ packet.
The verifier neither establishes actual physical $2I$ support,
native detector effects, a full-$SU(2)$ finite-outcome covariance,
nor a new energy, particle number or mass formula.
\end{remark}



\section{Finite-incidence descent versus constitutional realization}
\label{p2f:sec:v339-petersen-descent-scope}
\begin{proposition}[Rank-fifteen cochain descent and non-descending cell exclusion]
\label{p2f:prop:v339-projector-descent-certificate}
Use the five-core incidence map $B$ of
Proposition~\ref{p1:prop:v339-incidence-rank15-contrast} and
the rank-fifteen register projection $T$. Its restriction to
$\operatorname{im}T$ is injective and obeys
\[
 \|BM\|_{\ell^2(2I)}^2=30\|M\|_{\rm HS}^2,
 \quad M\in\operatorname{im}T.
\]
Consequently $\Pi_{15}=\tfrac1{30}BTB^*$ is a genuine
orthogonal projection on $\ell^2(2I)$, and the group-function
incidence role sees all fifteen oriented Petersen edge modes.
By contrast, for a fixed basis register $e_{ij}$ the rank-one
exclusion $L_{ij}=I-|e_{ij}\rangle\langle e_{ij}|$ on
$\mathbb R^{25}$ does \emph{not}, in general, preserve
$\ker B=(P\otimes Q)\oplus(Q\otimes P)$.
It therefore does not define a canonical descended operator
$\widehat L_{ij}$ satisfying $\widehat L_{ij}B=BL_{ij}$.
The separate $24/25$ single-register self-exclusion quantity
must not be substituted into the rank-fifteen physical export.
\end{proposition}
\begin{proof}
The singular-square value of $B$ on $Q\otimes Q$ is $30$;
$T$ is the orthogonal projection onto its trace-free component,
so $T^2=T=T^*$, $B^*B\,T=30T$, and
$(BTB^*/30)^2=BTB^*/30$. For non-descent choose a nonzero
matrix in $P\otimes Q$ with a nonzero $(i,j)$ entry.
Its image under $L_{ij}$ has a nonzero $Q\otimes Q$ component,
and cannot remain in $\ker B$. The same argument works for
$Q\otimes P$. Hence $L_{ij}$ cannot factor through the
quotient by $\ker B$ without extra marked data.
\end{proof}

\begin{remark}[Constitutional status: which observable is owned?]
\label{p2f:rem:v339-petersen-gp2-gp3-firewall}
The map $U$, the ordinary matrix commutator on
$\mathfrak{sl}(V)$, and $\Pi_{15}$ are exact, parameter-free
mathematical outputs of a supplied finite support and its
five-label incidence signature. If a GP2-frozen physical
constructor includes this support and its particular observable,
its whole generated image must be retained with the original
source, relation and test types. The finite proof by itself
provides no G0 occurrence, readout owner, action or dynamics.
In particular it does not discharge the conditional universal
GP3 multiplicity-one theorem for the unrelated $H_p$ of
Paper~II-B. The spinorial $\Spin(4)$ representation, the real
$A_5$ graph operator, and the 24-cell association algebra
remain distinct owners even when their dimensions happen to
match. An empirical distinction would require an identified
finite-support observable with the stated projector leverage,
not merely the existence of its finite geometric kernel.
\end{remark}


\section{Normalization and orthogonality of rank-two projective blocks}
\label{p2f:sec:v340-projective-binary-certificate}
\begin{proposition}[No rank-twenty-five register/kernel identification]
\label{p2f:prop:v340-projective-binary-orthogonality}
Let $\pi:2I\to A_5$ be the two-to-one quotient, and let
$\iota:\mathbb R[A_5]\to\mathbb R[2I]$ be normalized centre-even
pullback, $(\iota f)(g)=f(\pi(g))/\sqrt2$. If $\widehat B$
be the original $120\times25$ transporter incidence map and
$B_{\rm tr}$ its $60\times25$ projective version, then
\[
 \widehat B=\sqrt2\,\iota B_{\rm tr},\qquad
 \widehat B^{\mathsf T}\widehat B=2B_{\rm tr}^{\mathsf T}B_{\rm tr},
\]
\[
 \operatorname{im}B_{\rm tr}
 =\mathbf1\oplus\operatorname{im}\mathscr P_4,
 \quad \operatorname{rank} B_{\rm tr}=17,\quad
 \mathscr P_5B_{\rm tr}=0,
 \quad \operatorname{rank}\mathscr P_5=25.
\]
Under the existing nearest-shell character hypothesis,
$\iota\operatorname{im}\mathscr P_5$ is the
$25$-dimensional shell-walk zero-eigenspace and is orthogonal
to the image of all twenty-five transporter labels. It is therefore
impossible to identify the two spaces by counting the shared number
$25$. The marked $2T$ shell further decomposes that zero-eigenspace
as $15\oplus10$ from the branching
$\mathbf5\downarrow_{A_4}
=\chi_\omega\oplus\chi_{\bar\omega}\oplus\mathbf3$.
\end{proposition}
\begin{proof}
An element of a transporter cell contains both binary lifts of
each projective permutation; the displayed normalization is
immediate. The centred transporter identity from
Proposition~\ref{p1:prop:v340-transporter-hexagon-intersection}
identifies its nonconstant image with the central $\mathbf4$
isotypic block; its constant image has rank one. Schur
orthogonality puts the inequivalent $\mathbf5$ block in its
orthogonal complement. The restrictions to $A_4$ give eigenvalues
$-1/2,-1/2,0$ on $\chi_\omega,\chi_{\bar\omega},\mathbf3$,
so five right-multiplicity copies split as $10\oplus15$.
\end{proof}
\begin{remark}[Finite certification scope]
\label{p2f:rem:v340-ranktwo-proof-scope}
The accompanying \path{Paper00-RankTwoIncidence-v1_2.py} checks
the exact rational/integer incidence identities, Reye ranks,
character idempotents, signed double-cover polynomials and the
105-element orientation-sign stabilizers. This is a finite theorem
certificate, not a test of G0, GP3, dynamics or detector access.
\end{remark}



\begin{proposition}[Centre-even descent of class-tomography coefficients]
\label{p2f:prop:v342-class-frame-descent-scope}
Let $\pi:2I\twoheadrightarrow A_5$ and let $Z_g$ be the
trace-free four-contrast representation matrices of
Theorem~\ref{p1:thm:v342-class-frame-three-cycle-tomography}.
The coefficient map $g\mapsto Z_{\pi(g)}/\sqrt{30}$ on $2I$
induced from $g\mapsto Z_g/\sqrt{15}$ on $A_5$ is isometric
for the normalized class-frame reconstruction, because each
projective class member has exactly two binary lifts.  This
centre-even pullback retains the full rank-fifteen frame on the
three-cycle lifts and the same $5A/5B$ class-frame eigenvalues.
It does not identify either three-cycle orientation with a
physical chirality owner and cannot be promoted to a measurable
POVM by frame positivity alone.  In particular it provides no
whole-candidate GP3 comparison for the independently generated
fifteen-complex-dimensional momentum carrier $H_1$.
\end{proposition}
\begin{proof}
With the normalized rows $f_g=Z_g/\sqrt{15}$ and their double lifts
$f_{\widetilde g}=f_{\pi(\widetilde g)}/\sqrt2$, we obtain
\[
\sum_{\widetilde g\in\pi^{-1}(C)}
 |f_{\widetilde g}\rangle\langle f_{\widetilde g}|
 =\sum_{g\in C}|f_g\rangle\langle f_g|.
\]
The complete three-cycle frame follows from the positive
rank-fifteen expression for $\mathcal A_{3A}$ in
Theorem~\ref{p1:thm:v342-class-frame-three-cycle-tomography}.
The missing physical-ownership and GP3 hypotheses are explicit
in Remark~\ref{p1:rem:v342-class-frame-observable-firewall}.
\end{proof}


\begin{proposition}[Independent symmetry, matched readouts and GP3 limits]
\label{p2f:prop:v343-signed-owner-firewall}
For independently acting $A_{5,x}\times A_{5,p}$, the marked
factorwise central tensor operator $\mathcal C_{xp}$ of
Theorem~\ref{p1b:thm:v343-doubleodd-contrast} is invariant,
but the matched Born-even/odd decomposition is not invariant
without an intertwiner identifying $W_x\cong W_p$.
There is no nonzero invariant linear scalar bilinear
$W_x\times W_p\to\mathbb R$ under the independent group,
whereas under matched diagonal $A_5$ there are exactly four
independent invariant bilinears, one per irreducible summand.
The unique bilinear odd under outer fivefold-class interchange
is $\langle u,Lv\rangle$ up to scale.
The \emph{double-odd} requirement picks a one-dimensional
central tensor direction, but does not follow from Born exchange
alone.  Its coefficient in an action, any GP2 physical effect
realization and any GP3 whole-candidate multiplicity comparison
remain independent conditions.
\end{proposition}
\begin{proof}
Because neither factor module contains an $A_5$-fixed vector,
$\operatorname{Hom}_{A_{5,x}\times A_{5,p}}(W_x\otimes W_p,\mathbb R)=0$.
After matching the groups, Schur's lemma gives one scalar pairing
per self-dual multiplicity-one $A_5$ sector.  The outer odd
permutation exchanges the two three-dimensional sectors, so its
odd pairing line is generated by the contrast $L$ of
Proposition~\ref{p1:prop:v343-fivefold-signed-contrast}.
The coefficient-space and flip analysis in
Theorem~\ref{p1b:thm:v343-doubleodd-contrast} completes the
conditional statement.  No group-theoretic classification here
establishes physical sector occurrence.
\end{proof}

\printbibliography
\end{document}
